% =====================================================================
%  H. Brøchner
%  Om det Religiøse i dets Enhed med det Humane.
%  Et positivt Supplement til „Problemet om Tro og Viden“.
%  Kjøbenhavn: P. G. Philipsens Forlag, 1869.   [IV] + 211 pp.
%  Printer: Bianco Lunos Bogtrykkeri ved F. S. Muhle.
%
%  DIPLOMATIC TRANSCRIPTION (Danish), from the Royal Danish Library scan.
%  Copy digitised: KB shelfmark DA 1.-2.S 3 8°.
% =====================================================================
%
%  ---------------------------------------------------------------
%  1. PAGE MAP  (verified by eye, recon pass, 2 September 2026)
%  ---------------------------------------------------------------
%  The scan is 229 PDF pages, 250 ppi RGB JPEG, one image per page, with
%  no separate 1-bit mask. The map is UNIFORM over the whole book:
%
%      printed p. N   =   PDF page N + 11        (printed 1--210 = PDF 12--221)
%      Trykfeil leaf  =   printed p. 211 = PDF 222
%
%  Front matter, PDF 8 = title page [I], PDF 9 = title verso [II],
%  PDF 10 = Indhold [III], PDF 11 = Indhold IV. PDF 1--7 and 223--229 are
%  KB apparatus, boards, and blank flyleaves; there is no half-title leaf.
%
%  Verified by rendering the head band of every leaf from PDF 12 to PDF 222
%  and reading all 211 printed folios by eye. The offset holds with zero
%  exceptions. Two independent corroborations that no leaf is doubled or
%  skipped: the ABBYY folio sweep contradicts +11 nowhere, and the
%  binding-shadow side alternates verso-left / recto-right across all 211
%  leaves without exception. Three leaves carry no folio by design:
%  PDF 12 (p. 1, opening page of the treatise), PDF 16 (p. 5, opening page
%  of division I), and PDF 222 (p. 211, the Trykfeil leaf). Where the ABBYY
%  text layer appears to contradict the map -- notably PDF 57 and PDF 67 --
%  it has simply misread a perfectly legible numeral (46 and 56).
%
%  ---------------------------------------------------------------
%  2. WHAT IS TRANSCRIBED, AND HOW
%  ---------------------------------------------------------------
%  Diplomatic. 19th-century orthography exactly as printed: Christendommen,
%  Videnskaben, Theologie, høiere, aa, capitalised nouns. Printer's errors
%  are transcribed AS PRINTED, logged in a % comment at the site, and never
%  silently corrected -- including the one erratum the author himself lists
%  (see §5). Quote-balance counts therefore drift off zero on purpose.
%
%  Quotation marks: „…“, low open / high close, throughout. Checked on the
%  title page and on pp. 1, 17, 56, 123, 202. (This is the same practice as
%  the 1868 volume, and NOT the same-sort-at-both-ends practice found in
%  Nielsen's 1842 volume.) Em-dash: ---.
%    ONE EXCEPTION, and it follows the fount, not the sense: where the
%    quoted matter is the antiqua Latin, the book sets ANTIQUA GUILLEMETS
%    «…» round it -- „Deum pati“ is printed («Deum pati») on p. 41.
%    Verified on the image at 300 dpi. Kept as printed.
%
%  Emphasis: the book's only emphatic device inside a paragraph is
%  letterspacing (Sperrsatz), and it is used heavily -- whole phrases, not
%  just single words. All of it is rendered \emph{}. Bold (halbfette)
%  Fraktur occurs ONLY in display heads and in the top-level entries of the
%  Indhold; it never appears inside a paragraph.
%    The compositor's letterspacing is INCONSISTENT ACROSS A LINE BREAK,
%    and sometimes within a word: often only the first half of a word or
%    phrase is spaced (A a n d s-functioner p. 6; s a m m e  R e-sultat
%    p. 15; \emph{Tids}bestemmelsen p. 18; \emph{Erkjendelses}forhold
%    p. 32), and once the spacing runs past the word's end
%    (\emph{Livsforhold i Be}tydning, p. 32). Occasionally it starts
%    mid-word (Væsen\emph{stotaliteten}, p. 19). All of this is transcribed
%    AS PRINTED and flagged in a % comment at the site; it is not
%    regularised. Names are letterspaced Fraktur, never antiqua --
%    Kierkegaard, Grundtvig, Martensen, Feuerbach, Spinoza, Hegel.
%
%  Antiqua within the Fraktur: LATIN WORDS ONLY -- and the antiqua may
%  itself be letterspaced or not, page by page. „causa sui“ and „causa
%  finalis“ are letterspaced antiqua on p. 68 (verified at 300 dpi against
%  the tight antiqua of pp. 64 and 71); „conditio sine qua non“ on p. 59 is
%  half and half, „conditio“ tight and „sine qua non“ spaced; „modus“
%  (p. 83), „eo ipso“ (p. 89) and „implicite“ (p. 91) are tight. Spaced
%  antiqua is encoded \emph{\textit{…}}, tight antiqua \textit{…}.
%  Proper names stay Fraktur,
%  and are usually letterspaced instead. The decisive page is p. 123, where
%  „diligere Deum propter Deum … se ipsum“ is upright antiqua while
%  „(St. Bernhard)“, inside the same Latin sentence, is Fraktur. Not one
%  proper name was found set in antiqua anywhere in the book. The printed
%  Latin is UPRIGHT antiqua, not italic; it is rendered \textit{} here, for
%  consistency with the sibling volume `problemet-tro-viden' -- an
%  editorial choice, recorded as such.
%
%  Greek: real Greek occurs, always antiqua italic with accents. Neither
%  OCR witness reads any of it -- the ABBYY layer contains no Greek
%  anywhere in the file -- so every instance was caught by eye, page by
%  page. It is much commoner than the recon pass estimated: τέλος alone
%  runs to sixteen instances on pp. 30--36 and recurs at pp. 20, 23, 54,
%  65, 74 (twice), 75 and 185; p. 52 has ὁμοιωθῆναι θεῷ, p. 81 νόησις τῆς
%  νοήσεως and Νοῦς, p. 193 a longer word. Almost none of those pages was
%  predicted by any survey -- which is why every page is read for it. The
%  Greek enumerators α) β) γ) are antiqua italic too (Indhold, and pp. 31,
%  168, 177, 182).
%
%  The old id-est mark ɔ: (a reversed c followed by a colon) is used
%  constantly. It is typed as U+0254 + colon and rendered \reflectbox{c}.
%
%  NORMALISATION (the one departure from strict diplomacy, besides the
%  sort-level policy of TRANSCRIPTION-PLAYBOOK §2): Fraktur has a single
%  sort for I and J, so the print cannot distinguish them and the choice is
%  the editor's. It is resolved by the word -- I for I-words (Indhold,
%  Idee, Individet, „I denne“, never Jdee/Jndhold), J for J-words (Jfr.,
%  Jeg) -- following the house convention of `problemet-tro-viden'.
%
%  SHORT CENTRED RULES.  Settled in the verification pass of 3 September
%  2026, which swept all 210 printed pages mechanically for centred
%  horizontal rules (page images at 150 dpi, longest contiguous dark run per
%  scan line, tested for length, for centring on the measure, and for
%  isolation from type) and then read every hit on the image at 300 dpi.
%  Thirty-four such rules stand in pp. 1--210. The book uses them in two
%  distinct roles, and this edition treats the two differently:
%
%    (a) HEAD ORNAMENT -- a rule standing immediately above or immediately
%        below a display or run-in head. It belongs to the setting of the
%        head, not to the text, and is NOT transcribed. Sixteen of these:
%        p.   5  below the „I.“ division head
%        p.  30  between the „II.“ head and „1. Det Religiøses Væsen.“
%        pp. 43, 51, 56   above the heads „2.“, „3.“, „4.“
%        p.  64  below the „III.“ division head, and again above
%                „1. Det Absolutes Begreb.“   (two rules on the one page)
%        p.  97  above the run-in head „Ludvig Feuerbach.“
%        p. 126  above the run-in head „Den Hegelske Philosophie.“
%        pp. 128, 152, 164, 177, 182   above the display heads
%        p. 168  above the run-in head „α) Ophævelsen af Skylden.“
%        p. 202  below the „Slutning.“ head
%
%    (b) DIVISION-CLOSING RULE and MID-PAGE THOUGHT-BREAK -- every other
%        rule in the book. These ARE transcribed, uniformly, as
%        \begin{center}\rule{2cm}{0.4pt}\end{center}, with the printed
%        weight recorded in a % comment at the site. Eighteen of these:
%        closing a division or a sub-section, with the rest of the leaf
%        blank -- pp. 4 (end of the Indledning), 29 (end of I), 63 (end of
%        II), 89 (end of III.1), 141, 143, 201 (end of III) and 210 (end of
%        the Slutning, and of the book);
%        standing mid-page between two paragraphs, with no head adjacent --
%        pp. 19, 28, 42, 62, 97, 159, 175, 176, 190 and 208.
%
%  PRINTED WEIGHT.  The rule is a DOUBLE rule -- two parallel hairlines --
%  at pp. 4, 29, 63, 201 and 210: that is, at the close of the Indledning
%  and at the close of each of divisions I, II, III and the Slutning. Every
%  other rule in the book, head ornaments included, is a SINGLE hairline.
%  Each of the five doubles was checked on the image at 300 dpi at 2x. The
%  distinction is NOT reproduced typographically here -- one \rule serves
%  throughout -- but it is recorded at every site.
%
%  Two further findings of the sweep. A head does not always carry a rule:
%  pp. 90, 144, 165, 167, 190, 192, 193 and 195 have none. And no rule in
%  the book stands off centre: the „not centred“ rule reported at p. 175 in
%  the batch pass is centred to within 3 % of the measure, and is a
%  thought-break like the rest. (The FOOTNOTE separator rules are a
%  different thing altogether -- about one inch, set flush with the left
%  edge of the measure above each page's notes. They are part of the note
%  apparatus, are reproduced by LaTeX itself, and are never transcribed.)
%  The two rules on the Trykfeil leaf, printed p. 211, are that leaf's own
%  and are set as printed.
%
%  Words the compositor broke across a printed line or page are JOINED here,
%  with a discretionary hyphen and a comment character to eat the newline
%  („Cau\-“ then „%“) rather than with the printed line-break
%  hyphen: this edition reflows, so that hyphen is an artefact of the
%  setting and not part of the word. A SUSPENDED hyphen is different and is
%  kept as printed -- „Nødvendigheds- og Frihedsmomentet“.
%
%  NOT transcribed: signature numerals at the foot of the first page of
%  each gathering (1 on p. 1, 2 on p. 17, 7 on p. 97, …); the ornamental
%  two-line initials that open pp. 1, 5, 64, 97 and 202 (the letter is set
%  normally); the short centred rules that ornament a display or run-in
%  head, and the footnote separator rules (both above); and the
%  modern pencil marginalia of this copy (seen at pp. 16 and 112 at least),
%  which are a reader's, not the book's.
%
%  Footnotes are printed *) and **) and RESET PER PRINTED PAGE -- proved on
%  pp. 16 and 17, which each carry two. \thefootnote is therefore mapped
%  1 -> *), 2 -> **), and the counter is reset at the page marker of every
%  note-bearing page. (footmisc[perpage] is wrong here: it resets on the
%  TYPESET page, not the printed one.) A note may run over onto the
%  following page in the print (p. 6's **) does); it is set in full at its
%  call. There are more notes than the OCR sweep suggested: pp. 6, 14, 19,
%  24, 28, 33, 39, 40 and 41 all carry one that the sweep missed, and p. 22
%  carries none, contrary to it. Every page is checked by eye.
%
%  There are NO block quotations in this book. The long Feuerbach citations
%  of pp. 97--126 are run into the paragraph and marked by letterspacing
%  alone; p. 112 is a whole page of quotation without a single quotation
%  mark. They are not indented here either. There are no tables, figures or
%  displayed formulae. Displayed enumerations (1. 2. 3. …, hanging indent)
%  do occur, on pp. 1--4 and in the Slutning, pp. 202--210.
%
%  ---------------------------------------------------------------
%  3. THE HEADS, AND WHERE THE INDHOLD DISAGREES WITH THEM
%  ---------------------------------------------------------------
%  The book's own Indhold and its printed heads genuinely differ in EIGHT
%  places. Both readings are recorded; they are NOT normalised into
%  agreement. The disagreement is a datum. (The list ran to six until the
%  verification pass of 3 September 2026 collated every Indhold entry
%  against its printed head off the image; the two run-in heads of pp. 97
%  and 126 belong on it as well, and are added below. Every entry and every
%  page range on Indhold pp. [III]--IV was read again at 200 and 400 dpi in
%  that pass, and each divergence below was confirmed on both sides.)
%
%    p.   1  head: (none printed at all)   Indhold: „Indledning“ 1--4.
%            The division title used below is the Indhold's word; the body
%            simply begins. Editorial, and flagged at the site.
%    p.  90  head: „2. Det menneskelige Væsens Begreb …“
%            Indhold: „2. De menneskelige Væsens Begreb …“   (Det / De)
%    p.  97  head (run-in): „Ludvig Feuerbach.“
%            Indhold: „1) Ludvig Feuerbachs“   (enumerator, and genitive)
%    p. 126  head (run-in): „Den Hegelske Philosophie.“
%            Indhold: „2) den Hegelske Philosophies“   (likewise)
%    p. 159  head: „1. Forholdet mellem Nødvendigheds- og Frihedsmomentet
%            i det Onde.“   Indhold: the same without „i det Onde“.
%    p. 164  head: „2) Forsoningen, dens Kilde og Maal.“ standing on p. 164
%            Indhold: „2. Forsoningen, dens Kilde og Maal“ … 167--201.
%            Both the enumerator (2) / 2.) and the page differ.
%    p. 168  head: „α) Ophævelsen af Skylden.“
%            Indhold: „α) Skyldens Ophævelse“   (word-order variants)
%    pp. 190, 192  heads: „1. Den abstract metaphysiske …“, „2. Den
%            christelige Forsoningslære.“   Indhold: „1) den …“, „2) den …“
%
%  A printer's inconsistency INSIDE the Indhold, not listed in the Trykfeil:
%  entry „2. Forsoningen, dens Kilde og Maal“ is given as 167--201, but its
%  own first sub-entry „a) det i Forsoningen Virksomme“ is given as
%  165--167. Transcribed as printed. Re-read at 400 dpi in the verification
%  pass: both numerals are unambiguous, and every other page range in the
%  Indhold agrees with the print.
%
%  Two body heads have NO Indhold entry at all: „a) Opfattelsen af Synden og
%  af dens Virkning.“ (p. 193) and „b) Opfattelsen af Forsoningen og af
%  Ophævelsen af Syndens Virkninger.“ (p. 195). The Indhold runs straight
%  from „2) den christelige Forsoningslære“ 192--201 to „Slutning“. Both
%  heads are real -- centred, letterspaced regular Fraktur, like those of
%  pp. 165 and 167 -- and both are set in the body. They were missing from
%  the head table the batch agents worked from; a mechanical sweep of all
%  210 pages for centred display lines in the verification pass turned up no
%  third one.
%
%  Six rubric lines in the Indhold -- „Forberedende Begrebsudviklinger:“,
%  „Sammenligning af den udviklede Opfattelse med“, „Udvikling af
%  Problemerne:“, „Nærmere Udviklinger til Problemet om det Onde:“,
%  „Betingelser for Forsoningens Fuldstændighed:“, „Sammenligning af den
%  udviklede Opfattelse af Forsoningen med“ -- are printed in the CONTENTS
%  ONLY. The body has a lead-in sentence in each case instead. They are not
%  introduced into the body here.
%
%  Division extents: I. 5--29, II. 30--63, III. 64--201, Slutning 202--210,
%  Trykfeil 211. Corroborated by type-block measurement: pp. 29, 63, 201,
%  210 and 211 are the only short pages in the book.
%
%  ---------------------------------------------------------------
%  4. PRINTER'S DEFECTS AND DOUBTFUL READINGS
%  ---------------------------------------------------------------
%  Logged at the site in a % comment as they are found, and listed here.
%  Running quote-balance offset: 0 through p. 41.
%
%    p.  7   „hinandes“ for „hinandens“. Legible at 300 dpi, both OCR
%            witnesses agree, and „hinandens“ is set correctly on p. 5.
%            As printed. No effect on the quote balance.
%    p. 26   „I n h o l d“ (letterspaced) for „Indhold“. Verified on the
%            image, both witnesses agree. As printed.
%    p. 43   „Tilintetgjørrelse“, doubled r. As printed.
%    p. 68   „Asolutes“ for „Absolutes“, a dropped b. Verified at 600 dpi.
%            Not in the Trykfeil. As printed.
%    p. 78   a risen space printing as a thin rule inside „Immanents“.
%    p. 84   a full word-space standing inside „Gyldigheden“.
%    p. 87   „af en„ Physis i Gud“ “ -- the low quote is set tight to the
%            preceding word and loose from the quoted matter. As printed;
%            it does not disturb the balance.
%    p. 89   „ideal - real“ with a spaced hyphen; set „ideal-real“.
%    None of these disturbs the quote balance, which is still 0 at p. 91.
%    p. 62   The AUTHOR'S OWN erratum, and the only one: the page prints
%            „fra det“ where the Trykfeil leaf asks for „for det“.
%            Transcribed as printed; not applied. [pending -- p. 62 not
%            yet transcribed]
%
%  Rejected claims of a defect, recorded so they are not re-litigated
%  (playbook §2 rule 3: the burden of proof is on the claim):
%    p.  1   „ved hvilket det blev begrundet“ where „hvilke“ is expected --
%            a real word, so as printed.
%    p. 24   „som i deres Udfoldelse“ where „dens“ is expected -- likewise.
%    pp. 31, 35, 36  three ink specks that both OCR engines read as
%            punctuation (an accent over „med“; sub-baseline dots after
%            „Udtryk“ and „Distinction“). All sit off the baseline and the
%            sense needs no mark: transcribed as clean text, not as faults.
%    p.  4   an ink speck above the line after „Handlen,“ read as a quote
%            mark by OCR. Offset dirt; not transcribed.
%
%  ---------------------------------------------------------------
%  5. THE AUTHOR'S OWN ERRATA (printed p. 211, „Trykfeil:“)
%  ---------------------------------------------------------------
%  The errata leaf carries three blocks. Only the first belongs to THIS
%  book: p. 62, line 24, printed „fra det“, to read „for det“. It is
%  transcribed as printed at p. 62 and flagged there; it is not applied.
%  The other two blocks are errata for „Problemet om Tro og Viden“ (nine
%  items) and „Et Svar til Prof. R. Nielsen“, 2det Oplag (two items), and
%  are reproduced at the end of this file as printed. They are worth
%  carrying into `texts/brochner/problemet-tro-viden'.
%  Note the spelling: the heading is „Trykfeil:“, with -ei-.
%
% =====================================================================

\documentclass[12pt,a4paper]{book}
\usepackage[utf8]{inputenc}
\usepackage[T1]{fontenc}
\usepackage{libertinus}
\usepackage{libertinust1math}
\usepackage{textalpha}   % Greek (τέλος etc.), typed directly
\usepackage{graphicx}    % for \reflectbox (the old Danish »ɔ:« id-est mark)
\DeclareUnicodeCharacter{0254}{\reflectbox{c}}  % ɔ (U+0254) = reversed c
\usepackage[danish]{babel}
\usepackage[protrusion=true,expansion=true]{microtype}
\usepackage{setspace}
\usepackage[margin=1.2in]{geometry}
\usepackage{enumitem}
\usepackage{fancyhdr}
\setlength{\headheight}{28pt}
\addtolength{\topmargin}{-13.5pt}
\usepackage{hyperref}

\hypersetup{
  pdftitle={Om det Religiøse i dets Enhed med det Humane},
  pdfauthor={H. Brøchner},
  hidelinks
}

% The book prints *) and **), restarting on every printed page.
\renewcommand{\thefootnote}{%
  \ifcase\value{footnote}\or *)\or **)\or ***)\else \arabic{footnote})\fi}

\pagestyle{fancy}
\fancyhf{}
\fancyhead[LE,RO]{\thepage}
\fancyhead[RE]{\textit{Om det Religiøse i dets Enhed med det Humane}}
\fancyhead[LO]{\textit{\leftmark}}

\setstretch{1.3}

\title{\textbf{Om det Religiøse\\[2pt] i\\[2pt] dets Enhed med det Humane}\\[10pt]
  \large Et positivt Supplement til\\ „Problemet om Tro og Viden“}
\author{Dr.\ H.\ Brøchner\\[4pt]
  \small Professor i Philosophie}
\date{Kjøbenhavn: P.\ G.\ Philipsens Forlag, 1869}

\begin{document}
\maketitle
\thispagestyle{empty}
\newpage

% =====================================================================
%  INDHOLD  (printed pp. [III]--IV = PDF 10--11)
%  Read from the image, not from the OCR layer (playbook §6).
%  Top-level entries are bold Fraktur; the rubric lines are letterspaced;
%  the Greek enumerators α) β) γ) are antiqua italic.
% =====================================================================
\chapter*{Indhold}
\addcontentsline{toc}{chapter}{Indhold}
\markboth{Indhold}{Indhold}

\begingroup
\setstretch{1.15}
\begin{list}{}{%
  \setlength{\leftmargin}{2.2em}\setlength{\itemindent}{-2.2em}%
  \setlength{\itemsep}{2pt}\setlength{\parsep}{0pt}}

\item[] \hfill Side

\item \textbf{Indledning} \dotfill 1--4.

\item \textbf{I. Om Enhedsforholdet mellem den sande Erkjendelse og den
  sande Handlen} \dotfill 5--29.

\item \textbf{II. Om Tilstedeværelsen af det i sand Betydning Religiøses
  væsentlige Kjendemærker i den sande Erkjendelse og den sande Handlen,
  og om disses Enhedsforhold til hiint} \dotfill 30--63.

\item[] \hspace{1.2em} 1. Det Religiøses Væsen \dotfill 30--43.
\item[] \hspace{1.2em} 2. Erkjendelsens Forhold til det Religiøse efter
  dets ideale Bestemthed \dotfill 43--51.
\item[] \hspace{1.2em} 3. Villiens og Handlingens Forhold til det
  Religiøse efter dets ideale Bestemthed \dotfill 51--56.
\item[] \hspace{1.2em} 4. Menneskeaandens Forsoning i det paa den sande
  Erkjendelse hvilende ethisk-humane Liv \dotfill 56--63.

\item \textbf{III. De centrale ethisk-religiøse Problemer, opfattede fra
  Videns Synspunkt} \dotfill 64--201.

\item[] \hspace{1.2em} \emph{Forberedende Begrebsudviklinger:}
\item[] \hspace{1.2em} 1. Det Absolutes Begreb \dotfill 64--89.
% The Indhold reads „De menneskelige Væsens Begreb“; the printed head on
% p. 90 reads „Det menneskelige Væsens Begreb“. Both verified by zoom.
\item[] \hspace{1.2em} 2. De menneskelige Væsens Begreb, bestemt i
  Forhold til det Absolute og til Tilværelsestotaliteten \dotfill 90--97.
\item[] \hspace{1.2em} \emph{Sammenligning af den udviklede Opfattelse med}
\item[] \hspace{1.2em} 1) Ludvig Feuerbachs \dotfill 97--126.
\item[] \hspace{1.2em} 2) den Hegelske Philosophies \dotfill 126--128.
\item[] \hspace{1.2em} \emph{Udvikling af Problemerne:}
\item[] \hspace{1.2em} 1. Om Villiens Frihed \dotfill 128--143.
\item[] \hspace{1.2em} 2. Om det Ondes Væsen og Mulighed \dotfill 144--152.
\item[] \hspace{1.2em} 3. Det Ondes Selvopløsning som Forsoningens
  Mulighed \dotfill 152--159.
\item[] \hspace{1.2em} \emph{Nærmere Udviklinger til Problemet om det Onde:}
\item[] \hspace{1.2em} 1. Forholdet mellem Nødvendigheds- og
  Frihedsmomentet \dotfill 159--164.
\item[] \hspace{2.4em} a) i Virkeliggjørelsen af det Onde \dotfill 159--162.
\item[] \hspace{2.4em} b) i Ophævelsen af det Onde \dotfill 162--164.
% --- p. IV ---
% As printed: entry 2 is given 167--201 while its own sub-entry a) is
% given 165--167, and the printed head stands on p. 164. Not in the
% Trykfeil. Transcribed as printed.
\item[] \hspace{1.2em} 2. Forsoningen, dens Kilde og Maal \dotfill 167--201.
\item[] \hspace{2.4em} a) det i Forsoningen Virksomme \dotfill 165--167.
\item[] \hspace{2.4em} b) det, der skal naaes i Forsoningen \dotfill 167--168.
\item[] \hspace{2.4em} \emph{Betingelser for Forsoningens Fuldstændighed:}
\item[] \hspace{2.4em} α) Skyldens Ophævelse \dotfill 168--176.
\item[] \hspace{2.4em} β) Det Ethiske som Totalitetsbestemmelse for
  Menneskelivet \dotfill 177--182.
\item[] \hspace{2.4em} γ) Det individuelle Menneskelivs til Idealets
  Begreb svarende Evighedsbestemthed \dotfill 182--190.
\item[] \hspace{1.2em} \emph{Sammenligning af den udviklede Opfattelse af
  Forsoningen med}
\item[] \hspace{1.2em} 1) den abstract metaphysiske Opfattelse af
  Forsoningen \dotfill 190--192.
\item[] \hspace{1.2em} 2) den christelige Forsoningslære \dotfill 192--201.

\item \textbf{Slutning} \dotfill 202--210.

\end{list}
\endgroup

% =====================================================================
%  BODY
% =====================================================================

% The body of printed pp. 1--4 carries NO printed head whatever: the text
% begins directly, with a two-line ornamental initial. „Indledning“ is the
% Indhold's word for it, supplied here editorially.
\chapter*{[Indledning]}
\addcontentsline{toc}{chapter}{[Indledning]}
\markboth{Indledning}{Indledning}

% --- p. 1 ---
% p. 1 carries NO printed head of any kind: the text begins directly with a
% two-line ornamental initial N (set here as an ordinary letter, per the
% book's conventions). The signature numeral „1“ at the foot is not
% transcribed. „Indledning“ is the Indhold's word only, and is supplied
% editorially by the skeleton above.
Nærværende Skrift, hvis Udgivelse er bleven forsinket ved Omstændigheder,
uafhængige af Forfatterens Villie, knytter sig paa det Nøieste til den
historisk-kritiske Afhandling: Problemet om Tro og Viden, i hvilken der paa
flere Steder er henviist til det, og gjengiver, ligesom den, i alt Væsentligt
uforandret, Indholdet af en Universitetsforelæsning. Det danner det positive
Supplement til hiin Afhandling, hvis Resultat, efter Sagens Natur, væsentlig
maatte blive negativt, og idet det tager sit Udgangspunkt fra dette Negative,
søger det, forudsættende de Udviklinger, ved hvilket det blev begrundet, den
% „ved hvilket“ as printed (verified on the image at 300 dpi and confirmed by
% both OCR witnesses); one would expect „ved hvilke“. „hvilket“ is a word in
% the language, so it stands as printed and is NOT logged as a defect.
definitive Bestemmelse af Troens Forhold til Viden.

Det Resultat, af negativ Art, til hvilket alle de Udviklinger, der i det
tidligere Skrift ere blevne givne (i de „foreløbige Begrebsbestemmelser“, i
Kritiken af Theologiens Standpunkt og af Forsøget paa at tilveiebringe en
Forsoning ved en principiel Sondring), have ført hen, kan sammenfattes i
følgende Punkter:

% Displayed enumeration, printed with the numeral at the left margin and the
% text hanging: „1.“ and „2.“, both wholly on p. 1.
\begin{enumerate}[label=\arabic*., leftmargin=*, labelsep=1em,
                  topsep=0.35\baselineskip, itemsep=0.35\baselineskip]
\item Mellem Troen i den positive Religions Betydning ɔ: Troen som
  „Aabenbaringstro“, og Viden, bliver der, hvad enten Synspunktet tages i det
  positivt Religiøse eller i Videnskaben, ifølge deres forskjellige Indhold og
  forskjellige Kilde, en uophævelig Incongruents, en uforenelig Modsætning.
% „Incongruents“: the Fraktur I/J sort, normalised to I per the house
% convention (ABBYY reads it „Jncongruents“).
\item Det bliver derfor umuligt, at Troen (i den angivne Betydning) og Viden
  kunne forenes uden Splid i samme Bevidsthed, bevare deres Gyldighed for
  samme Personlighed.
\end{enumerate}

% --- p. 2 ---
Dette negative Resultat omsætter sig nu umiddelbart i et fra det kun formelt
forskjelligt positivt. Ere Modsætningerne uforenelige i Bevidstheden, saa
bliver det en Nødvendighed, at der \emph{vælges} mellem disse Modsætninger, af
% letterspaced: „v æ l g e s“ only; the following „disse“ is set tight.
hvilke en maa være Udtrykket for Menneskeaandens høieste Energie, i sig
indeslutte Maalet for dens høieste Stræben.

Dette Valg, hvis Nødvendighed under hiin Forudsætning er uafviselig, bliver
saa, med Hensyn til sin Charakteer, formelt at bestemme saaledes, at det,
ifølge den Fordring, der ligger i det menneskelige Væsens og den menneskelige
Selvbevidstheds Enhed, maa være et saadant, i hvilket Aandens
Grundvirksomheder ere i Harmonie i sig og med hverandre, i hvilket altsaa
Væsenet, udelt og ulemlæstet, er i Enhed med sig selv i sin Erkjenden og
Villen.

Idet denne Bestemmelse tages til Norm, bliver det tydeligt, hvilken Retning
Valget kan og maa tage, hvad der kan og maa blive Valgets Gjenstand.

Af de i det tidligere Skrift givne Udviklinger fremgaaer det nemlig, at Troen,
som Tro i de positive Religioners Betydning, \emph{ikke} kan vælges, naar den
ovenfor opstillede Fordring skal opfyldes. At den ikke kan det, ligger deels i
dens Forhold til det Praktiske, deels i dens Forhold til det Theoretiske. Den
% Both occurrences of „Theoretiske“ on this page are set TIGHT, not
% letterspaced --- checked at 300 dpi against „Praktiske“ on the line above,
% and spacing.py flags neither. The BATCH-AGENT calibration set lists
% p. 2 „T h e o r e t i s k e“; that entry does not check out on this page.
kan ikke vælges fordi den, i praktisk Henseende, ikke kan blive det formende
Princip for et ethisk Menneskelivs Heelhed, ikke kan forme Menneskenaturen.
Paa den positive Religions Standpunkt bliver nemlig den ethiske Fordrings
Kilde det fra Menneskevæsenet sondrede, supranaturale, transscendente
Aabenbaringssubject, hvis Villie er den absolute Vilkaarlighed, og denne
Villies Fordringer kunne altsaa ikke congruere med Menneskeaandens
Væsensfordring. Og naar den positive Religion er klaret til Aandsreligionens
høieste Form, saa kræver dens „ethiske“ Fordring Fornægtelsen af det
Naturlige, altsaa af et Fundamentalt i Menneskevæsenet, af et saadant, der
\emph{maa} gjøre sig gjældende i sin Væsentlighed. --- Sees der hen til det
Theoretiske, saa kan Troen ikke vælges fordi den ikke kan afslutte sig til en
Totalitet, fordi den, medens den ikke kan gjøre sig uaf\-%
% --- p. 3 ---
\setcounter{footnote}{0}% (no space: the word „uafhængig“ is broken across pp. 2/3)
hængig af Viden, ikke i sig kan optage det i Viden Væsentlige, ikke kan
gjennemskue og forklare Viden\footnote{Jfr. Problemet om Tro og Viden, S.\ 195;
210 f.}.
% printed mark: *) --- the only note on p. 3, set below a short rule.
% „Jfr.“: the Fraktur I/J sort; transcribed J here, as the abbreviation for
% „jævnfør“, following the sibling volume *Problemet om Tro og Viden*. The
% I-normalisation applies where I is the letter meant, not to this word.
% The signature „1*“ at the foot of p. 3 is not transcribed.
Viden maa da, som Væsensbestemmelse i den menneskelige Natur, uundgaaelig
hævde sig som berettiget ligeoverfor Troen og bringe Splid ind i den
Bevidsthed, der skulde udfyldes af denne. Et Videns Subordinationsforhold er
udelukket ved Modsætningens Art og ved Videns Væsen. Naar Troen vælges, kan
derfor Væsenets Harmonie og dets indre Forsoning ikke naaes.

Men naar det fordrede Maal ikke kan naaes gjennem Valget af Troen, saa viser
der sig, gjennem hvad der i det forrige Skrift blev udviklet, en Mulighed til,
at det kan naaes gjennem et Valg af Viden. Idet nemlig Viden vælges, saa
vælges det, der, ved at være Udtryk fra Menneskevæsenets Universalitet, har
Evnen til at udfolde sig som en i sig hvilende Totalitet med indre
Uendelighed. Idet Viden vælges, vælges det, i hvilket Væsenet i sin Enhed
bliver aabenbart for sig selv, det, der i den udelelige menneskelige Aand har
det ideelle Principat, og som derfor kan blive Normen for en sand ethisk
\emph{Livsheelhed}. Idet Viden vælges, vælges endelig det, der, som
indesluttende Erkjendelsen af Aandens \emph{hele} Liv og Livsudvikling, kan
gjennemskue de endelige Religionsformers Væsen og begribe dem i deres Udspring
af Menneskenaturen i dens Udfoldelse, og derfor ikke beholder en forstyrrende
Bevidsthedsskranke i dem, og som derfor ogsaa, idet det i de historiske
Religionsformer udskiller det Forgængelige og Endelige fra det Blivende og
Væsentlige, kan staae i Harmonie med og i sig optage dette Væsentlige, der
netop er det Religiøse efter dets sande og evige Betydning.

Hvorledes nu, fra dette Synspunkt af Nødvendigheden af et Valg og af et
saaledes bestemt Valg, en \emph{forsonende} Løsning af Problemet om Troens og
Videns Forhold bliver mulig; hvorledes Videns Ret og Gyldighed kan hævdes uden
at Troen efter dens Sandhed krænkes, uden at det
% --- p. 4 ---
af en Menneskeaandens Trang fordrede Religiøse fornægtes, bliver det de
efterfølgende Udviklingers Opgave at tydeliggjøre. Gangen i disse er bestemt
ved Opgavens Natur. \emph{Først} bliver det at godtgjøre, at der er et ved
Væsensenheden betinget Enhedsforhold tilstede mellem den sande Erkjendelse og
den sande Handlen, den ethiske Handlen. \emph{Derefter} skal det paavises, at
det i sand og evig Betydning Religiøses væsentlige Kjendemærker: Troen og
Hengivelsen, ere uadskillelige fra den sande Erkjendelse og den sande Handlen,
% A small tilted ink speck floats in the interline space just after
% „Handlen,“ here, above the ascender line and belonging to neither line's
% baseline. Both OCR witnesses read it as a quotation mark. Judged an ink
% speck / offset, not a printed sort, and so not transcribed.
at altsaa disse ere i Harmonie med det Religiøse og bære det i sig, saaat
Menneskelivet efter sin harmoniske Enhed med sig i sine Grundfunctioners
Totalitet kan vise sig som gjennemtrængt af det Religiøse. Og
\emph{endelig} skal Muligheden paavises af, at der ud fra Philosophiens
% „Og“ is set tight; only „e n d e l i g“ (broken „en=/delig“ across the
% line) is letterspaced.
Begreber af det concret bestemte Absolute, den høieste Vidensbestemmelse, og
af Menneskevæsenet i dets Forhold til Principet og den reale Tilværelse, kan
udvikles en Livsanskuelse, der kan tilfredsstille ethisk-religiøst, og naaes
til de væsentlige Bestemmelser af det ethiske og religiøse Forhold, om hvilke
det er blevet paastaaet, at de kun kunne udledes af et med Vidensprincipet
absolut uensartet Villies- og Troesprincip. Her ville altsaa de centrale
ethisk-religiøse Problemer: Problemerne om Villiens Frihed, om det ethisk
Ondes Mulighed og Virkelighed, og om Forsoningen og Restitutionen, komme under
Forhandling.

I Sammenhæng med disse Udviklinger vil tillige den antydede Livsanskuelses
Forhold til de Opfattelser, der frembyde virkelige eller tilsyneladende
Berøringspunkter med den, blive tydeliggjort, og Kritiken af dem given, deels
directe, deels indirecte.

\begin{center}\rule{2cm}{0.4pt}\end{center}
% RULE, p. 4: the Indledning closes with a short centred DOUBLE rule --- two
% parallel hairlines --- standing below the last line of type, the rest of the
% leaf blank. Verified on the image at 300 dpi at 2x; the same weight as the
% division-closing rules of pp. 29, 63, 201 and 210, and unlike every other rule
% in the book. Set with the single \rule this edition uses throughout.

\chapter*{I.\\ Om Enhedsforholdet mellem den sande Erkjendelse og den sande Handlen.}
\addcontentsline{toc}{chapter}{I. Om Enhedsforholdet mellem den sande Erkjendelse og den sande Handlen}
\markboth{I. Om Enhedsforholdet}{I. Om Enhedsforholdet}

% --- p. 5 ---
% p. 5 is the opening page of division I. The \chapter* head is in the skeleton
% and is not repeated here; the short centred rule below it is head ornament and
% is not transcribed at all (see the header, SHORT CENTRED RULES); the
% fragment begins at the first line of body text. p. 5 carries no folio (by
% design) and opens with a two-line ornamental initial G, set as an ordinary
% letter per the book's conventions.
%
% NOTE on letterspacing across a line break, observed six times in this batch
% (pp. 6, 7 fn., 13, 15, 17 twice): where a letterspaced word or phrase runs
% over a line ending, the compositor sometimes spaced only the part standing
% on the first line and set the continuation tight --- „A a n d s = / functioner“
% (p. 6), „B e v i d s t h e d s = / former“ (p. 6 fn.), „b e g g e  d e /
% Bestemmelser“ (p. 13), „s a m m e  R e = / sultat“ (p. 15), „V æ s e n s = /
% bestemmelse“ and „V i d e n s = / bestemmelse“ (p. 17). Each is transcribed
% exactly as spaced, so \emph{} sometimes ends inside a word. The opposite
% (spacing carried through the break) also occurs --- „l e d = / s a g e d e“,
% „E r = / k j e n d e n“ (p. 6), „V æ = / s e n“ (p. 8), „u i n t e r = /
% e s s e r e t“ (pp. 13/14) --- so the practice is genuinely inconsistent and
% is not normalised.
Grundforholdet mellem Erkjendelse og Villie som primitive Yttringsformer af
det menneskelige Væsen er blevet angivet i det tidligere Skrift (S.\ 134 ff.),
og til de der givne Bestemmelser kan jeg paa dette Sted vise tilbage. Den
Opfattelse, som der blev gjort gjældende, var den, at det menneskelige Væsen,
hvis udelelige Enhed er sat ved Subjectivitetens Begreb og afpræger sig i
Selvbevidsthedens Enhed, efter sit Grundbegreb er en Selvhedsenergie, og at
denne Selvhedsenergiens Bestemmelse er det gjennemgaaende Almene, der giver
sig sine særegne Virksomhedsformer i Erkjendelse og Villie. I disse særegne
Udtryksformer aabenbarer altsaa Væsenets almene Bestemthed sig, de blive
primitive Functioner af Selvhedsenergien, og medens deres Oprindelighed
indeslutter deres Eiendommelighed og Selvgyldighed, indeslutter deres Forhold
til Selvhedsenergiens almene, det Særegne iboende, Bestemmelse, at de ere
beherskede af Væsenets Enhed, at de reflectere hinandens Formbestemtheder, og
at derfor Væsensenergien bevarer sin Enhed i den dobbelte Fremtrædelsesform.
Ved det \emph{ligelige} Forhold, i hvilket Erkjendelse og Villie ere viste at
staae til Energiens Bestemmelse, er den Begrebsforvexling fjernet, der,
gjennem en vilkaarlig og uklar Udvidelse af Begrebet \emph{Drift},
identificerer denne med Energien og derved skaffer Villien en illusorisk
Forrang i Menneskevæsenet, uden saa at kunne forklare Erkjendelsens,
overhovedet Bevidsthedens,
% --- p. 6 ---
\setcounter{footnote}{0}% p. 6 carries TWO notes, printed *) and **)
Oprindelse, uden at kunne hæve den Forvirring, at „Driften“ selv skal sætte,
hvad Driften, i Ordets \emph{egentlige} Betydning, stræber at fastholde
\emph{eller} at fjerne, og uden at kunne gjøre Rede for den Kjendsgjerning, at
den første Yttring af Selvhedsenergien, som ubevidst organisk Formen,
vedbliver \emph{efterat} Villien er virkeliggjort, og vedbliver
\emph{uafhængig} af Villien. Gjennem Kritiken af denne ensidige og uklare
Opfattelse er saavel Erkjendelsens som Villiens Primitivitet hævdet, og idet
Spørgsmaalet om det ideelle Principat opkastedes, er dette blevet tillagt
Erkjendelsen ifølge dens Forhold til den Grundbestemmelse i det menneskelige
Væsen: \emph{Væsensuniversaliteten}, ved hvilken det kategorisk sondrer sig
fra Livets andre Udviklingsstadier\footnote{Jfr. Problemet om Tro og Viden,
S.\ 137; Et Svar til Prof.\ Nielsen, S.\ 12 f.}.
% printed mark: *) --- first of the two notes on p. 6, below a short rule.
% „Jfr.“ takes the J of the single Fraktur I/J sort, as the abbreviation for
% „jævnfør“ (contrast the lower-case „jfr.“ printed on p. 9).
Dette Erkjendelsens ideelle Principat aabenbarer sig ogsaa i
\emph{Bevidsthedens} Forhold til de særegne Aandsfunctioner. Som Functioner af
den menneskelige Selvhedsenergie ere de \emph{ledsagede af Bevidsthed}, og ere
netop kun derved \emph{Aands}functioner, men Bevidstheden staaer efter sit
% letterspacing: „A a n d s =“ at the line end is spaced, „functioner,“ on the
% next line is set tight. Verified at 300 dpi and confirmed by both witnesses.
Begreb i et andet Forhold til Erkjendelse og Tanke end f.\ Ex. til Villien;
den bliver det \emph{abstracte, formelle,} Udtryk for Erkjendelsen, det, der
kan udfoldes til en \emph{tilegnende Erkjenden}, til \emph{begribende}
Tænkning, til Aandens Gjennemsigtighed for sig i sit Indhold; og det, at den
træder i Forhold til de forskjellige Aandsfunctioner, idet den bliver Formen
for deres Væren for Selvet, viser netop hen til Erkjendelsens ideelle
Principat, og til Muligheden af en Selvets Gjennemsigtighed for sig gjennem
Erkjendelsen i sine forskjellige eiendommelige
Yttringsformer\footnote{Med Hensyn til den Beskyldning, en Kritiker har gjort
mod mig for Sammenblanding af det væsentlig Forskjellige, bemærker jeg kun,
hvad enhver eftertænksom Læser af det tidligere Skrift vel vil have seet, at
jeg ikke identificerer de \emph{umiddelbare Bevidstheds}former med den
udfoldede \emph{videnskabelige Erkjendelse}, men kun hævder, at det, der
træder frem i hine Umiddelbarhedens Former, deels er et saadant, der kan
udfoldes i klar Erkjendelse, deels er et \emph{resulterende} Umiddelbart, og
nærmere et Saadant, der er resulteret af en forudgaaende \emph{Erkjendelse}.
Ligesaalidt sammenblander jeg det \emph{Subjective} med det \emph{Objective},
men jeg hævder, at hint, som Udtryk for Subjectets \emph{Væsen}, har den samme
Gjennemskuelighed for Erkjendelsen ɔ: for en \emph{theoretisk} Erkjendelse, som
hint Umiddelbare.}.
% printed mark: **) --- the second note on p. 6. Its text RUNS OVER onto p. 7,
% where it fills the last six lines below the rule; it is set here in full at
% its call, and the page marker below stands at the beginning of p. 7's body
% text, as everywhere else.
% --- p. 7 ---
Naar dette Fundamentale: Væsenets udelte Enhed i Erkjendelse og Villie, og
Erkjendelsens ideelle Principat, allerede er vundet som Resultat i det
tidligere Skrift, saa bliver paa dette Sted den \emph{relative} Modsætning
nærmere at belyse, der maa fremtræde ifølge Grundfunctionernes Særegenhed og
ifølge Aandsudviklingens Betingelser; men gjennem hvilken, ifølge
Grundforholdet, Harmonien, den væsensbegrundede Enhed, maa virkeliggjøre sig
som beherskende Modsætningen.

\emph{Den relative Modsætning}, i hvilken Erkjendelse og Villie fremtræde i
deres Selvstændighed, bliver først at opfatte som en \emph{Begrebs}modsætning,
og derefter som en \emph{Facticitets}modsætning.
% „B e g r e b s modsætning“ and „F a c t i c i t e t s modsætning“: in both
% compounds only the first element is letterspaced, the second set tight, and
% both stand well inside their lines. Two lines below, „B e g r e b s m o d -
% s æ t n i n g e n“ is spaced throughout.

\emph{Begrebsmodsætningen} mellem de to primitive Functioner er allerede
angivet i „Problemet om Tro og Viden“ (S.\ 135 f.). Den viste sig som
umiddelbart fremtrædende i Formbestemthedens Differents; men Formdifferentsen
saaes, ikke at constituere en abstract Modsætning, idet de Bestemmelser, der
udtrykte den, viste hen paa hinanden og udviklede sig af hinanden, saaledes at
Grundfunctionerne reflecterede hinandes Formbestemtheder.
% PRINTER'S DEFECT: „hinandes“ as printed, for „hinandens“. A dropped n, not a
% scan artefact: the word is fully legible at 300 dpi, both OCR witnesses read
% „hinandes“, and „hinandens“ is set correctly on p. 5 three lines from the
% foot. Transcribed as printed, not corrected.
Derved blev det saa muligt, at under Formdifferentsen Indholdets
\emph{væsentlige} Enhed kan bevares: en Enhed, der fordres derved, at det
samme udelte Væsen giver begge Former deres Indhold. Naar saaledes i den
ethiske Beslutning Erkjendelsen af Pligten omsættes i en Villen af Pligten,
naar altsaa det, der som Erkjendelsesgjenstand har Mulighedens Form, sættes i
Virkelighedens Form i Jeget, saa er dette \emph{bestemte} Indhold i begge
Former \emph{væsentlig} uforandret, og maa være det allerede af den Grund, at
ellers Pligtopfyldelsen i Handlingen \emph{ikke} vilde være en Opfyldelse af
dette bestemte, som Pligt Erkjendte; og i den forudgaaende
% --- p. 8 ---
Erkjendelse er Mulighedens Almeenbestemmelse saaledes punktualiseret i Forhold
til den Existerende og Erkjendelsens Gjenstand saaledes bragt i Forhold til
det erkjendende \emph{Væsen}, saaledes seet som en Væsensbestemmelse, der
fordres sat som Realitetsbestemmelse, at Omsætningen i Villiesindhold ikke kan
ophæve Indholdets væsentlige Enhed.

At Differentsen mellem Erkjendelse og Villie ikke forvandler sig til en
\emph{abstract} Modsætning viser sig ogsaa indirecte naar vi belyse Forsøget
paa at hævde en saadan Modsætning gjennem Modsætningen mellem den
„theoretiske“ og den „praktiske Bevidsthed“, mellem den „theoretiske“ og den
„praktiske Erkjendelse“ (Tænkning).

Med Hensyn til denne Theorie berøre vi kun, hvad vi andensteds have paavist,
den Uklarhed i Anvendelsen af Kategorierne, i hvilken den gjør sig skyldig,
idet Begreberne: Bevidsthed, Erkjendelse, Tænkning, bruges iflæng uden klar
Begrændsning, og Erkjendelsen, der forholder sig til Handling, blandes sammen
med selve Villiesacten, Villiesbeslutningen. Vi undersøge, idet vi see bort
fra denne forvirrede Anvendelse af Begreberne, og idet vi sondre hvad der er
blevet kaldt praktisk Erkjendelse fra den umiddelbare „praktiske Bevidsthed“,
der deels er en Bevidsthed om Subjectets Tilstand i praktisk Henseende, om
Subjectets Villiesbeslutning og om Motivet, det factisk Bevægende, deels en
Bevidsthed om det ethisk Aprioriske, om den praktiske Fordring, --- hvorledes
den søgte Modsætning mellem en theoretisk og en praktisk Erkjendelse bliver at
bestemme.

Det vil da her strax være indlysende, at der ikke kan gives to væsenssondrede
Arter af Erkjendelse, to absolut uensartede Erkjendelsesformer, da isaafald
Erkjendelsens Bestemmelse vilde forvandle sig til et rent indholdsløst Navn og
i sin Dobbelthed ikke kunne være det samme erkjendende Væsens Eiendom. En
Modsætning, der udspringer af det samme Væsens Grund, maa i sig bevare den
væsentlige Enhed. Og derfor bliver det at fastholde, at Erkjendelsen,
\emph{som} Erkjendelse, i enhver Form, og i enhver Anvendelse maa bevare sin
Væsensbestemthed.
% --- p. 9 ---
Modsætningen mellem en „theoretisk“ og en „praktisk“ Erkjendelse bliver da en
blot relativ Modsætning mellem en Erkjendelse, der er bestemt ene ved Tankens
Love og som har theoretisk Indsigt til Resultat, og en Erkjendelse, der er
bestemt ogsaa ved Villiesmotiver (Villieslove), og som har Handling,
Beslutning, til Resultat.

Betegnelsen af den praktiske Erkjendelse som bestemt ved Villiesmotiver viser
deels hen til denne Erkjendelses Forudsætninger og Udgangspunkter
(Præmisser), deels til dens Paavirkning ved Villien under dens Fremadskriden.
Dens Udgangspunkter ere praktiske Incitamenter: deels Drifttilskyndelser,
deels Virkninger af et ethisk Apriorisk (Samvittighedsgrunde). Fra disse
Udgangspunkter, der staae som et Givet, skrider saa den „praktiske“
Erkjendelse frem paa samme Maade som den theoretiske, der drager Consequentser
af sine Præmisser, deels af uopløste, deels af gjennemsigtige. ---
Erkjendelsens Paavirkning under dens Fremadskriden kan udgaae fra en klarere
eller dunklere bevidst Villie, der enten supplerer ufuldstændige Præmisser,
eller corrigerer ensidige Theorier, eller hæmmer den sande Erkjendelse. I
første Tilfælde giver Villien et Udgangspunkt, der, forat benyttes, ikke maa
staae i Strid med det theoretisk Erkjendte. Forholdet danner her en Analogie
til det, hvor i theoretisk Raisonnement en ubeviist gængs Mening (Fordom)
benyttes. I andet Tilfælde skyder Villien et Correctiv ind, som Erkjendelsen
vedkjender sig, og som danner Udgangspunktet for en ny Slutning. Ogsaa her
frembyder det theoretiske Gebet fuldkommen tilsvarende Analogier (jfr.
Problemet om Tro og Viden S.\ 144 f.). --- I sidste Tilfælde er Forholdet mere
compliceret. Under Erkjendelsens Fremadskriden indvirker Villien nemlig deels
saaledes, at den drager Instantser frem, der kunne forhindre en Conclusion,
deels saaledes, at den spatierer Udfoldelsen af Præmisserne fra Conclusionen,
Beslutningen, hvorved Erkjendelsens Klarhed kan fordunkles, og modsigende
Instantser, Skingrunde og Paaskud, kunne faae Præg af Gyl\-%
% --- p. 10 ---
dighed. Til denne Hæmmen dannes paa det theoretiske Gebet Analogien af
Fordommes Indvirkning paa Erkjendelsen.

Den „praktiske Erkjendelses“ \emph{Resultat} er deels Antagelsen af en Lære,
der forholder sig til Subjectets Handling, til dets praktiske Liv, deels en
Slutning i Retning af Handling, der omsættes i Beslutning. Er Resultatet
Antagelsen af en Lære, saa kommer denne Antagelse frem som en Slutning, der
fremgaaer med Nødvendighed af sine Præmisser, men blandt disse Præmisser
kunne, efter det ovenfor Udviklede, saadanne findes, der ere skudte ind af
Villien. Er Resultatet Beslutning som Beslutning til den bestemte Handling,
saa er den forudgaaende Erkjendelsesact, til hvilken Beslutningen i egentlig
Forstand altid viser tilbage, ligesom den Erkjendelse, af hvilken Antagelsen
af den praktiske Lære afhang, og ligesom den Erkjendelse, der uddrog
Consequentser af praktiske Præmisser, i sin Form ikke væsentlig forskjellig
fra den theoretiske; men dette vilde være umuligt hvis en absolut
Uensartethed adskilte „praktisk“ og „theoretisk“ Erkjendelse fra hinanden. Thi
Erkjendelsen er kun hvad den er ved sin Form; men ved Form forstaaes da
naturligviis noget Andet og Mere end den paa et vilkaarligt Indhold
\emph{udvortes} anvendte tomme syllogistiske Form.

Resultatet med Hensyn til den Erkjendelse, der forholder sig til det
Praktiske, bliver altsaa det, at den, ifølge sin Form og ifølge sit Forhold
til Erkjendelsens almene Begreb, ligesom ogsaa ifølge det Praktiskes
(Udgangspunkternes og Resultatets) Forhold til Menneske\emph{væsenet}, maa
% letterspacing: „Menneske“ is set tight, only „v æ s e n e t ,“ is spaced.
% Verified at 300 dpi (zoom on the line).
blive væsenseet med den Erkjendelse, der betegnes som theoretisk. Al
Erkjendelse er \emph{som} Erkjendelse theoretisk. At Erkjendelsen kan anvendes
i praktisk Retning og i denne Anvendelse paavirkes af Villien, beviser Intet
for en væsenssondret praktisk Erkjendelse: thi en Villiens Enhed med
Erkjendelsen træder \emph{overalt} frem, dens Vexelvirkning med Erkjendelsen
bliver ogsaa aabenbar paa den uomtvistet theoretiske Erkjendelses Gebet, og
den theoretiske Erkjendelse kan \emph{som saadan} blive Gjenstand for Villien.
% --- p. 11 ---
At Forholdet er et saadant, vil blive endnu mere indlysende naar de
\emph{Charakteermærker} undersøges, ved hvilke den paastaaede abstrakte
Modsætning mellem theoretisk og praktisk Erkjendelse \emph{beskrives}: de
% „abstrakte“ with k as printed here; the book's usual form is „abstract“
% (see pp. 8, 13, 15). Both witnesses agree, and the k is plain at 300 dpi.
ville alle vise sig, fra den relative Differents at føre tilbage til den
væsentlige Enhed.

Naar først Modsætningen charakteriseres saaledes, at den \emph{theoretiske}
Bevidsthed (Erkjendelse) er \emph{upersonlig}, den \emph{praktiske
personlig}, saa viser denne Modsætning, som abstract og exclusiv Modsætning,
sig allerede af den Grund som uholdbar, at den theoretiske Erkjendelse, som
Udtryk for det Væsens Enhed, hvis Grundbestemmelse er Personlighed, ikke blot
overhovedet maa være en Personlighedsyttring, men til sit endelige Maal maa
have en Erkjendelse af Personligheden i dens Sammenhæng med Tilværelseshelheden
og i dens Underordning under Tilværelsesprincipet. I denne Erkjendelse af det
personlige Selv vil da ogsaa Erkjendelsen af den objective reale Tilværelse,
af den Legemlighedens Verden, til hvilken Personligheden staaer i Forhold
gjennem sin Realitetsside og sine Realitetsbetingelser, vise sig som havende
„personlig“ Betydning. Hiin Bestemmelse af den theoretiske Erkjendelse som
upersonlig indeslutter den Consequents, at den \emph{theoretiske Erkjendelse,
skjøndt primitiv Yttring af Selvet, ikke kan blive Selverkjendelse}, idet den
faaer sin Grændse i det Selv, der er selve denne Erkjendelses Kilde. Vilde
man, forat hævde Paastanden om den theoretiske Erkjendelses upersonlige
Charakteer, lægge Eftertrykket paa, at ogsaa Erkjendelsen af Selvet, som
theoretisk Erkjendelse, maa blive i det Almenes Form, og at derfor det
personlige Selv gjennem den dog kun kan erkjendes paa upersonlig Maade, saa
vilde derimod være at gjøre gjældende, at det Almene i denne Erkjendelse ikke
fastholdes i den oprindelige Abstractionens Form, men bestandig mere og mere
punktualiseres til Concretion, saa at det blivende „Upersonlige“ i den
theoretiske Erkjendelse kun vilde betegne det i det Almenes Punktualisation
blivende Relative, --- og dette Relative er tilmed et
% --- p. 12 ---
saadant, der ikke blot er blivende i den theoretiske Erkjendelse, men ogsaa er
og maa være det i Villiesbeslutningerne, netop fordi disse ere bevidste og
Aandsyttringer, og som saadanne ere hævede over den Enkelthedens Form, der er
Naturtilværelsens og den naturlige Driftbestemtheds.

Den upersonlige, theoretiske Erkjendelses Gjenstand bestemmes som det
\emph{objectivt Gyldige}, den personlige praktiske Erkjendelses som det
\emph{subjectivt Gyldige}, og der gjøres en abstract Modsætning mellem det
objectivt Almeengyldige og det subjectivt Almeengyldige. Naar imidlertid hiin
modsigende Grændse for den theoretiske Erkjendelse fjernes, saa viser den sig
ogsaa at omfatte det, der, som et „subjectivt Gyldigt“ er sat i Modsætning til
det objectivt Gyldige, og at omfatte ikke blot hvad der, ifølge sin
Sammenhæng med den erkjendte Menneskenatur, har Gyldighed i sin Almindelighed
for alle Subjecter, men ogsaa hvad der, som i sin Concretion bestemt ved
samvirkende almene Grunde, har Gyldighed for dette bestemte Subject.

Kun ifølge dette Forhold kan der blive Tale om et „subjectivt
Almeengyldigt“, medens den abstracte Modsætning mellem det Subjective og det
Objective maa føre til, at lade den \emph{individuelle} Subjectivitet som
saadan bestemme det Gyldige, som den \emph{vil} bestemme det, og saaledes at
lade de subjective Grunde blive til Eet med individuelt tilfældige Grunde,
Drifttilskyndelser o.\ desl., at lade „Villiesgrunde“ blive saadanne, som en
vilkaarlig Villie gjør til Grunde.

Bestemmes saa Modsætningen mellem de to Erkjendelsesarter saaledes, at den
theoretiske Erkjendelse er bevæget af \emph{Nødvendighed}, medens den
praktiske „\emph{grunder i Frihed}“, idet hiin, som fortabt i det Almene,
tvingende bestemmes af de uimodsigelige Vidensgrunde, denne, som betinget ved
Selvhensyn, hviler paa Samvittighedsgrunde, saa viser ogsaa denne Modsætning
sig som dialektisk. Den theoretiske Erkjendelse er behersket af Nødvendighed
hvor Slutningens Forudsætninger ere fuldt gjennemsigtige, hvor Beviset faaer
Stringents ved klare og fuldstændige Præ\-%
% --- p. 13 ---
misser. Men ligesaa fuldt er den „praktiske Erkjendelse“ bestemt med
Nødvendighed hvor disse Betingelser ere tilstede, hvor Forudsætningerne ere
givne klare og fuldstændige. Hvad der tilveiebringer Skinnet af en abstract
Modsætning, er \emph{Forskjellen i Concretion} mellem den praktiske
Erkjendelses Gjenstand og visse Sphærer af det Theoretiske. Paa det concrete
Ethiskes Gebet kan der vedblive at være en tilsyneladende uophævelig
Differents netop paa Grund af Præmissernes Ufuldstændighed, paa samme Maade
som indenfor det Theoretiske, paa de concretere Gebeter. Differentsen kan
holde sig, fordi de mangfoldige Forudsætninger ikke ere klarede og derfor
partielt suppleres paa individuel Maade. Saaledes bliver, under samme
Forudsætninger, Nødvendigheden den samme. Og den Nødvendighed, der bestemmer
den theoretiske Erkjendelse, er, naar Erkjendelsen er \emph{begribende} og som
saadan \emph{tilegnende}, et Udtryk for \emph{selve det erkjendende Væsens
Lovmæssighed}, og omsætter sig saaledes i \emph{Frihed}. I den ufrie
Nødvendigheds Form fremtræder den theoretiske Erkjendelse kun hvor det ved
Sandsningen Givne, i dets Fremmedhed, er bestemmende. Men til dette Stadium
svarer det Stadium i den praktiske „Erkjendelses“ Udvikling, hvor det
Selvhensyn, der bestemmer den, er Hensynet til et uklart Driftagtigt, ved
hvilket Villien bringes ind under Naturnødvendigheden. Kun som Udtryk for et
Væsentligt har Subjectets Trang, der bestemmer hiin „Erkjendelse“,
Berettigelse, men som Udtryk for Væsenet har denne Trang paa een Gang den
sande Nødvendigheds og Frihedens Præg, medens Driftbestemthedens saakaldte
„Frihed“ er ufri Nødvendighed ligesom Sandsningens Nødvendighed. Baade i den
theoretiske og i den praktiske Erkjendelse fremtræde altsaa \emph{begge de}
% letterspacing: „b e g g e  d e“ ends the line and is spaced;
% „Bestemmelser,“ opening the next line is set tight. Verified by zoom.
Bestemmelser, der skulde constituere den skarpe Modsætning; det virkelige
Forhold bliver kun uklart derved, at Individet, der handler „vilkaarligt“,
ikke bliver sig bevidst, netop i denne Handling at være ufrit.

Naar endelig Modsætningen bestemmes paa den Maade, at den theoretiske
Erkjendelse skal være \emph{uinter}\-%
% --- p. 14 ---
\setcounter{footnote}{0}% p. 14 carries one note, printed *)
\emph{esseret, contemplativ, ansvarsløs}, medens den praktiske er
% the letterspacing of „u i n t e r = / e s s e r e t“ IS carried across the
% page break here; both halves are spaced on the image.
\emph{interesseret}, bevæget af \emph{Frygt og Haab}, Gjenstand for
\emph{Ansvar}, saa forholder det sig med denne Modsætning paa samme Maade som
med de forrige. Den theoretiske Bevidsthed er \emph{uinteresseret}, naar
„interesseret“ tages i Egoismens Betydning; den er derimod i høieste Grad
interesseret, naar Interesse tages i Betydning af en Stræben hen mod et
væsentligt Formaal, og dette Formaal for den theoretiske Erkjendelse er netop
Selvets Klarhed over sig selv, dets Væren for sig selv som selvbevidst efter
sin totale Concretion. At Aanden theoretisk finder Hvile i et hvilketsomhelst
Resultat, til hvilket Consequentserne af Forudsætningerne føre, har kun
forsaavidt sin Rigtighed, som den momentant kan standse i dette Resultat; men
i det usande eller kun halvt sande Resultat er der bestandig en Stimulus, der
atter vækker Erkjendelsens Interesse for en ny Undersøgelse, og atter driver
den fremad, netop fordi det usande Resultat ikke kan komme i Samklang med
eller omfatte den Erkjendelses-Totalitet, som Erkjendelsen efter sit Væsen maa
stræbe at udforme\footnote{Den theoretiske Erkjendelses Interesse for sit
Resultat træder iøinefaldende frem i dens Suppleren af Inductionsslutningens
ufuldstændige Præmisser.}.
% printed mark: *) --- the only note on p. 14, below a short rule.

Den theoretiske Erkjendelse skal være \emph{contemplativ}, gaaende op i
Begrundelsen af det Objective; men den er det kun i den ved de tidligere
Bestemmelser begrændsede Betydning. Fra den beskuende, selvforglemmende
Fordybelse i det Almene og Objective søger den tilbage til det Punktuelle og
til den Selverkjendelse, til hvilken en uimodstaaelig indre Trang driver den,
og idet den stræber hen mod sit Maal, er der i Erkjendelsens \emph{Kamp} for
Tilegnelsen af Virkeligheden, i den \emph{sande} Erkjendelses Kamp mod det
Usande, et Tilsvarende til Villiens Conflict med Existentsen, til dens Kamp
for sin Realisation, til den gode Villies Kamp mod det Onde.

Den theoretiske Erkjendelse skal endelig være \emph{uden}
% --- p. 15 ---
\emph{personligt Ansvar}, den praktiske ledsaget af Bevidstheden om et
saadant. Ogsaa denne Modsætning er uholdbar i den Form, i hvilken den
opstilles. At der med Hensyn til Bevidstheden om Ansvar bliver en Forskjel
mellem \emph{Villiesacten} og Tankevirksomheden, naar der kun sees hen til, at
hiin griber bestemmende ind i Realiteten, medens denne ikke gjør det, ialfald
ikke paa bevidst Maade, forstaaer sig af sig selv. Men derfor falder
Erkjendelsesvirksomheden ikke udenfor Ansvaret. Erkjendelsesacten, der er
underkastet \emph{Sandhedsideens Norm}, har i denne et „Bør“, i Forhold til
hvilket der er en Bevidsthed om Ansvar. Naar Ansvaret overhovedet viser hen
til, at det, for hvilket Individet er ansvarligt, er Individets egen Gjerning,
at det fremgaaer af dets Væsen, og at det staaer i Forhold til en høiere, mere
omfattende Lov, saa bliver Ansvaret fælleds for det, der fremgaaer af
Erkjendelsens og Villiens Virksomhed, og det maa saameget mere blive det, naar
Erkjendelsens normale Virksomhed, som det i det Efterfølgende vil vise sig, er
Betingelsen for Villiens normale Virksomhed. Erkjendelsens Sammenhæng med
Ansvaret bliver altsaa ikke blot saaledes at opfatte, at Villiens Forhold til
Erkjendelsen kan bringe denne ind under Ansvaret, men saaledes, at
Erkjendelsen \emph{i og for sig}, ved sit Forhold til Individets Væsen og til
Sandhedens høiere Lov, gaaer ind derunder.

Ogsaa gjennem Betragtningen af de Charakteermærker, ved hvilke en abstract
Modsætning mellem „theoretisk“ og „praktisk Erkjendelse“ (Bevidsthed) og
derigjennem af Erkjendelse og Villie, skulde hævdes, føres vi saaledes til det
\emph{samme Re}sultat som ovenfor: at den \emph{abstracte Modsætning} viger
% letterspacing: „s a m m e  R e =“ closes the line spaced, „sultat“ opens the
% next line tight.
Pladsen for en \emph{væsentlig Enhed}, der \emph{behersker} den
\emph{relative} Differents, og \emph{maa} gjøre det, paa Grund af de primitive
Væsensyttringers Forhold til Væsensenheden. Dette Resultat vil nu faae en ny
Bestyrkelse, naar vi betragte \emph{Facticitetsmodsætningen} mellem
Erkjendelse og Villie, den Modsætning, i hvilken de, under Aandslivets
individuelle Udvikling, kunne træde til hinanden.

En \emph{factisk} Modsætning fremtræder paa Erkjendelsens
% --- p. 16 ---
\setcounter{footnote}{0}% p. 16 carries TWO notes, printed *) and **)
og Villiens ufuldkomne Udviklingstrin, deels under Formen af en ensidig
Erkjendelses- eller Villiesudvikling; deels under Formen af en Villiens
Ligegyldighed mod eller Uafhængighed af Erkjendelsen; deels under Formen af
Villiesactens Modstrid imod denne.

Den \emph{første} af disse Former for Divergentsen finder sin
Forklaring\footnote{Jfr. Problemet om Tro og Viden, S.\ 136 f.} i det
Sporadiske, der træder frem og maa træde frem paa de underordnede
Udviklingstrin, men som paa Udviklingens høieste Trin, paa det Trin, hvor
Erkjendelsen og Villien ere fornuftbestemte i deres concrete Indhold, maa
sammensluttes i Enheden og Harmonien. Og Enheden er ikke blot Udviklingens
Maal, men under den divergerende Udvikling vises der hen til den derigjennem,
at den ensidige Erkjendelsesudvikling ikke blot er mangelfuld med Hensyn til
Villie, men ogsaa mangelfuld med Hensyn til selve Erkjendelsen, ligesom den
ensidige Villiesudvikling ikke blot er mangelfuld med Hensyn til Erkjendelse,
men ogsaa med Hensyn til Villie, saa at altsaa det Sondredes Sammenhøren og
gjensidige Betingen bliver tydelig. Den ensidig udviklede Erkjendelse, der
% A modern pencil bracket runs down the margin of this copy from „Den“ here to
% „befriet Villie.“ four lines below. A reader's mark, not part of the
% edition; not transcribed.
ikke er forenet med en Personlighedens Villiesudvikling og Villiesfasthed, vil
altid vise sig at være en abstract og uigjennemført Erkjendelse; den ensidigt
udviklede Villie, der ikke er forenet med Erkjendelsens Rigdom og Klarhed, vil
vise sig at være en endnu ikke gjennemluttret, og derfor ikke fra det
Instinctmæssige, Fanatiske og Voldsomme, og derved Egoistiske, befriet Villie.

I den \emph{anden} Form for Differentsen gjør ligeledes\footnote{Sammesteds,
S.\ 142 ff.} en Virkning af det Sporadiske sig gjældende, som hvor der handles
ud fra en blot Drifttilskyndelse eller i blind Lidenskab, eller ogsaa har
Beslutningen en kun ufuldstændig Erkjendelse til Forudsætning, og der handles
ud fra den concentrerede Personligheds Umiddelbarhed. Med Hensyn til det
sidste Tilfælde er i det tidligere Skrift (S.\ 144) den Form
% --- p. 17 ---
\setcounter{footnote}{0}% p. 17 carries TWO notes and RESTARTS at *)
paavist, under hvilken i en saadan Beslutning en forudgaaende Erkjendelse, der
er „gaaet tilbage til umiddelbar Væren i Aanden“, er virksom\footnote{En
grundigere psychologisk Iagttagelse vil vise, at Begrebet om en Villiens
Bestemmethed ud fra en til umiddelbar Væren tilbagegaaet Erkjendelse ---
Erkjendelse taget efter Begrebets videste Betydning --- ogsaa faaer en
Anvendelse ved Villiens første Yttringer i Barnet, idet disse Yttringers
„Udprægethed“ bestandig Skridt for Skridt følger Forestillingslivets
Udprægethed.}. Det Forhold, der er mellem Beslutningen og den ufuldstændige
Erkjendelsesforudsætning, har sin nøiagtige Parallel indenfor det Theoretiske
i den Slutning, der, med et ufuldstændigt empirisk Grundlag, en ikke udtømt
Erfaring, til Forudsætning, udsiger det Almene, ikke hypothetisk, men
kategorisk.

I den \emph{tredie} Form for Divergentsen gjør deels\footnote{Problemet om Tro
og Viden, S.\ 144.} en Virken af det ethisk Aprioriske sig gjældende i
Villiens Bestemmen af sig i Modsætning til den ensidige Theorie, deels en
Villiens Modvirken af det Sande, gjennem en Paavirken af Erkjendelsen, og
gjennem en Forhindren af Erkjendelsens Realisation. Med Hensyn til den
Virksomhed, det ethisk Aprioriske udøver i Bevidstheden, er der i „Problemet
om Tro og Viden“, S.\ 145, viist hen til Parallelen med den Virken, det
Aprioriske udøver indenfor det Theoretiske, og det er godtgjort, at det
\emph{ethisk} Aprioriske, som \emph{Væsens}bestemmelse, og som Udtryk for det
\emph{udelte} Væsen, for hvilket ogsaa \emph{Viden} er Udtryk, ogsaa maa være
\emph{Videns}bestemmelse og være gjennemskueligt for Erkjendelsen, saa at dets
% letterspacing: „V æ s e n s =“ and „V i d e n s =“ close their lines spaced;
% „bestemmelse“ opens the following line tight in both cases.
Virken altsaa er en Virken af det, der finder sin Plads indenfor Viden, gaaer
ind i Systemet af dens Grundbegreber. --- Hvad Villiens Strid mod det
erkjendte Sande angaaer, saa vil denne vise sig, kun at finde Sted i Forhold
til en \emph{ufuldkommen} Erkjendelse, altsaa at begrændses til de
Udviklingsstadier, hvor det Sporadiskes Magt er uovervunden. En skarpere
psychologisk og ethisk Iagttagelse vil stadfæste, at den klare,
\emph{concrete} Erkjendelse aldrig er magtesløs imod Villien, at det er
umuligt at handle imod
% The signature numeral „2“ at the foot of p. 17 is not transcribed.
% --- p. 18 ---
en saadan \emph{actuel} Erkjenden. Kunde Villien modstaae hvad der var saaledes
erkjendt som \emph{Væsenets} Bestemmelse, som dets Fordring, vilde den handle
blot i Kraft af det Væsenløse, eller Væsenet maatte være spaltet i to heterogene
Dele; men hint er selvmodsigende, og en saadan Spaltning i Væsenet, ifølge
hvilken Villien skulde blive et Irrationelt, Erkjendelsen et Rationelt, er en
tom Forestilling, der ikke bliver mindre tom, om der end istedenfor en Spaltning
postulatsviis tales om en „Defecthed“ i Væsenet. I saadanne Tilfælde, der synes
at tale for Villiens ubetingede Magt til at modstaae Erkjendelsen, vil
Erkjendelsen deels vise sig som saaledes ufuldstændig, at den i sin Abstraction
ligesom giver Plads for den modstridende Handling; deels vil det vise sig, at
den, som endnu ikke helt udfoldet, ved en, klarere eller dunklere bevidst,
Villiesact, der paavirker Tankens Retning, bliver bragt udenfor Bevidsthedens
Focus, rykket ud i en Tidsafstand for Bevidstheden, og derved fordunklet og
hæmmet, saa at den kan paralyseres ved en modstridende illusorisk Erkjendelse,
en Undskyldning, et Paaskud; og først gjennem en saadan Spatieren bliver den
modstridende Handling mulig. Den klare Erkjendelse af det Uethiske i en
Handling, af dens Modsætning til en forudgaaende bedre Erkjendelse, indtræder
først \emph{efter} Handlingen. Her faaer \emph{Tids}bestemmelsen en afgjørende
Betydning.
% letterspacing: „T i d s“ is spaced and „bestemmelsen“ runs on tight, within one
% and the same line. Transcribed as printed; not regularised.

Medens saaledes den factiske Divergents, der paa de ufuldkomne Udviklingstrin
kan finde Sted mellem Erkjendelse og Villie, ikke viser sig at være af en saadan
Art, at den ophæver den indre Enhed, saa finder denne indre Enhed, i hvilken
Erkjendelse og Villie staae som Udtryk for det samme Væsen, sit Afpræg i den
Parallelisme mellem Erkjendelsens og Villiens almene Udviklingsstadier, der
træder frem, trods den uensartede Udvikling, de to Aandsfunctioner kunne have
hos den Enkelte. Ligesom Erkjendelsen fra sin første Form, Sandsningen, udvikler
sig, gjennem den uvilkaarlige Dannen af Almeenforestillinger, til en bevidst og
fremskridende Udsondring af det Almene i Begrebet, saaledes udvikler Villien sig
fra den første, umiddelbare, enkelte Driftbestemt%
% --- p. 19 ---
\setcounter{footnote}{0}%
hed, gjennem en Bestemmethed ved uvilkaarligt dannede, ved en spontan fremkommen
momentan og partiel Reflexion frembragte, \emph{praktiske Almeenforestillinger},
som dem, der gjøre sig gjældende i den blivende Tilbøielighed, til en
Bestemmethed ved bevidst dannede Grundsætninger, Maximer, af hvilke saa en
Livsanskuelses Enhed kan fremgaae.

Den Harmonie og Enhed, til hvilken der vistes hen under selve Divergentsen,
træder saa utilsløret frem paa det Erkjendelsens og Villiens Udviklingstrin, paa
hvilket den menneskelige Aand er naaet til at blive sig klar efter sin concrete
Væsensbestemmelse. Idet Menneskeaanden her i Erkjendelsen bliver sig sit Væsens
Indhold bevidst, udtrykker den dette Erkjendelsens Indhold i Villiens
Bestemthed: Villien gjør det Erkjendelsesindhold, i hvilket Væsenet bliver
\emph{sig} aabenbart, til \emph{sit} Indhold, realiserer Væsenet i sin Form, i
Beslutningen, og den gjør det ifølge den i det individuelt-almene, ideelt-reale
Væsens Enhed liggende Nødvendighed, ifølge den Væsenets Lov, der lader det
aabenbare sig i sin \emph{totale} Formbestemthed, og i denne aabenbare sig i
Harmonie med sig selv\footnote{Jvfr.\ Problemet om Tro og Viden, S.\ 142.}.
% printed mark: *) --- the only note on p. 19, below a short rule.
I Beslutningen er Villien bestemt ved Væsenets Nødvendighed, men idet denne
Nødvendighed er dens \emph{eget} Væsens Nødvendighed, er den \emph{selv}%
bestemmende ved Frihed.
% letterspacing: the line closes with the spaced „s e l v =“ and „bestemmende“
% opens the next line tight. As printed.
Forat Væsenet kan realisere sig i Villieshandling, i dette Begrebs egentlige
Betydning, maa det først være blevet sig aabenbart i Erkjendelsen; men Villiens
Selvbestemmen hævdes gjennem hint Forhold til Væsen\emph{stotaliteten}.
% letterspacing: the spacing begins mid-word, at the „s“ --- „Væsen“ is tight and
% „s t o t a l i t e t e n“ is spaced, all on one line. As printed.

\begin{center}\rule{2cm}{0.4pt}\end{center}
% A short centred rule stands here on p. 19, between the two paragraphs; it is a
% thought-break in the running text, not a head. A SINGLE hairline (checked at
% 300 dpi at 2x), like every rule in the book but the five doubles of pp. 4, 29,
% 63, 201 and 210.

Søge vi, efterat have belyst det almene Forhold mellem det Theoretiske og det
Praktiske, den nærmere Bestemmelse af \emph{Erkjendelsens sande Form}: af den
Form, der er dens Udviklings Maal, saa bliver den at finde ud fra Bevidsthedens
Grundforhold. I det Forhold, der i Menneskeaanden er givet mellem
Selvbevidsthed og Bevidsthed, ligger
% --- p. 20 ---
Nødvendigheden af en Stræben efter at bringe alt Bevidsthedsindhold ind under
Selvbevidsthedens Enhed, at ophæve Bevidsthedsgjenstandens Fremmedhed og at
gjøre den gjennemsigtig for Aanden i en begribende Erkjenden, i hvilken Selvet
har sin ideelle Udfoldelse og sin Selvaabenbarelse, idet det ideelt tilegner sig
Virkeligheden. Først i en saadan begribende Erkjenden, i hvilken
Bevidsthedsindholdet for Aanden danner en Enhed, svarende til Selvbevidsthedens
Enhed, er den Modsigelse hævet, der ligger i Bevidsthedsforholdet: at
Gjenstanden er i min Bevidsthed som \emph{min} Bevidstheds Bestemmelse, og dog
efter sit Indhold er et for mig Fremmedt, og først i den er ligeledes den
Modsigelse hævet, at Aanden, der i Selvbevidstheden paa uendelig Maade forholder
sig til sig selv og er i Enhed med sig selv, tillige er fremmed for sig selv ved
det fremmede Indhold, Selvet omfatter.
% letterspacing: only „m i n“ is spaced (the first „min“ four words earlier is
% tight, in the same line); „Bevidstheds Bestemmelse“ opens the next line tight.
Men naar en saadan Tilbageføren af Bevidsthedsindholdet til gjennemsigtig Enhed
med Selvbevidstheden er Erkjendelsens, ved det psychologiske Grundforhold
bestemte, Fordring, saa bliver en saadan Tilbageføren og en saadan Enhed kun
mulig paa Betingelse af en Erkjendelse af det Princip, der er Selvets og
Tilværelsens Enhedsprincip, og Erkjendelsens sande Form bliver da en
principbestemt Erkjendelsens Enhed, i hvilken Selvet, idet det erkjender sig som
organisk Led indordnet i den principbeherskede Totalitet, kan erkjende sig efter
sin Concretion. En saadan Erkjendelsens Form, i sit Væsen liig den, Philosophien
efter sit Begreb skal give, maa, ifølge sit Forhold til Bevidsthedens
oprindelige Fordring, ifølge sit universelle Princip, ifølge den Tankens Art,
den forudsætter, ifølge den Bevidsthed, den indeslutter om Tankens Væsen og
Grundforhold, ifølge sin Forbinden af det Aprioriske og det Empiriske, hævde sig
Pladsen som Erkjendelsens τέλος.
% Greek, antiqua italic with accent, read off the image; neither OCR witness
% reads it (tesseract „rédos“, ABBYY nothing).

Søge vi \emph{Handlingens sande} Form, da maa denne bestemmes ud fra Villiens
Grundforhold til den selvbevidste Aands Væsen. Idet Mennesket er selvbevidst
Væsen og som selvbevidst er villende, faaer den Virken, der hos Dyret, som
ubevidst og ufri Driftfuldbyrdelse, kun er en
% --- p. 21 ---
Gjerning, Charakteren af \emph{Handling}, ɔ: af en af \emph{Villie} fremgaaende
\emph{bevidst} Virken. Nærmere bliver Handlingens Natur saaledes at bestemme, at
den ikke blot er ledsaget af Bevidstheden om, \emph{at} der handles, men at den
er bestemt ved Bevidstheden om, \emph{hvorfor} der handles. Handlingen som
Handling, i Ordets egentlige Betydning, er rettet mod et Maal, et
\emph{Øiemed}, er ledet ved dette, er en bevidst Virken for dette Øiemeds
Opnaaelse. Handlingens Maal, det bevidst villede Øiemed, kan, seet fra det
Objectives Synspunkt, vise sig som et umiddelbart eller som et middelbart Maal,
som det, der staaer directe i Forhold til Villien, villes for sin egen Skyld,
eller som det, der kun gjennem et Andet staaer i Forhold til den, villes for et
Andets Skyld, som Middel til Opnaaelsen af Øiemedet. Men naar vi, med Hensyn til
det Objective, tale om en Villen af det for dets egen Skyld, saa er dette dog
kun at forstaae i Sammenligning med det objective Middel, ikke saaledes, at
Subjectets Forhold til sig selv udelukkes, at den Handlendes subjective
Interesse ved det Objective glemmes. Handlingens Subject forholder sig, som
Subject, paa ethvert Punkt til sig selv; idet det bestemmer det Objective som
Øiemed, bestemmer det det som \emph{sit} Øiemed, og forholder sig \emph{til sig
selv} i Virkeliggjørelsen af Øiemedet, vil \emph{sig selv}, idet det vil det
Objective. Gjennem Relationen til det Objective føres vi derfor tilbage til
Relationen til Subjectet, og at vi føres tilbage dertil, ligger i selve
Subjectivitetens Væsen, i dens Grundbestemmelse: at være for sig i sit Andet.
Handlingens Øiemed bliver da overalt Subjectet selv, og dette gjælder ogsaa om
den Handling, gjennem hvilken Subjectet hengiver sig til det Høiere; thi gjennem
Selvhengivelsen vindes den høiere Selvbekræftelse.

Men naar den angivne Bestemmelse gjælder for \emph{al} Handlen, saa falder
indenfor den Modsætningen mellem den \emph{sande} og den \emph{usande} Handlen.
% letterspacing: „a l“ closes its line spaced, „Handlen“ opens the next tight.
Handlingens Sandhed bestemmes ved Villiens (Villiesbestemthedens) Sandhed, men
Villiens Sandhed bestemmes ved, om det villende Subject, idet det vil \emph{sig
selv}, vil sig selv \emph{efter sin Sandhed},
% --- p. 22 ---
ɔ: efter sin \emph{Væsensbestemmelse}, og nærmere, efter sin \emph{fulde}
Væsensbestemmelse. Den \emph{sande Handling} bliver altsaa den, i hvilken
Villien gjør Virkeliggjørelsen af det Væsen, for hvilket den er et primitivt
Udtryk, til Øiemed, i hvilken den altsaa virker som \emph{Væsensvillen}.

I Modsætning til denne sande Handlingens Form, som vi anticiperende betegne som
den \emph{ethiske} Handlen, bliver, idet vi endnu see bort fra den uethiske
Handling, hvis Væsen senere vil blive bestemt, den \emph{usande} Handlingens
Form den \emph{endnu ikke} ethiske, ufuldkomne Handling, der fremgaaer af
Subjectets Villen af sig selv efter sin particulære Umiddelbarhedsbestemthed,
efter sin Endelighed og Tilfældighed, efter sin phænomenale Existents. En saadan
Villen bliver en Subjectets Villen af sig selv efter sin ubeherskede
Sandselighedsside --- Sandselighed tages her i Begrebets videste Omfang --- som
en ubevidst egoistisk Villen af Nydelsen. I denne Handlingens Form, og i den
Villie, af hvilken den fremgaaer, ligger en indre Modsigelse, der aabenbarer
dens Usandhed. Naar Selvet, som sandseligt bestemt, vil Nydelsen, vil det
\emph{sig} gjennem Nydelsen, idet det i Nydelsen vil sin Selvbekræftelse som
sandseligt. Men idet Selvet vil Nydelsen, vil det det, der er
\emph{uafhængigt} af Selvets Villen, og som efter sin Natur er
\emph{foranderligt}, altsaa det, der \emph{ikke} kan beherskes af Villien og i
sin Incommensurabilitet med Selvets almene Væsen \emph{ikke} kan fastholdes af
Villien. Og idet Selvet vil sig i Nydelsens uforenede Mangfoldighed og væsenløse
Foranderlighed, vil det sig selv i en Bestemthed, den sandselige Udvorteshedens
Bestemthed, i hvilken dets Enhed ophæves, det vil sig altsaa selv i det, i
hvilket det, medens det søger en Selvbekræftelse, \emph{taber sig selv}; det vil
sig som Ikke-Selv. Den indre Modsigelse hæves kun, naar Villien, der som
Væsensyttring, efter sit Begreb maa være Væsensvillen, bliver en Villen af det,
der i Sandhed \emph{kan} villes, fordi det er Eet med den Villendes Væsen, en
Villen af dette, i hvilken Selvet bekræfter sig efter sin Enhed og blivende
Væsentlighed, efter sit Væsens Uendelighed. Til en saadan Villen af sig selv
% p. 22 carries NO footnote --- the page was on the brief's list of note-bearing
% pages, but the image shows an unbroken type block with no rule and no note.
% --- p. 23 ---
vises der paa Villiens ufuldkomne Udviklingstrin hen som til det Maal, ved
hvilket Villien finder Hvile, idet den virkeliggjør, hvad der uklart tilstræbtes
paa hine lavere Udviklingstrin. Villiens og Handlingens Maal, der bestemmes ved
Selvet, og i hvilket Selvet forholder sig til sig selv, maa fra Begyndelsen af
skjult være bestemt ud fra \emph{Selvets Enhed} --- en Bestemmen, der røber sig
gjennem det Almenes Fremarbeiden --- og idet Selvet bliver sig klart efter sit
Væsen, maa det \emph{bevidst} og \emph{frit} være bestemt ud fra \emph{Selvets
indholdsrige Enhed}, saa at Selvet efter sin \emph{totale} Væsensbestemthed,
efter sin sande Almindelighed og Uendelighed, bliver Villiens og Handlingens
modsigelsesfrie τέλος. Men en saadan Selvets Villen af sig efter sin Sandhed,
efter sit Væsens Uendelighed, er den \emph{ethiske} Villen.
% Greek: a second τέλος, at the head of a line on p. 23. It is NOT in the recon
% list (which gives pp. 20, 65, 185, 193); read off the image, antiqua italic
% with the accent. ABBYY drops the word entirely; tesseract reads „zéloc“.

For denne ethiske væsentlige Villen er i det nærmest Foregaaende en dobbelt
Bestemmelse antydet, der staaer i nøieste indre Sammenhæng. Det Ethiske
constitueres, idet Selvet vil sig \emph{efter sit Væsen}, og idet det, hvad der
er Betingelsen herfor, drager sig ud af sin umiddelbare Væren med dens Ufrihed,
Endelighed og Foranderlighed, og ved \emph{Frihed} bestemmer sig \emph{som Aand}
efter sin \emph{indre Uendelighed}. Paa det Stadium, der ligger forud for det
Ethiske, vilde Selvet det, der som endeligt er et Forsvindende og som er
uafhængigt af Villien; det vilde altsaa \emph{sig} i hvad der er et Forsvindende
og et for Selvet Fremmedt, d.\ v.\ s.\ det tabte sig selv. Idet det bliver sig
det Villedes Forgængelighed og den deri liggende Modsigelse bevidst, træder
Modsætningen frem mellem det forgængelige Endelige og det Uendelige og
Ubetingede; der viser sig et \emph{Valg}, og idet Villien, i Fortvivlelsen om
det Endelige, frigjør sig fra det og vælger det Uendelige, bryder det Ethiske
igjennem; men det Uendelige vælges netop, idet Selvet vælger sig selv efter sit
Væsen; thi det vil sige: efter \emph{Evighedsbestemmelsen} i sig, efter sin
\emph{evige Gyldighed}. I denne \emph{Valgets Beslutning}, der er en
\emph{Frihedens Selverhverven}, faae Modsætningerne i Valget en ny Belysning og
en ny
% --- p. 24 ---
\setcounter{footnote}{0}%
Bestemmelse: Bestemmelsen af det \emph{Gode} og det \emph{Onde}: Begreber, som
først høre hjemme paa det ethiske Gebet, og i det ethiske Valg vælges det Gode.

Idet Selvet vælger sig selv efter sit Væsen og gjennem Valget vil sit Væsen, har
det altsaa gjort en \emph{Uendelighedens} Bevægelse: dets \emph{umiddelbare}
Væren efter dets \emph{Natur}bestemthed er omsat i en \emph{Aandens} Væren ved
\emph{sig selv}, ved sin \emph{Frihed}.
% letterspacing: „N a t u r“ is spaced and „bestemthed“ runs on tight in the same
% line; „ved s i g / s e l v“ carries the spacing across the line break.
Idet Selvet, gjennem hint uendeliggjørende Valg, vælger og vil sig selv, vil det
sig \emph{absolut}, ɔ: \emph{som et Absolut}, og \emph{kun sig selv}, ɔ:
\emph{sit Væsen}, kan Selvet ville absolut; men denne Villen er kun mulig, idet
det i Selvet vil det \emph{Absolute}. \emph{Det Selv}, der vælger sig selv ved
den ethiske Beslutnings absolute Villen, har en \emph{dobbelt} Betydning: det er
det \emph{ideale, almene Selv}, Selvets absolute Væsen, der vælges, men det
vælgende Selv, der vælger det som \emph{sit} Væsen, er det \emph{individuelle
Selv}, der i Valget af hint Selv vælger sig i sin Individualitet. Dermed er
sagt: det \emph{overtager} sin, forud for det ethiske Valg liggende Udvikling,
men idet det overtager den i Valget af det ideale Selv, overtager det den
saaledes, som den viser sig i dettes Lys. Det vil med andre Ord sige: det fælder
\emph{Angerens} ethiske Dom derover, det \emph{forvandler} sig selv principielt
i den ethiske Beslutning i Kraft af sit Grundvæsen. I denne angrende Vælgen
samler Selvet sig i hele sin endelige Concretion, og idet den i Valget ligesom
tager sig ud af Endeligheden, er det i den absoluteste Continuitet med den og
har sig selv til Opgave i sin ethisk valgte concrete Virkelighed, der skal
gjennemarbeides med, gjennemtrænges af det Almene, saa at det virkelige Selv
bestemmes ved det ideale Selv, det har i sig, og paa een Gang udtrykker det
Individuelle og det Almeenmenneskelige\footnote{Jfr.\ Udviklingerne i
\emph{Kierkegaards} Enten --- Eller, II.}.
% printed mark: *) --- the only note on p. 24, below a short rule. The name
% „K i e r k e g a a r d s“ is letterspaced Fraktur, not antiqua.

For den \emph{ethiske Villen}, hvis Væsen og Opgave saaledes er bestemt, viser
sig nu, saavel i dens første Constitueren af sig, som i deres Udfoldelse i
Handlingens con%
% „deres“ (for „dens“) is what the page prints; transcribed as printed.
% --- p. 25 ---
crete Bestemthed, den \emph{sande Erkjendelse} som væsentlig og nødvendig
\emph{Betingelse}.

Den umiddelbare Virken af det ethisk Aprioriske, ved hvilken en saadan
Erkjendelsesforudsætning kunde synes at erstattes, viser sig nemlig som
utilstrækkelig. At hint Aprioriske har en Virksomhed, der aabenbarer sig i
Samvittigheden, er uimodsigeligt; men forat denne Virksomhed kan hæves over den
blot indhyllede (implicite) Tilstedeværen i det kun Instinctmæssiges og
Driftagtiges Form, saa forudsætter den en Erkjendelsesudvikling, ligesom det
logisk Aprioriske forudsætter en saadan forat frigjøres for det Hæmmende og
Fordunklende. I hiin første Virken arbeider det ethisk Aprioriske umiddelbart
sammen med Drifttilskyndelser, virker selv under Drifttilskyndelsens,
f.\ Ex.\ under Sympathiens, den naturlige Kjærligheds Form; men er ikke udklaret til
ethisk Frihed. Derfor bliver hint ethisk Aprioriske i sin umiddelbare Virken
utilstrækkeligt til Regulativ for Handlingen. I den umiddelbare Tilskyndelse,
der endnu ikke er gjennemklaret af Erkjendelsen, bliver det bestandig usikkert,
hvorvidt det Tilskyndende er et i Sandhed Apriorisk, og som saadant en sand
Væsensbestemmelse, eller om det ikke er et Usandt, en Yttring af et Uberettiget,
et Phænomenalt og Tilfældigt. Og idet det Aprioriske \emph{som Apriorisk} har
det Abstractes Præg, saa bliver det, naar det ikke udfoldes til Concretion ved
Erkjendelsen, utilstrækkeligt til Bestemmelsen af den concrete Handling. Idet
der skal handles, vil Subjectet komme til at søge Handlingens Bestemthed
udenfor den abstracte Norm i sin ikke gjennemklarede naturbestemte
Umiddelbarhed, i hvilken det ikke er sig bevidst efter sit Væsen, og der vil da
blive en Delthed og Dobbelthed i Villien, idet det Almene og det Enkelte, Normen
og Handlingens concrete Indhold, ikke kan villes som Eet. Denne modsigende
Dobbelthed kan kun forsvinde naar Væsensfordringen udfolder sig i concret
Bestemthed for Erkjendelsen.

Forskjelligheden i det, der til de forskjellige Tider, paa de forskjellige
menneskelige Udviklingstrin, har været anseet
% --- p. 26 ---
for det sædeligt Rette, stadfæster Utilstrækkeligheden af den umiddelbare
Tilskynden ved det Aprioriske. Man tage som Exempel blot den Maade, paa hvilken
Følelsen af Ærefrygt for det Guddommelige eller den umiddelbare Pietetsfølelse
mod Forældre har fundet sit Udtryk.

I Kraft af det ethisk Aprioriske vil der kunne være en umiddelbar Bevidsthed om
en ethisk Forpligtelse, som en Forpligtelse mod et Høiere; men den klare
Bevidsthed om \emph{det Forpligtende} er betinget ved den sande Erkjendelse, og
ligeledes er den \emph{sikkrede} Bevidsthed om Forpligtelsens concrete Indhold
betinget ved en saadan Erkjendelse, ved den concrete Erkjendelse af det
Forpligtende og af den Forpligtedes \emph{Væsensbestemthed} og \emph{væsentlige
Inhold}.
% Printer's defect: the page prints „I n h o l d“ (letterspaced) for „Indhold“ ---
% the „d“ is simply absent. Positively verified on the image at 200 dpi, and both
% OCR witnesses read „Inhold“. Transcribed as printed, not corrected.
Instinctmæssig kan der leves i Overensstemmelse med det Sædelige, idet det
objectivt Sædelige: Sæder og Tradition, umiddelbart bestemme det handlende
Individ, der i et Pietetsforhold underordner sig under dette Objective som under
en høiere Magt, i hvilken det føler et Beslægtet, og hvis Lov det derfor følger,
ikke som en fremmed Lov. Instinctmæssig kan den umiddelbare Bevidsthed gribe det
Gode, handle i Sædelighedens Aand, idet den, ubevidst paavirket af det
omsluttende sædelige Livselement, anerkjender et Alment og Ideelt som
Handlingens Norm, opgiver det Egoistiske og i den enkelte Handling handler med
Resignationens Villie, ledet ved sin principielle Bestemmethed. Men Sæder og
Tradition, den omsluttende sædelige Livsatmosphære, kunne, hvad enten de ere
Resultanten af det Aprioriskes umiddelbare Virken eller af en forudgaaet
Erkjendelse, der ligesom er nedfældet i umiddelbar Væren, kun være den sikkre
Norm naar de atter fra den umiddelbare Form, i hvilken et Tilfældigt og Fremmed
kan klæbe ved dem, kunne frigjøres i levende Erkjendelse: den klare, frie,
sikkrende Bevidsthed om det Sædelige og den indre Sammenhæng i det kan kun naaes
gjennem en saadan, Umiddelbarhedens Tilfældighed udskillende Erkjendelse, og
bestemtere gjennem den Erkjendelsens Art, ved hvilken det Erkjendte kan faae
indre Enhed og Nødvendighed, alt%
% --- p. 27 ---
saa ved den principbestemte Erkjendelse, hvis Væsen i det Foregaaende er
udviklet. En saadan Erkjendelse har, i sin fra de menneskelige Udviklingsvilkaar
uadskillelige Ufuldstændighed, i sin relative Abstraction, den væsentlige
Sandhed og Muligheden til ligesom at voxe ud fra sit Centrum og at udfylde sine
Former uden at de forskydes. Med den til Forudsætning faaer den ethiske Villie,
netop fordi Væsenet i hiin Erkjendelse er klart for sig selv, sin sande Form, og
i sin Frigjorthed fra endelig Hildethed, fra det uklare Driftagtige, der i sig
er det Foranderlige og det, der, selv i sin Retning mod det Gode, ikke kjender
Grændse, faaer den den sande Besluttethed, Fasthed og Renhed, der er betinget
ved det Villedes Væren for Erkjendelsen. Subjectets concrete Erkjendelse af sit
Væsen og af Væsenets sande Fordring faaer her sin afgjørende praktiske
Betydning.

Efter det Udviklede bliver det klart, at Erkjendelsens Betydning for det Ethiske
allerede træder frem hvor det ideale Selv \emph{med Bevidsthed} skal gjøres til
Villiens Gjenstand, hvor et \emph{Absolut} og \emph{Ubetinget} forkynder sig og
skal sondres ud fra det Tilhyllende og Fordunklende. Men endnu klarere træder
den frem naar det Godes Idee skal realiseres i \emph{Concretion} af den sædelige
Villie, udformes i en Livsheelheds Virkelighed, naar, hvad der bliver identisk
hermed, Selvets Villen af sig selv efter sit ideale Væsen skal naae til den
ethiske Handlings concrete Virkelighedsform, til det menneskelige Væsens
Virkeliggjørelse ved den frie, ved den med sig selv og sit Væsen i Sandhed
overensstemmende menneskelige Handlen. Thi kun ved den sande og concrete
Erkjendelse af Selvets Væsen og af dets Grundforhold bliver den frie Handlen i
sin Bestemthed \emph{fornuftbestemt} Handling, udfolder det Ethiske sit
Fornuftindhold \emph{klart og ublandet}. Idet det \emph{ethisk Aprioriske} i
dets Abstraction mødes med et concret Givet, med et \emph{Empirisk}, saa kan i
denne Møden, der skal føre til en Bestemmen af dette ved hiint, til en Udfolden
af den ethiske abstracte Væsensbevidsthed til væsentlig
Virkelighedsbevidsthed, kun den sande Erkjendelsens Form, der anerkjender det
% --- p. 28 ---
\setcounter{footnote}{0}%
Aprioriske og virkeliggjør dets væsentlige Enhed med det Empiriske, altsaa en
Erkjendelse af den ovenfor bestemte Art, være i sidste Instants veiledende. Kun
ved den kan Selvets Villen af sit Væsen blive en Villen af Væsenet efter dets
fuldstændige concrete Bestemthed; kun ved den kan det i sin Handlen virkelig
„ville Eet“\footnote{Jfr.\ Afsnit III (om Idealets Realisation).}. ---
% printed mark: *) --- the only note on p. 28, below a short rule. „III“ is set
% in antiqua, as all figures in this book are; not italicised here.

Et Udtryk for denne Erkjendelsens Betydning for Handlingen er det, at
\emph{Ethiken}, det Ethiskes videnskabelige Erkjendelse, maa faae sin
\emph{væsentlige} Betydning for dette Ethiske. Er den \emph{videnskabelige}
Erkjendelse Erkjendelsens klareste og fuldstændigste Udfoldelse, saa maa det,
der er sagt om den sande Erkjendelse, umiddelbart faae sin Anvendelse paa den.

Idet den ethiske Handlen af den Grund er bestemt som Handlingens sande Form, at
Villien, af hvilken den fremgaaer, er en \emph{Væsens}villen, en Villen af
Selvet efter dets Væsensbestemthed, saa faaer den ethiske Handling sin
\emph{Norm} i \emph{Menneskevæsenet}, sin \emph{Lov} i Væsenets egen
Virksomhedslov, og dens Opgave bliver at udtrykke den totale virkelige
Personlighed.
% letterspacing: „V æ s e n s“ closes spaced and „villen“ runs on tight, in the
% same line. As printed.
Herved vises der atter tilbage til \emph{Modsætningen til det positivt
Religiøse}. Denne Modsætning viser sig her deels som en Modsætning mellem
Væsenets iboende Lov og den supranaturale, vilkaarlig bestemte Handlingslov,
deels som en Modsætning mellem et ethisk Liv, der omfatter den totale, reale
Existents, og et Liv, der, som lydende den supranaturale Fordring, flygter bort
fra det Naturlige, forholder sig negativt til det gjennem Askesen, og udtrykker
Væsenets Aandsside abstract spiritualistisk.

\begin{center}\rule{2cm}{0.4pt}\end{center}
% A short centred rule stands here on p. 28, between „spiritualistisk.“ and the
% closing summary: a thought-break, not a head rule. A SINGLE hairline.

Resultatet af de i det Foregaaende givne Udviklinger sammenfatte vi i følgende
Sætninger:

\begin{enumerate}[label=\arabic*., leftmargin=*, labelsep=1em,
                  topsep=0.35\baselineskip, itemsep=0.35\baselineskip]
\item Den sande Erkjendelse, der vælges i Valget mellem \emph{Viden} og
  \emph{Tro}, bliver Betingelsen for den klart be%
% --- p. 29 ---
vidste Villen af det Ethiske, og for dettes concrete Udfoldelse til en Heelhed,
  der omfatter den virkelige Personligheds Liv.
\item Det personlige Liv kan gjennemklares af den sande principbestemte
  Erkjendelse, og denne Erkjendelse kan, idet den skrider fremad i Concretion,
  vinde en saadan Punktualitet, at den concrete ethiske Handling i den kan finde
  sin \textit{ratio sufficiens}.
% antiqua Latin on p. 29, upright in the print; set \textit{} per the house
% convention taken from the 1868 volume.
\item Den sande Enhed af Theorie og Praxis, af den „theoretiske“ og den
  „praktiske“ Bevidsthed, er tilstede i dette Forhold, i hvilket Subjectets
  Bevidsthed om sin Handling og sine „Villiesgrunde“ er i Harmonie med dets
  theoretiske Erkjendelse af sit Væsen og af Væsenets Fordring; --- og den er
  \emph{kun} tilstede i dette.
\item Ifølge dette Forhold mellem Erkjendelse og ethisk Handlen maa, da ingen
  Adsplittelse i incommensurable Former finder Sted indenfor Erkjendelsen, alt
  sandt Ethisk være et rationelt Ethisk, ligesom al sand Ethik en rationel
  Ethik. Et antirationelt Ethisk er en Selvmodsigelse, en antirationel Ethik en
  Uting.
\end{enumerate}

\begin{center}\rule{2cm}{0.4pt}\end{center}
% p. 29 is the last page of division I and is a short page: the type block ends
% after item 4 and a short centred rule closes the division. The rule is printed
% as a DOUBLE rule (two thin lines); set here with the same single \rule the
% skeleton uses for the other short rules in this book.

\chapter*{II.\\ Om Tilstedeværelsen af det i sand Betydning Religiøses væsentlige Kjendemærker i den sande Erkjendelse og den sande Handlen, og om disses Enhedsforhold til hiint.}
\addcontentsline{toc}{chapter}{II. Om Tilstedeværelsen af det i sand Betydning Religiøses væsentlige Kjendemærker}
\markboth{II. Om Tilstedeværelsen}{II. Om Tilstedeværelsen}

% --- p. 30 ---
\setcounter{footnote}{0}%
% p. 30 carries the division head II (already in the skeleton), then a short
% centred rule (not transcribed, per the book's convention for rules under
% display heads), then the bold centred sub-head below. The text opens with an
% enlarged initial capital D (about one and a half lines); set normally.
\section*{1. Det Religiøses Væsen.}

Det Udtryk for det religiøse Forhold, der i de „foreløbige
Begrebsbestemmelser“ i „Problemet om Tro og Viden“ er brugt til at betegne
det efter dets ideale Bestemthed, bliver her vort Udgangspunkt, og gjennem
Analysen af det ville vi kunne naae til den concretere Udfoldelse af det
Religiøses Begreb.

Vi bestemte der det religiøse Forhold som „\emph{det absolute
Forhold}\footnote{Hvis det for „omvendte“ Dialektikere skulde synes
modsigende at tale om et \emph{absolut Forhold}, fordi Forhold udtrykker
\emph{Relation}, behøve vi vel blot at minde om det Absolutes Forhold til sig
selv, og om Aandens Uendelighedsforhold til sig selv i Selvbevidsthedens
Concretion.} \emph{til det Absolute}“, altsaa ved et Udtryk, der ikke blot i
Formen nærmer sig til S.\ Kierkegaards almindelige Betegnelse af den endnu
ikke christelige Religiøsitets Stadium, hvilket han charakteriserer ved „den
absolute Retning mod det absolute τέλος“, men som allerede er brugt af ham
hvor han sammenstiller det religiøse absolute Forhold med de relative
Forhold, der hos ham ogsaa omfatte de ethiske Forhold, i hvilke hiint ikke
skal udtømme sig eller blive concret. Medens vi saaledes tilegne os hans
% printed mark: *) --- the only note on p. 30, below a short rule.
% Greek: τέλος, antiqua italic with acute, read off the image; the ABBYY layer
% carries no Greek at all. The second quotation, „den absolute Retning mod det
% absolute τέλος“, is NOT letterspaced, though the first one is.
% --- p. 31 ---
Begrebsbetegnelse, vil det imidlertid gjennem de efterfølgende Udviklinger
vise sig, hvorledes denne Betegnelse i vor Opfattelse faaer en anden Betydning
end den hos Kierkegaard havde, allerede hvor han taler om det forchristelige
religiøse Stadium, og endnu bestemtere hvor han taler om det paradox
Religiøse.

Analysere vi Bestemmelsen af det religiøse Forhold som det absolute Forhold
til det Absolute, saa fører Analysen os til følgende Consequentser, gjennem
hvilke Begrebet faaer sin concretere Bestemthed.

a) \emph{Med Hensyn til det Absolute}, der er Forholdets Object, indeslutter
Grundbestemmelsen Følgende:
% the enumerators a)--h) on pp. 31--41, and the Greek α) β) on p. 31, are
% antiqua italic in the print; set plain here, as in the Indhold.

α) Det „Absolute“, der er Gjenstand for det religiøse Forhold, kan kun være
et Absolut \emph{for os}, naar det er \emph{vort} Absolute ɔ: naar det staaer
i et \emph{Væsensenheds}-Forhold til os. I Væsensmodsætningen vilde
umiddelbart ligge en Bestemmelse, der er uforenelig med Absolutheden. Mod
denne Væsensenhed strider ikke den af Kierkegaard saa stærkt accentuerede
ethiske „absolute Distinction“ mellem Gud og det skyldige Menneske; thi
Skyldens Bestemmelse ligger, som det nærmere vil blive udviklet under III,
indenfor Immanentsforholdet, og gjør derfor intet definitivt Brud paa
Væsensenheden.
% letterspacing: „V æ s e n s e n h e d s =“ closes the line spaced and
% „Forhold til os.“ opens the next tight. The capital F shows this is a real
% hyphenated compound, not a broken word. Transcribed as printed.
% „med Absolutheden“: a small speck of ink stands over the e of „med“ (both
% OCR witnesses read an accent). It is dirt, not a printed sort; „med“.

I den positive Religions Opfattelse af Forholdet mellem Gud og Mennesket er
der en Dobbelthed tilstede. Paa den ene Side accentueres, ud fra den
umiddelbare religiøse Trang, Ligheden: Mennesket sættes som skabt i Guds
Billede; men paa den anden Side fører den af Forestillingsformen bestemte
Reflexion til at accentuere Modsætningen mellem den transscendente Gud og
Mennesket, og denne Modsætning fæstner sig som en uophævelig (jfr.\ Afsnittet
III, om Forsoningen).

β) Det religiøse Forholds Gjenstand kan kun være det Absolute for os, naar
det er \emph{det i Sandhed Uendelige}. Kun det i Sandhed Uendelige kan absolut
villes, absolut søges gjennem det \emph{hele} menneskelige Væsens Trang, og
kun hvad vi ville absolut, absolut stræbe hen imod, er for
% --- p. 32 ---
os et Absolut. En Villen af et endeligt Bestemt er en endelig Villen,
Forholdet til det et endeligt, relativt Forhold. Denne Bestemmelse: at det
Absolute maa være det i Sandhed Uendelige, er uadskillelig fra den, at det maa
være i Væsensenhedsforhold til os. Thi det kan kun være det i Sandhed
Uendelige naar det ikke ved en forestillet, altsaa endelig bestemt,
Transscendentses Skranke er adskilt fra os.

Kun disse mere formelle Bestemmelser af det Absolute ere i nærværende
Sammenhæng nødvendige; de concretere Begrebsbestemmelser ville blive givne
under III, Cap.\ 1.

b) \emph{Med Hensyn til Forholdets Subject: Mennesket}, følger af
Bestemmelsen: det absolute Forhold, at Subjectet maa være tilstede i Forholdet
efter sit Væsens \emph{indre Uendelighed}, d.\ v.\ s.\ efter sin \emph{hele}
Væsensbestemthed, efter sit \emph{universelle} Væsens Totalitet. Det religiøse
Forhold kan altsaa ikke være et Forhold, der forholder sig til en enkelt Side
af den menneskelige Aand eller til en enkelt Grundfunction, selv om denne
tænktes at indeslutte de andre i sig som Momenter; hvorved de netop kun
\emph{i den Bestemthed}, de kunde have som Momenter, vilde træde ind i
Forholdet; men det maa være et Forhold, i hvilket Menneskeaandens
Grundfunctioner, som i sig totaliserede, ligeligt træde ind. Ligesom det
altsaa bliver en Ensidighed ved den \emph{Hegelske} Religionsphilosophie, at
den opfatter det religiøse Forhold udelukkende som et \emph{Erkjendelses}forhold,
saaledes bliver det ligesaafuldt en Ensidighed, at bestemme det som et blot
\emph{Villies}forhold, eller at betegne det som det, man med et noget uklart
Udtryk har kaldt: et \emph{Livs}forhold, i \emph{Modsætning} til et
Erkjendelsesforhold. Det religiøse Forhold er et \emph{Livsforhold i Be}tydning
af et Forhold, der omfatter \emph{hele Menneskelivet}; men Menneskelivet,
Aandslivet er Erkjendelsens Liv ligesaavel som Villiens og Følelsens; det
opfattes derfor paa ensidig og ufuldkommen Maade, naar det kun opfattes som
Erkjendelsen modsat. Men idet Forholdet omfatter Menneskelivets Totalitet,
faaer paa den anden Side Villiesforholdet en ligesaa afgjørende og primitiv
Betydning, og det
% letterspacing on p. 32 breaks in the middle of compounds and runs past the
% word: „E r k j e n = d e l s e s“ spaced with „forhold“ tight; the same for
% „V i l l i e s forhold“, „L i v s forhold“; and „L i v s f o r h o l d  i
% B e =“ is spaced while „tydning“ on the next line is not. Transcribed as
% printed; not regularised.
% --- p. 33 ---
\setcounter{footnote}{0}%
kan træde i Forgrunden hvor Forsoningen \emph{søges}. Villien i Forholdet til
det Absolute er, naar Forholdet er i Sandhed, den religiøs bestemte,
\emph{ethiske} Villie, der handlende omdanner Existentsen ved Forholdet. Ved
denne Villie bliver Forholdet et \emph{Samvittigheds}forhold. At Forholdet er
et \emph{Personlighedens Forhold}, viser ligemeget tilbage til Erkjendelse og
Villie\footnote{Jfr.\ forrige Afsnit.}: hvad der er en Grundyttring af
Væsenet, er med det Samme en Grundyttring af Personligheden, ellers vilde
Personlighedsbestemmelsen staae i et blot tilfældigt Forhold til Væsenet, hvad
der umiddelbart er selvmodsigende.
% printed mark: *) --- the only note on p. 33, below a short rule.

Det, hvorefter Aanden uendeligt stræber gjennem det saaledes opfattede
absolute Forhold, er \emph{Forsoningen}, og det religiøse Liv bliver et Liv
for Aandens Forsoning \emph{i sig og med sin Væsensgrund}. Dette er endnu mere
en \emph{metaphysisk} Bestemmelse af det religiøse Forhold, og den naaes
allerede gjennem Begrebet af Aanden som det efter sit Væsen Uendelige, der
forat finde Hvile i sig selv stræber at virkeliggjøre denne Uendelighed i sig,
men først virkeliggjør den gjennem det absolute Forhold til det i Sandhed
Uendelige, til sin og Tilværelsens Væsensgrund. Den specifisk
ethisk-religiøse Bestemmelse af Forholdet: Forsoningens Bestemmelse som
Forsoning af \emph{Skylden}, vil først senere blive udviklet.

Idet den \emph{hele} Menneskeaands ligelige Gaaen ind i det religiøse Forhold
saaledes er accentueret, blive de afvigende Opfattelser at berøre, der
fremhæve den ensidigt begrændsede Bestemmelse af et \emph{Livsforhold}. Vi
kunne her tage Hensyn deels til \emph{Grundtvig}, deels til \emph{Martensen}.

Hos den \emph{Første} lægges der et Eftertryk paa \emph{Livs}forholdet i
Betydning af et \emph{Umiddelbarheds}forhold, der bliver saaledes opfattet, at
Umiddelbarheden faaer den sandselige Nærværelses æsthetiske Charakteer og
kommer til i sig at indeslutte det Magiskes Moment. I Opfattelsen af det
„levende Ords“ Virkning, af Sacramentets Magt, giver
% signature numeral „3“ at the foot of p. 33 (first page of the gathering);
% not transcribed.
% --- p. 34 ---
her Aandens Frihedsforhold til Aanden Plads for en magisk Virken, og medens
den religiøse Bevidsthed, der ved et Postulat afskærer Reflexionen paa et
bestemt Punkt, ved Symbolet, søger Hvile i en urokkelig, ubevægelig
Objectivitet, ombytter den det i Sandhed religiøse Forhold med et æsthetisk,
digterisk Forhold, i hvilket Personligheden selv ikke har gjort
Uendelighedens Bevægelse, gjennem den sande Inderliggjørelse, i hvilket den
ikke kjender den uendelige Resignation, men forholder sig til det
Guddommelige uden at kjende til Forholdets Alvor, uden at mindes, at Tvivlens
Mulighed i hvert Øieblik bliver det, der skal overvindes, forat Forholdet kan
blive i Sandhed, og at Uvishedens Angest bliver den Grund, paa hvilken
Troesforholdets seierrige Vished skal hvile. Idet der, med halvbevidst Sky
for Tanken, uklart tales om et Livsforhold, i hvilket Tankens Liv ikke
optages, bliver det paa dette Standpunkt overseet, at et Forhold gjennem
Erkjendelse, gjennem Tanke, netop fordi Erkjendelsen er Væsenets Udtryk, fordi
den er Realisationsform for Væsensforholdet, i sand Betydning bliver et
Livsforhold, at det fører til \emph{den} umiddelbare Enhed, der kan tænkes paa
Aandslivets Gebet.

Hos \emph{Martensen} træder der en Bestræbelse frem for at bringe \emph{hele}
Menneskelivet ind i det religiøse Forhold: \emph{Troen} bestemmes som
\emph{Enhed af religiøs Følelse, Erkjendelse og Villie}. Men ifølge den
underordnede og beherskede Stilling, han lader Erkjendelsen indtage i Forhold
til Troen; ifølge den Opfattelse, han gjør gjældende af \emph{Villien} som
„det afsluttende Moment i den religiøse Bevidsthed“, bliver denne Stræben ikke
realiseret, og Erkjendelsesforholdet trænges tilbage. Idet det religiøse
Gudsforhold adskilles fra det spekulative og det æsthetiske, der bestemmes som
Forhold „paa anden Haand“ medens det religiøse bliver et
\emph{Existentialforhold}, en „Bevidsthed om Gudsforholdet, der er Eet med den
personlige Leven og Væren i dette Forhold“, saa er i denne Adskillelsens Form
det i Erkjendelsesforholdet Væsentlige trængt tilside; den eensidige Sondring
mellem Upersonligt og Personligt er
% --- p. 35 ---
gjort gjældende ogsaa for den sande, concrete Erkjendelse, i hvilken Selvet
søger en Tilfredsstillelse for sin dybeste, i dets Personligheds Væsen
begrundede Trang. Og denne eensidige Tilbagetrængen af Erkjendelsesmomentet
røber sig ogsaa i de nærmere Udviklinger, i hvilke det religiøse Forhold
charakteriseres som et \emph{helligt} Forhold, et \emph{Samvittigheds}forhold;
thi gjennem disse Bestemmelser skal det udtrykkes, at Religionsforholdet er et
Forhold til den „\emph{personlige}“ Gud, i \emph{den} Opfattelse af Begrebet,
der hæver det over det „Logiske“, og gjør \emph{Villien} til Organ for
Gudsforholdet.

c) Med Hensyn til de \emph{relative} Forhold følger af Bestemmelsen af det
religiøse Forhold som det \emph{absolute}, at dette Forhold ikke kan falde
udenfor hine, men maa omfatte dem, gjennemtrænge dem og gjøre dem absolute ved
Enheden med sig. De relative Forhold maae i denne Enhed --- en Enhed, der er
betinget ved, at det Ethiske bliver et Altomfattende --- faae en høiere
Betydning, saa at de kunne blive adæquate Udtryk for det Religiøse, der har
sin concrete Udfoldelse i dem.
% unlike a) on p. 31 and b) on p. 32, only „r e l a t i v e“ is letterspaced
% here; „Med Hensyn til de“ is set tight. As printed.

Denne Bestemmelse faaer sin skarpere Belysning ved en Sammenligning med S.\
\emph{Kierkegaards} Opfattelse af Forholdet mellem det absolute religiøse
τέλος og de relative. Idet han søger Udtrykket for den Religiøsitetens Form,
der endnu ikke er det paradox Religiøses, bliver for ham den religiøse
Existentses \emph{første} Udtryk \emph{den absolute Retning (Respect) mod det
absolute} τέλος, \emph{handlende udtrykt i Existentsens Omdannelse}. I
Forholdet til det absolute τέλος skulle \emph{alle} Endelighedens Momenter
eengang for alle handlende være nedsatte til hvad der maa opgives i Forhold
til den evige Salighed; kun under Betingelse af denne undtagelsesløse
Opgivelse skal den Existerende i Virkeligheden forholde sig til det absolute
τέλος, ville det absolut. Som det absolute τέλος for den Villende, der absolut
vil stræbe, skal det være saa abstract, at det æsthetisk seet, er den
fattigste Forestilling; den Villende skal end ikke ville have Noget at vide om
dette τέλος, Andet end at det er til; thi
% between „Udtryk“ and „den absolute Retning“ there is a small round speck
% below the baseline; tesseract reads it as a full stop, ABBYY reads nothing.
% It is dirt, not a printed sort: no punctuation transcribed.
% --- p. 36 ---
\setcounter{footnote}{0}%
saasnart han faaer Noget at vide om det, skal han allerede begynde at sinkes i
Farten. Det absolute τέλος skal da ikke i Ro kunne besiddes i Existentsen;
hele Tiden, hele Existentsen bliver Forventningens Tid. Mediationen skal ingen
Plads faae her, fordi den Existerende er forhindret fra at faae et saadant
Fodfæste udenfor Existentsen, at han fra det kan mediere hvad der desuden ved
at være i Vorden unddrager sig Færdigheden. I Forhold til det absolute τέλος
skal det, at der medieres, netop betyde, at det absolute τέλος nedsættes til
et relativt τέλος. Det skal ikke være sandt, at det absolute τέλος bliver
concret i det relative; thi Resignationens absolute Distinction skal hvert
Øieblik sikkre det mod al Fraterniseren. Det absolute Forhold skal kunne
fordre Forsagelsen af \emph{alle} de relative. Men derfor skal dog den, der
forholder sig til det absolute τέλος, meget godt kunne være i de relative, men
det vil netop være for forsagende at indøve det absolute Forhold: det Forhold,
der hviler paa den absolute Distinction, Forholdet til det, der kun vindes ved
absolut at vove Alt\footnote{\emph{Kierkegaard}: Afsluttende Efterskrift, S.\
295 ff.}.
% between „Distinction“ and „skal“ a round speck sits below the baseline
% (tesseract reads a comma there, ABBYY a comma too). Dirt, not a printed
% sort: no punctuation transcribed.

I den saaledes nærmere udviklede første Bestemmelse for det religiøse Forhold
ere allerede hos Kierkegaard de Consequentser antydede, der afgjørende komme
frem hvor tilsidst, indenfor Christendommen, det Religiøses Fordring løftes
saa høit, at den fortærende forbrænder al Endelighed og med den alle de
ethiske Forhold, der, naar de sees fra denne Høide, gaae ind under
Endelighedens Bestemmelse. Det religiøse Forhold faaer her den yderste
Abstraction: det kan blive den Magt, der med Uendelighedens Kraft drager den
Existerende ud af Endeligheden, bestandig giver ham Retningen mod det Evige;
men det kan ikke vinde Skikkelse i den Existerende, ikke forme hans concrete
Virkelighed; thi, som jeg et andet Sted\footnote{Problemet om Tro og Viden,
S.\ 221.} har paavist: det bliver kun en \emph{Fordring}, at Individet paa een
Gang skal kunne forholde sig absolut til det absolute τέλος og relativt til de
relative,
% printed marks on p. 36: *) and **), below a short rule. The name in the
% first note is letterspaced Fraktur: „K i e r k e g a a r d“.
% --- p. 37 ---
i hvilket hiint ikke maa blive concret. Men ved denne Abstraction, der
uundgaaeligt hæfter sig til det religiøse Forhold, faaer det netop
\emph{Relativitetens} Charakteer og taber Muligheden til at udtrykke det
udelte menneskelige Væsens høieste Energie, og i det Endelige, som det skal
kunne fordre opgivet, er ogsaa det, der, som et Væsentligt, er et Uopgiveligt:
de ethiske Forhold som Totalitet. Et Correctiv mod hiin Opfattelse ligger
allerede deri, at hos Kierkegaard det ethiske Grundforhold kun bliver til som
et Villiesforhold til \emph{det Evige}, og at Bestemmelsen af det Evige paa
det Ethiskes Gebet ikke er til at sondre fra Bestemmelsen af det paa det Trin
af det Religiøse, vi have charakteriseret. Fastholdes dette, saa kan det
Ethiske ikke i den Forstand opfattes som Gjennemgangsstadium, at der skulde
kunne gives et ethisk Liv, der endnu ikke var religiøst, men det Ethiske
bliver kun ethisk ved fra første Færd at være religiøst, og netop derfor
staaer det Religiøse i et saadant Forhold til det Ethiske, at det altid maa
finde sit Udtryk i dette, at der altsaa ikke kan tænkes en ved det Religiøse
berettiget Opgivelse af det Ethiskes Totalitet, af alle de ethiske Forhold;
men at Forsagelsen af et ethisk Forhold bestandig maa være i Kraft af et
høiere og mere omfattende Ethisk, at altsaa den religiøse Forsagelse bestandig
maa bæres af ethisk Handling.

d) Den første, væsentlige Bestemmelse for det religiøse Forhold blev altsaa
den, at det var et \emph{absolut, uendeligt} Forhold. Og i Menneskeaandens
Begreb, i dens Væsensbestemmelse, ligger \emph{Nødvendigheden} af et saadant
Forhold for Menneskeaanden. Uendeligheden er dens Væsen, dens sande
Livselement; kun i et Uendelighedsforhold kan den derfor komme til Hvile, og
det saaledes, at den ikke blot efter en enkelt af sine Grundfunctioner, men
efter deres Totalitet, staaer i et absolut Forhold ɔ: i et
Uendelighedsforhold. Ogsaa naar Aandsindholdet er omsat i \emph{Følelsens}
Form, kan Forholdets Absoluthed bevares, idet Følelsen, i totaliseret Enhed
med Erkjendelse og Villie, bliver den Aandens intensive Umiddelbarhedsform,
hvis Indhold har den uudfoldede Muligheds uendelige Rigdom, og som i hvert
% --- p. 38 ---
Moment kan blive bevægelig og gjennemklares i Erkjendelse og Villie.

e) Hiin første Bestemmelse af det religiøse Forhold, som det absolute Forhold
til det Absolute, faaer nu en \emph{concretere} Form idet
\emph{Existents}begrændsningen, \emph{Endeligheds}bestemtheden ved det
Menneskevæsen, der træder ind i det religiøse Forhold, accentueres. Paa dette
Punkt blive dog kun de almindelige Bestemmelser, til hvilke dette Synspunkt
fører, at angive; de nærmere Consequentser, nemlig med Hensyn til
Tidsforholdets Betydning for det ethisk-religiøse Liv, for Skylden og dens
Ophævelse, ville blive paaviste i sidste Afsnit.

De nye Bestemmelser, der træde til, og som udtrykke, at det absolute Forhold
til det Absolute er \emph{den efter sit Væsen uendelige}, men i
\emph{Endelighedsbestemthed existerende}, Menneskeaands absolute Forhold til
det Absolute, er \emph{Troens} og \emph{Hengivelsens} Bestemmelser, af hvilke,
naar de sondres, den første nærmest viser hen til Erkjendelsen, den sidste til
Villien, men saaledes, at Troens Erkjendelse efter sin Sandhed maa bestemmes
som værende i indre, harmonisk Enhed med Villien, Hengivelsens Villie som
staaende i samme Enheds-Forhold til Erkjendelsen.

f) \emph{Troens} Bestemmelse træder frem idet den endelige Bevidsthed ud fra
sin Endelighedsbestemthed skal gribe det Uendelige og Evige i Vishedens Form
og fastholde det som det, der gjennemtrænger og bestemmer Endeligheden. Troen
er et Erkjendelsens Postulat, en Slutning, incommensurabel med sine
Præmisser: den overspringer, i Forholdet til det evigt Nærværende, den
begrundende Erkjendelses Mellemled og griber det som dette Nærværende; den
forvandler det Tilkommendes Uvished til et Nærværendes Vished. I dette Forhold
viser Troen sig paa een Gang at staae i Forhold til Væsenets
Endelighedsbestemthed og til dets Uendelighedsbestemthed. Troen er kun som den
Vished, der fødes af Uvisheden, der ligesom vikler sig ud af denne, og gjennem
Uvisheden, der i hvert Øieblik skal beseires, røber Endeligheden sig. Men det,
at Visheden kan vindes gjennem
% modern pencil marginalia on p. 38 (a tick above „idet“ and a short vertical
% stroke in the right margin); not transcribed.
% --- p. 39 ---
\setcounter{footnote}{0}%
Uvisheden, at denne \emph{kan} beseires, viser Aandens Uendelighed. Det, der
berettiger til Troens Postulat; det, der udfylder Afstanden mellem dens
Præmisser og dens Conclusion, det er \emph{Menneskeaandens
Væsensbestemmelse}, der med det \emph{Aprioriskes} Nødvendighed viser hen paa
\emph{det Absolute} som det, uden hvilket Bevidstheden vilde tabe Tilholdet i
sig selv. Det er Væsenets Uendelighed, der gjennem Endelighedens Omslutning
tvinger Bevidstheden til at anerkjende det Uendelige og
Evige\footnote{Umiddelbart ligger heri endnu kun en saadan Anerkjendelse, der
kan forenes med en Sindelagets Bortvendthed, som naar der om Dæmonerne siges
i det N.\ T., at de troe paa Gud og frygte. Men naar det Uendelige skal
anerkjendes som det i Sandhed Uendelige, saa er denne Anerkjendelse betinget
ved en Tro, der anerkjender det Uendelige og Evige i dets Gyldighed som
Menneskelivets Ideal, ved en Tro, der er i Enhed med Hengivelsen. I hiint
Forhold, i hvilket Troen staaer ved Siden af en Sindelagets Bortvendthed,
bliver det Guddommelige, gjennem \emph{Modsætningen} til Selvets abstracte
Uendelighed, kun et abstract Uendeligt, i Virkeligheden altsaa kun et
Endeligt, med den negative Uendeligheds Form. Den Uendelighedens Hvile, som
Aanden søger, fordi Uendeligheden er dens Væsensbestemmelse, findes derfor kun
hvor Tro og Hengivelse ere i Enhed ɔ: hvor Troen er efter sit Væsens
Fuldstændighed.}. Og gjennem \emph{Væsenets} Uendelighed virker
\emph{Principets} Uendelighed, ved hvilken hiin har sin Væren. Hiint Troens
Postulat har derfor paa een Gang menneskelig og guddommelig Oprindelse, er et
Enhedsudtryk for den menneskelige og den guddommelige Aands Virken.
% printed mark: *) --- the only note on p. 39, below a short rule.

Som et saadant Erkjendelsens Postulat er Troen virksom idet i det ethiske Valg
Menneskets Leven bestemmes ved det Uendelige og Evige, og som et saadant er
den virksom idet Livsuendeligheden fastholdes i den enkelte Handling, der
udfylder Menneskelivet, og, som \emph{ethisk}, stadfæster hiint primitive
Valgs Constitueren af Aandslivets Charakteer.

I det ethiske Valg og den Valget bekræftende ethiske Handling aabenbarer sig
den \emph{sande Form for Enheden af Erkjendelse og Villie} i Troesforholdet.
Valget og den Handling, der hviler paa Valget, er paa een Gang
% --- p. 40 ---
\setcounter{footnote}{0}%
Udtrykket for den Erkjendelsens Energie, med hvilken den gjennembryder den
Uendeligheden tilhyllende Endelighed og med sit Væsens Magt griber
Uendeligheden, og for Villiens Beslutning. Den gjennem Erkjendelsen over sig
selv klare Villie kræver som det Maal, i hvilket den, som Væsensvillen, kan
komme til Hvile, det Uendelige, der er Væsensuendelighedens Betingelse; den
kræver det, forat kunne være sig selv\footnote{Villiens mulige Modstand mod
det Guddommelige vil blive belyst under Afsnit III.}. Og ligesom Villien
fordrer det Absolute forat hvile i det, saaledes finder ogsaa Følelsen først
Hvile i det; og den bevarer dets Absoluthed, idet den i hvert Moment af sin
intensive Concentration ligesom skjuler Erkjendelsens Erindring og
Spændkraften til Villie i sig.
% printed mark: *) --- the only note on p. 40, below a short rule.
% modern pencil marginalia on p. 40: a diagonal stroke across „tilhyllende“
% and underlinings under „Formen af Resignation“ and „af en Selvhævdelse“.
% Not transcribed.

g) \emph{Hengivelsens} Bestemmelse træder frem idet den enkelte endelige
Villie skal gjennembryde Endelighedsbestemtheden, i hvilken den umiddelbart
har sit Liv, forat ville uendeligt, og i denne Villen, der er en Villen af sit
Væsens sande Uendelighed, gaae ind i det absolute Forhold, sætte sig i Enhed
med det i Sandhed Uendelige, lade sig bestemme af det som Led i den Totalitet,
det Aandssamfund, som det behersker.

Hengivelsen er først \emph{Valgets} Hengivelse, den vælgende Villies
afgjørende principielle Opgivelse af Endelighedsbestemtheden forat kunne gribe
Uendeligheden. Men det constituerende Valgs Øieblik skal udfoldes til at
omfatte et \emph{helt} Liv, i hvis enkelte Handlinger det til blivende
Beslutning formede Valg gjentager sig, gjennem Stadfæstelsen af Villiens
første Opgiven af det endelige Villiesindhold og af sin Selviskhed. I Valget
og i Gjentagelsen af Valget faaer Hengivelsen paa den ene Side Formen af
Resignation, af Afdøen fra Umiddelbarheden (S.\ Kierkegaard), af
Inderliggjørelse, af „Selvtilintetgjørelse“ i Betydning af, det Selviskes
Ophævelse, paa den anden Side, idet Maalet er en Selvets Gjenvinden af sig
selv efter sin Evighedsbestemthed og i concret Virkelighed, af en
Selvhævdelse efter Selvets
% --- p. 41 ---
\setcounter{footnote}{0}%
Sandhed, men da en saadan Gjenvinden, en saadan Selvhævden, kun er mulig
gjennem Enheden med det Guddommelige, der er Væsenets Princip, faaer
Hengivelsen Formen af en Hengivelse til det Guddommeliges Virken («\textit{Deum
pati}»), en Hengivelse, der ikke blot er den endelige Villies Luttring, men,
som dens Uendeliggjørelse, Udtrykket for dens høieste Energie: den bliver en
Liden, der er Eet med den høieste Handlen\footnote{Hos S.\ \emph{Kierkegaard}
sættes den religiøse Lidelse og Selvtilintetgjørelse i Modsætning til den
ethiske Selvhævdelse, i hvilken den Existerende, ved at „angre igjennem“ skal
gjenfinde sig selv. Denne Modsætning hænger sammen med Opfattelsen af det
ethiske Stadium som et Gjennemgangsstadium. Opfattes det Ethiske i dets
ovenfor paaviste oprindelige Enhed med det Religiøse, bliver allerede den
ethiske Selvhævdelse en Selvhævdelse ved Enheden med Selvets absolute Princip,
og Modsætningen mellem det religiøse og det ethiske Forhold hæves.}.
% printed mark: *) --- the only note on p. 41, below a short rule.
% antiqua on p. 41: «Deum pati», set upright antiqua inside French guillemets
% (the guillemets are antiqua too). Latin --> \textit{} per the house
% convention; the guillemets are kept as printed. The name „Kierkegaard“ in
% the note is Fraktur and letterspaced, as everywhere in this book.

Sin \emph{definitive} Bestemmelse faaer Hengivelsen, idet Subjectet, der
hengiver sig til det Guddommelige, bestemmer sit Liv i
Endelighedsbestemtheden som \emph{Skyld}, idet Hengivelsen altsaa bliver Eet
med Bevidsthedens angrende Anerkjenden af Skylden, og med dens Stræben efter
dens Forsoning i Gud ved en guddommelig Virken. Denne Bestemmelse bliver her
kun berørt; dens nærmere Udvikling vil blive given under Afsnittet III,
ligeledes Paavisningen af dens Forhold til det Christeliges Bestemmelser.

I Hengivelsen, i hvilken Villiens Virken primitivt træder frem, er, som det
fremgaaer af de tidligere Udviklinger, denne Villie betinget ved og i Enhed
med Erkjendelsen. Og ligesaafuldt som Hengivelsen under det Uendelige maa
bestemmes som Villiens Maal, maa den ogsaa bestemmes som Erkjendelsens og
Følelsens; thi naar den erkjendende Aand, gjennem Opgivelsen af sin
Endelighed, finder Hvile i Hengivelsen til det Uendelige, finder Aanden som
følende Hvile i den Hengivelse, der er den absolute Afhængighedsfølelses.

h) I det Foregaaende er det allerede antydet, hvorledes
% --- p. 42 ---
Troens og Hengivelsens Moment i det religiøse Forhold staae i nøieste
Forbindelse med hinanden. Deres Begreber forudsætte hinanden og de kunne, i
deres fulde og sande Betydning, ikke tænkes uden hinanden. Troen forudsætter,
som det blev viist (S.\ 39 Anm.), Hengivelsen forat kunne være Tro paa det i
Sandhed Uendelige; den forudsætter den forat kunne voxe sin Væxt, og den bærer
Hengivelsen i sig idet den som en Griben og Fastholden af det Evige, der i
Troens Vished løftes ud af Tvivlens Uvished, er en Aandens Sighengiven til den
evige Sandhed, i hvilken Menneskelivet som Heelhed finder Hvile. Og ligesom
Troen forudsætter Hengivelsen, saaledes forudsætter Hengivelsen Troen netop
forat kunne være Hengivelse til det \emph{sande} Guddommelige ɔ: for i Sandhed
at kunne være Hengivelse; og den Hengivelsens Handling, ved hvilken Selvet
vinder sit Liv gjennem Opgivelsen af sit Livs Selviskhed, er kun ved Troen paa
det Evige, til hvilket Livet saaledes kan hengives, at det kan gjenvindes i sin
Sandhed og Evighed.
% letterspacing on p. 42: „s a n d e“ only --- „Guddommelige“ that follows it is
% tight. Verified at 260 dpi.

i) Sammenfatte vi de givne Bestemmelser af det religiøse Forholds Væsen, saa
bliver altsaa Religionen, efter sin ideale Betydning, det absolute Forhold til
det Absolute som det med os i Væsensenhedsforhold staaende sande Uendelige. Den
bliver et Forhold, der som begrundet i selve Menneskevæsenet omfatter hele
Menneskelivet, alle Aandens Grundvirksomheder. Den bliver, idet Mennesket sees
som Enhed af Uendelighed og Endelighed, et Troens og Hengivelsens Forhold; et
Forhold, gjennem hvilket den i Endeligheden hildede, i Skylden bundne,
Menneskeaand finder Forsoning i sig og med sit Ophav.
% the enumerators i) on p. 42 and a)--e) on pp. 44--51 are antiqua italic in the
% print, exactly like a)--h) on pp. 31--41; set plain here, following the
% convention already established in the batch pp. 30--41 and in the Indhold.

\begin{center}\rule{2cm}{0.4pt}\end{center}
% A short centred rule stands here on p. 42, between „med sit Ophav.“ and the
% paragraph on the positive religions; a thought-break in the running text, not
% a head. Same treatment as the rules on pp. 19 and 28. A SINGLE hairline.

Forholdet mellem Religionen i den angivne Betydning og de \emph{positive}
Religioner er i Grundtrækkene antydet i de „foreløbige Begrebsbestemmelser“ i
det tidligere Skrift; det vil fremgaae bestemtere af de efterfølgende
Udviklinger.

Den ovenfor udviklede Opfattelses Forhold til S.\ \emph{Kierkegaards}
Bestemmelser af det Religiøses Væsen er
% --- p. 43 ---
\setcounter{footnote}{0}%
allerede tildeels berørt. Det viser sig ved første Øiekast, at den har
væsentlige Tilknytningspunkter i hans Bestemmelse af det Religiøsitetens Trin,
der endnu ikke er det paradox Religiøses; men en indgribende Differents traadte
frem ved Forskjellen i Opfattelsen af det Religiøses Forhold til det Ethiske,
af det absolute Forholds Forhold til de relative. Fra \emph{Kierkegaards}
Bestemmelse af det Religiøse som det \emph{paradox} Religiøse maa hiin
Opfattelse selvfølgelig fjerne sig endnu mere afgjørende. Hvad der her skærper
Differentsen, er Modsætningen mellem Immanentsen og den abstracte
Transscendents, mellem Anskuelsen af Menneske\emph{væsenets} Uforstyrrelighed
og Menneskets \emph{Væsens}forvandling ved Synden, mellem Erkjendelsens
Hævdelse i det religiøse Forhold og dens Tilintetgjørrelse ved Paradoxet, det
absolut Ubegribelige, mellem Anerkjendelsen af det Reales væsentlige Gyldighed
og væsentlige Enhed med det Ideale i Tilværelsens Princip, i Tilværelsens
Heelhed og i Menneskevæsenet\footnote{Jfr. det følgende Afsnit (III).}, og
Christendommens absolut spiritualistiske Opfattelse, der hos Kierkegaard bliver
aabenbar i de negative Consequentser, med hvilke han slutter. Med disse
Modsætninger vil ogsaa en væsentlig Differents i Opfattelsen af Forsoningen
vise sig at staae i Forbindelse.
% PRINTER'S DEFECT, p. 43: „Tilintetgjørrelse“, with a doubled r, for
% „Tilintetgjørelse“. Both r's are unambiguous at 480 dpi and both OCR
% witnesses read them; the same page sets „Uforstyrrelighed“ correctly, and
% „Virkeliggjørelsen“ is correct on p. 41. Transcribed as printed. No effect on
% the quote balance.
% letterspacing across the line break, p. 43: the compositor spaces only the
% SECOND half of „Menneske-væsenets“ (the break falls after „Menneske=“) and
% only the FIRST half of „Væsens-forvandling“. Transcribed as printed, not
% regularised. „Kierkegaard“ in the last sentence is tight, though the same
% name is letterspaced twice above --- also as printed.
% printed mark: *) --- the only note on p. 43, below a short rule. The Roman
% numeral „(III)“ in it is antiqua, as all figures and Roman numerals in this
% book are; not italicised here.

\section*{2. Erkjendelsens Forhold til det Religiøse efter dets ideale Bestemthed.}
% Bold centred head in two lines, standing mid-page on p. 43 above a short
% centred rule (rule not transcribed). Wording verified against the image; it
% agrees with the Indhold exactly.

Forat Erkjendelsens Forhold til det Religiøse i dette Begrebs ideale Betydning
kan vise sig som et Enhedsforhold, maa det godtgjøres, at Erkjendelsen ved sit
Forhold til Menneskevæsenet og ved den Form, i hvilken den udtrykker Væsenet,
kan blive Mediet for et Uendelighedsforhold, et absolut Forhold, at den kan
have absolut Betydning for Menneskeaanden, og at den i sig indeslutter det
Troens og
% --- p. 44 ---
Hengivelsens Moment, der viste sig at være et væsentligt Kjendemærke for det
Religiøse.

a) Betingelsen for, at Erkjendelsen for Menneskeaanden kan blive Formen for et
absolut Forhold, er den, at den kan bevare Gjenstandens sande Uendelighed i
Forholdet. At den maa kunne bevare Uendelighedsbestemtheden, ligger forsaavidt
allerede i Erkjendelsens Form, som denne er bestemt ved det Almene, og dette
Almene, hvis Væren som Alment for den menneskelige Bevidsthed er betinget ved
Væsenets indre Uendelighed, staaer i Væsenssammenhæng med det Uendelige. Men
forat gjennem denne Erkjendelsens Form det Uendelige kan bevares som det i
Sandhed Uendelige, maa Erkjendelsens Almindelighedsform være i concret Enhed
med Begrebets andre Grundformer. Forat den kan være dette, forudsættes, at
Erkjendelsen i Forholdet til det Absolute ikke staaer i abstract Modsætning til
hvad der ligesaafuldt er en Grundfunction af Aanden, at den ikke bliver en
ensidig, fra det Supplerende og Totaliserende isoleret Form for Aandens
Virksomhed. Som saadan maatte den f.\ Ex. uundgaaelig opfattes naar den ved en
principiel absolut Uensartethed sondredes fra Villien; thi i denne Sondring
vilde Villien tilegne sig Enkelthedsbestemmelsen som sin Grundform og lade
Erkjendelsens Almindelighedsform blive den abstracte Almindeligheds, medens
Villiens Enkelthedsform, i sin Sondring fra det Almene, netop kun blev den
umiddelbare Naturenkelthedsform. I denne abstracte Modsætning vilde ingen af
Aandsfunctionerne blive opfattet efter sin Sandhed, ɔ: som Yttring af den
udelte, totale Menneskeaand, og ingen af dem vilde komme til i Sandhed at
forholde sig til det Uendelige, fordi det Uendelige ikke i sin sande
Uendelighed vilde kunne være for nogen af dem.

Med Hensyn til Forholdet mellem de primitive Aandsfunctioners Formbestemtheder
er det Fornødne udviklet i forrige Afsnit og i det tidligere Skrift, S.\ 134
ff. Der er paaviist, i hvilket Forhold Erkjendelsens Almindelighedsbestemthed
staaer til det Enkelte og Concrete, og hvorledes Villiens
Enkelthedsbestemthed, netop fordi \emph{Villiens} Enkelthed er
„Aandsenkelt\-%
% --- p. 45 ---
hed“, udfolder det Almenes Bestemthed af sig. Af hvad Natur Villiens
Enkelthedsbestemthed er, viser sig allerede gjennem den Kjendsgjerning, at
Villien i Organismen har sin Grændse mod den naturbestemte Realitets Enkelthed
og Umiddelbarhed, og Villiens psychologiske Forhold til den Driftbestemthed, af
hvilken den udfolder sig under Aandslivets Udvikling, og i hvilken den ubevidst
omsættes som indvirkende paa Selvets Naturside, dets Organisme, f.\ Ex. naar en
Villiesact fremkalder en momentan eller en mechanisk fortsat Bevægelse, viser
hen til en Villiens Grændse som Villie, ved hvilken dens Enkelthedsform
bestemmes. Vilde man imidlertid forat hævde den absolute Uensartethed lade
Villiens Almene stille sig som et absolut uensartet fra Erkjendelsens Almene,
saa vilde en saadan Fordobling af det Almene i to absolut uensartede og dog ved
samme Fælledsbestemmelse betegnede Arter af Alment umiddelbart indeslutte en
Selvmodsigelse, en Meningsløshed. Af hvad der i det Foregaaende er udviklet
fremgaaer det altsaa, at for den Erkjendelse, der, som væsentligt Udtryk for
den i sig uadsplittede totale Menneskeaand, ikke kan fastholdes i abstract
Sondring fra hvad der ligesaafuldt er et primitivt Udtryk for Aanden, og derfor
ikke ensidig bestemmes ved en enkelt Begrebets Grundform, kan det Absolute være
Gjenstand efter sin Sandhed, som det i Sandhed Uendelige, --- og dets
Uendelighed kan bevares under den relative Abstraction; det kan altsaa være
tilstede for Erkjendelsen i den Form, i hvilken det er Gjenstand for Religionen
efter dens ideale Betydning, i hvilken det er Gjenstand for Tro som det, der
fordres af Menneskeaanden, som det Aandens Maal, i hvilket Væsenets Krav finder
Tilfredsstillelse, i hvilket Aanden først finder Hvile.

b) At Erkjendelsen kan have \emph{absolut} Betydning for Menneskeaanden, altsaa
en saadan Betydning, som den maa kunne have forat kunne udtrykke det religiøse
Forhold, viser sig først derigjennem, at Erkjendelsen, ifølge sit Forhold til
Menneskevæsenet, maa være \emph{Øiemed i og for sig}. Den bliver Øiemed idet
den er en Virkeliggjørelse af en Grundbestemmelse i Menneskeaanden, og
nærmere: af den efter
% letterspacing p. 45: „a b s o l u t“ (not „Betydning“), and the whole phrase
% „Ø i e m e d  i  o g  f o r  s i g“; the second „Øiemed“, two lines below,
% is tight. Both verified at 260 dpi.
% --- p. 46 ---
Begrebet principale Bestemmelse i Aanden; men dette vil sige: den er Øiemed som
Væsenets Virkeliggjørelse i en af dets Grundformer. Men idet den er Form for
det universelle Væsens Virkeliggjørelse, er den Formen for den individuelle
Aands Udvidelse til, ideelt at optage den uendelige Tilværelses Heelhed; den er
den Energie, der ideelt sammenknytter den individuelle Aand med denne
Tilværelsens Heelhed. Derved viser Erkjendelsens absolute Betydning sig fra et
nyt Synspunkt; thi idet den virkeliggjør Aandens Universalitetsbestemthed, og
ideelt sammenknytter Aanden med Tilværelsens Heelhed, sammenknytter den Aanden
med Tilværelsens Princip, det Absolute, der, som saadant, er Erkjendelsens
Princip og bevægende Magt, det, paa hvis Erkjendelse al sand Erkjendelse
hviler. Erkjendelsen faaer altsaa, fra dette Synspunkt, absolut Betydning fordi
den virkeliggjør det i Menneskevæsenet anlagte Aandsforhold til det Absolute,
og idet denne Virkeliggjøren af Menneskevæsenets dybeste Forhold sees som kun
værende ved det iboende Absolutes Virken, saa viser Erkjendelsens absolute
Betydning sig atter i en ny Belysning, der gjør det endnu klarere, hvorledes
Aanden kan forholde sig religiøst i Erkjendelsen og til Erkjendelsen.

Fra de fremhævede Synspunkter er Erkjendelsen endnu directe kun seet efter sin
Betydning i \emph{theoretisk} Henseende, gjennem hvilken den allerede viser sig
som selvgyldigt Øiemed for det Væsen, der i den virkeliggjør sig efter en
Grundbestemmelse i sig, og som netop kun, idet det forholder sig til den som
Øiemed, kan være nærværende i den med Aandens hele Magt, med den Kjærlighed til
Erkjendelsen, ved hvilken alene denne bliver frembringende og skabende,
besjælet af Ideen. Men idet nu Menneskevæsenet, der i Erkjendelsen
virkeliggjør sig og virkeliggjør sin ideelle Enhed med Tilværelsen og dens
Princip, som et Væsen, der omfatter en \emph{Totalitet} af Functioner, træder
ind som activt \emph{Led} i en Tilværelsens Heelhed og bliver sig bevidst som
saadant, saa bliver Erkjendelsens absolute Betydning ogsaa at see fra
\emph{Handlingens} Synspunkt. Hvorledes den fra
% p. 46: modern pencil marginalia --- a short vertical stroke and underline
% after „al sand Erkjendelse hviler.“, and a stroke before „Men idet nu“. Not
% transcribed. The speck above the line after „Ideen.“, which tesseract reads
% as an opening quote, is offset dirt and show-through; not transcribed.
% letterspacing p. 46: „t h e o r e t i s k“ (not „Henseende“), „T o t a l i t e t“,
% „L e d“ (which stands at the head of its line, „activt“ before it tight), and
% „H a n d l i n g e n s“ (not „Synspunkt“). „sige:“ in line 1 is merely a
% loosely justified line-end, not spaced; likewise „de primitive
% Aandsfunctioners“ on p. 44.
% --- p. 47 ---
dette Synspunkt bliver at opfatte, fremgaaer af det i forrige Afsnit Udviklede.
Kun ved Erkjendelsen bliver den Handling mulig, i hvilken Individet virker
efter sin absolute Bestemmelse, der er Væsenets Virkeliggjørelse ved Villiens
Frihed. Erkjendelsen viste sig at blive Forudsætning for en saadan Handling,
deels ved sin klarende, frigjørende, luttrende Indvirkning paa Villien, deels
som bestemmende den sande Villies Indhold. Den bliver fra dette Synspunkt, hvor
den viser sig som Function i en organisk Livshelhed, medens den virker som
Betingelse, som Middel for det Ethiske, tillige at see som ethisk
\emph{Øiemed} og som virkende til sin egen Realisation idet den virker til
Realisationen af den Livsenhed, i hvilken den har sin Sandhed. Og idet
Erkjendelsen udøver den samme luttrende og klarende Indvirkning, som den
udøvede paa Villien, paa de andre Aandens Virksomheder, paa Følelseslivet og
Phantasielivet, saa viser den sig ogsaa i Forholdet til dem i sin Betydning som
realiserende det totale Menneskevæsens indre Enhed og Harmonie, og dette
Synspunkt fører saaledes tilbage til det første, fra hvilket Erkjendelsen saaes
som en Realisation af Menneskevæsenet.

c) At der i Erkjendelsen er et \emph{Troens} Moment, ikke i den Jacobiske
Betydning af en umiddelbar Vished om Væren, men i en Betydning, i hvilken det
religiøse Troesforholds Grundbestemmelse er bevaret, viser sig paa
\emph{ethvert} Trin af Erkjendelsens Udvikling. Dette Troesmoment er skjult
tilstede i de lavere Erkjendelsesformer overalt hvor et Alment udsiges paa
Grundlag af det dermed incongruente Enkelte. Hvor en almen Lov udsiger i Kraft
af det ufuldstændigt foreliggende Enkelte, hvor Inductionens ufuldstændige
Præmis altsaa gives Gyldighed som en fuldstændig, der bringes det Almene frem,
idet det Empiriske suppleres i Kraft af et Apriorisk, altsaa i Kraft af en
\emph{Væsenets} Fordring paa Enhed og Lovmæssighed. Og hvad der saaledes
allerede viser sig, hvor i en lavere Sphære et \emph{relativt Alment} træder
frem, det viser sig endnu tydeligere, hvor det i \emph{høieste Forstand Almene},
den \emph{totale Tilværelses}
% letterspacing p. 47: „Ø i e m e d“ (not „ethisk“); „T r o e n s“;
% „e t h v e r t“, with „Trin“ on the next line tight --- as printed;
% „V æ s e n e t s“, with the following „For=“ tight (checked at 480 dpi, the
% F-o-r are close-set); „r e l a t i v t  A l m e n t“, spaced through the line
% break; and „h ø i e s t e  F o r s t a n d  A l m e n e“ and
% „t o t a l e  T i l v æ r e l s e s“ as two runs, the „den“ between them tight.
% The run stops at the page break: „universelle Princip“ on p. 48 is tight.
% --- p. 48 ---
universelle Princip, gribes af Erkjendelsen ud fra de ufuldstændige
Forudsætninger. Thi idet der gaaes ud over Erfaringens uundgaaelige
Ufuldstændighed, saa gribes den ikke af Erfaringen udtømte Tilværelses
Enhedsprincip ikke som det, hvis Vished Erfaringen betrygger, men som det, der
naaes ved et nødvendigt Postulat ud fra den individuelle Bevidsthed, som uden
det vilde tabe Tilholdet i sig selv, blive usikker paa sin egen Enhed, forvirres
i sig selv. Dette Postulat af Principet, der som Tilværelsens Princip er
Menneskeaandens Livsprincip, træder ogsaa frem i den Erkjendelse, der er
Forudsætningen for det primitive ethiske Valg, ved hvilket Selvet constituerer
sig som ethisk Selv, og ligesom altsaa den Erkjendelse, paa hvilken dette Valg
hviler, indeslutter Troens Moment, saaledes er dette ogsaa tilstede i den
Erkjendelse, for hvilken den enkelte ethiske Handlings ufuldstændige Præmisser
faae Gyldighed som fuldstændige, som en Handlingens \textit{ratio sufficiens}.

Deri, at den er en af selve Væsenets Kilde med Nødvendighed fremgaaende
Fordring, og en Fordring, der indeslutter den umiddelbare Anerkjendelse af det
Fordrede, stemmer den Tro, som Erkjendelsen indeslutter, overens med den Tro,
der er bleven bestemt som det Religiøses Særkjende. Men den Tro, der er tilstede
i selve Erkjendelsen, maa have det \emph{Rationelles} Charakteer. Idet den
fordrer det Absolute, hviler dens Fordring paa det tænkende menneskelige
Væsens Bevidsthed om sig selv, den kommer ikke i Strid med denne, men
stadfæster den. Troen som rationel Tro bliver en Tro paa det, der er høiere end
Individet, idet det er Principet, der begrunder de individuelle Aander og hele
Tilværelsen; men som, idet det bestemmes som høiere, med det Samme bestemmes
som væsensligt, netop fordi det er Menneskevæsenets og Tilværelsens
\emph{Princip}. Og da det Væsen, fra hvilket den religiøse Troesfordring
udgaaer, er en udelelig Enhed, saa ligger deri uimodsigeligt, at denne Fordring,
for i Sandhed at udtrykke Væsenets Krav, ikke kan staae i Strid med nogen sand
Væsensyttring, altsaa ikke med Tanken, med Erkjendelsen. Troen i det i sand For\-%
% antiqua on p. 48: „ratio sufficiens“, upright antiqua in the print, verified
% at 480 dpi; set \textit{} per the house convention.
% letterspacing p. 48: „R a t i o n e l l e s“ (not „Charakteer“) and
% „P r i n c i p“ in „Menneskevæsenets og Tilværelsens P r i n c i p.“
% p. 48 carries modern pencil strokes before „Thi“ (line 2) and after
% „Bevidsthed,“ (line 7). Not transcribed.
% --- p. 49 ---
stand Religiøses Betydning kan derfor ikke være en Tro, der principielt staaer
Viden imod; den vilde ved denne principielle Disharmonie vise sig, ikke at have
sin Kilde i det Væsentlige. Denne Tro maa, ifølge Væsenets Enhed, staae i
principiel Harmonie med Viden; men ifølge Menneskeaandens
Endelighedsbestemthed maa Troen, som \emph{foregribende Resultatet}, faae en
\emph{blivende} Betydning som et for Menneskeaanden Nødvendigt. Den er Formen
for den uundværlige Vishedens Anticipation; den er den væsensbegrundede
Fordring, der har selve Væsenets Nødvendighed; den er derfor ikke som en usikker
Viden, men som en Slutning uden Præmisser, der dog har uimodsigelig Vished,
fordi Præmisserne ere skjulte i Væsenet selv. Men netop derfor, netop fordi
Troen henter sin Vished fra denne Væsenets Kilde, kan den ikke komme i Strid med
hvad der klart og gjennemsigtigt udspringer af samme Kilde. Hvor derfor det, der
sættes som en Troens Bestemmelse, uimodsigeligt strider mod Viden, og dette gjør
det, naar det ikke blot strider mod denne eller hiin Vidensbestemmelse, men mod
det, der gjør Viden til Viden, mod \emph{Videns Princip, Begrundelsens og
Nødvendighedens} Princip, der kan det ikke have Sandhed. Men i en saadan Strid
mod selve Videns Princip staaer f.\ Ex. Troen paa Miraklet, naar Miraklet tages
i den Betydning, i hvilken den positive Religion tager det, ikke i den udviskede
og forfalskede Betydning, i hvilken det tages af en halvtheologisk omvendt
Speculation.

Naar „Troen“ paa de positive Religioners Trin optræder med Kravet paa ubetinget
Gyldighed, saa maa den enten støtte dette Krav paa sin guddommelige Oprindelse
eller paa sit Udspring fra hvad der er det Inderste og Væsentligste i
Menneskeaanden. Støtter den Kravet paa sit Udspring fra det transscendente
Guddommelige --- og kun naar det Guddommelige opfattes saaledes, falde de to
Synspunkter udenfor hinanden --- saa maa Troen med sin overnaturlige Oprindelse
og sit overnaturlige Indhold sættes som hævet over al Viden og med det Samme som
staaende i
% letterspacing p. 49: „f o r e g r i b e n d e  R e s u l t a t e t“ (spaced
% through the line break, „foregri=/bende“), „b l i v e n d e“ (with
% „Betydning“ tight), and the run „V i d e n s  P r i n c i p ,
% B e g r u n d e l s e n s  o g  N ø d v e n d i g h e d e n s“ --- the second
% „Princip,“ that closes the phrase is tight. All verified at 260 dpi.
% The signature numeral „4“ at the foot of p. 49 is not transcribed.
% --- p. 50 ---
\setcounter{footnote}{0}%
Modsætning til al Viden, men idet Viden sees som et Udtryk for den
Menneskenatur, der har sit Udspring fra Gud, saa bliver der ikke blot en Splid
mellem det Guddommeliges modsatte Aabenbaringsformer: i Troen og i Tanken, men
der bliver en uophævet og uophævelig Splid i Menneskeaanden, der skal fornægte
en Væsensyttring forat anerkjende en Virkning, udgaaende fra den samme Kilde,
fra hvilken Menneskevæsenet er udsprunget. Vises der saa, forat hævde
Nødvendigheden af „Aabenbaringens“ Anerkjendelse, hen til Personlighedens
Fordring, til en Norm i Personligheden, saa er herved det sidste af de ovenfor
angivne Synspunkter sat i Stedet for det første. Men naar saa Troen hviler paa
Væsenets uafviselige Krav, saa gjælder atter det ovenfor Sagte, at dette Væsens
Enhed maa bevare Harmonien mellem Væsensyttringerne, og der fordres da en
Harmonie mellem Troens Indhold og det i Viden Væsentlige. Saasandt der har været
og bestandig vil kunne være falske Former af det positivt Religiøse, Former af
Lære og Cultus, der paa det Dybeste krænke den sædelige
Bevidsthed\footnote{Det er her tilstrækkeligt at minde om de forasiatiske
Religioner.}, saasandt Samvittigheden kan være hildet og uklar i Individet,
saasandt fordres en Norm, der ikke kan søges i et med Viden absolut Uensartet,
men maa søges i Videns negative Correctiv, i det Videns Veto, der maa beholde
sin Gyldighed, fordi, naar Viden fornægtes, det Regulativ tabes, ved hvilket der
kan sondres mellem den sande, i Væsenet begrundede Tro og det, der kun er et
Udtryk for endelig Attraa og Begjær, for et Uvæsentligt og Particulært i
Menneskenaturen. Videns Forhold til Troessætningen kan være et saadant, at Viden
maa standse i et \textit{non liquet}, men Videns Forhold kan aldrig blive det,
at den maa godkjende, hvad der uomtvistelig modsiger dens \emph{Princip},
tilintetgjør den \emph{som} Viden.

d) I Erkjendelsen, saaledes som den er efter sin Sandhed, er endelig det
\emph{Hengivelsens} Moment tilstede, der betegnede det i Sandhed religiøse
Forhold. I al sand Erkjendelse er der en Tankens Luttring, en Opgiven af
individuel
% antiqua on p. 50: „non liquet“, upright antiqua, verified at 480 dpi. It is
% not in any OCR-derived list --- neither witness marks it --- and it is a
% second Latin tag on this stretch besides „ratio sufficiens“ (p. 48).
% letterspacing p. 50: „P r i n c i p“ and „s o m“ in „modsiger dens
% P r i n c i p, tilintetgjør den s o m Viden.“, and „H e n g i v e l s e n s“.
% spacing.py's two „RUN“s on this page („Dybeste krænke sædelige“, „Udtryk
% for“) are FALSE POSITIVES: both lines are merely loosely justified, and the
% type is tight at 480 dpi. Nothing on those lines is emphasised.
% printed mark: *) --- the only note on p. 50, below a short rule.
% --- p. 51 ---
Hildethed, Menen og Fordom, med eet Ord: af Tankens Selviskhed; der er i den en
Selvhengivelse til Sandhedens absolute Magt. Og denne Hengivelse, der er
tilstede i al sand Erkjendelse, den viser sig klarest i det Forhold, i hvilket
Erkjendelsen, ifølge Væsenets Troesfordring, hengiver sig til det Væsenets
Trang tilfredsstillende Uendelige. I denne Hengiven, der er betinget ved en
Væsenets indre Nødvendighed, er Erkjendelsen i Enhed med Villien, og idet den
selv i Hengivelsen uendeliggjøres, klarer den den Villie, hvis Valg bestemmes
ved den samme Væsenets Magt.

e) Resultatet af det Udviklede bliver altsaa, at Erkjendelsen i sin Sandhed
staaer i et væsentligt Enhedsforhold til det Religiøse i dets ideale
Bestemthed, til den religiøse Tro, der er Udtrykket for Menneskevæsenets
uendelige Fordring. Under Menneskeaandens Endelighedsvilkaar bliver der vel
ligesom et Mellemrum mellem Tro og Viden, men et Mellemrum, der maa kunne
gjennemløbes af Viden. En principiel, uophævelig Splid bliver en Umulighed,
fordi den Tro, der ophævede Videns Væsensbestemmelse, netop derved vilde vise
sig som ikke stemmende med selve Troens Væsen. Under sine historiske Formers
Vexlen har Troen en Udvikling, og dens Væsenssammenhæng med Viden viser sig
deri, at denne Troens Udvikling ikke er løsreven fra Videns, ikke er uden den
frigjørende Viden.
% p. 51: a short centred rule stands between this paragraph and the sub-head
% below --- a rule under/above a display head, and so not transcribed, as at
% p. 30. Modern pencil marginalia on p. 51: „maa kunne gjennemløbes“ is
% underlined and a „?“ stands in the right margin. Not transcribed.

\section*{3. Villiens og Handlingens Forhold til det Religiøse efter dets ideale Bestemthed.}
% Bold centred head in two lines, p. 51. Wording verified against the image at
% 260 dpi; it agrees with the Indhold exactly.

Ved Bestemmelsen af Villiens og Handlingens Forhold til det Religiøse ville de
samme Synspunkter være at gjøre gjældende, som i det Foregaaende bleve
fremhævede med Hensyn til Erkjendelsen.

a) Der bliver da først at udhæve, at Muligheden til et absolut Forhold til det
Absolute gjennem Villien er tilstede, idet Villien, i sin Harmonie med
Erkjendelsen, og i sin Gjennemklarethed ved det Almene, har Evnen til at be\-%
% The signature numeral „4*“ at the foot of p. 51 is not transcribed.
% No letterspacing anywhere on p. 51 --- checked line by line at 260 dpi; the
% one word spacing.py flagged („være“, 1.19) sits in a loosely justified line.
% --- p. 52 ---
vare sin Gjenstands sande Uendelighed --- en Evne, der maatte benægtes, naar man
lod en absolut Uensartethed adskille Villien fra Erkjendelsen, eller naar der,
paa Grundlag af en uklar Forestilling om Væsenets „Defecthed“, statueredes et
\emph{blivende} Irrationelt i Villien. Det Absolute er tilstede som absolut i
Villiens Form som \emph{det Gode}, og Menneskets Villen af det Gode er i Enhed
med Erkjendelsen af det. Idet denne Villen af det Absolute som det Gode
udpræger sig i Handling, faaer denne Handling, hvad der i det Efterfølgende
nærmere vil blive udviklet, Skikkelse af den ved Villiens ubetingede
Underkastelse under det Absolutes Fordring beherskede Virkeliggjørelse af det
Gode i det menneskelige Væsens Form, saaledes at denne Virkeliggjøren ved den
ubetingede Underkastelse bliver Udtrykket for Realisationen af Væsensenheden
med det Guddommelige (ὁμοιωθῆναι θεῷ), og derfor Formen for et absolut
Villies-Forhold til det Absolute.
% GREEK on p. 52: „ὁμοιωθῆναι θεῷ“, antiqua italic, broken „ὁμοιω-/θῆναι“
% across the line; rough breathing on the initial omicron, no acute;
% circumflex on the eta; omega with iota subscript and circumflex in „θεῷ“.
% The printed sort for theta is the script ϑ; transcribed θ, as elsewhere in
% this edition. Read off the image at 480 dpi --- neither OCR witness reads
% any of it, and p. 52 is not in the recon pass's Greek list.
% The parenthesis „( hvad der i det Efterfølgende nærmere vil blive udviklet)“
% is printed with plain commas, not parentheses: tesseract's brackets are the
% modern pencil strokes that stand over „hvad“ and after „udviklet“ in this
% copy. Not transcribed. Further pencil marks on p. 52: a dash in the left
% margin beside „lutes Fordring“ and an underline under „det“ two words later.

b) Den sande Villens og den sande Handlens \emph{absolute Betydning} for
Menneskeaanden viser sig umiddelbart, idet Villien som \emph{Væsens}villen
betinger Menneskevæsenets Anlægs Udfoldelse i real Virkelighed, og constituerer
det sande Forhold til det Absolute. Som ovenfor paavist gaaer dette dobbelte
Synspunkt sammen til en uadskillelig Enhed.

Naar der i Udviklingen af det Ethiskes Tilbliven i Individet brugtes det
Udtryk, at Selvet i det ethiske Valg vælger og vil \emph{sig absolut}, saa
vistes der hen til, at denne Vælgen kun er mulig derved, at Selvet staaer i
\emph{Væsensenhedsforhold til det Absolute}, thi kun ifølge et saadant Forhold
kan det vælges absolut. Hint Udtryk maa tages som Udgangspunkt, fordi, naar vi
bestemte Valgets Gjenstand som det Godes absolute Idee, Valget af denne som det,
der skal virkeliggjøres i \emph{menneskelig} Handling, vilde blive liig med
Valget af det \emph{menneskeligt Gode}, af det \emph{idealt Menneskelige}. I
dette er den absolute Idee for Mennesket i \emph{dets Væsens Bestemthed}, og
idet det virkeliggjøres ved den frie Villie, gives der Ideen Virkelighed.
Udenfor Bestemmelsen af det idealt Menneskelige
% letterspacing p. 52: „b l i v e n d e“ (not „Irrationelt“); „d e t  G o d e“
% at its first occurrence only --- the two that follow are tight;
% „a b s o - / l u t e  B e t y d n i n g“, spaced through the line break;
% „V æ s e n s“ with „villen“ tight, i.e. only the first half of
% „Væsensvillen“; „s i g  a b s o l u t“; „V æ s e n s = / e n h e d s f o r h o l d
%  t i l  d e t  A b s o l u t e“; „m e n n e s k e l i g“ with „Handling“ tight;
% „m e n n e s k e l i g t  G o d e“; „i d e a l t  M e n n e s k e l i g e“; and
% „d e t s  V æ s e n s  B e s t e m t h e d“. The last line of the page repeats
% „det idealt Menneskelige“ TIGHT --- as printed, not regularised.
% --- p. 53 ---
som den menneskelige Handlings Opgave kunne vi ikke komme. Vi kunne kun realisere
det Absolute, der er \emph{vort} absolute Væsen, kun virkeliggjøre det Ideal, der
er \emph{vort} Ideal. Fra Bestemmelsen af den ethiske Opgave som
Virkeliggjørelsen af det Godes Idee som saadan føres vi strax tilbage til den
angivne Grændse derved, at det Virkeliggjørende maa være congruent med det, der
skal virkeliggjøres; det Absolute kan da kun \emph{positivt} realiseres af os,
forsaavidt det udtrykker sig i vort Væsens Bestemthed. Men denne Begrændsning af
det positive Udtryk for Ideens Virkeliggjørelse ved Villien, udelukker ikke, men
forudsætter netop en anden negativ Virkeliggjørelsesform, i hvilken Ideen gjør sig
gjældende \emph{som saadan}. See vi det menneskelige Individ indordnet som Led i
en principbestemt organiseret Totalitet, saa følger det vel af dets Bestemthed som
enkelt organiske Led, at det positivt kun kan udtrykke Principet gjennem
Individualitetens eget Væsens virksomme Udfoldelse, men i Negationens
(Selvhengivelsens) Form udtrykker det Heelheden som Heelhed, Principet som
Princip, ifølge Heelhedens og Principets Tilstedeværen i det i Mulighedens
abstracte Form. Gjennem den Bestemmelse for det ethiske Valg og den ethiske
Opgave, der var vort Udgangspunkt, vises der altsaa hen til Forholdet til det
Absolute som til det, ved hvilket Selvet i sin Idealitet kan blive sig selv
absolut Opgave, og som faaer Nærværelse i det ethisk udfoldede Selv; og den
ethiske Villens og Handlens absolute Betydning viser sig altsaa fra et dobbelt
Synspunkt.

Til det samme Resultat føres vi, naar vi gaae ud fra det andet Udtryk for det
ethiske Valg, at Selvet i det vælger \emph{sig efter sin Væsensbestemthed}. Valget
af Selvet efter dets Væsen er nemlig et Valg af Væsenet \emph{efter dets totale
Concretion}, altsaa efter dets Bestemthed \emph{som organisk Led i en Tilværelsens
Totalitet}. Idet Selvet erkjender sig i denne Bestemthed, er det fra Begyndelsen
af i Forhold til det Princip, der er den totale Tilværelses styrende Livsprincip.
Jegets Selvhævden i det ethiske Valg er kun ved Principets Gjennemvirken i Jeget,
% letterspacing p. 53, all confirmed at ~480 dpi: „v o r t“ twice;
% „p o s i t i v t“; „s o m  s a a d a n.“ (spaced through the line break);
% „s i g  e f t e r  s i n  V æ s e n s b e s t e m t h e d“;
% „e f t e r  d e t s  t o t a l e  C o n c r e t i o n“; and „s o m“ at the end
% of a line followed by „o r g a n i s k  L e d  i  e n  T i l v æ r e l s e n s
%  T o t a l i t e t.“ --- one emphasis carried across the break, not two.
% „Men“ (spacing.py 1.29) is a loosely justified line, not letterspacing.
% „Livsprincip.“ is a period, not the comma tesseract reads.
% A modern pencil dot stands in the left margin beside line 2. Not transcribed.
% --- p. 54 ---
og idet Jeget hævder sig selv, men hævder sig i sin Inordnethed i Totaliteten, er
dets Selvhævdelse kun som Selvhengivelse til Principet. Og hvad der gjælder om
Jegets Forhold i det ethiske Valg, gjælder om dets Forhold i den enkelte ethiske
Handling, for hvilken dette Valg danner Grundlaget. Idet det Ethiske vinder
Concretion og den ethiske Villen udformer sig som den ved den sande Erkjendelse
klarede Villen af Væsenet efter dets Fuldstændighed, er paa ethvert Punkt
Forholdet sat til den Totalitet, i hvilken Selvet som ethisk frit indordner sig,
og til Totalitetens Princip. Det Forhold, der i det ethiske Valg sattes, idet det
ethisk Aprioriske først fik Gyldighed for Villien, bevares idet dette Aprioriske
vinder Fylde i den concrete Handling. Det dobbelte Synspunkt for det Ethiskes
absolute Betydning samles saaledes atter til Enhed.

Naar det Forhold til det Absolute, der er constitueret i det ethiske Valg, sættes
som bevaret overalt i den ethiske Handling, saa fører dette til en, fra et andet
Synspunkt allerede ovenfor dragen, væsentlig Consequents for Opfattelsen af de
\emph{særlige} ethiske Forhold. Deraf, at der i disse Forhold i deres Enkelthed
paa ethvert Punkt er et Forhold til det gjennemvirkende Absolute, følger, at de
særlige ethiske Forhold, der i deres Isolation vilde være at bestemme som blot
relative Forhold, blive til Forhold med absolut Betydning, til Udtryk for det
absolute Forhold. Det absolute Forhold sees som omfattende alle de ethiske
Forhold, som iboende i dem, og som gjennem dem udfoldende et concret Indhold og
vindende Virkelighedsskikkelse. Derved faae de særlige ethiske Forhold en høiere
Betydning, og idet Subjectet opfatter dem i denne, forholder det sig religiøst til
det Ethiske, der nu i sin Articulation faaer absolut Gyldighed, bliver absolut
Øiemed for det.

Kun et andet Udtryk for det Samme er det, naar der sagdes, at det Ethiske aldrig
som Heelhed kan opgives for et religiøst τέλος, der som et Absolut skulde hæve sig
over alt det Ethiske. Det kan det ikke, fordi det er Væsensudtryk, fordi det ved
Forholdet til det Absolute har absolut
% GREEK on p. 54: „τέλος“, antiqua italic, acute on the epsilon, final sigma.
% Read off the image at ~460 dpi; tesseract renders it „7é4o0ç“ and the ABBYY
% layer drops it altogether. p. 54 is not in the recon pass's Greek list.
% letterspacing p. 54: only „s æ r l i g e“, at its FIRST occurrence
% („de særlige ethiske Forhold“); the second occurrence eight lines below is
% tight, and „ethiske Forhold“ is tight in both. As printed, not regularised.
% --- p. 55 ---
Betydning, og fordi det staaer i uopløselig Sammenhæng med det Religiøse. I
Virkelighedscollisionen mellem det Ethiskes forskjellige Former kan det Ethiske
fra en Sphære vige for det Ethiske fra en anden, der fremtræder med det høiere
Almenes Ret, men idet der resigneres paa et enkelt af de ethiske Forhold, der, som
Led i den ethiske Totalitet, efter deres Væsen ere absolute Forhold, er
Resignationen paa det enkelte Forhold i Kraft af den ethiske Totalitet, der i
Resignationen \emph{integrerer sig selv}. Resignationen paa det enkelte ethiske
Forhold bliver altsaa i Kraft af selve det Ethiske; der bliver intet absolut
Forhold udenfor dette, men det Absolute bliver allestedsnærværende i det Ethiske.

c) I det saaledes opfattede Ethiske, Udtrykket for den sande, væsentlige Villen og
Handlen, er Troesmomentet allerede viist at være tilstede. Det er viist, at det
ikke blot træder frem i det ethiske Valgs første Griben af det Absolute, men
overalt i Valgets Stadfæstelse i Gjentagelsen, i enhver ethisk Beslutning, i
hvilken Individet bestemmer sig til Handling, trods Erkjendelsens
Ufuldstændighed og Virkeliggjørelsens, Udfaldets, Usikkerhed, og det træder
ligeledes frem, idet Individet, under den ufuldstændige Realisation af den ethiske
Fordring, i det iboende, gjennemvirkende Princip, med hvilket det er i Enhed
gjennem Kjærlighedens Selvhengivelse, anerkjender det ethisk Fuldstændiggjørende.

d) Hengivelsens Moment i det Ethiske er allerede berørt, idet den negative Form
for Virkeliggjørelsen af det absolute Gode angaves, og idet det vistes, at Selvets
Hævdelse af sig, som en Hævdelse af Selvet efter dets Væsens
\emph{fuldstændige} Bestemthed, efter hvilken det er Moment i en Heelhed,
behersket af Principet, blev Et med en Kjærlighedens og Resignationens Hengivelse
til det Princip, ved hvilket Selvets Væsen er, og i hvilket det først gjenfinder
sig efter sit Væsens Sandhed. Dette gjælder med Hensyn til det første ethiske Valg
og til hver enkelt ethisk Handling. Idet nu den ethiske Selvhævdelse er
Selvhævdelsen gjennem Selvhengivelse, er \emph{Kjærligheden} det Ethiskes Princip,
og idet Selvet sees i Forhold til det Absolute og til den aan\-%
% letterspacing p. 55: „i n t e g r e r e r  s i g  s e l v.“;
% „f u l d s t æ n d i g e“ (with „Bestemthed“ tight);
% „K j æ r l i g h e d e n“ --- verified letter by letter at ~490 dpi.
% The section letters „c)“ and „d)“ are antiqua, as „a)“/„b)“ on pp. 42--52.
% Modern pencil marginalia on p. 55, none transcribed: a dot between „i“ and
% „det“ in „allestedsnærværende i det Ethiske“ (tesseract turns it into a
% guillemet); a caret after „sit Væsens Sandhed.“ (tesseract turns it into a
% closing quotation mark --- there is NO quotation here); and a vertical stroke
% in the left margin beside „Idet nu den ethiske Selvhævdelse“.
% --- p. 56 ---
delige Livsheelhed, i hvilken det er indordnet, bliver Kjærlighedens Gjenstand det
Absolute, Principet, og Mennesket \emph{som} Menneske (Næsten). I
Kjærlighedsforholdet til Mennesket er Kjærlighedsforholdet til det Absolute
aabenbaret: i det ethisk religiøse Forhold til Menneskeheden er der et Forhold til
det Absolute, der er dens Livsgrund. I denne Kjærlighed til det Absolute og til
Mennesket \emph{som} Menneske er Selvkjærligheden forklaret og Menneskets inderste
Væsen aabenbaret. „Kun det Menneske er Noget, der elsker Noget, og hvad et
Menneske er, det skal kjendes paa hvad han elsker“. (Feuerbach).
% The Feuerbach quotation, checked at ~490 dpi: low „ opening, high “ closing,
% and the full stop stands AFTER the closing mark, not before it. „(Feuerbach).“
% runs on as the last, short line of the same paragraph; it is not set off.
% Quote balance for this batch: 1 „ and 1 “ --- balanced.

e) Saaledes bliver Resultatet med Hensyn til Villiens Forhold til det Religiøse
det samme som med Hensyn til Erkjendelsens. Forholdet mellem Villien i dens
Sandhed: den i ethisk Handling virksomme Væsensvillen, og det Religiøse efter dets
ideale Bestemthed, er et Enhedsforhold, og maa være det, fordi en Væsenets
Grundfunction ikke kan modsige, hvad der udtrykker Menneskevæsenets dybeste Krav,
dets uendelige Fordring.
% A short centred rule stands between this paragraph and the sub-head below ---
% a rule under/above a display head, and so not transcribed, as at pp. 30 and 51.

\section*{4. Menneskeaandens Forsoning i det paa den sande Erkjendelse hvilende ethisk-humane Liv.}
% Bold centred head in two lines, p. 56, verified word for word at ~490 dpi; it
% agrees with the Indhold. The line break falls after „sande“; „ethisk=humane“
% is the ordinary Fraktur hyphen, set here as „-“.

Naar det menneskelige Liv er udfoldet i harmonisk Enhed af den sande Erkjendelse
og den sande Handlen; naar Erkjendelsen i dens Sandhed i sig bærer Troen efter
dette Begrebs ideale Betydning; naar den ethiske Handling, i hvilken det
menneskelige Væsen efter dets Concretion finder sin Virkeliggjørelse, som ethisk
Handling har det Religiøses sande Præg, og naar endelig Menneskevæsenet gjennem
Erkjendelse og Handling staaer i Forhold til en Tilværelse og til et
Tilværelsesprincip, fra hvilke al dualistisk Adsplittelse er fjernet, saa ere de
Betingelser tilstede, uden hvilke Aanden ikke kan finde Hvile og Forsoning. Hvad
der umiddelbart fremgik af Menneskeaandens Væsen, var en Stræben efter
Virkeliggjørelsen af dens indre Uendelighed, som en
% letterspacing p. 56: only „s o m“, twice, in „Mennesket s o m Menneske“ ---
% in both places the neighbouring „Mennesket“/„Menneske“ is TIGHT, checked at
% ~490 dpi. The compositor spaced the little word alone; transcribed as printed.
% No footnote and no Greek on p. 56.
% --- p. 57 ---
sand Uendelighed, og denne Virkeliggjørelse kan kun finde Sted paa de Vilkaar, at
Aandens Virksomheder ikke adsplittes, ikke følge uensartede Love, og at Aanden
ikke i Forholdet til Tilværelsen og til Principet hæmmes ved en Dualismes
Endelighedsskranke. Men i den sande Erkjendelse af det Absolute, af
Tilværelsestotaliteten og af Selvet, i den sande ethiske Villen, der ligelig
accentuerer Selvets Selvhævdelse og dets Indordning i Totaliteten, dets Aandsside
og dets Naturside, i det sande religiøse Forhold, der, som det absolute Forhold,
gjennemtrænger de relative og i dem vinder Concretion, er denne concrete
Uendelighed tilstede, og idet Erkjendelse, Sædelighed og Religiøsitet ere
harmonisk forbundne, saa er den Betingelse tilveiebragt for at Aandens
oprindelige Stræben kan komme til Hvile, for at Aandens Forsoning med sig selv i
sin Virkelighed kan finde Sted, der maa mangle, saalænge Aanden er deelt af
modstridende Magter, saalænge Bevidstheden er adsplittet i sig selv, i ethvert
Øieblik sønderreven af Tvivlen om Erkjendelsens Sandhed og om Handlingens
Gyldighed. Først ved hiin harmoniske Virkeliggjørelse af alle Aandslivets Kræfter
bliver ogsaa Menneskelivet i Sandhed et \emph{humant} Liv efter dette Begrebs
fulde Betydning, og Humanitetens Charakteer aabenbarer sig i den
\emph{universelle} Kjærlighed til Menneskeheden, til det Menneskelige \emph{som}
menneskeligt.

Naar der paa dette Punkt af Udviklingen, hvor \emph{Skylden}, Bruddet i det
Ethiske, og Muligheden af Restitutionen trods dette Brud, endnu ikke er omhandlet,
tales om Forsoningen, saa bliver selvfølgelig deels dette Begreb at tage i en
abstractere, mere metaphysisk Betydning, i den Betydning, det maa faae, idet der
gaaes ud fra Uendeligheden som Væsensbestemmelse, og idet den Form søges, under
hvilken den af Menneskevæsenets Grund fremgaaede Uendelighedsstræben kan finde sin
Tilfredsstillelse; deels bliver Forsoningens Virkelighed endnu kun
\emph{hypothetisk}, idet kun de almene Betingelser angives, \emph{uden hvilke den
ikke bliver} tænkelig. Men naar i det Efterfølgende Forsoningen vil blive udviklet
i concretere Bestemthed som \emph{Skyldens For\-%
% letterspacing p. 57: „h u m a n t“ (Liv tight); „u n i v e r s e l l e“;
% „s o m“ alone in „det Menneskelige s o m menneskeligt“, the same idiom as on
% p. 56; „S k y l =/d e n“, spaced through the line break, with „Bruddet“ tight;
% „h y p o t h e t i s k“; „u d e n  h v i l k e  d e n  i k k e  b l i v e r“
% with „tænkelig.“ tight; and „S k y l d e n s  F o r =“ running over onto
% p. 58, where „s o n i n g,“ is spaced too. „For=“ is TIGHT six lines above
% (spacing.py's „For-“, 1.20, is a false positive) --- checked at ~470 dpi.
% A modern pencil bracket stands before „Først ved hiin“; not transcribed.
% --- p. 58 ---
soning}, saa maa denne concretere Bestemmelse faae hiin abstractere til
Forudsætning og Betingelse, den vil ikke ophæve den, men udfylde den.

I den Betydning, i hvilken der paa dette Punkt af Udviklingen tales om
Forsoningen, betegner dette Begreb først \emph{Aandens Forsoning i sig}, som
begrundet i dens indre Harmonie, i alle Aandsfunctioners indre Enhed, i Væsenets
udelte Virken. Denne, paa den sande Erkjendelse hvilende bevidste Harmonie, og den
Forsoning, den indeslutter, adskiller sig væsentlig fra den \emph{græske
Bevidstheds} umiddelbare Skjønhedsharmonie med dens illusoriske Forsoning. I den
græske Aands Skjønhedsharmonie vare Modsætningerne kun paa illusorisk Maade bragte
til Enhed, netop fordi Foreningens Form var Skjønhedens Umiddelbarhed, den
\emph{umiddelbare} Forbindelse af det \emph{oprindeligt Forskjellige}. Saasnart
denne for Anskuelsen værende Harmonie skulde bestemmes af Tanken, skiltes
Modsætningerne ad, og den græske Tankes Bestemthed ved Anskuelsens Form lod de
skilte Modsætninger fæstnes i Sondringen uden at kunne tilveiebringe en høiere
Forsoning. Derfor svæver den græske Bevidsthed mellem den umiddelbare
Skjønhedsenhed og den uforsonede Modsætning, der uundgaaelig afløser hiin, og fordi
hiin Skjønhedsenhed i sin Umiddelbarhed ikke har Aandslivets sande Form, saa savner
Aanden i den den dybere Inderliggjørelse, og netop derfor den indre Uendelighed,
der er det i Sandhed Ethiskes Betingelse; og ligeledes mangler Muligheden til, at
Livets concrete Indhold kan udfolde sig for den ethiske Bevidsthed som en
principbestemt Enhed. Idet den Enhed, vi have bestemt som det i Sandhed humane
Livs Enhed, som en frit erhvervet, for Erkjendelsen gjennemsigtig Enhed, der hviler
paa en anden Bestemmelse af Grundforholdet i det Menneskelige og i Tilværelsen, er
betrygget mod Adsplittelsen, indeslutter den Muligheden til den Forsoning, som den
græske Aand, ifølge sin Bevidsthedsform, maatte savne.

Til Bestemmelsen af Menneskeaandens Forsoning i sig knytter sig saa Bestemmelsen
af dens \emph{Forsoning med sit
% letterspacing p. 58: „s o n i n g,“ (the run-over from p. 57);
% „A a n d e n s  F o r s o n i n g  i  s i g“ with „først“ tight;
% „g r æ s k e  B e v i d s t h e d s“ with „umid=/delbare“ tight;
% „u m i d d e l b a r e“ with „Forbindelse af det“ tight, then
% „o p r i n d e l i g t  F o r s k j e l =/l i g e.“; and
% „F o r s o n i n g  m e d  s i t“ on the last line, running over onto p. 59,
% where „g u d d o m m e l i g e  P r i n c i p;“ is spaced --- one emphasis
% carried across the page turn, so the \emph{} spans the page marker below.
% „Til Bestemmelsen af Menneskeaandens Forsoning i sig“ on the penultimate line
% is TIGHT, despite the loose justification. „maatte savne.“ ends with a period
% (the Fraktur square dot), not the comma tesseract reads.
% --- p. 59 ---
guddommelige Princip}; og denne Bestemmelse af Forsoningen viser sig at staae i
nødvendig Sammenhæng med hiin første, idet først i Aandens Forsoning med sit
iboende Princip dens indre Forsoning kan blive fuldstændig, ligesom, omvendt,
Forsoningen med Principet, der er den udelte Aands Forsoning, kun kan fuldbyrdes
idet Aanden vinder den indre Forsoning. Forsoningen med Principet finder Sted idet
Væsensenheden bliver bevidst og finder sit --- negative og positive --- Udtryk;
den finder Sted idet Menneskets Forhold til det Absolute bliver et absolut
Forhold, gjennem Troens og Hengivelsens Ubetingethed. For den forsonede Bevidsthed
er det Absolute den gjennemvirkende Livsmagt, der, ifølge Væsensenheden, ubetinget
kan beherske den menneskelige Tanke, bestemme den menneskelige Villie, og i hvilken
Selvet i høieste Betydning kan gjenfinde sig selv. Hvorledes denne Gjenfinden
bliver mulig trods Skylden, bliver det det efterfølgende Afsnits Opgave at vise.

Menneskeaandens Forsoning bliver endelig at opfatte som dens \emph{Forsoning med
Tilværelsen og i Tilværelsen}. Den faaer denne Bestemmelse, idet Mennesket bliver
sig bevidst som indordnet i en organisk Tilværelse, hvis Princip er
Menneskevæsenets Princip, hvis styrende Love ere Fornuftens ubrødelige Love, og
idet det i Kjærlighed sammenslutter sig med den Livsheelhed, i hvilken det udgjør
et Led.

Naar det nu af det Foregaaende er blevet klart, at det, der er \textit{conditio}
\emph{\textit{sine qua non}} for Aandens sande Forsoning, for dens Forsoning i
Virkeligheden, er tilstede i den Religiøsitetens Form, som vi kunne kalde
\emph{Humanitetens Religiøsitet}, saa vil det ogsaa blive klart, at det ikke kan
findes indenfor de positive Religioner.

Hvad der forhindrer dette, er med eet Ord \emph{Endeligheds}begrændsningen i det
positivt Religiøse. Den Stræben efter det i Sandhed Uendelige, der er Særkjendet
for det i ideel Betydning Religiøse, er ogsaa den bevægende Magt i de historisk
fremtraadte positive Religionsformer; men
% ANTIQUA LATIN on p. 59: the whole of „conditio sine qua non“ is antiqua, and
% within it „s i n e  q u a  n o n“ is ALSO letterspaced while „conditio“ is
% tight --- verified at ~490 dpi. Latin --> \textit{} per the house convention;
% the letterspacing is recorded by the extra \emph{} round the spaced half.
% Other letterspacing p. 59: „g u d d o m m e l i g e  P r i n c i p;“ (the
% run-over from p. 58); „F o r s o n i n g  m e d  T i l v æ r e l s e n  o g
%  i  T i l v æ =/r e l s e n.“; „H u m a n i t e t e n s  R e =/l i g i ø s i t e t“;
% and „E n d e l i g =/h e d s“ with „begrændsningen“ TIGHT --- the same
% half-spaced compound as „Aands“functioner on p. 6. Not regularised.
% A pencil dot stands before „den menneskelige Villie“; not transcribed.
% --- p. 60 ---
\setcounter{footnote}{0}%
denne Stræben faaer, som det paa et andet Sted\footnote{Jfr. Problemet om Tro og
Viden, S.\ 33.} er viist, sin Endelighedsbegrændsning ved den theoretiske og
praktiske Aands eiendommelige Formbestemthed paa disse Religioners Stadium.
% printed mark: *) --- the only note on p. 60, below a short rule. The asterisk
% stands after „Sted“, before „er viist“.

At Menneskeaanden, paa de positive Religioners Trin, maa savne Betingelserne for
Forsoningen, viser sig fra de forskjellige Synspunkter, fra hvilke ovenfor
Forsoningens Begreb opfattedes.

Menneskeaanden kan paa de positive Religioners Trin \emph{ikke} naae den
\emph{indre Forsoning}, Aandens harmoniske Enhed med sig selv. Den kan det ikke,
fordi Aandsfunctionerne ere i Strid med hinanden og i sig selv, og fordi denne
Strid er en saadan, der ikke kan hæves. Erkjendelse og Villie, der begge fastholdes
i en Endelighedsbestemthed, maae blivende divergere; mellem Viden og Tro bliver der
en uløselig Modsigelse, og ligeledes mellem den ved den sande Erkjendelse af
Væsenets Lov bestemte Handlen, og den Handlen, der er Opfyldelse af det
supranaturale Princips vilkaarlige Bud. Dette finder ogsaa sit Udtryk deri, at den
positive Religions Endelighedsbestemthed gjør det religiøse Forhold, Forholdet til
det transscendente Guddommelige, til et \emph{særligt} Forhold, der sondrer sig fra
de andre væsentlige, og som saadanne berettigede, menneskelige Forhold, og negerer
dem, idet det gjør Krav paa Aandens hele Magt. Den Forsoning, Aanden skulde naae i
det exclusive Gudsforhold, vilde, hvis der kunde tænkes en Forsoning hvor et
Væsentligt søges negeret, ialfald ikke kunne blive en Forsoning \emph{indenfor
Virkeligheden}, ikke den i Virkelighedens Concretion existerende Aands Forsoning i
sig selv. Det vilde blive en Forsoning i en Tro, der som forholdende sig til det
Supranaturale taber Viden af Syne, og i en Handlen, i hvilken det Ethiske er
absorberet i et abstract Asketisk eller i et vilkaarligt Statutarisk.

Menneskeaanden kan paa de positive Religioners Trin ikke heller naae den fulde
Forsoning \emph{med sit Princip}. Dette viser sig hvad enten vi see hen til
Menneskeaandens
% letterspacing p. 60, checked letter by letter at ~500 dpi: „i k k e“ IS spaced
% in „Trin i k k e naae den i n d r e  F o r s o n i n g“, but „naae den“ between
% them is TIGHT --- an emphasis the compositor set in two pieces. The parallel
% sentence twelve lines below reads „ikke heller naae den fulde Forsoning“ ALL
% tight, with only „m e d  s i t  P r i n c i p.“ spaced, which is why the first
% „ikke“ is a positive reading and not a justification artefact.
% Also spaced: „s æ r l i g t“ (Forhold tight) and
% „i n d e n f o r  V i r k e l i g h e d e n,“.
% --- p. 61 ---
\setcounter{footnote}{0}%
Vilkaar, eller til det Guddommeliges Bestemthed. See vi, med Hensyn til
Menneskeaandens Vilkaar, bort fra det allerede Berørte, Aandsfunctionernes
Modsætning, der maa forhindre Aanden fra at træde ind i det i Sandhed absolute
Forhold, saa bliver der en særlig Hindring for et saadant i det uigjennemklarede
Driftagtige, der ifølge Erkjendelsens og Villiens Endelighedsbestemthed, ifølge
Villiens og Troens Divergents fra Erkjendelsen, bliver tilbage paa dette
Standpunkt. Thi det Driftagtige som saadant, der er bevaret i Frygtens og Haabets
Affect i den Form, i hvilken den træder frem i Religionerne, skjuler, idet det
søger sin Tilfredsstillelse, ubevidst en Egoismens Rest, der kun kan fjernes ved
Driftens Gjennemklaring ved den sande Erkjendelse, ved Tilveiebringelsen af Aandens
indre, for den selv gjennemsigtige Enhed. --- Sees der hen til Opfattelsen af det
Guddommelige paa de positive Religioners Trin, saa er det viist\footnote{Jfr.
Problemet om Tro og Viden, S.\ 35 f.}, hvorledes Forestillingsformen bringer en
blivende Endelighedsbestemthed ind deri. Og denne Endelighedsbestemthed ved det
transscendente Guddommelige medfører saa uundgaaeligt, at Menneskets Forhold til
det Guddommelige ikke kan blive i Sandhed absolut. Den ubetingede, for al Egoisme
frigjorte, Selvhengivelse maa mangle; thi den kan kun findes hvor Mennesket
forholder sig til det i Sandhed Uendelige. Paa de positive Religioners Trin kan
ikke heller det sande Subjectivitetsforhold finde sin Plads. Subjectet er her i det
religiøse Forhold bestemt ved et for Subjectet Andet, der altsaa, skjøndt selv
bestemt som Subject, som Personlighed, paa Grund af Transscendentsen, af
Væsenssondringen, bliver et for Subjectet Objectivt. Kun naar Forholdet er bestemt
ved Immanentsen, bliver Subjectivitetsforholdet i fuldeste Betydning et
Subjectivitetsforhold.

Menneskeaanden kan endelig, paa de positive Religioners Trin, ikke naae Forsoningen
\emph{med og i Tilværelsen}. Hvad der i denne Henseende er det Forhindrende, er
først og fremmest Forestillingen om \emph{Underet}. Ved denne Fore\-%
% printed mark: *) --- the only note on p. 61, below a short rule. The asterisk
% stands after „viist“ and before the comma.
% letterspacing p. 61: only „m e d  o g  i  T i l v æ r e l s e n.“ (with „ikke
% naae Forsoningen“ tight) and „U n d e r e t“ --- neither of which spacing.py
% flagged; its two flags („mangle“, „med“) are loose justification.
% „Immanentsen“ is an I-word, set with I per the house convention.
% A pencil stroke stands between „Guddommelige“ and „paa“ (tesseract turns it
% into a semicolon); not transcribed.
% --- p. 62 ---
\setcounter{footnote}{0}%
stilling, der, som det paa et andet Sted er viist\footnote{Jfr. Problemet om Tro og
Viden, S.\ 35.}, endeliggjør det Guddommelige, og bringer en Splid ind i Gud,
bringes der nemlig ogsaa en Splid og en Modsigelse ind i Tilværelsen, og denne
Tilværelsens indre Modsigelse: dens ved Underet satte Væsenløshed, der dog atter i
hvert Moment ligesom maa tilbagekaldes\footnote{Sammesteds, S.\ 79 ff.}, om end kun
for paany at sættes, forhindrer Aanden fra at finde Hvile i denne Tilværelse, der
viser sig som i et forvirrende Dobbeltlys for Tanken og Villien. Og vilde den
religiøse Bevidsthed ud af denne Dobbelthed gribe Forestillingen om det Naturliges
Væsenløshed, saa vilde Aanden vises hen til at søge sit Livs Maal gjennem en Flugt
fra den naturlige Tilværelse, gjennem en Fornægtelse af den; men om en Forsoning i
denne, til hvilken Menneskevæsenets Naturside dog bestandig \emph{maa} drage det
tilbage, kunde der ikke blive Tale.
% printed marks on p. 62: *) and **), below a short rule. The first stands after
% „viist“, the second after „tilbagekaldes“.

Fra det Synspunkt, der her er gjort gjældende, bliver det altsaa det
\emph{foreløbige} Resultat, at det, der er \emph{uundværlig} Betingelse for
Forsoningens Mulighed, kun er tilstede i den Religiøsitetens Form, der er det
Humanes høieste Udtryk. Forsaavidt altsaa dette foreløbige Resultat ikke omstødes
ved den \emph{definitive} Bestemmelse af Forsoningens Betingelser, vil denne Form
af Religiøsitet vise sig som Virkeliggjørelsen af Religionens Begreb, som det
adæquate Udtryk fra det i de positive Religioner Tilstræbte.
% ERRATUM, p. 62 line 24: the page prints „fra det i de positive Religioner
% Tilstræbte“; the Trykfeil leaf (printed p. 211) asks for „for det“.
% Transcribed AS PRINTED, per the editorial stance. In this copy a modern hand
% has struck „fra“ through in pencil and written „for.“ in the left margin ---
% a reader applying the author's own erratum. Neither mark is transcribed.

\begin{center}\rule{2cm}{0.4pt}\end{center}
% A short centred rule stands here on p. 62, between the two paragraphs (not
% under a head), as at pp. 19 and 28. A SINGLE hairline.

Ved dette Resultat vises der altsaa hen paa vor sidste og vanskeligste Opgave:
gjennem en Behandling af de centrale ethisk-religiøse Problemer: Problemerne om
Villiens Frihed, om det Ondes Mulighed og Virkelighed og om det Ondes Ophævelse i
Forsoningen, at hævde det angivne Standpunkts Livsbetydning, dets
Virkelighedsbetydning. Ud fra et Videns Udgangspunkt skal der søges naaet til en
Løsning af disse Problemer, der tilfredsstiller ethisk-religiøst, og den givne
Opfattelse af
% letterspacing p. 62: „f o r e l ø b i g e“ (Resultat tight);
% „u u n d =/v æ r l i g“ (Betingelse tight); „d e f i n i t i v e“; and „m a a“
% alone at the end of the line „... Naturside dog bestandig m a a“, with „drage
% det tilbage“ on the next line TIGHT. The last is the one doubtful call on the
% page: it sits at a line end, where a compositor may space to justify, and the
% line is already loosely spaced. Set as printed, on the strength of the visible
% inter-letter gaps in „m a a“ and of the same half-executed pattern at pp. 6,
% 15, 59 and 60; the alternative reading (justification artefact, no emphasis)
% is recorded here rather than silently adopted.
% spacing.py's „gangspunkt“ (1.98) is the tight „Udgangspunkt skal der søges“
% of the last paragraph, not letterspacing.
% --- p. 63 ---
det Menneskeliges Grundforhold saaledes gjøres concretere og fordybes, at den, idet
den opfylder den vidende Aands Krav, opfylder den sædelige og religiøse Aands
dybeste Fordring. Medens Bruddet i Menneskeaandens Liv gjennem Skylden anerkjendes,
skal altsaa det i Sandhed absolute Forhold til det Absolute vise sig at kunne
bevares, og den sande og fulde Forsonings Mulighed hævdes, og hævdes som det, der
kan blive Gjenstand for Erkjendelsen. Forsoningen, hvis Mulighed vistes ud fra
Menneskeaandens almene Væsensbestemthed, skal hævdes ogsaa ud fra dens concreteste
Virkelighedsbestemthed.

\begin{center}\rule{2cm}{0.4pt}\end{center}
% p. 63 is the last page of division II and is a short page: the type block ends
% after eleven lines and a short centred rule closes the division, exactly as at
% p. 29. The rule is printed as a DOUBLE rule (two thin lines) and is genuine
% type, not show-through: it is crisp and black, while the verso of p. 64 shows
% through this leaf only as grey mirrored text. Set here with the same single
% \rule the skeleton uses for the other short rules in this book.
% No letterspacing, no footnote and no Greek anywhere on p. 63.

\chapter*{III.\\ De centrale ethisk-religiøse Problemer, opfattede fra Videns Synspunkt.}
\addcontentsline{toc}{chapter}{III. De centrale ethisk-religiøse Problemer, opfattede fra Videns Synspunkt}
\markboth{III. De centrale Problemer}{III. De centrale Problemer}

% --- p. 64 ---
% p. 64 opens division III. Order on the page: the bold centred division head
% (already in the skeleton), a short centred rule (not transcribed), an
% ornamental two-line initial F (set normally), the introductory paragraph
% below, a second short centred rule (not transcribed), then the bold centred
% sub-head "1. Det Absolutes Begreb."
% Letterspacing on this page is asymmetric across the parallel phrases:
% "Begrebet om det Absolute" is spaced throughout, but in the second member
% only "det menneskelige Væsen" is (the words "og Begrebet om" are set close).
% Verified at 300 dpi. Likewise "Det Absolutes Begreb" has the spacing on
% "Absolutes" alone. Transcribed as printed; not regularised.
Forudsætningen og Grundlaget for Løsningen af disse Problemer fra det angivne
Synspunkt maa blive Udviklingen af de hinanden gjensidig betingende Begreber:
\emph{Begrebet om det Absolute}, og Begrebet om \emph{det menneskelige Væsen}.
Bestemmelsen af det Absolutes Begreb maa have et Fodfæste i Bestemmelsen af
det Menneskelige, som det, der er den endelige Tilværelses høieste Form,
særlig i Bestemmelsen af den Subjectivitetens Form, der frigjort kommer
tilsyne i den menneskelige Selvbevidsthed. Det Menneskelige bliver da fra
dette Synspunkt \textit{principium cognoscendi} for det Guddommelige. Men naar
Menneskevæsenet definitivt skal bestemmes efter sin Eiendommelighed og sine
Grundforhold, saa maa Begrebet om det Absolute som Menneskevæsenets
\emph{Princip} blive Betingelsen, saaledes at det, der som \textit{principium
essendi} har Prioriteten, definitivt ogsaa faaer den som \textit{principium
cognoscendi}, idet det for Tanken, der naaer til det fra Endeligheden, bliver
det, der, som hvilende i sig selv og havende Vished i sig selv, forklarer det,
der har Væren ved det. Det \emph{Absolutes} Begreb bliver altsaa det
fundamentale.
% antiqua Latin on p. 64: principium cognoscendi (twice), principium essendi.
% All three are set in UPRIGHT antiqua at normal letterspacing --- the
% calibration case for the letterspaced antiqua found on p. 68 below.

\section*{1. Det Absolutes Begreb.}

For den concrete Bestemmelse af det Absolutes Begreb bliver den endelige
Tilværelse som Heelhed Grundlaget, idet
% --- p. 65 ---
det Absolute skal opfattes som denne Tilværelses Princip, og saaledes
bestemmes, at det forklarer den af det begrundede Tilværelses Kjendsgjerning,
medens det erkjendes som det, der, som Princip, forklarer sig selv, idet det
er sin egen Grund og hviler i sig selv.

At det Absolute opfattes som Tilværelsens Princip, faaer deels den Betydning,
at det opfattes som \emph{Værens}princip for Tilværelsen ɔ: som Tilværelsens
Udspring, dens Grund, deels den, at det opfattes som Tilværelsens
\emph{Maal}, dens Lov og dens Norm, og som Tilværelsens τέλος er det Absolute
atter deels bestemt som det teleologisk Virkende, deels som det ethisk
Bestemmende (som ethisk Ideal). Men disse forskjellige Synsmaader for
Opfattelsen af Principet blive at see i deres indre Enhed.
% letterspacing broken across the line break: "V æ r e n s -" ends the line
% spaced, "princip" opens the next close-set. As printed.
% Greek on p. 65: τέλος, antiqua italic with accent. Neither OCR witness reads
% it; caught by eye at 300 dpi.

Idet det Absolute skal bestemmes som \emph{Værensprincip}, træder først
Forholdet frem mellem det i Tilværelsen \emph{Ideale og Reale}. Den
Tilværelse, der umiddelbart er given for den menneskelige Bevidsthed og som
afgiver Udfyldningen for Selvets i Tiden første Bevidsthed om sig selv, er den
reale, naturlige Tilværelse. Denne Tilværelse skal forklares ved det Absolute,
og dens Værens Udspring søges i det Absolutes Bestemmelser. Men
Realitetsbestemmelsen viser tilbage, ud over sig selv, til det Ideale.
Fastholdes Realitetsbestemmelsen i sin Isolerethed, og i den Form, der er
nødvendig til Forklaringen af Virkelighedens concrete Mangfoldighed, saa fører
den os til Begrebet om \emph{Atomerne}, som de oprindelige selvstændige
Realia, der i deres Uforanderlighed ere det Tilgrundliggende i al
Tilværelsens Mangfoldighed. Men netop dette Begreb, der umiddelbart
indeslutter en Modsigelse derved, at Atomet som selvstændigt, evigt og
uforanderligt sættes i en i sig afsluttet absolut Væren, medens det i sin
Mangfoldighed og sin endelige Qualitetsbestemthed sættes som et Relativt,
fører, naar de Forhold, i hvilket Atomet indtræder, gjennemtænkes, ud over det
kun Reale til et Ideelt. Forat Tilværelsesphænomenerne kunne forklares, maa
det statueres, at Atomerne indtræde i bestandig vexlende Forbindelser, og at
de gjennem gjensidig
% "Realia" is Fraktur, not antiqua --- checked at 300 dpi; the book sets only
% Latin phrases in antiqua, and this naturalised plural is not one.
% signature "5" at the foot of p. 65 (first page of the gathering): not
% transcribed.
% --- p. 66 ---
Bevægelsesmeddelelse blive Led i den Aarsagernes og Virkningernes uendelige
Kjæde, der omfatter al real Tilværelse. Men i Atomets Natur som afsluttet i
sig og uforanderligt kan ingen Tilskyndelse til saadanne, en Forandring
indesluttende, Forhold tænkes at ligge; thi selv om Atomerne udrustes med
oprindelige Kræfter, saa vil Atomets Uforanderlighed betinge Kraftvirkningens
Uforanderlighed. Den som Kjendsgjerning ubestridelige bestandige Forandring i
den reale Tilværelse maa da uundgaaelig tænkes at have sin Kilde udenfor
Atomerne som enkelte; men dette vilde, da paa dette Standpunkt Tanken om en
supranatural Factor maa bortfalde, være i Totaliteten af det reale Værende, i
Naturen som Heelhed, saaledes at altsaa Forandringen blev til ved en
Heelhedsvirksomhed. Men saasnart Totalitetens Bestemmelse bringes ind, som
Udtryk for en virkelig Enhed, ikke blot for Aggregatets tilsyneladende Enhed,
der ikke vilde kunne forklare Kjendsgjerningen, saa gaaes der ud over det
\emph{kun} Reale, og sættes et Ideelt; thi Totaliteten \emph{som Tota}litet er
\emph{ideelt} bestemt idet den er og virker i sine Dele, sine Momenter,
\emph{som Totalitet}, altsaa gjennem dem forholder sig til sig selv, og i
Selvfordobblingen i dette Forhold, der er Idealitetens Grundform, hævder sig
selv. Forandringens Kjendsgjerning viser altsaa hen til Totaliteten med dens
ideelle Enhed, for hvilke Atomerne --- der som den reale Tilværelses
oprindelige Principer maatte sættes som selvstændige og uforanderlige ---
blive tjenende Momenter. --- Til det samme Resultat føres vi ogsaa gjennem
Betragtningen af Causalitetsforholdet, der, trods Atomernes Selvstændighed og
Uforanderlighed faaer Gyldighed for dem, idet Atomerne, som meddelende
hverandre Bevægelse, opfattes som fungerende Led i en Aarsagernes og
Virkningernes Kjæde, om end det derved satte Forhold bliver et ligesom
overfladisk, der ikke berører Atomets Indre. Aarsagsbegrebet reduceres da her
til Begrebet om den exclusiv reale Aarsag, Tingsaarsagen, den endelige Aarsag.
Den saaledes bestemte Aarsag maa have Tilskyndelsen til Virksomhed udenfra;
dens Virken er Aarsag til et Andet idet den er Virkning af et
% p. 66 carries NO footnote, contrary to the OCR-derived hint list; checked at
% 250 and 300 dpi. It carries modern pencil marginalia (a bracket before
% „Den“, a squiggle under „som Totalitet“): not transcribed.
% "kun Reale": only "kun" is spaced, "Reale" is not. "som Tota-" is spaced at
% the line end and "litet" is not. Both as printed.
% --- p. 67 ---
Andet: den er Aarsag og Virkning med Hensyn til et Forskjelligt. Idet nu det
Paavirkede, ifølge den begrebsnødvendige Sammenhæng mellem Action og
Reaction, sees som tilbagevirkende, fører Aarsagsbegrebet til Begrebet om en
\emph{Vexelvirkning}; men paa det angivne Standpunkt maa
Vexelvirkningsforholdet, ligesom Aarsagsforholdet, blive et overfladisk
Forhold, i hvilket Vexelvirkningens Led, selv om Virksomhedens Substrater
forestilles i indre Enhed med Virksomheden, ikke skulle gaae op: de skulle
bevare en Væren som givne udenfor Vexelvirkningen, skulle kun med en Side af
deres Existents gaae ind i det bestemte Vexelvirkningsforhold. Saasnart
imidlertid en Vexelvirkning er statueret, hvilket nødvendigviis maa skee
saasnart Aarsagens og Virkningens Begreber statueres, om end Virksomhedens
Oprindelse nok saa meget bliver uforklarlig ifølge Forudsætningen --- saa
indeslutter den Consequentser, der uundgaaelig føre ud over de endelige
Forhold og de endelige Tingsaarsager, til en Opfattelse af Aarsagsbegrebet,
svarende til det Totalitetsbegreb, der traadte i Stedet for Atombegrebet som
Begrebet om det i Tilværelsen Primitive. Hvad der kan gaae ind i
Vexelvirkningsforhold til et Andet, viser ved selve dette Forhold hen til en
væsentlig Sammenhæng saavidt Vexelvirkningen naaer. Men naar nu enhver Ting
overhovedet kun er, hvad den er, i sine Relationer, naar en Fastholden af den
som en Ting i og for sig, som relationsløs Ting, gjør den til en indholdsløs
Abstraction, saa maa Tingen efter \emph{hele} sit Væsen gaae ind i en
omfattende activ Relation, i en \emph{total Vexelvirkning}. Og idet
Vexelvirkningen saaledes omfatter den hele reale Tilværelse, vises der gjennem
den omfattende Vexelvirkning hen til en omfattende Væsenssammenhæng, og den
enkelte Ting, den enkelte reale Aarsag, der paa ethvert Punkt paa een Gang er
betingende og betinget, forudsætter sin i Væsenssammenhæng med den staaende
Betingelse ligesom den forudsættes af denne som væsensforbunden Betingelse.
Idet nu ingen af de i Væsenssammenhæng staaende reale Ting, som gjensidig
forudsættende hinanden, bliver \emph{primitiv} Aarsag, viser Cau\-%
% "Consequentser" is the printed form, verified at 300 dpi.
% signature "5*" at the foot of p. 67: not transcribed.
% --- p. 68 ---
salitetsforholdet, gjennem de betingede endelige Tingsaarsagers
Vexelvirkning, tilbage til et alles Væsen Sammenfattende, dem alle Betingende
og deres Vexelvirkning Bestemmende som primitiv Aarsag. I
Underordningsforholdet til denne primitive Aarsag, der virker gjennem dem som
\emph{iboende}, blive de endelige reale Aarsager \emph{secundære} Aarsager,
\emph{Mellemaarsager}, Aarsager i \emph{afledet} Betydning. Som Principet for
de væsensforbundne Reales virksomme Totalitet, for den Totalitet, der, som
værende i sine Momenter, er sig selv, bliver den primitive Aarsag \emph{ideel}
Aarsag; den bliver Aarsag i Forhold til sig selv (\emph{\textit{causa sui}})
og Øiemed for sig selv (\emph{\textit{causa finalis}}). Ved sit Forhold til de
betingede Aarsager hævder den ideelle Aarsag sig som det \emph{Ubetingede} ɔ:
\emph{det sig selv Betingende}, det, der sætter sine egne Betingelser:
Systemet af de reale Aarsager. Omfattede af Totaliteten, som dennes Led, blive
disse i en anden Betydning end i den endelige Vexelvirkning paa een Gang
Aarsag og Virkning: Aarsagernes og Virkningernes Række bøies tilbage i sig
selv i et saadant Kredsløb, i hvilket de enkelte Led, idet de som
Totalitetsmomenter forholde sig til hverandre, ikke i Forhold til et
Forskjelligt, men i Forhold til det Samme, nemlig til Totaliteten, paa een
Gang ere Aarsager og Virkninger, Betingelser for det Ubetingede, men som
Betingelser satte af dette, der gjennem dem forholder sig til sig selv, og er
\emph{Princip for sin egen teleologiske Virkeliggjørelse} idet det er
\emph{Princip for al Virkelighed}.
% antiqua Latin on p. 68 --- causa sui, causa finalis --- is itself
% LETTERSPACED, unlike principium cognoscendi/essendi on p. 64, which is set
% close in the same upright antiqua. Compared side by side at 600 dpi; the
% difference is unmistakable. Encoded \emph{\textit{...}}: \textit{} for the
% antiqua sort (the house convention), \emph{} for the letterspacing.

Naar saaledes den reale Tilværelse, selv idet den ensidigt opfattedes fra
Realitetens Synspunkt, med Nødvendighed viste tilbage til et Ideelt som det
Forenende og Begrundende, saa bliver altsaa dens Forklaring at søge i et
\emph{ideelt} bestemt Absolut; men forat dette kan være \emph{den reale
Tilværelses Grund}, maa i selve det Asolutes Idealitetsbestemmelse
Tilknytningen for det Reale paavises, saaledes at det ideelle Princip sees som
\emph{idealt-realt} Princip. Denne Sammenhæng mellem Idealitetens og
Realitetens Bestemmelse i det Absolute viser sig først
\emph{\textit{in abstracto}}
% PRINTER'S DEFECT, p. 68: "Asolutes" for "Absolutes" --- the b is a dropped
% sort. Verified at 600 dpi: A, long-s, o, l, u, t, e, s, with no b and no
% gap. Transcribed as printed; not in the author's Trykfeil list.
% "in abstracto" is antiqua and letterspaced, like causa sui above.
% --- p. 69 ---
derigjennem, at selve Idealitetsbestemmelsen indeslutter Forholdet \emph{til
det Andet}, gjennem hvilket det Ideale er i Forhold til sig selv. Er altsaa
det Ideales Andet, dets Modsætning, \emph{det Reale}, saa maae vi sætte denne
Bestemmelse, fra hvilken vi \emph{ikke kunne abstrahere}, og som vi saae som
visende tilbage til det Ideelle, som \emph{væsentlig for selve det Absolute},
som fordret ved dette, og som oprindelig Bestemmelse ved det. Men naar vi paa
denne Maade, gjennem \emph{Andethedens} Begreb, der indeslutter
\emph{Modsætningens Nød}vendighed, komme til Bestemmelsen af det Reale, saa er
denne Modsætning mellem Ideelt og Reelt endnu \emph{kun formel}, det Reales
\emph{concrete Bestemthed} er endnu ikke forklaret, og kun naar ogsaa den
forklares af det ideelle Princips Grundbestemmelser, er Principet i
Virkeligheden Forklaringsgrund for den reale Tilværelse. Forklaringen af det
Reales væsentlige Bestemmelser finde vi, idet vi accentuere, at det ideelle
Princip, som saadant, idet det som \textit{principium sui} gjennem sit Andet
hævder sig selv, er \emph{Energie}. Energiens Bestemmelse, der ligger i
Idealitetsbestemmelsen, forudsætter, ligesom denne, \emph{det Andet}, gjennem
Forholdet til hvilket Energien hævder sig som Energie. Denne den ideelle
Energies indre Modsætning, der er begrundet i dens \emph{Væsen}, er det
\emph{Passive}, men dette i Energien begrundede Passives Bestemmelse betegner
ikke den absolute Passivitet, men en Passivitet, der har Activitetens Mulighed
i sig, der er som en slumrende Virksomhed, som en Virksomhed, der, idet den er
gaaet til Ro i Passiviteten, bærer Energien i sig som real Mulighed. En
Passivitet af denne Art fordres af Energien, er Betingelse for dens
Virkelighed, og er kun ved Energien, tænkes kun ved den. Men en saadan
Passivitet, en saadan skjult Energie, tilhyllet under Inertien, er
Grundpræget i det \emph{Materielle}, i det, som Mechanismen bestemmer som
\emph{død Væren, som passiv Materie}. Dødt, passivt i absolut Forstand kan et
Værende ikke være, thi al Væren fordrer, forat \emph{hævde} sig i sin
\emph{Væren, Virken og Virksomhed}. Men Begrebet: \emph{passiv Materie}, kan
faae sin Gyldighed i den Betydning, at Activiteten kun er tilstede som Mulig\-%
% antiqua Latin on p. 69: principium sui, set close, NOT letterspaced. (The
% batch brief placed this phrase on p. 71; it stands here.)
% "Modsætningens Nød-" is spaced at the line end and "vendighed" is not.
% As printed.
% --- p. 70 ---
hed medens dens Virkelighed er betinget ved hvad der falder udenfor det
saaledes bestemte Materielle. Og dettes Begreb falder sammen med det
\emph{Totalitetsløses}, idet Totaliteten, ogsaa som \emph{relativ Totalitet},
ved sit ideelle Forhold til sig selv, ved sin Væren for sig selv, umiddelbart
viser hen til Energiens Bestemmelse, det Passive altsaa kun kan være det, der
mangler Totalitet. Men den i Bestemmelsen af det Totalitetsløse liggende
Mangel af det ideelle Forhold til sig indeslutter atter
\emph{Udvorteshedens} Bestemmelse, en Væren, bunden af dennes Former,
behersket af \emph{Rums- og Tidssondringen}. Som bestemt ved disse Former
existerer det Materielle \emph{som} materielt i de objectivt enkelte
Stofelementers Sondring, og disses Væren som isolerede er ikke blot for den
abstraherende Tanke, der seer bort fra den concrete Tilværelses Betingelser,
og skiller Elementet fra Totaliteten; men den har sin relative Gyldighed i det
Objective, Reale, ifølge selve Grundforholdet i Principet.

Saaledes viser der sig i selve Principets Idealitetsbestemmelse en Tilknytning
for Realitetsbestemmelsen i Tilværelsen. Hvorledes nu denne
Realitetsbestemmelse beherskes af det Ideelle, \emph{idealiseres}, men
\emph{uden at tabe sin Væsentlighed}, og hvorledes \emph{det Reale} som
Bestemmelse \emph{i det Absolute} forholder sig til \emph{det Reale i
Tilværelsen}, vil fremgaae af Udviklinger i det Efterfølgende.

Paa dette Punkt bliver den fremsatte Opfattelse nærmere at belyse gjennem
Sammenligninger. Fra saadanne Opfattelser, der fra Begyndelsen af postulere
Idealitet og Realitet som primitive Bestemmelser i det Absolute, men \emph{kun
fordre} denne Forening eller indskrænke sig til at substituere Realitetens
Begreb for den ved Idealiteten fordrede Andethed, adskiller den angivne
Opfattelse sig derved, at den ud fra Idealitetsbestemmelsen deducerer det
Reales \emph{eiendommelige} Bestemmelser, og at den derfor kan føre til
Erkjendelsen af en Enhed, i hvilken det Reale beherskes af det begrundende
Ideelle, uden at den tvinges til at postulere et modsigende forenende Medium
mellem de, ved Mangelen paa en Begrebssammenhæng, i deres Eiendommelighed op\-%
% "idealiseres," is spaced, the following "men" is not, and "uden at tabe sin
% Væsentlighed," is spaced again. As printed.
% --- p. 71 ---
rindeligt sondrede, og i deres umiddelbare Uensartethed uforenelige
Modsætninger.

Til Forestillingen om en \emph{Skabelse} staaer den angivne Opfattelse i det
Forhold, at den fra denne Forestilling tilegner sig den Sandhed, \emph{at det
Reale er begrundet i det Ideelle, Aandelige og bestemt ved dette;} medens den
seer det Usande i Skabelsesforestillingen deri, at denne ikke lader det Reale
være af væsentlig Betydning for det Ideales Væren, ikke lader det ideelle
Absolute i dets Ubetingethed betinge sig selv gjennem det Reale, at den ikke
anerkjender dette Reales Væsens-Enhed med det Ideale. Den fastholder, at det
Ubetingede kun er ubetinget som det, der gjennem et System af reale
Betingelser betinger sig selv; at det Absolute er \textit{principium sui},
\textit{causa sui}, idet det er Princip for et Andet, og at dette Andet, som
Principets Andet, kun kan være dette ideelle Princips \emph{egen} Bestemmelse.

Allerede naar vi gik ud fra Realitetssiden, og bestemtere: fra den abstract
opfattede Realitetsside i Tilværelsen, vistes vi altsaa hen til
Idealitetsbestemmelsen i det Absolute, og den immanente Udvikling af dette
Begreb førte til concretere Bestemmelser. Men disse kommer \emph{directe} frem
idet Tilværelsen sees fra de \emph{høieste Tilværelsesformers} Synspunkt, fra
det ved det blot Reale uforklarlige \emph{Levendes} Synspunkt, og nærmere: fra
det \emph{Sjæleliges og det Aandeliges Synspunkt}. Idet Livet faaer sit
høieste Udtryk i Menneskelivet, og Menneskelivet bliver \emph{selvbevidst}
Liv, saa viser Selvbevidstheden, idet dens \emph{Grund} skal søges, tilbage
til det Absolutes Bestemmelse som \emph{absolut Subjectivitet}. Til denne
Bestemmelse er allerede ovenfor viist hen gjennem Bestemmelsen af Principet
som \textit{causa sui}, som det Ubetingede, der idet det forholder sig til sig
selv gjennem Totaliteten af de reale Betingelser, staaer i et Frihedsforhold
til sig selv. Af denne Principets Bestemthed er \emph{Selvbevidsthedens}
Bestemmelse et Afpræg, men i Ende\emph{lighedens Form}. Fra forskjellige
Synspunkter viser denne Selvbevidsthedens Endelighedsbestemthed sig. Den viser
sig derved, at Selvbevidsthedens Bestemmelse indeslutter og for\-%
% antiqua Latin on p. 71: principium sui, causa sui (twice) --- all set close,
% NOT letterspaced, unlike p. 68.
% "Ende-lighedens Form": here the spacing starts only in the SECOND half of
% the divided word --- the mirror image of the usual case. As printed.
% --- p. 72 ---
udsætter Bevidsthedens, og derigjennem Jegets Adskillelse fra Objectet og de
andre Jeger som det begrændsende Andet. Den viser sig derved, at
Selvbevidstheden er i sin Virkelighed betinget ved den \emph{individuelle
Organisme}, og saaledes bliver en \emph{Individualitetens} Bestemmelse. Og den
viser sig derved, at Selvbevidsthedens Forsigværen, idet den hviler paa et
concret, men individuelt begrændset, Grundlag, er en Selvets \emph{abstracte}
Fordobbling, en Forsigværen i den abstracte Almindeligheds Form.
Selvbevidstheden, saaledes som vi kjende den i \emph{Mennesket}, og det er:
saaledes som vi \emph{overhovedet} kjende den, viser sig saaledes som den
begrændsede Aandsindividualitets Form. \emph{Umiddelbart} kunne vi derfor ikke
overføre Selvbevidsthedens Bestemmelse paa det Absolute uden at endeliggjøre
dette; Begrebet maa undergaae en nødvendig \emph{Metamorphose}, hvilket skeer
derved, at de endeliggjørende Bestemmelser fjernes, og disse fjernes idet
Begrebet om den \emph{absolute Subjectivitet} betegner det Absolute som det,
der, som frigjort fra Selvbevidsthedens endelige Begrændsning, kan være
\emph{Princip for Selvbevidstheder}. Vi kunne give Begrebet om den absolute
Subjectivitet dets Indhold idet vi gaae ud fra Forholdet mellem den
fornuftbestemte objective Tilværelse og det endelige Subjective. Hvad der i
dette Sidste træder frem i Reflexionens Form, i de relativ abstracte
Begrebsbestemmelsers Relationer, træder i det Objective frem som det reale
Concretes objective Sammenhæng; hvad der træder frem i det efter sit Væsen
universelle, i Individualitetens Form existerende Ideelle i Jegets abstracte
Almindeligheds Form, har sit concrete Tilsvarende i det Objective i den reale
Totalitet. Men i denne reale Totalitets \emph{Enhedsprincip} maa dette
objective Reale, der i sig indeslutter de objective Relationers Heelhed, være
reflecteret i en ideel Reflexions Form, der i en dobbelt Henseende skiller sig
fra Formen i det Menneskelige. Idet nemlig denne ideelle Reflexion ikke hviler
paa et \emph{individuelt} Concret, der idet det udelukkede det øvrige
Concrete, kun kunde give den Negativitetens og den relative Abstractions
% I/J: "Jegets", "Jeger" are J-words (Jeg); "Individualitetens",
% "Individet" etc. are I-words. Per the recorded normalisation.
% --- p. 73 ---
Form, men idet den er Form for Principet, der bærer det Concretes Heelhed i
sig som sin, ved det satte, Bestemmelse, saa har Principet i den en
Forsigværen i den absolut \emph{concrete Almindeligheds Form}, og er, ikke
blot i Abstractionens Form reproducerende, men i Virkelighedens Form sættende,
producerende, har ikke Modsætningen, som abstract, formel Modsætning i sig,
som concret, reel Modsætning udenfor sig, men omfatter i sig i uendelig
concret Forsigværen den i den reale Tilværelses Mangfoldighed udfoldede
Andethedens Modsætning. I det Absolute som absolut Subjectivitet see vi derfor
den al Tilværelse og med den alle de endelige Aandsexistentser omfattende,
bærende, gjennemvirkende Magt, i hvis uendelig concrete, energiske
Forsigværen Selvbevidsthedernes Kilde er, og i hvis evige Væren Rummets,
Tidens og Stoffets Uendelighed er sluttet sammen til Idealitetens levende
indholdsrige Enhed. I denne Sammenfatten til Idealitetens indholdsrige Enhed
er det Absolute \emph{absolut vidende Subjectivitet}, som en Viden af
\emph{sig i Alt}.

I det saaledes bestemte Absolute søge vi Forklaringen af den endelige
Selvbevidstheds Væsen, idet vi benytte det ovenfor udviklede Forhold mellem
Idealitetsbestemmelsen og Realitetsbestemmelsen til Udgangspunkt. Idet
Tilværelsen blev opfattet som den ved det ideelle Princip bestemte
\emph{universelle Totalitet}, saa ligger i dette Begreb Muligheden af
\emph{relative Totaliteter} indenfor den universelle, idet dennes Enhed, som
articulerende sig, sætter \emph{særegne} Bestemmelser, der, netop som
\emph{Totalitetsbestemmelser}, \emph{selv} kunne have \emph{Totalitetens}
Bestemmelse i sig. Dette Begreb om den relative Totalitet, der udtrykker det
\emph{Organiskes} Bestemmelse, og danner Mellemleddet mellem det
Totalitetsløses, det Uorganiskes, Begreb og den universelle Totalitets, finder
sit \emph{adæquateste} Udtryk i \emph{Menneskenaturen}, der ved sin indre
\emph{Universalitet} har Muligheden til at blive Grundlaget for den
Forsigværen af det Ideelle, der er i \emph{Selvbevidstheden, i den frigjorte
Almindeligheds Form}. Ifølge Væsenets \emph{Universalitet}
% "adæquateste": the second sort is a d, not a b --- compared at 600 dpi with
% the certain d of "finder" three words earlier; both OCR witnesses read d.
% --- p. 74 ---
har den individuelle Menneskesjæl Selvbevidsthedens Mulighed, og gjennem den
Muligheden til en Bevidsthed om Tilværelsen og dens Princip, men idet
Menneskeindividet er \emph{relativ} Totalitet, faaer den Almindelighed, der er
forsigværende i Bevidstheden, \emph{Abstractionens}, om end \emph{bestandig
integrable} Form, netop fordi i den relative Totalitet, som organisk Led i den
universelle, denne og de øvrige Led kun kunne være tilstede i
\emph{Mulighedens} Form, ikke i den udfoldede Realitets Form, i hvilke de kun
ere \emph{som} universel Totalitet, og som disse bestemte realt udformede Led.
Derfor bliver Selvbevidstheden et \emph{abstract} Afbillede af den absolute
Subjectivitets uendelig concrete Forsigværen: den har en
Endelighedsbestemthed i sig, der forklares ved Principet, men ikke umiddelbart
kan overføres paa Principet.

Gjennem Bestemmelsen af \emph{absolut Subjectivitet} viser, som allerede
berørt, det Absolute sig som det, der, idet det griber ud over det Objective,
forholder sig til sig selv \emph{og hviler i sig selv som det Ubetingede, der
er sin egen Grund}. Og i denne Bestemmelse af det er det ogsaa, at
Sammenhængen mellem hvad man har villet sondre som en \emph{blot
logisk-metaphysisk og en ethisk} Bestemmelse af det skal vise sig. Den ethiske
Bestemmelse maa have sit Fundament i den metaphysiske og kommer frem som
ethisk-metaphysisk Bestemmelse. Med Hensyn hertil faaer Sammenhængen mellem
den \emph{dobbelte} Bestemmelse af \emph{det Absolute som Tilværelsens} τέλος
sin Betydning, idet Bestemmelsen af det som τέλος deels viser hen til det
\emph{Organiskes}, deels til det \emph{Ethiskes Synspunkt}, og netop i den
\emph{organiske} Opfattelse af Tilværelsen ligger Udgangspunktet for den
Bestemmelse, i hvilken den \emph{menneskelige Frihed} grunder.

I Bestemmelsen af det Absolute som \emph{absolut Subjectivitet} og som
\emph{ideal-realt Princip} ligger Bestemmelsen indesluttet af det, der sætter
Realiteten som i dens reale Væren bestemt ved det Ideelle. Idet Eftertrykket
her falder paa Realitetens Væren \emph{som} Realitet, som værende i
% Greek on p. 74: τέλος twice, antiqua italic. Neither OCR witness reads it,
% and no list mentioned this page; caught by eye.
% "deels til det Ethiskes Synspunkt": the d of the second "deels" is
% letterspaced and its bowl is thin, which made it look like r/k at 250 dpi.
% Compared at 600 dpi with the certain "deels" one line above; both OCR
% witnesses read d. Read as "deels".
% --- p. 75 ---
den reale Existentses Grundformer, bliver dens Bestemmethed ved det Ideelle
\emph{teleologisk} Bestemmethed, og det Absolute sees fra dette Synspunkt som
teleologisk virkende. Det sees som τέλος for den Udvikling i den reale
Tilværelse, gjennem hvilken denne lader Idealiteten faae Virkelighed i sig.
Men i denne Bestemmelse er Tilknytningspunktet givet for en høiere. I den
teleologiske Omsætten af Realitetens Bestemmelse i Idealitetens faaer Begrebet
af den \emph{relative Totalitet} sin Betydning: den bliver Mellemleddet,
gjennem hvilket Omsætningen, Tilbageførelsen, realiseres. Men den relative
Totalitets Begreb finder først sin \emph{Afslutning} i Begrebet om den
\emph{ved Væsensuniversaliteten bestemte} relative Totalitet ɔ: i den
\emph{selvbevidste} menneskelige individuelle Aandsexistents. I denne bliver
det Ideelles teleologiske Virken \emph{bevidst og fri Virken}: en bevidst
Idealiseren, i fremskridende Tankeudvikling, af det optagne Reale, ɔ: en
\emph{theoretisk} Virken, og en bevidst fri Formen, under fremskridende
Villiesudvikling, af Realiteten ved Idealiteten ɔ: en \emph{praktisk} Virken.
Den Bestemmelse: \emph{med Frihed at give Idealiteten Virkelighed i det
Reale}, er det \emph{Ethiskes abstracte} Udtryk, men et Udtryk, der endnu ikke
sondrer det fra den „praktiske“ og „poietiske“ Virksomhed i Almindelighed.
Forat udtrykke det Ethiskes \textit{differentia specifica} maa den abstractere
Bestemmelse faae en concretere Form. Den ethiske Handlen er væsentlig en
Virkeliggjørelse af det Ideelle i \emph{Personlighedens Realitetsside}, --- i
den ved Naturenergien bestemte Drift, der faaer det Ideelles Form, ---
gjennem en \emph{fri Beslutning}, og Virkeliggjørelsen af det Ideelle udenfor
Personligheden er kun ethisk ved at være Udtrykket for Personlighedens indre
Handlen; uden denne Sammenhæng bliver det Ethiske til et Historisk, eller til
et Praktisk i almindelig Betydning. Idet Realitetens Bestemmen ved Idealiteten
bliver sat gjennem ethisk Beslutning og Handling, og Villien, for at kunne
realisere denne Bestemmen, maa gaae tilbage til sin egen Uendelighed, og
derved til \emph{sit Princip}: det Absolute, saa bestemmer det Absolutes
\emph{logisk-metaphysiske} Begreb sig til et \emph{ethisk-metaphysisk}:
% Greek on p. 75: τέλος, antiqua italic. Latin: differentia specifica,
% antiqua, set close (not letterspaced).
% --- p. 76 ---
det Absolute opfattes under Bestemmelsen: \emph{det Godes Idee}.

Det \emph{abstracte} Grundlag for denne Ideebestemmelse er \emph{Idealitetens
Forhold til sig selv som sin egen Energies Maal}. Men dette abstracte Grundlag
faaer en concretere Skikkelse, idet det Ideales Forhold til sig selv sees som
et Forhold \emph{gjennem det Reale}, og idet det \emph{Reale} opfattes efter
sin \emph{concrete} Væsensbestemthed. Det Absolute sees som \emph{det Gode i
metaphysisk} Forstand, idet det sees som det, til hvilket hele den reale
Tilværelse stræber hen som til \emph{det ideale Maal}, ved hvilket selve denne
Stræben er; idet det sees som det, der i Væsensmeddelelse giver det Værende en
Væren, der kun er Væren ved Forholdet til det Ideale. Det sees som \emph{det
Gode i ethisk-metaphysisk} Betydning, idet det sees som den \emph{ethiske
Handlings Maal og Norm}, som det, der i Individualitetens Form skal
virkeliggjøres gjennem den frie Handlen; som det, der, som et ubetinget
Gyldigt for Villien, skal beherske den og udfylde den. Det, at Forholdet til
den \emph{menneskelige Villie} reflecterer sig i det Absolutes Forhold til sig
selv, er Betingelsen for, at det Absolute sees som \emph{det Godes Idee i
ethisk} Bestemthed. Dette Forhold træder frem hvad enten vi bestemme det
Absolute som det, der er Normen, Loven for den menneskelige Villie, som det,
der skal udtrykkes gjennem den, eller vi bestemme det som det, der i sin evige
Selvvirkeliggjørelse begrunder de menneskelige Villier, i hvilke det giver sig
Skikkelse. Ved dette Forhold til Villien ligger det nær at overføre dennes
Bestemmelse paa det Guddommelige. Det er imidlertid i det tidligere Skrift
(S.\ 165 f.) paaviist, hvilken Metamorphose der maa foregaae med dette Begreb,
ligesom med Videns, forat det kan faae Anvendelse paa det Evige og Uendelige.
Ligesom vi i dette søgte en Selvbevidsthedens Grund, saaledes søge vi ogsaa i
det en Villiens Grund, men vi søge Villiens Grund i et fra Endeligheden
Frigjort.

Med Bestemmelsen af Villie træder Bestemmelsen af det Absolute som
\emph{absolut Personlighed} frem. Hvad der
% --- p. 77 ---
\setcounter{footnote}{0}%
fra \emph{Theologiens} Side søges hævdet gjennem denne Bestemmelse, er
allerede tidligere paaviist\footnote{Problemet om Tro og Viden, S.\ 68 f.}.
Det er paa den ene Side Forestillingssondringen fra det Menneskelige, paa den
anden Side Ubundetheden ved den almene Nødvendighedslov. Hvad der fra
\emph{Philosophiens} Side tilstræbes, naar dette Begreb anvendes, er deels det
Samme; men deels er det ogsaa det: i det Guddommelige at faae et til det
personlige Livs Fordringer svarende Villiesideal; deels endelig det: at holde
den evigt færdige guddommelige Personlighed ude fra Verdens Vorden og
Udvikling.
% printed mark: *) --- the only note on p. 77, below a short rule. Its text is
% plain Fraktur, neither letterspaced nor in quotation marks.

Naar der i det Guddommelige fordres et til det personlige Livs Krav svarende
Villiesideal, saa forstaaes denne Fordring nærmere saaledes, at der, eftersom
det Gode skal virkeliggjøres i Individets personlige Liv, og det personlige
Liv har Enkelthedens Form, i det Godes Ideal skal være en tilsvarende
Enkelthedens Form, og at det skal opfattes som guddommelig Person med
Villiesbestemmelser, der udtrykke Personlighedens med Vidensbestemmelserne
uensartede Væsen. Men idet Fordringen forstaaes saaledes, saa er deels den
allerede tidligere i sin Falskhed paaviste abstracte Modsætning mellem
Erkjendelsens Almindelighedsbestemthed og Villiens Enkelthedsbestemthed gjort
gjældende; deels er det overseet, at Gjennemførelsen af Bestemmelsen af den
\emph{absolute} Personlighed, ved at gjøre Gud til det \emph{absolut
singulære} Væsen, netop vilde tilintetgjøre Bestemmelsen af Gud som Villiens
Ideal; thi det absolut singulære Væsens Bestemthed kan ikke have Gyldighed for
noget andet Væsen end netop for dette Væsen selv. At recurrere til en
Combination af Almagtens Bestemmelse med Villiesidealets og at opfatte den
guddommelige Villie som den almægtigt bydende Villie, hvis Bud, som Almagtens
Bud, skulde blive Love for den endelige Villie, vilde ikke kunne hjælpe; thi
Almagtens \emph{vilkaarlige} Bud kunne ikke blive det absolute Ideal for
Villiens \emph{væsensbestemte} Frihedsforhold.

Saavel i Theologiens som i de forskjellige philosophiske
% "Saavel i Theologiens som i de forskjellige philosophiske" is set close
% throughout --- unlike the letterspaced "Theologiens"/"Philosophiens" earlier
% on the page. Checked at 300 dpi.
% --- p. 78 ---
Systemers Anvendelse af Personlighedsbestemmelsen paa det Guddommelige er den
nødvendige \emph{Metamorphose}, som den maa undergaae, idet den overføres paa
det Absolute, enten helt glemt eller kun inconsequent gjennemført. Den er
allerede tildeels belyst gjennem de tidligere Udviklinger, og i det
Efterfølgende vil det blive angivet, hvilken Betydning \emph{Hellighedens},
\emph{Retfærdighedens} og \emph{Kjærlighedens} Bestemmelser, der sættes som den
absolute Personligheds Prædicater, under Forudsætning af hiin Metamorphose,
kunne faae i Anvendelsen paa det Absolute.
% The three nouns are letterspaced but the „og“ between them is set close;
% checked at 200 dpi. Compare p. 82, where the „og“ inside a run IS spaced.

Naar der i Bestemmelsen af det Absolute som Personlighed lægges et særligt
Eftertryk paa, at ved denne Bestemmelse Gud som den evigt Værende skal holdes
ude fra den naturlige Tilværelses og den endelige Aands Vorden og Udvikling,
--- hvad dog, ifølge Sagens Natur, intet af de philosophiske Systemer, der have
accentueret hiin Bestemmelse, har kunnet gjennemføre, --- saa vil det i det
nærmest Efterfølgende blive paaviist, hvorledes der ud fra den givne
Bestemmelse af \emph{absolut Subjectivitet} vil kunne naaes til, idet
\emph{Immanents}forholdet statueres, at frigjøre det Absolute fra Vordens
Ufuldkommenhed og til at hævde det i dets Absoluthed, medens de Modsigelser
undgaaes, der ere uadskillelige fra Bestemmelsen af absolut Personlighed. Dette
vil fremgaae af en Undersøgelse af \emph{Pantheismens Væsen}, af dens
\emph{Sandhed} og af dens \emph{Ensidighed} i den almindelige Opfattelse.
% PRINTER'S DEFECT, p. 78: in „Immanentsforholdet“ a thin full-height vertical
% stroke stands between the two m's, plainly visible at 600 dpi --- a risen
% space or quad from the letterspacing, taking ink. Not a sort; not transcribed,
% but logged here. (ABBYY read the word „JMjMAnentsforholdet“ because of it.)
% Letterspacing: „Immanents“ is spaced and „forholdet“ is not --- the
% compositor's usual habit of dropping the spacing part-way through a compound.
% „af dens“ and „og af dens“ are set close between the spaced nouns (600 dpi).

Ved Navnet \emph{Pantheisme} klæber der ligesom en Plet, og den umiddelbare
Bevidsthed, navnlig den religiøse Bevidsthed, har en instinctmæssig Sky derfor.
Hvad der skræmmer tilbage, er, paa den ene Side, Individets Modstræben mod at
gaae op i det Absolute: en Opgaaen, der, selv naar det Absolute bestemmes som
absolut Subjectivitet, viser sig for Bevidstheden som en Fortaben i det
Selvløse, og paa den anden Side er det det, at det Guddommelige, ved at skulle
omfatte al Tilværelse, synes at nedværdiges, at drages ind i det Endeliges
Ringhed, at gaae sammen med det Forgængelige, endog med det Onde. Men hvad det
\emph{Sidste}
% --- p. 79 ---
angaaer, saa har den pantheistiske Opfattelse aldrig, selv ikke der, hvor den
optraadte paa Aandens lavere Udviklingstrin, som i Orienten, været saaledes at
forstaae, at Gud skulde være Alt, i Betydning af Indbegrebet, Collectivsummen
af alle de endelige, enkelte Ting. Naar den opfattede Gud som Alt, saa
forstodes derved, at Gud var Tingenes blivende Væsen, ved hvilket det Endelige
kun var som forsvindende Accidentser, der evigt vexle medens Væsenet bliver. De
skulde fremgaae af Substantsen og gaae tilbage deri, og gjennem deres
Forsvinden skulde de vise dens Ophøiethed og den endelige Værens Usandhed som
endelig. Og naar dette Væsen skulde finde et Udtryk i det Endelige, saa blev
det det Herligste i hver Art af det Endelige, der blev Billedet, Symbolet, for
det: det Fortrinlige staaer Væsenet nærmest og brugtes derfor til dets
Betegnelse. Sagdes der da: Gud er Alt, saa mentes der altsaa ikke: Gud er de
endelige Ting; thi disse have som endelige, som Accidentser, ingen Væren: i
deres Forgængelighed aabenbares deres Væsen, og Substantsen viser sig som det
Blivende i dem. Men ved dette: Gud er Alt, udtryktes det, at Gud er
Tilværelsens almindelige Væsen, det absolut Almene, der er de enkelte
Existentser iboende, men saaledes, at der ikke tilkommer disse som enkelte en
virkelig Realitet. Selv paa det Standpunkt altsaa, hvor Aanden endnu ikke
vidste at sondre sit Væsen fra Naturens, var Pantheismen ikke en Lære om Guds
Identitet med Summen af det Endelige, men Læren om Guds Identitet med den
almene Naturmagt, overfor hvilken intet Endeligt har Bestaaen, intet
Ufuldkomment Ret til at kaldes værende. Gud endeliggjøres altsaa ikke, drages
ikke ned i det Ringe, men det Endelige ophæves ligeoverfor det Guddommelige:
Verden, ikke Gud forsvinder. Og hvad der gjælder, naar Gud opfattes som den
almindelige Naturmagt, Natursubstantsen, det gjælder ogsaa om de indenfor
Philosophiens Gebet fremtraadte høiere Former af Pantheismen, der bestemme Gud
som \emph{Subject}, som \emph{Aand}, men i hans Væsen som absolut Aand lade de
endelige Aander optages og ophæves som Momenter. Ogsaa i disse
% --- p. 80 ---
Former, hvis Ophavsmænd, ved en \emph{vilkaarlig} Begrændsning af Pantheismens
Begreb, søge at holde dem udenfor dette, drages ikke Gud ned i det Endelige,
men det Endelige forsvinder i Gud. Og hvad Guds Forhold til \emph{det Onde}
angaaer, saa vil det i det Efterfølgende blive forsøgt at vise, at der, om end
de historisk givne pantheistiske Systemer ere blevne staaende ved det Dilemma:
enten at ophæve det Ondes Virkelighed eller at gjøre det til Bestemmelse ved
det Guddommelige, er en Mulighed til, \emph{indenfor en pantheistisk
Anskuelse} at fastholde det Ondes Virkelighed, uden at lade det ophæve
Bestemmelsen af det Absolute som det Godes Ideal.
% „Ophavsmænd“, not „Ophavsmand“ (ABBYY) or „Ophavsmend“ (tesseract): read off
% the image at 400 dpi.

Hvad det \emph{første} af de ovenfor fremhævede Motiver til en Sky for
Pantheismen angaaer: Frygten for at lade Enkeltvæsenet gaae op i det
Guddommeliges Enhed, saa viser det sig, ved en nærmere Betragtning, at denne
Sky er betinget, ikke ved noget Almindeligt og Gjennemgaaende i
Menneskenaturen, men ved noget under bestemte historiske Forhold udviklet
Eiendommeligt. Naar Inderen længes efter Befrielse fra Enkelttilværelsens
Skranker, og for at naae den lader sig knuse under Sivas Vogn, eller opsluge af
Ganges's Bølger, eller nedsænker sin Sjæl i det formløse Intet, saa er hiin Sky
ikke tilstede. Naar persiske Digtere udvikle en pantheistisk Mystik, saa er
Hengivelsen til det Uendelige, Ene, en Kilde til Begeistring. Og gaae vi fra
Orienten til Grækenland, saa finde vi, hos det Folk, hvis Eiendommelighed var
en Udprægen af den skjønne Individualitet, at saagodtsom alle dets
eiendommelige Tænknings største Repræsentanter, ogsaa saadanne, hvis Systemer
ikke have Pantheismens Charakteer, uden Savn og uden Attraa lade Individet
forsvinde ligeoverfor den evige guddommelige Tanke, der besjæler Universet.
Naar Eleaterne sætte det Ene, hvad enten det nu bestemmes som Gud eller som det
Værende, som Tilværelsens Heelhed, saa finde de netop Ro og Tankens Hvile i
denne Tanke, og Parmenides bliver Digter, idet han udtaler den. Naar
Aristoteles søger en Udødelighed for Menneskeaanden, saa er det kun for den
% --- p. 81 ---
Aandens høieste Form, der er i individualitetsløs Enhed med den evige
guddommelige i en νόησις τῆς νοήσεως hvilende Tanke. Naar Stoikerne i deres
Stræben efter Enhed lade den Aristoteliske Νοῦς blive Verdens iboende Kraft, og
lade denne fornuftige Kraft gaae sammen med Materien, saa lade de indenfor
denne pantheistiske Enhed Individerne staae beherskede af en Verdenslov, som de
villigen følge, og under hvilken de bøie sig med Resignationens Ro. Og naar
Nyplatonikerne stræbe henimod den ekstatiske Enhed med det guddommelige Ene,
saa er det saa langt fra, at det, at denne Enhed ophæver den individuelle
Existents, har noget Afskrækkende for dem, at det tværtimod er Maalet for hele
Livets brændende Higen. Og indenfor Christendommen gjælder det Samme om
Mystikernes Liden af Gud, om deres Stræben efter en Forening, i hvilken Sjælen
„gjør sig enfoldig“ for ikke at have nogen Bestemmelse, der skiller fra det
Guddommelige.
% GREEK on p. 81, two instances, both antiqua italic with accents, both read off
% the image at 600 dpi; neither OCR witness reads either (tesseract gives
% „»yonyouç re voyoswc“ and „Noùg“, ABBYY gives blanks). The first is
% νόησις τῆς νοήσεως, with the perispomeni on τῆς clearly resolved; the second
% is Νοῦς, capital nu, perispomeni over the upsilon, final sigma. Neither is
% letterspaced. p. 81 was on no list of Greek pages.
% The „6“ at the foot of the page is the gathering signature; not transcribed.

Det er væsentlig i den moderne Tid, at vi finde den ofte i sygelige Former
fremtrædende Higen efter Individualitetens Bevaren, der har den
instinctmæssige Afsky for Pantheismen til Følge, og den har, i den umiddelbare
Bevidsthed, sin Grund i den for den moderne Tid egne Fremhæven af
Subjectiviteten, der fra først af antager hiin Form, idet det spiritualistisk
bestemte Individ opfattes som uendeligt i dets Løsrevethed fra dets
Naturbetingelser. Og efterat et saadant Eftertryk er lagt paa Individualitetens
Bevaren, saa træder Attraaen derefter netop stærkest frem i de Tider og hos de
Folkeslag, hvor Individet, idet dets Endelighed uendeliggjøres ved Egoismens
Interesse, bliver Centrum i Tilværelsen, og hvor det Uendelige skydes bort som
ubegribeligt eller ligegyldigt. Oplysningsperioden var saaledes Pantheismen
fjendsk, og hos de romanske Nationer har Antipathien mod den været stærkest.

Naar saaledes hiin instinctmæssige Sky for Pantheismen har viist sig, ikke at
være begrundet i en Nødvendighed i Menneskenaturen, men i en bestemt, ved
særegne historiske Forudsætninger betinget Udvikling, saa kunne vi altsaa,
upaa\-%
% --- p. 82 ---
virkede af den, undersøge Pantheismens \emph{Væsen og Sandhed}. Og naar vi da
ikke begrændse Pantheismens Begreb til de \emph{ufuldkomne og ensidige} Former
af den, saa maae vi erkjende den pantheistiske Opfattelse som \emph{væsentlig
og nødvendig for Philosophien}. Et pantheistisk Element maa der være i al
Philosophie, fordi ingen Philosophie kan tænke det Uendelige som værende
adskilt fra og udenfor det Endelige, men, naar det Uendelige skal bevare sin
Uendelighed, maa tænke det Endelige som værende i det og kun havende Sandhed i
det; og fordi al Philosophie søger Erkjendelsens Enhed og Sammenhæng og kun er
Philosophie saavidt den forfølger denne. Idet Tanken giver sig ifærd med den
Opgave, at finde Fornuften i Tilværelsen, saa maa alt det, hvortil Tanken
naaer, danne en Enhed, en Sammenhæng, i hvilken intet Brud paa Immanentsen
findes. Thi enhver høiere Tanke, til hvilken der stiges op, er netop kun den
høiere derved, at den omfatter de lavere, at den, som den rigere og fyldigere,
har dem i sig som sine Momenter. Og hvad der gjælder om enhver høiere Tanke,
gjælder ogsaa om den høieste: den kan ikke for Tænkningen sondre sig fra de
Tanker, der indesluttes i den som dens Momenter og finde deres sidste
Begrundelse i dens omfattende Heelhed. Men denne høieste Tanke, Ideens Tanke,
er netop Tankens Udtryk for det Absolute, for Gud. Det Absolute altsaa, der
omfatter og i sig bærer de enkelte Tanker, den fornuftige Tilværelses Tanker,
maa tænkes i Væsens-Enhed med selve denne Tilværelse. Idet Tænkningen finder
Fornuft i Tilværelsen og i Fornuften har det Absolute, søger den Intet udenfor
Universet, udenfor Enheden af Tilværelsen, den legemlige og aandelige. I denne
Forstand er altsaa al Philosophie pantheistisk: den maa fastholde det Absolutes
Immanents; den kan ikke tænke sig det udenfor den Tilværelse, hvis Princip det
er; den kan ikke tænke sig en abstract Transscendentses Sondring af det
Guddommelige; thi derved vilde Uendeligheden tabes. Ved sin egen Natur føres
Tanken ikke bort fra Immanentsens Enhed, men den kan, ved Consequentser af
ensidige Forudsætninger (f.\ Ex.
% Emphasis on p. 82: three runs in which the connecting „og“ / „for“ ARE
% letterspaced with the rest (400 dpi) --- unlike p. 78, where the „og“ between
% spaced nouns is set close. Transcribed as printed in both places.
% „Væsen og Sandhed“ breaks across the line at „Sand-“/„hed.“; the continuation
% „hed.“ measures 1.26 on spacing.py and is visually near-tight, so the spacing
% may in fact stop at the break, as it does at pp. 6 and 15. Set here as one
% emphasised phrase; flagged in case a later pass wants it as „\emph{Væsen og
% Sand}hed“.
% --- p. 83 ---
Aristoteles, Descartes), og ved fremmede Paavirkninger føres til at bryde
Enheden, og indenfor den pantheistiske Opfattelse kan Tankernes Enhed faae en
forskjellig Form. Men al virkelig Philosophie er som saadan pantheistisk i den
angivne Betydning, som fastholdende det Absolutes Immanents og den endelige
Værens Usandhed, naar den som endelig Væren adskilles fra det abstract
transscendente Guddommelige.

For at bestemme Pantheismens \emph{sande} Form, udsondre vi de Hovedformer af
en \emph{ensidig} pantheistisk Opfattelse, hvis Mulighed er given i
Begrebsforholdene.

Som den \emph{første} ensidige Form fremhæve vi den, i hvilken Grundbegrebet er
Begrebet om \emph{Naturmagten}, nærmere opfattet som den mægtige ene
Natursubstants i Forhold til de blot forsvindende Accidentser. I denne Form af
Pantheismen kommer Ensidigheden \emph{umiddelbart} tilsyne, idet
Aandsexistentsen, den menneskelige Selvbevidsthed, skal sættes i Forhold til
den Substants, der ikke forklarer den, og i dette Forhold gaae ind under den
forsvindende Accidentses Bestemmelse; og \emph{middelbart} viser Ensidigheden
sig, idet Opfattelsen af det Accidentelle som det blot Forsvindende, Magtløse,
virker tilbage paa Opfattelsen af det Absolute, Substantsen, der, naar den kun
sætter et Forsvindende og kun aabenbarer sig som Magt over det blot Magtesløse,
selv bliver det Tomme og det Afmægtige.

Som den \emph{anden} ensidige Form af Pantheismen opfatte vi den, i hvilken det
Absolute faaer Bestemmelsen af Substantsens uendelige Enhed, der i sig
sammenknytter Tanke og Udstrækning. Her er vel, gjennem Tankens Bestemmelse,
Aandens Bestemmelse bragt ind i Substantsen, det Absolute; men idet Substantsen
i sin Enhed skal forene de \emph{kun disparate} Attributer, bliver den
\emph{som} Substants det rent Bestemmelsesløse, det, i hvis forskjelsløse Enhed
Alt kun forsvinder, og dens Accidentser, Modificationerne, komme til at indtage
samme Stilling til den som hiin Natursubstantses Accidentser. Selvbevidsthedens
Form bliver ikke saaledes klaret i den enkelte tænkende \textit{modus}, at den,
tilbagevirkende, kunde bestemme det Absolutes Begreb; Selvbevidstheden bliver
% ANTIQUA LATIN, p. 83: „modus“ is set in upright antiqua against the Fraktur
% and is NOT letterspaced (400 dpi). Set \textit{} per the house convention.
% „som“ in „bliver den som Substants“ is letterspaced (400 dpi); so is „som“ on
% p. 84 („det Absolute, som ideelt“) and on p. 88 („bestaaer som Rum“). The book
% emphasises this qua-„som“ repeatedly.
% The „6*“ at the foot of the page is the gathering signature; not transcribed.
% --- p. 84 ---
i Virkeligheden kun et postuleret Accessorium til Naturcausaliteten, ligesom
Udtrykket for dennes Selvfornemmen.
% „Accessorium“ and „Imponderabile“ (p. 88) are set in FRAKTUR, not antiqua ---
% checked at 400 dpi. Only „modus“ (p. 83), „eo ipso“ (p. 89) and „implicite“
% (p. 91) are antiqua in this batch.

Som den \emph{tredie} ensidige Form opfatte vi endelig den, i hvilken det
Absolute bestemmes som \emph{Subjectivitet}, men \emph{saaledes}, at i denne
Subjectivitets evige Proces Naturen, hvis Væreform er det Usande, bliver et
evigt ophævet Moment, og de endelige Aander, hvis Individualitet opfattes som
kun hvilende paa en forsvindende Naturgrund, blive kun forsvindende Former i
den absolute Enheds evige Bevægelse.

For Ensidigheden i disse Former bliver Correctivet at søge i den
\emph{concrete} Bestemmelse af det Absolute som \emph{absolut Subjectivitet},
og det saaledes, at Correctivet anbringes \emph{indenfor} den pantheistiske
Opfattelse.

Opgaven bliver, med Hensyn til det \emph{Absolute}: at hævde et
Transcendentsens Moment \emph{indenfor} Immanentsen; med Hensyn til \emph{det
ved det Absolute Begrundede}, at sikkre den reale Tilværelse og den endelige
Aand en relativ Selvstændighed.
% „Transcendentsens“ here has ONE s, and „Transscendentsen“ two lines below has
% two; both readings confirmed at 400 dpi. The book varies; not normalised.

Denne dobbelte Opgave løses, idet Consequentserne drages af en
Begrebssondring, der i det Foregaaende kun har været berørt, fordi det der
nærmest gjaldt om overhovedet at hævde Gyldigheden af det Ideale, af den
Bestemmelse, der er fælleds for Totaliteten og for det Princip, ved hvilket
Totaliteten constitueres. Denne Begrebssondring er Sondringen mellem
\emph{Tilværelsestotaliteten og Tilværelsesprincipet}.
% „Gyldigheden“ is printed with a word-space-sized gap inside it
% („Gyldig heden“) --- a loose space in a loosely justified line, not
% letterspacing: the letters within each half are tight (400 dpi). Transcribed
% as one word.

Mellem det som ideelt-realt bestemte Princip og den ved det begrundede
Tilværelsesheelhed fremtræder den Sondring, ved hvilken et Transcendentsens
Moment, i den Betydning, i hvilken Transscendentsen overhovedet kan have
Gyldighed for Philosophien, bringes ind, idet Eftertrykket falder paa
\emph{Andethedsmomentet}: det Reale, og paa \emph{det Reales Værens Form}.
Idet, som det ovenfor blev viist, det Absolute, \emph{som} ideelt, i sit Begreb
indeslutter \emph{Andethedens} Bestemmelse, som Energie Energiens Modsætning:
det Passive, saa maa det begribes som det, der i sin evige Sætten af sig sætter
det ved Andethedens og Passivitetens
% --- p. 85 ---
Bestemmelser Charakteriserede, og sætter det som værende i Realitetens
\emph{væsensgyldige}, \emph{begrebsnødvendige For}mer, medens det tillige
tænkes som det, der, i dette Forhold, evigt er i Enhedsforhold til sig selv.
Der sættes altsaa paa een Gang et dobbelt Forhold: et \emph{Evigheds}forhold og
et \emph{Tids}forhold, en \emph{Realitetens Sammenfattethed i det concrete
Ideelles Enhed} og en \emph{Realitetens Existents i
Udvorteshedssondringen}, og ved dette dobbelte Forhold, ved denne dobbelte
Bestemmelse af Realiteten, er Muligheden af den ovenfor stillede Opgaves
Løsning betinget.
% Three partial letterspacings on this page, all as printed and all confirmed at
% 400 dpi: „begrebsnødvendige For“ is spaced and „mer,“ on the next line is not;
% „Evigheds“ and „Tids“ are spaced and „forhold“ close in both compounds.
% The „og“ joining the two long runs is set close here, unlike p. 82.

Den Udvorteshedssondring, der, som væsentlig for det Reale, er uadskillelig fra
Tilværelsen, gjør sig gjældende som gjennemgaaende Bestemmelse i \emph{alle}
Former af de reale Enkeltexistentser. De primære Enkeltexistentser, de reale
Grunddele, blive, som bestemte ved Idealitetens og Energiens Modsætning,
materielle sondrede Elementer, selvløse (totalitetsløse) og passive i den
ovenfor angivne Betydning. Og idet indenfor den reale Tilværelse det Ideales
Bestemmen aabenbarer sig deri, at dens materielle Elementer sammensluttes til
det Organiskes relative Totaliteter, i hvilke et Livsprincip med Energiens og
Idealitetens Bestemmelse er virkeliggjort, saa bevarer det materielle Grundlag
for dette Livsprincips Virkeliggjørelse, under Bestemmetheden ved det Ideale,
sin væsentlige Væreform.

I de materielle Elementer og i de, i Tidsbestemthed og Rumssondring, i et
materielt Grundlag realiserede relative Totaliteter, kommer
Endelighedssondringen i den reale Tilværelse til sin Ret: den reale
Tilværelses Led faae ved den deres relative Selvstændighed og kunne gjøre deres
individuelle Existents gjældende. Den relative Selvstændighed er betinget ved
Existentsformen, derved, at hvert enkelt Led er traadt ud i Adskilthed og
virker ud fra denne sondrede Væren. Da al Væren er Virken, saa bliver al
individuel Væren --- og denne har kun Virkelighed ved Realitetssondringen ---
en Virken ud fra et individuelt Actionspunkt, og en saadan Virken er det
abstracte Tilgrundliggende i
% --- p. 86 ---
Frihedens Bestemmelse, ogsaa i Begrebet om Friheden i prægnant Betydning, om
den menneskelige Frihed, om Villiens Frihed, der er Udtryk for det universelle
Væsens Idealitet og er en Energie, der med Bevidsthed kan gjøres gjældende mod
selve det Absolute.
% NO FOOTNOTE on p. 86. The OCR-derived hint list names this page; the image
% shows an unbroken type block down to the foot, with no rule and no note.

Men naar saaledes Endelighedssondringen kommer til sin Ret, og de reale Leds
relative Selvstændighed gjør sig gjældende; saa vises der, ifølge det
oprindelige Væsensforhold, igjennem Sondringen tilbage til Enheden, igjennem
den relative Selvstændighed til Gjennemvirketheden ved det universelle Princip.

De materielle Stofelementer, hvis Væren udtrykker den største Modsætning til
Principets Enhed og Idealitet, mangle i deres selvløse Væren Energiens
Actualitet og Oprindelighed og ere bestemte ved Inertien, saaledes at de kun
yttre den i dem slumrende Activitet i Forhold til et Andet og ifølge en Impuls
udenfra, der vækker den slumrende Virksomhed tillive. Men dette Forhold til det
Virksomheden vækkende \emph{Ensartede}, der selv i sin Enkelthed ikke har
Virksomhedens Kilde i sig, viser tilbage til \emph{de materielle Elementers
Indbegreb}, hvilket, som allerede tidligere udviklet, bliver at bestemme som en
\emph{Totalitet}, der som saadan har Energiens Bestemmelse og Betingelserne for
Virksomhedens Vedvaren i sig.

De enkelte Led i den reale Tilværelse give deres individuelt bestemte Væren
Udtryk i en individuelt bestemt Virken, men denne individuelle Væren er dragen
ind i Heelhedens Sammenhæng og bestemt ud fra denne, og det enkelte Leds
Virken, der actualiseres ved de andre Leds Virken, bliver omsluttet af
Heelhedens Nødvendighed. Det enkelte Udgangspunkt for en individuel Virken
bliver derfor Gjennemgangspunkt for en universel Virken, og dets Virken
omsættes, gjennem den af alle de enkelte individuelle Virksomheder
resulterende Totalitetsvirkning, i Totalitetsprincipets Virken. Idet det af det
Enkeltes Virken Resulterende faaer Bestemmelsen af
\emph{Totalitetsvirksomhed}, saa udtrykker denne Bestemmelse, at det
Resulterende efter
% The two „Virken,“ that spacing.py flags on this page are NOT letterspaced ---
% the line is loosely justified. Checked at 400 dpi.
% --- p. 87 ---
Væsensforholdet er det Begrundende, Primitive; det Øiemed, ved hvilket det
Enkeltes Virken bestemmes.

Naar Tidens og Rummets Uendelighed er Formen for den reale Tilværelses Væren,
saa er Evighedsforholdet til sig selv Formen for det ideale Princips Væren, og
i dets Forsigværen er det Reale, Tilværelsesheelhedens Stofside, idealiseret.
Men Principets Idealitet vilde blive abstract, hvis ikke Evighedsbestemmelsen i
sig reflecterede Tidsbestemmelsen, og derved Idealiteten reflecterede det
Reales Virkelighedsform. I det Evighedens Kredsløb, i hvilket det Reale
fremgaaer af Principet, for at give dets Idealitet Fylde og Virkelighed, er
Tiden og Realitetsbestemmelsen overhovedet saaledes reflecteret, at den bevarer
sin Virkelighedsbetydning. Det Reale i Gud er en Virkelighed, ikke et
Phantasiebillede af en„ Physis i Gud“, og i den reale Tilværelses evige Enhed i
Gud er den Enhed i den naturlige Tilværelse, ved hvilken den i \emph{hvert
Øieblik} aabenbarer sin iboende Energie, optagen som Moment. Det Reale, der i
hvert Tidselement er sammensluttet til en Heelhed og i denne ideelt bestemt,
omsættes gjennem denne Mellembestemmelse i hvert Tidselement med bevaret
Realitet i det absolute Ideelle. Netop paa denne Omsætning i hvert Tidselement
(Tidsdifferential) falder her Eftertrykket. Totaliteten, gjennem hvilken den
finder Sted, bliver Mellembestemmelsen mellem det Reales Væren i
Udvorteshedsformen og dets evige Enhed i det Absolute: den Mellembestemmelse,
gjennem hvilken denne evige Enhed faaer Concretion ved den deri bevarede Reflex
af Udvorteshedsbestemtheden.
% PRINTER'S DEFECT, p. 87: the opening low quote is set TIGHT against the
% preceding word and LOOSE from the quoted matter --- „af en„ Physis i Gud““, as
% transcribed. Confirmed at 400 dpi; both OCR witnesses reproduce the fault.
% The quoted matter is Fraktur and is NOT letterspaced. The quote balance is
% unaffected: one „ and one “.

Forholdet i det Absolute mellem det Reale, der selv som idealiseret ikke maa
tilintetgjøres, men beholde sin Væsentlighed, og det Ideale, kan modtage en
Belysning gjennem Analogier fra det Endelige, gjennem psychologiske og
physiologiske Kjendsgjerninger. I den besjælede Organisme er
Rumsudstrækningen idealiseret, idet i Selvfølelsens Enhed Rumspunkternes
Udvorteshed er hævet; en Affection af en enkelt Deel af Organismen fornemmes i
det ideelt concentrerede Selv; det physiske rumsbestemte Indtryk for\-%
% --- p. 88 ---
vandles til et ideelt Moment i Selvfornemmelsen. Men Organismens Væren i Rummet
er derved ikke tilintetgjort og det idealiserlige Rum bestaaer \emph{som} Rum
for Fornemmelsen og Sandsningen. --- Materien idealiseres i Organismen, idet
den fornemmende Organisme i sin normale Tilstand for sig er et Imponderabile;
men Organismen ophører derved ikke at være materiel, at være Tyngdens Love
underkastet. --- Bevidstheden knytter de forskjellige Forestillingselementer
sammen i en Forestillingsenhed, til hvilken der svarer en Tidsenhed af
Tidselementerne, i hvilken Successionens Udvorteshed er hævet; medens dog af
Forestillingselementerne, der, som knyttede til en Nerveact, fordre hver sit
Tidselement, i hvert Øieblik kun eet kan være for Bevidsthedens Focus,
Bevidsthedens Affection ved de enkelte Elementer altsaa maa have Successionens
Form. Men idet Tiden saaledes sammenfattes til en Tidsenhed, og i denne
Sammenfatten, i hvilken Successionen er concentreret og derved idealiseret,
taber sin Udvorteshed, vedbliver den at have sin reale Gyldighed for
Bevidstheden. --- Saaledes vise disse Kjendsgjerninger hen paa en Bevaren af
det Reales Grundbestemmelser, medens disse idealiseres. Og naar vi nu, med
disse Kjendsgjerninger som Tilknytningspunkter, ved en Analogieslutning
overføre hvad der gjælder om det Individuelle paa Tilværelsestotalitetens
absolute Princip, saa ville vi, --- idet vi opfatte dette som det, der efter sit
Begreb sætter det Reale, der er dets \emph{egen} Væsensbestemmelse, som et
\emph{Væsentligt og Gyldigt}, som et i \emph{væsensgyldige Existentsformer}
værende Tilværelsessystem, og med det Samme see det som det, der, idet det
hævder sin Idealitet, sammenslutter dette reale Tilværelsessystem i den
Idealitetens indholdsrige Enhed, ved hvilken den er begrundet, --- faae Rum, Tid
og Materie, det Reales Grundbestemmelser, som væsensgyldige for det Absolute,
men tillige som saaledes idealiserede, at en Evighedens levende Nærværenhed i
sig concentrerer Tidens Uendelighed, at det Absolutes Livsenhed i sig
concentrerer i et Idealitetens og Subjectivitetens Punkt det grændseløse Rums
Udvorteshed og Materiens Uendelighed.
% „Idealitetens“ twice, not „Identitetens“: ABBYY reads „Identitetens“ in both
% places and is wrong; the image at 400 dpi shows I-d-e-a-l-i-t-e-t-e-n-s.
% --- p. 89 ---
Den Opfattelse, vi have gjort gjældende af det Absolutes Forhold til den i sine
væsensgyldige Former med relativ Selvstændighed existerende reale Tilværelse,
indeslutter, i Begrebssondringen mellem det absolute Princip og
Tilværelsesheelheden og i Bestemmelsen af det idealt-reale Princips
Evighedsforhold til sig selv, et \emph{Transscendentsens Moment}; men saaledes,
at der ikke ved det er gjort et Brud paa Immanentsen; thi Væsenssammenhængen er
bevaret, og Tilværelsen er kun opfattet som Totalitet i Kraft af Principets
Iboen. Ifølge det oprindelige Væsensforhold til den reale Tilværelse bliver det
Absolute, idet det bestemmes som \emph{Aand}, \textit{eo ipso} at bestemme som
\emph{Naturprincip}, (\emph{Naturidee}), ved hvilket Udtryk det da betegnes som
Princip for den naturlige Tilværelses ideelle Enhed, for dens universelle
Totalitet.
% ANTIQUA LATIN, p. 89: „eo ipso“ is upright antiqua and is NOT letterspaced ---
% compare the letterspaced Fraktur „Aand“ immediately before it, on the same
% line (600 dpi). So this book's Latin is not uniformly letterspaced: cf. the
% letterspaced Latin at p. 68 against the close-set Latin at pp. 64, 71, 83 and
% 91.
% The letterspacing of „Naturprincip, (Naturidee)“ runs through the comma and
% the parentheses; set here with the emphasis on the two words only.

Ved det Absolutes Evighedsforhold til sig selv er det, under sin
Væsenssammenhæng med den i Realitetens Grundformer værende reale Tilværelse,
unddraget Vorden. Dets Forhold til det enkelte Endelige, der, seet fra dettes
Synspunkt, er i Tidsbestemthed, bliver, idet det optages som Moment i det
Absolutes Evighedsforhold, en ideal-real Bestemmelse i den concrete Evighed, og
som saadan frigjort fra Udvortesheden. Og med det Absolutes Frigjorthed for
Vorden staaer dets Frihed for endeliggjørende Begrændsning i Forbindelse. Ved
Idealiseringen af det Reale, der netop stadfæster dette i dets Væsensgyldighed
som Realt, er det Absolutes Uendelighed hævdet som concret Uendelighed, og i
dets af Immanentsforholdet omsluttede Transscendentsforhold er den Begrændsning
fjernet, der ikke undgaaes i den almindelige Bevidstheds paa
\emph{Forestillingen} om Transscendentsen hvilende Opfattelse, og ikke heller i
de philosophiske Forsøg, der under Benævnelsen: logisk \emph{Theisme}, skjule
den theologiske \emph{Theisme}. --- I det Efterfølgende vil, gjennem Udviklingen
af det Ondes Problem, det Absolutes Frigjorthed fra Endelighedens
Ufuldkommenhed trods Immanentsen blive belyst fra et nyt Synspunkt.
% „ideal-real“: the hyphen is printed with a space on each side („ideal - real“),
% where „ideelt-realt“ (p. 84) and „idealt-reale“ (six lines above, same page)
% are set tight. Neither word is letterspaced --- the letters are close; only the
% hyphen is loose. Transcribed as a plain compound; the spaced hyphen is logged
% here rather than reproduced.

\begin{center}\rule{2cm}{0.4pt}\end{center}
% p. 89 is the last page of the sub-section „1. Det Absolutes Begreb“ and is a
% short page: the type block ends and a short centred rule closes the
% sub-section, exactly as at pp. 29 and 63. Confirmed at 600 dpi. Unlike p. 29
% this one is a SINGLE rule, not a double.

% --- p. 90 ---
% The printed head is bold Fraktur on two lines, the „2.“ standing to the left
% of the first line. Read at 600 dpi: „Det menneskelige Væsens Begreb“, D-e-t.
% The book's own Indhold reads „De menneskelige Væsens Begreb“ --- one of the six
% head/Indhold divergences in this book. Printed head in the body; the Indhold's
% reading recorded here. Do not normalise.
\section*{2. Det menneskelige Væsens Begreb, bestemt i Forhold til det Absolute og til Tilværelsestotaliteten.}

Efterat det Absolutes Begreb er udviklet i dets Sammenhæng med Begrebet om
Tilværelsens Heelhed, bliver det Opgaven, at bestemme det menneskelige Væsens
Begreb i Tilknytning til hine Begrebsudviklinger. I Bestemmelsen af
Menneskevæsenets Begreb er Bestemmelsen af den menneskelige \emph{Livsopgave}
indesluttet, og fra den og fra hine Begrebsbestemmelser maa Forsøget paa en
Løsning af de afgjørende ethisk-religiøse Problemer udgaae.

Det \emph{almene} Begreb, der i det Foregaaende er vundet til Bestemmelsen af
det menneskelige Væsen, er Begrebet om den \emph{relative Totalitet}: det
Begreb, der betegner det Levende, det Organiske, overhovedet, og paa den ene
Side adskiller det fra Begrebet om det totalitetsløse Elementære, paa den anden
Side fra Begrebet om den universelle ved Principet constituerede Totalitet, den
reale Tilværelses Heelhed. Indenfor den relative Totalitets almene Begreb
stiller sig saa Menneskevæsenet som særegen Bestemmelse fra det Levendes lavere
Former: fra Planten og Dyret. Plantelivet har til sin Grundbestemmelse
\emph{Selvhedens} Bestemmelse, det Organiskes \emph{abstracte
Almeenbestemmelse}, der angiver Idealiteten, Totalitetsprincipets Forhold til
sig selv gjennem sit Andet, men det mangler Concentrationen i et
\emph{sjæleligt Selv}. En saadan Concentration naaes i Dyrelivet; men
Dyrelivets høieste Bestemmelse er den \emph{Selvfornemmelse}, i hvilken det
individuelt-almene Selv bestandig kun er værende for sig i en Enkelthedens
Bestemthed, i umiddelbar Enhed med denne enkelte Bestemmelse og ved dette
Forhold røber sit Væsens Begrændsethed. Først i Menneskelivet er
\emph{Selvbevidstheden} Grundbestemmelsen: det efter sit ideelle Væsen almene
Selv er som værende for sig i den \emph{frigjorte} Almindeligheds Form, og
denne Selvbevidsthedens Bestemmelse er ikke en accessorisk Bestemmelse, en
Egenskabsbestemmelse ved et realt forudsat Subject, eller rettere: Substrat,
men \emph{Væsensbestemmelse i} det Menneskelige. Den udtrykker selve det ideale
Grundvæsen, der
% „accessorisk“, not „accessorist“ (tesseract): read at 600 dpi.
% A small ink speck stands between the letterspaced „i“ and „det“ in
% „Væsensbestemmelse i det Menneskelige“; tesseract read it as a hyphen. It is
% not a sort and is not transcribed (cf. the specks rejected at pp. 4, 31, 35,
% 36). The letterspacing stops after „i“; „det Menneskelige“ is close-set.
% --- p. 91 ---
er for sig værende i sin Realitetsbestemthed. Men forat Selvbevidstheden skal
blive mulig, maa den particulære Væsensbestemmelse, der betegner Dyrelivet,
have givet Plads for \emph{Væsensuniversalitetens}; kun gjennem den bliver det
Almenes Væren for sig \emph{som alment}, der er Selvbevidsthedens Grundform,
mulig, og kun ved den bliver \emph{Tænkning og Villie} mulig. Og idet
Selvbevidstheden saaledes viser tilbage til Væsensuniversaliteten, viser den
tilbage til \emph{Væsenets Uendelighed}: kun det efter sit Begreb
\emph{uendelige} Væsen er \emph{selvbevidst} Væsen.

Denne Sammenhæng mellem \emph{Selvbevidsthedens} Bestemmelse og det
\emph{Uendelige} har først \emph{Leibnitz} antydet idet han satte
Selvbevidstheden i Forbindelse med Erkjendelsen af de „evige Sandheder“, der
have deres Udspring fra det Guddommelige. Først gjennem Muligheden til
Erkjendelsen af disse evige Sandheder med deres Nødvendighed og Almindelighed
har Selvets Perception Reflexibilitetens Mulighed; først ved den er, med andre
Ord, Selvbevidstheden mulig. I hiin Erkjendelse, i hvilken Gjenstanden har
Uendelighedens Mærke, forholder nemlig Bevidstheden sig ikke til noget
Fremmedt, bindes ikke af det Particulære, men forholder sig i Gjenstanden til
sig selv, gjennem det Almene til sit eget ved det Almene bestemte Væsen.
% „Leibnitz“ is a proper name set in FRAKTUR and letterspaced, as the book's
% habit is (400 dpi) --- no antiqua name has turned up in this batch.

Idet Selvbevidstheden er Menneskets \emph{Væsensudtryk}, Formen for Væsenets
Selvaabenbarethed, saa indeslutter Selvbevidstheden et Forhold til
Menneskevæsenets \emph{Grund}, og da det Væsen, der gjør Selvbevidstheden
mulig, er det \emph{universelle} Væsen, saa er den Menneskevæsenets Grund, til
hvilken Selvbevidstheden har et Forhold, umiddelbart bestemt som
\emph{Tilværelsens universelle Grund} (Princip), og gjennem Forholdet til den
er et Forhold sat til \emph{den reale Tilværelses Totalitet}. I
Menneskeaandens Selvbevidsthedsforhold til sig selv indesluttes altsaa
\textit{implicite} det dobbelte Forhold: til det Absolute og til den
\emph{Tilværelses Totalitet}, der er begrundet i det Absolute. Dette dobbelte
Forhold maa træde frem for Bevidstheden under Selvbevidsthedens Udfoldelse til
concret Selverkjen\-%
% ANTIQUA LATIN, p. 91: „implicite“ is upright antiqua against the Fraktur and
% is not letterspaced (400 dpi). A third Latin word in this batch, on no list.
% The „(Princip)“ after the spaced run is close-set, as is „den“ before
% „Tilværelses Totalitet“ in the following sentence.
% --- p. 92 ---
delse; det maa vise sig i concret Enhed med
Selvbevidsthedsforholdet hvor den abstracte Selvbevidsthed gjennem den
begribende Erkjenden er udfoldet til fornuftbestemt Erkjenden af det, der som
indre Mulighed ligger i det ideelt altomfattende Menneskevæsen. ---

Ud fra Begrebet om Selvbevidstheden, som Væsensudtryk for det Menneskelige, og
som indesluttende Selvets Forhold til det Absolute og til Tilværelsesheelheden,
maa den \emph{menneskelige Livsopgave} bestemmes. See vi hen til den Væsenets
Art, som Selvbevidstheden aabenbarer, saa bliver Opgaven i dens Almindelighed
at bestemme ved den \emph{Enhed af Uendelighed og individuel Væren}, der er
udtrykt i Selvbevidstheden, og den bliver saaledes: \emph{en Virkeliggjørelse
af Væsenets Uendelighed i den givne Individualitets Form}. Og idet det
menneskelige Individ, som oprindelig anlagt til Selvbevidsthed og
Selvbestemmelse, er anlagt til Personlighed, kan Opgaven med andre Ord
bestemmes som \emph{Personlighedens Virkeliggjørelse gjennem Udfoldelsen af
Væsenets indre Uendelighed}. Men Personligheden maa da ikke opfattes som Udtryk
for en enkelt Side eller en enkelt Virksomhedsform af Menneskevæsenet, men som
\emph{Enhedsudtrykket} for det \emph{totale} menneskelige Væsen i Indbegrebet
af dets Grundfunctioner.
% EMPHASIS, p. 92 (300 dpi): the run breaks mid-phrase — „Enhedsudtrykket“ is
% spaced across its line break (Enheds-/udtrykket, both halves), „for det“ is
% close-set, „totale“ spaced again, and „menneskelige Væsen i Indbegrebet af
% dets Grundfunctioner“ close-set. Transcribed as printed, not regularised.

\emph{Den almene} Opgave maa, som Opgave for et Selv, der aabenbarer sig i en
Flerhed af primitive Yttringsformer, articulere sig i en Flerhed af
væsensforbundne, \emph{særlige} Opgaver, men det almene Udtryk for den maa
bevare sin Gyldighed for disse.

Saaledes bliver den menneskelige \emph{Erkjendelses Opgave}: Uendelighedens
Virkeliggjørelse i den med Selverkjendelsen til Enhed forbundne Erkjendelse af
Tilværelsen og dens Princip. Ved en saadan Erkjendelse, i hvilken Selvet ideelt
optager den fornuftbestemte, uendelige, reale Tilværelse, og sammenslutter sig
i Enhed med sit og Tilværelsens begrebne Princip, virkeliggjøres Selvets indre
Uendelighed og i Uendelighedsforholdet til sig selv virkeliggjøres Friheden.

% --- p. 93 ---
Den menneskelige \emph{Villies} Opgave bliver Selvets ethiske Villen af sig
selv efter sin concrete Uendelighed. I den ethiske Villen hævder Selvet sig
selv efter sin Uendelighed og gjør sin Frihed gjældende idet det er
bestemmende ud fra sin sande Uendelighed. Men Betingelsen for en saadan Villen
er, at ingen dualistisk endeliggjørende Sondring adskiller Villien fra Aandens
andre Virksomhedsformer; thi kun saaledes kan Villien ville Væsenet efter dets
totale Bestemthed. Ligesom Erkjendelsen paa sin Side maa bevare Universalitetens
Charakteer, saaledes, at den ikke, i Forhold til noget Aanden Tilhørende,
bliver et blot afspeilende Gjemmested, men virkelig kan gjøre Aandens Indhold
gjennemsigtigt for den, saaledes maa ogsaa Villien bevare denne Universalitetens
Charakteer, og kunne have alle Aandens væsentlige Virksomhedsformer til sit
Indhold. Og dette vil, med andre Ord, sige: den ethiske Opgave maa blive en
altomfattende Opgave, omslutte Menneskelivet som Totalitet, saaledes, at hele
Menneskelivet, alle Aandsyttringerne, gaae ind under det ethiske Ansvar.

Til disse Bestemmelser af Opgaven naae vi idet vi tage vort Udgangspunkt fra
den i Selvbevidstheden aabenbarede \emph{Væsensuendelighed}. Men idet vi nu
accentuere det i Selvbevidstheden liggende oprindelige \emph{Forhold til
Principet og Tilværelsen}, saa fører dette os til et concretere Udtryk for
Opgaven.

Idet vi bestemte Erkjendelsens Opgave, naaede vi til Begrebet om den udfoldede
Selvbevidstheds Enhed med Erkjendelsen af Tilværelsen og af Principet. I denne
Erkjendelse bliver Selvet sig bevidst \emph{i sin Indordning som organisk Led i
den af det concret bestemte Princip beherskede Tilværelse}, og idet Selvet, der
gjennem Bevidsthedens Selvfordybelse finder det Absolute som det, i hvilket
først dets Selvbevidsthed fuldbyrdes, i \emph{Selvhengivelsen til det Absolute}
har Bevidstheden om \emph{Væsensenheden} med det, faaer det Absolutes
Væsensbestemmelse Gyldighed for Selvet og for den af Principet satte Totalitet.
Men Principets Væsensbestemmelse var \emph{det}
% --- p. 94 ---
\emph{Ideales og Reales concrete Enhed, der udelukkede al dualistisk Sondring}:
denne Bestemmelse faaer altsaa, idet Immanentsforholdet ikke brydes, sin
Gyldighed for det ved Principet Værende.
% EMPHASIS across the p. 93/94 joint: the spaced run begins with „det“ as the
% last word of p. 93 and runs to „Sondring“ on p. 94; the colon is close-set.
% Set as two \emph{} groups so the page marker can stand between them.

Naar da Selvet er det, af ingen Dualisme adsplittede, i sine Grundfunctioners
væsentlige Enhed sig aabenbarende, og naar det som integrerende Led er indordnet
i den ideelt-reale Tilværelsesorganisation, saa ligger heri, at dets Opgave maa
blive at udvikle den udelte Personlighed i dens indre Harmonie, at udvikle den
efter Væsenets Fuldstændighed som naturligt-aandeligt, i uadskillelig Enhed af
det Ideale og det Reale, og at udvikle den som eiendommeligt Led i den ved
Principet bestemte Tilværelsens Heelhed, ved at gjennemtrænge dens
Eiendommelighed med det Almene og gjøre den til et eiendommeligt Udtryk for det
Almeengyldige. Men idet Menneskeselvet bliver sig bevidst i sin Indordnethed i
Tilværelsestotaliteten, bliver det sig tillige, ifølge sit Væsens
Universalitet, bevidst i sin \emph{relative Selvstændighed som integrerende
Led}, og dets Opgave bliver da, at ville sig i denne Selvstændighed og at gjøre
den gjældende gjennem hiin harmonisk-fuldstændige Udvikling af sit Væsens
Eiendommelighed. Mennesket maa ville sit Væsen som dette Concrete, i hvilket
det reale Naturgrundlag ikke skal negeres, men optages i Aanden, gjennemformes
med Bevidsthed, saaledes at altsaa Villien faaer Indhold og Concretion ved
Individets Realitetsbestemthed, medens den har sit Regulerende i det
Aprioriske, der paa det Ethiskes Gebet faaer samme Gyldighed som paa det
Intellectuelles. Men idet det vil sig \emph{ethisk} i sin individuelle
Eiendommelighed og relative Selvstændighed, er denne Villen i Enhed med dets
Villen af sig som Led i den omfattende Heelhed, som underordnet Principet, og
ved sin Villen gjør det saaledes en Uendelighedens Bevægelse, eviggjør sin
Handlen medens den bevarer Tiden og Virkeligheden som sin Skueplads. ---

Saaledes opfattet bliver altsaa Livsopgaven en \emph{Uendelighedsopgave}, idet
Selvets Villen er en Villen af sig efter sit Væsens Uendelighed, efter sit
Uendelighedsfor\-%
% --- p. 95 ---
hold til Tilværelsen og dens absolute Princip, efter sin Betydning som ideelt
Moment i Aandens evige Livsenhed. Men idet Livsopgaven bliver en
Uendelighedsopgave, en Evighedsopgave, faaer den den nærmere Bestemmelse af en
Uendelighedsopgave, der skal løses i Endeligheden, af en Evighedsopgave, der
skal løses i Tiden, under Existentsbetingelser, Endelighedsbetingelser. Herved
falder der et Eftertryk paa det \emph{ethiske Valg som Selvets Valg i et
Øieblik af Tiden}, og dermed paa den \emph{individuelle frie Villies Valg} til
Forskjel fra den metaphysiske Nødvendigheds evige Bestemmen; og der falder et
Eftertryk paa den Betydning, Tiden faaer for Løsningen af Evighedsopgaven og
for dens Forfeilen, paa den Splid i Selvet, der opstaaer ved Selvets dobbelte
Forhold, til det Evige og til Tidsbetingelserne, og paa den Skyld, Selvet i
dette Forhold erkjender at have paadraget sig.

Fra dette \emph{Synspunkt} vil da \emph{Skyldens} Bestemmelse \emph{ikke} blive
forflygtiget til et \emph{metaphysisk} Moment; den vil blive \emph{i Sandhed
ethisk} Bestemmelse med afgjørende \emph{Virkeligheds}betydning, og den vil,
idet \emph{Idealets For}dring hævdes, blive seet som \emph{almeen} Bestemmelse i
det endelighedsbetingede Menneskelige. Det ethiske Individ, der anerkjender den
\emph{uendelige} Fordring, maa bøie sig under Følelsen af \emph{Skylden},
anerkjende \emph{Angerens} Nødvendighed og Trangen til en \emph{supplerende
Bistand} forat Skylden kan forsones, Restitutionen skee. Men idet Tilværelsens
Heelhed sees som baaren af det immanente Princip, og dens \emph{Harmonie som
Heelhed} derfor ikke kan tænkes som forstyrret, ligesom eiheller
Enkeltexistentsens \emph{Væsen} kan tænkes som forvandlet, saa maa Individet
ikke søge den supplerende Hjælp udenfor Tilværelsen, men det maa søge de
forsonende Kræfter i denne Tilværelse: i de andre supplerende Enkeltvæsener,
med hvilke Individet slutter sig sammen i Kjærlighed, i gjensidig
Væsensmeddelelse, og, definitivt, i det Princip, der er Individet selv iboende
og gjennemvirkende, og som slutter det sammen med de andre Individer i Væsenets
Enhed. Det vil i det Efterfølgende blive viist, hvorledes
% EMPHASIS, p. 95 (300 dpi): two more mid-word breaks in the letterspacing —
% „Virkeligheds“ is spaced and „betydning“ close-set in the same word; and
% „Idealets For-“ is spaced while „dring“, after the line break, is close-set.
% As printed in both cases.
% --- p. 96 ---
Redintegrationen definitivt bliver at søge gjennem den absolute Selvhengivelse
af det i Modstanden mod det Ethiske efter sin reale Grund actualiserede Selv
til dette gjennemvirkende Princip, og hvorledes gjennem Redintegrationen
Livshelhedens Ethiceren bliver mulig.

Idet Livsopgaven er saaledes bestemt, at Personligheden i sin frie, harmoniske
Selvudarbeidelse bestandig skal leve sig ind i Heelheden og saaledes
virkeliggjøre Forholdet til dens Princip, saa er hermed det Samme udtrykt, der
tidligere er hævdet: \emph{at det religiøse Forhold bliver det afsluttende}, i
hvilket de særlige Forhold først finde deres Fuldendelse. At det religiøse
Forhold maa blive det, gjennem hvilket Livsopgavens Løsning fuldbyrdes, ligger
i Menneskevæsenets Grundforhold; thi Menneskevæsenet grunder i det
Guddommelige, og kan som selvbevist Væsen kun komme til Virkeliggjørelse, idet
dette Forhold, ved hvilket Væsenet er, bliver til for Væsenets Bevidsthed, og
bliver saaledes til for det, at Gudsforholdet gaaer i Enhed med Selvets Forhold
til sig selv i Selvbevidstheden.
% PRINTER'S DEFECT, p. 96, l. 15: „selvbevist“ for „selvbevidst“ — the d is
% simply absent from the sort-string; confirmed at 300 dpi and by both OCR
% witnesses. Transcribed as printed.
% Note also „Livshelhedens“ (single e) four lines above, where the book
% elsewhere sets „Heelhed“. As printed.

Men naar saaledes det Religiøse atter har viist sig som det Afsluttende, saa er
det af tidligere Udviklinger klart, at den Religiositet, der hviler paa Videns
og den ethiske Villies Grund, og har sit Indhold i det \emph{hele}
Menneskelivet omfattende Ethiske, bliver en saadan, der ikke sondrer sig ud som
et exclusivt Forhold, men at dens Cultus bliver Handling, dens Ideal det, at
hele Livet \emph{som ethisk} Liv er Cultus.

For det ethisk-religiøse Selv bliver det Absolute bevidst i det
\emph{Ethiskes} Bestemthed; thi Selvets Væren bestemmer Formen af det Absolutes
Væren for Selvet: som Mennesket er, er dets Gud. Religionen har sin Kilde i
selve Menneskevæsenet: dens Indhold er betinget ved dette Væsens Aabenbarethed
for sig selv: dets Gudsbevidsthed holder Skridt med dets Selvbevidsthed. Men
Religionen har kun sin Kilde i Menneskevæsenet fordi det grunder i det
Absolute, og i Selvet, der er sat ved det Absolute, hvis successive Aabenbaring
for Bevidstheden hviler paa Væsens\-%
% --- p. 97 ---
\setcounter{footnote}{0}%
meddelelsen til Menneskeaanden, er det Absolute ogsaa i det religiøse Forhold
det Gjennemvirkende og det primitivt Virkende: i \emph{sidste} Instants søge vi
altsaa Religionens Kilde i det.

\begin{center}\rule{2cm}{0.4pt}\end{center}
% RULE, p. 97, the first of the two on this page: a short centred SINGLE hairline
% standing mid-page between „…Religionens Kilde i det.“ and „Den i det
% Foregaaende udviklede Opfattelse…“, with no head adjacent to it --- a
% thought-break in running text, like the rules of pp. 19, 28, 42 and 62. A small
% ink blot sits on its middle; it is not a second rule. Transcribed.

Den i det Foregaaende udviklede Opfattelse vil blive klarere belyst i sin
Eiendommelighed gjennem en Sammenligning med Theorier, med hvilke den har
iøinefaldende Berøringspunkter, medens den dog i et Afgjørende adskiller sig
fra dem. Med Hensyn til Opfattelsen af det Absolute, af Religionen og af
Forholdet mellem det Ethiske og det Religiøse vil Sammenligningen med
\emph{Feuerbachs} Theorie være oplysende; med Hensyn til det ethiske Livs
Forhold til Virkelighedsbestemmelserne vil Sammenligningen med den
\emph{Hegelske} Synsmaade være det.
% p. 97 opens with an ornamental two-line initial D (set normally here). The
% SECOND of the page's two rules stands immediately below this paragraph and
% immediately above the run-in head „Ludvig Feuerbach.“: it is head ornament and
% is NOT transcribed, as at pp. 126, 128, 152, 164, 168, 177 and 182. (The
% page's first rule, above this paragraph, is a thought-break and IS
% transcribed --- see the note at the head of the page.) It too is a single
% hairline. The signature numeral 7 at the foot is not transcribed.

\textbf{Ludvig Feuerbach.} --- Allerede i Forordet til „Problemet om Tro og
Viden“ er Feuerbachs Standpunkt i Almindelighed charakteriseret, og de
væsentlige Punkter, i hvilke jeg afviger fra det, antydede. Det vil
imidlertid, for Sammenligningens og Kritikens Skyld, være nødvendigt,
udførligere at gaae ind paa en Fremstilling af hans Opfattelse i dens
Sammenhæng og i dens indre Udvikling, og det vil saameget mere være det, som de
Udtalelser, der hertillands ere fremkomne om Feuerbach, saagodtsom uden
Undtagelse\footnote{Som en saadan Undtagelse maa jeg navnlig fremhæve en
Afhandling, der for mange Aar siden blev trykt i \emph{L. Helwegs}
Fjerdingsaarsskrift for Literatur og Kritik (1845, 3. Hefte, S. 1--17) og hvis
Forfatter var min afdøde Ven, den som pædagogisk Forfatter bekjendte Adjunct
\emph{Chr. Fenger Christens}. Men denne Afhandling ligger i Tiden forud for de
Skrifter af Feuerbach, i hvilke han har gjennemført sin Theorie til dens sidste
Consequentser.} røbe ikke blot et høist ufuldstændigt Kjendskab til hans
Skrifter, men en ringe Indsigt i det for hans Theorie Charakteristiske.
% HEAD, p. 97: the printed form is a run-in halbfette (bold) Fraktur head,
% „Ludvig Feuerbach.“ followed by an em-dash — not letterspaced; checked at
% 300 dpi against the close-set roman of „Allerede“ beside it. The Indhold
% instead reads „1) Ludvig Feuerbachs“. Printed head kept in the body.
% FOOTNOTE p. 97: printed mark *). Its two names are Fraktur and letterspaced,
% hence \emph{} and not \textit{}. The dash in „S. 1--17“ is a short (en) rule,
% distinctly shorter than the em-dash of the run-in head.
% --- p. 98 ---
Hvad der vanskeliggjør Opfattelsen og Fremstillingen af Feuerbachs Anskuelse af
det Religiøse, ligesom af hans philosophiske Anskuelse overhovedet, er deels
hans Philosopherens aphoristiske Form, navnlig paa hans sidste
Udviklingstrin, deels den bestandige Bevægelse, der finder Sted i hans
Tænkning, og som lader ethvert nyt Skrift betegne et nyt Stadium i en
Udvikling, hvis Begyndelses- og Slutningspunkt ligge i en uhyre Afstand fra
hinanden, og tilhøre to heelt forskjellige Aandsretninger. Kun idet denne
rastløse Udviklings drivende Kraft findes, og Udviklingens Heelhed overskues,
bliver den rette Forstaaelse mulig.
% p. 98, l. 3: „hans Philosopherens aphoristiske Form“ as printed — awkward but
% read the same way by both witnesses and confirmed at 300 dpi. The pencil
% marginalia and the pencil underline under „heterogent“ are modern and are not
% transcribed.

\emph{Feuerbach} begyndte som Theolog, med lidenskabelig Stræben efter at
gjennemtrænge Troens Indhold med Tanken. Den Form af Theologien, der drog ham
til sig, var derfor den speculative Theologie; ved den førtes han over til
\emph{Hegels Religionsphilosophie}, og idet han opfattede denne i Logikens
Belysning, blev han snart fremmed for Theologien: det Guddommelige blev for ham
\emph{Fornuftvæsenet}, Fornuftens Enhed det Evige, i hvilket de forgængelige
Individer opløses. Trods begyndende Tvivl fastholdt han endnu den Hegelske
Philosophies Grundanskuelse: Anskuelsen af Tanken som i og for sig værende
Virkelighed og Sandhed, saaledes at det \emph{Tænkte som saadant} er virkeligt,
Tanke og Væren identiske. Men snart fik Tvivlen om denne Anskuelses Sandhed
Magt, og medens han fjernede sig fra den, omformedes hans Opfattelse af det
Guddommelige og af Religionen. \emph{Indirect} havde den Hegelske Philosophie
anerkjendt Naturens Væsentlighed og Sandhed, ved at sætte den som nødvendigt
Moment i Ideens evige Proces; men denne Anerkjendelse blev atter ligesom
tilbagekaldt idet Naturen \emph{som} Natur blev et kun Usandt, Forsvindende.
Feuerbach fastholdt hiin Anerkjendelse; han tillagde først „det Sandselige“
Sandhed og Væsentlighed \emph{indenfor Naturens Gebet}, og saa \emph{absolut
Realitet}, idet „det Væsen, der modsættes Sandseligheden som et heterogent
Væsen, ikke er Andet end Sandselighedens abstracte Væsen“. Det Primitive blev
nu for ham ikke længere Aanden, men Na\-%
% --- p. 99 ---
turen, ikke længere Tanke, men Væren: en Materiens Fremgaaen af Aanden erklærer
han for utænkelig; altsaa maa Materien gaae forud for Aanden, Naturen være
Aandens Grund. „Det sande Forhold mellem Tænken og Væren er kun dette: Væren er
Subject, Tænken Prædicat. Væren er ud af sig selv og ved sig selv, --- Væren
har sin Grund i sig selv, fordi kun Væren har Mening, Fornuft, Nødvendighed,
Sandhed, kortsagt: er Alt i Alt. Væren er fordi Ikke-Væren er Ikke-Væren ɔ:
Intet, Meningsløshed. Værens Væsen som Væren er Naturens Væsen“. Og Væren,
Naturen, har nu for Feuerbach ikke Betydning af en almeen Enhed, men af
\emph{Indbegrebet af det Sandseligt-Enkelte}, af de „ganske bestemte,
individuelle, sandselige Ting og Væsener“, der blive det ene Virkelige; medens
det Almene kun bliver en subjectiv Fælledsbetegnelse, ligesom „Aanden“ kun
bliver Sandsernes Enhed.

Denne Udvikling af Feuerbachs philosophiske Grundanskuelse fra
\emph{Idealisme} til en eiendommelig formet \emph{Realisme}, der ved
Anerkjendelsen af Menneskevæsenets oprindelige Activitet og ved Opfattelsen af
de Aandsphænomener, i hvilke denne lægger sig for Dagen, skiller sig fra den
almindelige Sensualisme og Materialisme, maatte blive afgjørende for hans
\emph{Opfattelse af det Religiøse}. „Gud var min første Tanke“ siger han,
„\emph{Fornuften} min anden, \emph{Mennesket} min tredie og sidste Tanke.
Guddommens Subject er Fornuften, men Fornuftens Subject er Mennesket“. Den
transscendente Guds Existents uden for Mennesket forvandler sig først for ham
til en \emph{Fornuftexistents} ɔ: en Existents i Fornuften, og Guds Væsen
identificeres paa \emph{pantheistisk} Viis med \emph{Fornuftens Væsen}; saa
bliver den religiøse Forestilling seet som en \emph{ubevidst Projection} af det
menneskelige Selv, der i Gud objectiverer sit eget almene Væsen som et fra sig
forskjelligt; og tilsidst bliver, --- idet det \emph{Almenes} Realitet paa
\emph{nominalistisk} Maade benægtes, og Fornuftvæsenet, der var identificeret
med Menneskeslægten som Naturens Afslutning, absorberes af de enkelte
Individer, --- den religiøse Forestilling kun et
% --- p. 100 ---
\emph{Speilbillede af disses Attraa}. I denne Proces giver \emph{Pantheismen
Plads for Atheismen}, Theologien opløses, ligesom før Metaphysiken, i
\emph{Anthropologie}, Gud forsvinder, og kun \emph{Mennesket} bliver tilbage
som det Sande og Virkelige, og det \emph{religiøse Forhold}, som den
\emph{positive} Religion sætter, bliver \emph{Sandhedens yderste Modsætning}.
Kun i overført Betydning tales der endnu om Religion som noget Gyldigt.

Hvad der i Udviklingen af Feuerbachs philosophiske Standpunkt, der betinger
Udviklingen af hans Opfattelse af det Religiøse, er det bevægende Motiv, er
\emph{Interessen for Virkeligheden}. Og medens denne Interesse, som Interesse
for den reale \emph{Naturvirkelighed}, det concrete Menneskelivs Grundlag,
leder hans Polemik mod den abstracte Idealisme, leder den som Interesse for det
\emph{virkelige Menneskelivs ethiske Fordring} hans Polemik mod det positivt
Religiøse.

De Skrifter, der betegne Udviklingsstadierne i Feuerbachs \emph{eiendommelige}
Opfattelse af det Religiøse, ere i det Væsentlige følgende: 1) Philosophie und
Christenthum; das Wesen des Christenthums; (jfr. P. Bayle og de samtidige
kritiske Afhandlinger) 2) Ueber das Wesen des Glaubens im Sinne Luthers; das
Wesen des Religion; Vorlesungen über das Wesen der Religion; 3) Theogonie;
Ueber Gottheit, Freiheit und Unsterblichheit. --- Af disse Skrifter maa der
udførligst tages Hensyn til \emph{das Wesen des Christenthums}, i hvilket
Feuerbachs Opfattelse af det Christelige er udviklet i Sammenhæng. Til dette
Skrift, der netop betegner Overgangen i hans philosophiske Grundanskuelse og
kun forstaaes naar de Stadier, mellem hvilke det danner Overgangen, klart
opfattes, er Skriftet: Ueber Philosophie und Christenthum, den polemiske
Indledning, Afhandlingen: vom Wesen des Glaubens im Sinne Luthers
Efterskriften. I Afhandlingen om Religionens Væsen er det, der i W. d.
Christenthums blev udviklet med særligt Hensyn til Christendommen, udført med
Hensyn til Naturreligionerne, og disses Forhold til Aandsreligionerne er
bestemt. Forelæsningerne over Religionens Væsen er en
% PRINTER'S DEFECTS in the list of Feuerbach's titles, p. 100, both confirmed
% at 300 dpi and read alike by both witnesses, both transcribed as printed:
%   „das Wesen des Religion“ for „der Religion“;
%   „Unsterblichheit“ for „Unsterblichkeit“ (the sort is plainly the same h as
%   in the preceding „lich“, not a k).
% The German titles are set in Fraktur, not antiqua, so none of them takes
% \textit{}; only „das Wesen des Christenthums“ at the foot is letterspaced.
% --- p. 101 ---
videre Udførelse af den sidstnævnte Afhandlings Grundtanker;
\emph{Theogonien} og de dermed sammenhørende Partier af F.'s sidste Værk:
Gottheit, Freiheit und Unsterblichheit, en Udvikling af Religionens Genesis og
af dens Sammenhæng med det Pathologiske i Mennesket: en Sammenhæng, som
allerede Afhandlingen om Religionens Væsen stærkt betoner. F.'s Polemik mod
\emph{Theologien} træder deels frem paa forskjellige Steder i de nævnte
Skrifter, navnlig i de under 1) anførte, deels i hans Kritiker; saaledes i hans
„Beiträge zur Kritik des modernen Afterchristenthums“.

\emph{Grundtrækkene af Feuerbachs Opfattelse af Religionen}, saaledes som den
er given i „\emph{das Wesen des Christenthums}“, ere følgende.

\emph{Da Religionen kun findes hos Mennesket, maa den være begrundet i det, der
er Menneskevæsenets eiendommelige} Kjendemærke. Dette er, almindeligt udtrykt,
Bevidstheden, men Bevidsthed i stræng Forstand er Bevidsthed om sin Slægt, sin
Væsenhed. Dyret er kun Gjenstand for sig som Individ, Mennesket er ogsaa
Gjenstand for sig som Slægt; derfor har det Evne til Tænken, til Viden; det kan
gjøre andre Ting til Gjenstand efter deres væsentlige Natur. Religionen i
Almindelighed, \emph{som identisk med Menneskets Væsen}, er identisk med
Menneskets Selvbevidsthed, med dets Væsensbevidsthed; og naar nu tillige
Religionen er Bevidsthed om \emph{det Uendelige}, saa bliver den dette som
Menneskets Bevidsthed om \emph{sit uendelige Væsen}. Et virkelig endeligt Væsen
har ikke den fjerneste Anelse, endsige Bevidsthed, om et uendeligt Væsen; thi
\emph{Væsenets Skranke er ogsaa Bevidsthedens Skranke}. Bevidsthed om det
Uendelige forudsætter det bevidste Væsens Uendelighed, derfor viser Religionen
tilbage til \emph{Menneskevæsenets Uendelighed}.

Mennesket er Intet uden Gjenstand: gjennem Gjenstanden bliver Subjectet
aabenbart; men den Gjenstand, til hvilken det menneskelige Subject
\emph{væsentlig og nødvendig} forholder sig, den \emph{indre og udvalgte}
Gjenstand, er ikke Andet end dette Subjects \emph{eget}, men
\emph{objectiverede}
% From p. 101 the long Feuerbach citations begin. They are run into the
% paragraph and marked by letterspacing alone, with no quotation marks; they
% are neither indented nor set as a quotation environment.
% --- p. 102 ---
Væsen. En saadan Gjenstand er den \emph{religiøse} Gjenstand; i den bliver
Mennesket derfor kun sig selv bevidst. Ved den sandselige Gjenstand er
Bevidstheden om Gjenstanden adskillelig fra Selvbevidstheden, ved hiin falde de
umiddelbart sammen: Gjenstanden er Menneskets aabenbare Væsen, dets sande
objective Jeg, dens Magt over Mennesket er dets Væsens egen Magt.

\emph{Mennesket kan ikke blive sig et Væsen bevidst, der er høiere end dets
eget}. Hvilken Gjenstand det end bliver sig bevidst, saa bliver det sig
bestandig sit eget Væsen bevidst; det kan ikke bekræfte et Andet uden at
bekræfte sig selv. Idet dets Væsens Kræfter, der constituere Væsenet, ere
Realiteter, saa er enhver Virken af dem, altsaa ogsaa den, ved hvilken de
udtrykke den absolute Gjenstand, den umiddelbare Bekræftelse og Stadfæstelse af
dem selv; enhver Yttring af dem maa være ledsaget af Bevidstheden om dem som
Fuldkommenheder. Hvis Mennesket kunde blive sig selv ɔ: sit Væsen, bevidst som
endeligt ligeoverfor et andet Væsen, der vidstes som uendeligt, saa vilde det i
Bevidstheden begrændse sit Væsen ɔ: negere det: \emph{dets Bevidsthed vilde
altsaa naae ud over dets Væsen}. Men Bevidstheden er kun Væsenets
Selvobjectivering, altsaa dets Selvbekræftelse, dets Realisation af sig selv;
den er netop Selvbekræftelsens høieste Form; hiin Negation er altsaa en
Umulighed. Men hvor faaer da Distinctionen mellem det Endelige og det Uendelige
sin Plads? I Distinctionen mellem \emph{Individ og Slægt}. \emph{Slægtens Væsen
er Individets absolute Væsen}; ud over det kan Individet kun \emph{indbilde}
sig at kunne naae. Ethvert Væsen er uendeligt for sig selv; dets Skranker
existere kun for et andet Væsen udenfor og over det: Væsenets Maal er
Forstandens Maal.

Hvad der gjælder om Bevidstheden overhovedet i Forhold til Gjenstanden, gjælder
selvfølgelig om \emph{de enkelte Aandens Organer}. I det, der er væsentlig
Gjenstand for et enkelt Aandens Organ, forholder Organet sig til sig selv,
objectiverer det sig selv. Saaledes er \emph{Fornuftens}
% --- p. 103 ---
\setcounter{footnote}{0}%
\emph{Gjenstand den for sig objective Fornuft, Følelsens Gjenstand} den for sig
objective \emph{Følelse}. Derfor bliver \emph{Formforskjellen} mellem den
religiøse og den philosophiske Bevidsthed af saa \emph{afgjørende} Betydning:
\emph{netop Formen}, der er Udtryk for Organets Eiendommelighed, bestemmer
\emph{Religionens Væsen}. Det Absolute er \emph{ikke} Religionens Gjenstand
\emph{som saadant}, men \emph{saaledes som} det er Gjenstand for
Gemyttet\footnote{„Gemyttet“ stiller F. paa den ene Side i Modsætning til
\emph{Fornuften} som det Individuelle til det Almeengyldige; paa den anden Side
i Modsætning til „Hjertet“ som det subjectivt vilkaarlige, skrankeløst
Attraaende til den naturmæssige, sin Skranke anerkjendende Følelse og Drift.}
og Phantasien. Hvad den transscendente Speculation over det Religiøse sætter
som det blot Secundære, som Middel og Organ, det har netop Betydningen af det
Primitive, af Væsenet, af Gjenstanden selv.
% FOOTNOTE p. 103: printed mark *). This page is NOT on the OCR-derived list of
% note-bearing pages — found by eye. „Fornuften“ inside the note is Fraktur and
% letterspaced across its line break (For-/nuften, both halves), hence \emph{}.

Men naar det nu siges, at Mennesket i Bevidstheden om den religiøse Gjenstand
forholder sig til sit eget objectiverede Væsen, saa at \emph{Bevidstheden om
Gud er Menneskets Selvbevidsthed}, saa bliver dog, forat det \emph{religiøse
Standpunkt} kan betegnes i sin Eiendommelighed, en nærmere Bestemmelse at
tilføie som væsentlig: nemlig den, at det religiøse Menneske \emph{ikke er sig
dette Forhold til sig selv bevidst}, og at Mangelen paa en saadan Bevidsthed
netop constituerer Religionens eiendommelige Charakteer. \emph{Religionen er
Menneskets første og indirecte Selverkjendelse}. Ligesom Barnet seer sit Væsen:
Mennesket, udenfor sig, saaledes at Mennesket som Barn er sig Gjenstand som et
\emph{andet} Menneske, saaledes henlægger ogsaa, under den historiske
Udvikling, Mennesket først sit Væsen udenfor sig før det finder det i sig; dets
eget Væsen bliver det først Gjenstand som et \emph{andet Væsen}. Den historiske
Udvikling i Religionerne bestaaer derfor deri, at det, der gjaldt de tidligere
Religioner som noget Objectivt, nu anskues som noget Subjectivt, d. v. s. at
det, der anskuedes og tilbades som Gud, nu erkjendes for noget Menneskeligt.
For den
% --- p. 104 ---
sildigere Religion er den tidligere Afgudsdyrkelse: Mennesket har i den tilbedt
sit eget Væsen. Men hver enkelt Religion undtager naturligviis sig selv fra
denne Dom; fordi den har en anden Gjenstand, et andet Indhold, indbilder den
sig at være hævet over de Love, der constituere Religionens Væsen; indbilder
den sig, at dens Gjenstand, dens Indhold, \emph{er over}menneskeligt. Men
Philosophien gjennemskuer Religionens for den selv skjulte Væsen; den viser, at
Modsætningen mellem Gud og Menneske er fuldkommen illusorisk, at den ikke er
Andet end Modsætningen mellem det menneskelige Væsen (Slægten) og de
menneskelige Individer, saa at følgelig \emph{enhver}, ogsaa den høieste,
Religions Gjenstand og Indhold ere aldeles menneskelige.

Naar der indrømmes, at Gud af Mennesket \emph{uundgaaeligen} maa bestemmes ved
menneskelige Prædicater, og at altsaa Bevidstheden forsaavidt i Gud forholder
sig til et menneskeligt Bestemt; men tillige siges, at Gud efter sit Væsen
\emph{i og for sig} er uendeligt høiere end Mennesket, saa forstyrrer denne
Adskillelse mellem hvad Gud \emph{er for Mennesket} og hvad han \emph{er i og
for sig} ikke blot Religionens Fred, men er i og for sig ubegrundet og
uholdbar. I denne Distinction sætter Mennesket sig ud over sig selv, over sit
Væsen, sit absolute Maal; men det er en Illusion. Hiin Distinction har kun
Betydning hvor en Gjenstand virkelig \emph{kan vise sig} anderledes for mig end
den gjør det, eller hvor Gjenstanden er given umiddelbart sandseligt, og derfor
netop ogsaa for andre Væsener end Mennesket; men den har det ikke hvor en
Gjenstand, der udelukkende er for Mennesket, viser sig saaledes for mig som den
viser sig \emph{efter mit absolute Maal}, \emph{som den maa} vise sig for mig.
Svarer Gudsforestillingen til \emph{Slægtens Maal}, saa falder hiin Distinction
bort og Forestillingen er absolut; thi Slægtens Maal er Menneskets absolute
Maal, Lov og Kriterium. Dette fastholder ogsaa \emph{enhver} Religion med
Hensyn til \emph{sin} Gudsforestilling: hiin Distinction er skeptisk og
religiøs. Religionen opgiver sig selv naar den opgiver Guds Væsen; den er ingen
Sandhed mere naar den giver Afkald paa
% EMPHASIS, p. 104: „er over“ is letterspaced and „menneskeligt“, in the same
% word, close-set — the third mid-word break in this batch. As printed.
% --- p. 105 ---
Besiddelsen af den \emph{sande} Gud. Hvad der for Mennesket har Betydning af
det Iogforsigværende, hvad der for det er det høieste Væsen, ud over hvilket
det ikke kan forestille sig noget Høiere: det er netop for det det
\emph{guddommelige} Væsen. Ved denne guddommelige Gjenstand kan der ikke mere
spørges: hvad den er i og for sig. At spørge, om Gud i og for sig er saadan som
han er for mig, er det Samme som at spørge, om \emph{Gud er Gud}; det er at
hæve sig over sin Gud, at gjøre Oprør imod ham. Betvivler man først
\emph{Prædicaternes objective Sandhed} og gjør dem til Anthropomorphismer, saa
maa man consequent ogsaa betvivle disse Prædicaters \emph{Subjects objective
Sandhed}, ogsaa gjøre Subjectet, altsaa Guds Existents, overhovedet Troen paa
en Gud, til en Anthropomorphisme, en fuldkommen menneskelig Forudsætning.
\emph{Subjectets Nødvendighed ligger kun i Prædicaternes Nødvendighed}; thi
Subjectet er kun det personificerede, det existerende Prædicat. Prædicatet er
det egentlig Primitive. \emph{Visheden om Guds Existents afhænger af Visheden
om Guds Qualitet}; kun Prædicatets Realitet er Garantien for Existentsen. For
den Christne er kun den christelige Guds, for Hedningen den hedenske Guds
Existents en Vished. \emph{Gud er Menneskets Væsen, anskuet som absolut
Sandhed}, men han er Menneskets Væsen, \emph{opfattet i en bestemt Qualitet},
der, ved at være Sandheden, bliver den høieste Existents, thi hvad Mennesket
anskuer som sandt, anskuer det umiddelbart som existerende.

\emph{Naar de guddommelige Prædicater ere menneskelige Prædikater, saa maa
altsaa ogsaa det guddommelige Subject være et menneskeligt Subject}. Men her
maa et Phænomen fremhæves, der charakteriserer Religionens inderste Væsen. Jo
menneskeligere det guddommelige Subject er i Væsenet, desto større er
tilsyneladende Differentsen mellem Gud og Mennesket, desto mere nægter
Theologien, Reflexionen over Religionen, Identiteten, og nedsætter det
Menneskelige, saaledes som det \emph{som saadant} er Gjenstand for Menneskets
Bevidsthed. Grunden
% PRINTER'S DEFECT, p. 105: „Prædicater ere menneskelige Prædikater“ — the same
% word spelled with c and then with k within one sentence, confirmed at 300 dpi
% and read alike by both witnesses. Transcribed as printed.
% --- p. 106 ---
hertil er ligefrem den, at da det \emph{Positive} i Anskuelsen af det
\emph{guddommelige} Væsen netop er de \emph{menneskelige} Væsensbestemmelser, saa maa
Anskuelsen af det \emph{Menneskelige}, saaledes som det \emph{directe} er Gjenstand for
Bevidstheden, \emph{kun} være en \emph{negativ}.
% p. 106: the compositor spaces „kun“ and „negativ“ but leaves „være en“ between
% them unspaced; confirmed at 600 dpi. Transcribed as printed, not regularised.
Forat berige Gud maa Mennesket gjøres fattigt; det har \emph{sit Væsen} i Gud,
hvorledes skulde det saa have det i sig selv. Mennesket negerer kun med Hensyn
til sig selv hvad det sætter i Gud. Religionen abstraherer fra Mennesket, fra
Verden; men den kan kun abstrahere fra Phænomenet, ikke fra Væsenet, og derfor
maa Religionen ubevidst sætte alt det i Gud, som den bevidst har negeret, naar
det er et \emph{Væsentligt} og \emph{Sandt}, et Saadant altsaa, fra hvilket den ikke
kan abstrahere, som den ikke \emph{overhovedet} kan negere: saaledes Sandselighed,
Fornuft, Godhed, moralsk Virksomhed.
% p. 106: „og“ between the two spaced words is itself unspaced (ratio 0.90);
% the same holds at pp. 108 and 114, though not at pp. 111, 115, 117, 119.
Gjennem Gud faaer Mennesket dette, der, ved at sættes i Gud, Aandens væsentlige
Gjenstand, allerede indirecte er udtalt som en væsentlig Bestemmelse ved den
menneskelige Natur, i uendelig rigere Maal tilbage. Religionens Hemmelighed er
den, at Mennesket objectiverer sit Væsen, men atter gjør sig til Object for
dette objectiverede, til et Subject forvandlede Væsen. Vel har Mennesket i
Religionen Gud til Øiemed, men Gud har intet Andet til Øiemed end Menneskets
Vel: den guddommelige Virksomhed bliver Middel for dette, altsaa har Mennesket
i Religionen kun sig selv til Øiemed i og gjennem Gud. Og hvad der, ligesom ved
en religiøs Repulsionskrafts Act, sættes som en Virksomhed af Gud, Menneskets
udsondrede Subjectivitet, fra hvem den religiøse Bevidsthed maa lade enhver
Impuls udgaae, det bliver atter ved en religiøs Attractionskrafts Act, idet den
guddommelige Virksomhed skal virke paa Mennesket, i Mennesket, blive Princip
for Menneskets gode Sindelag, for dets gode Handlinger, gjort til en i sit
Væsen menneskelig Virksomhed. Religionens Udviklingsgang bestaaer derfor
nærmere deri, at Mennesket bestandig tilegner sig Mere af det, hvad det før
tilskrev Gud. Hvad der paa det tidligere Trin var
% --- p. 107 ---
Religion, er det ikke længere paa det sildigere: hvad der paa hint var
Atheisme, gjælder paa dette for Religion.

Men naar saaledes Religionen er bestemt som Menneskets indirecte
Selverkjendelse og som dets ubevidste Objectiveren af sit Væsen, saa maa det
Spørgsmaal besvares: i hvilket Forhold den \emph{religiøse} Bevidsthed om Væsenet
staaer til den Bevidsthed, der er tilstede i andre Aandslivets Former, med
andre Ord: \emph{hvad af Menneskets Væsen} der væsentlig giver den religiøse
Bevidsthed Indhold.
% p. 107: the letterspacing stops at „Væsen“; „der væsentlig giver den religiøse
% Bevidsthed Indhold“ is set tight. Transcribed as printed.
Svaret paa dette Spørgsmaal vil fremgaae af en Betragtning af Religionens
Charakteer og væsentlige Standpunkt.

\emph{Religionens væsentlige Standpunkt} viser sig umiddelbart som det
\emph{subjective, praktiske}. Religionens Øiemed er Individets Vel, dets Frelse,
dets \emph{Salighed}. Menneskets Relation til Gud er kun dets Relation til sin
Frelse. \emph{Gud er den realiserede Salighed eller den ubegrændsede Magt til at
virkeliggjøre Menneskets Salighed}. Gud som Gud ɔ: saaledes som han er
Gjenstand \emph{for Religionen}, er væsentlig \emph{kun} Gjenstand for Religionen,
ikke for Philosophien, for Gemyttet, ikke for Fornuften, for Hjertets Trang,
ikke for Tankens Frihed, kortsagt: et Væsen, der \emph{ikke} udtrykker det
\emph{theoretiske, objective}, men det \emph{praktiske, subjective} Standpunkts
Eiendommelighed. Dette bliver ogsaa tydeligt derved, at Religionen knytter
Forbandelse og Velsignelse, Fordømmelse og Salighed, til sine Lærdomme, til
Troen paa dem. Den appellerer altsaa ikke til Fornuften, men til Gemyttet,
Lyksalighedsdriften, til Frygtens og Haabets Affect. Tvivlen, den theoretiske
Friheds Princip, bliver her en Forbrydelse, Tvivlen om Gud den høieste
Forbrydelse. Men hvad jeg ikke kan eller tør betvivle uden at føle mig
foruroliget i mit Gemyt, det er ikke Theoriens Sag, men Samvittighedens, ikke
et Fornuftvæsen, men et Gemytsvæsen.

Naar nu Religionens Standpunkt er det praktiske eller subjective, saa ligger
deri, at \emph{kun} det \emph{praktiske}, efter sine bevidste, physiske eller
moralske, Øiemed handlende Menneske, der alene betragter Verden i Relation til
disse Øiemed,
% --- p. 108 ---
ikke i og for sig, gjælder den for det \emph{hele} væsentlige Menneske, og det er
i Relation til Mennesket i denne Bestemthed, at Religionens høieste Væsen
bestemmes som et \emph{Gemyttets} Væsen, som et Væsen, der tilfredsstiller
Gemyttets og Phantasiens Trang. Alt det, der ligger bag den praktiske
Bevidsthed, men er den væsentlige Gjenstand for Theorien --- for \emph{Theorien} i
den oprindeligste og almindeligste Betydning, i Betydning af objectiv Anskuelse
og Erfaring, Fornuft, Videnskab overhovedet --- forholder sig ikke væsentlig
til den religiøse Bevidsthed. Religionen mangler, idet den er udelukkende
praktisk, subjectiv, den sande Anskuelse af Universet, af det virkelige
Uendelige, den sande Bevidsthed om Slægten. Og idet nu Naturens og Menneskets
sande almindelige Væsen som et saadant, der ikke kan negeres, gjør sig
gjældende for Aanden, saa søges vel hiin Mangel erstattet i Gud, i et særskilt,
personligt Væsen: Naturens og Menneskets kun for det theoretiske Øie
tilgængelige Væsensbestemmelser lægges uvilkaarligt over i Gud, men deels
saaledes, at de kun udgjøre den religiøse Bevidstheds yderste
Tilknytningspunkt, „Religionens mathematiske Punkt“; deels saaledes, at de
omformes, og i denne Omformning tilegnes af Gemyttet og Phantasien. Idet
Religionen antager Menneskehedens og Naturens tilsyneladende, overfladiske
Væsen for deres sande, indre Væsen, og forestiller sig deres sande esoteriske
Væsen som èt andet Væsen, metamorphoserer den under Personificationen
Menneskets og Naturens sande Væsen til en \emph{overmenneskelig} og
\emph{overnaturlig Personlighed}, ubegribelig for Erkjendelsen, virkende gjennem
\emph{Underet} og gjennem dette tilfredsstillende det religiøse Individs praktiske
Trang i Modsigelse med Fornuftloven.
% p. 108: „èt andet Væsen“ — the accent is printed and is a GRAVE (è), confirmed
% at 600 dpi on the crop; 19th-c. Danish „èt“ = the numeral one, against the
% article „et“. Not the sporadic scan dirt described in the batch brief.

Religionens oprindelige Sondring af det Guddommelige og Menneskelige er
uvilkaarlig og naiv; den identificerer oprindelig ligesaa umiddelbart Gud med
Mennesket som den adskiller ham fra det; den bevæger sig i Phantasiens
quantitative Forskjelle, gjør ikke Forskjellen til en qualitativ og væsentlig.
Men idet Reflexionen vaagner, saa bliver Ad\-%
% --- p. 109 ---
% (no space: the word „Adskillelsen“ is broken across pp. 108/109)
skillelsen forsætlig og udstuderet, netop for at fortrænge den Bevidsthed om
Identiteten, der begynder at dæmre i den religiøse Bevidsthed. Den efter
Begrebet første Maade, paa hvilken Reflexionen sikkrer sig Gud som et fra
Mennesket forskjelligt Væsen, er ved \emph{Beviserne for Guds Tilværelse}. Disses
Øiemed er nemlig just at sondre Gud ud fra Mennesket, ved at hævde ham en
Existents som en Ting i og for sig, der ikke blot skal have Væren i vor Tro,
vor Følelse, vor Tanke, men ogsaa have en herfra forskjellig real Væren, en
Tilværelse \emph{udenfor} os, og dette: \emph{udenfor}, maa tages \emph{egentligt}; thi
ellers bliver selve Existentsen en uegentlig Existents. En saadan Existents
udenfor os vilde imidlertid ikke være Andet end en \emph{sandselig} Væren. Men
Guds Væren har ikke den sandseligt-reale Værens Eiendommelighed: at jeg
afficeres uvilkaarligt af den, at den er, selv naar jeg ikke er, ikke tænker
den, ikke føler den; Gud er kun til for Mennesket, idet det hæver sig over det
sandselige Liv, idet det føler, tænker og troer Gud. Men idet han \emph{kun} er
Gjenstand for \emph{Mennesket}, saa vil dette simpelthen sige: Gud er kun, idet han
føles, tænkes og troes. Den Existents, man vil hævde Gud, vilde blive en real
Væren, der dog ikke var real: en \emph{aandelig} Væren kalder man den som Udflugt.
Men aandelig Væren er netop kun at være følt, at være tænkt, at være troet.
Guds Væren er altsaa en modsigende Mellemting mellem sandselig Væren og tænkt
Væren, en sandselig Væren, der, ved at berøves alle Sandselighedens
Bestemmelser, reduceres til en vag Existents overhovedet, og derved ophæver sig
selv, da der til Existents hører fuld, bestemt Realitet. En \emph{nødvendig Følge
af denne Modsigelse er Atheismen}. Guds Existents, der ikke betegner indre
Realitet, Sandhed, men ydre Existents, har en empirisk Existentses Væsen, men
uden at have dens Kjendetegn; den er i og for sig en Erfaringssag, og dog i
Virkeligheden ingen Gjenstand for Erfaringen. Den opfordrer selv Mennesket til
at søge sig i Virkeligheden, og idet det finder Erfaringen i Modsigelse med de
Forestillinger, den har vakt, er det fuldkommen berettiget til
% --- p. 110 ---
at nægte denne Existents. Kun Sandsningen, ikke Fornuften, kan give Beviset for
empirisk Existents, og i Beviserne for Guds Tilværelse har Væren netop
Betydningen af en saadan Existents. For Religionen bliver Troen paa Guds
Existents uden Hensyn til den indre Qualitet, det aandelige Indhold, tilsidst
uundgaaelig Hovedsagen. Guds Existents i og for sig sætter sig fast som religiøs
Sandhed, og til den kan, som Erfaringen viser, de raaeste og vanhelligste
Forestillinger knytte sig.

Paa Religionens Standpunkt er det \emph{Indbildningskraften}, Stedet for den
fraværende, fra Sandserne unddragne, men dog efter Væsenet sandselige
Existents, der løser Modsigelsen i en Existents, som paa eengang skal være
sandselig og usandselig. Kun Indbildningskraften bevarer for Atheisme; thi i
den har Existentsen sandselige Virkninger. Hvor Guds Existents er en levende
Sandhed, en Indbildningskraftens Sag, der troes ogsaa paa Tilsyneladelser af
Gud. Hvor derimod den religiøse Indbildningskrafts Ild slukkes, og hvor de
sandselige Virkninger eller Phænomener, der nødvendigviis ere forbundne med en
i og for sig sandselig Existents, ophøre, der bliver Existentsen til en død og
selvmodsigende Existents, der redningsløs hjemfalder til Atheismens Negation.

Den samme Modsigelse, der gjør sig gjældende i Forestillingen om Guds
Existents, gjør sig gjældende i Bestemmelsen af \emph{Guds Væsen overhovedet}.
Alle de positive Bestemmelser ved dette ere menneskelige Bestemmelser; men idet
Bestemtheden i Bestemmelserne negeres, gjør Reflexionen dem til
ikke-menneskelige. Men ved denne Negation opløses Begrebet; der bliver kun en
negativ, indholdsløs, selvmodsigende Forestilling tilbage; thi netop den
Bestemthed er negeret, der gjorde Bestemmelsen til hvad den er, og kun dette er
et realt Begreb; det derimod, der skal gjøre den til et væsentlig Andet, er
objectiv Intet, subjectiv blot en Indbildning.

Dette udfører nu Feuerbach i „das Wesen der Christenthums“ med Hensyn til de
christelige Hoveddogmer.
% PRINTER'S ERROR, transcribed as printed: „das Wesen der Christenthums“ for
% „des Christenthums“. Confirmed at 600 dpi; both witnesses read „der“. The
% same error recurs at pp. 113 and 117, so it is the book's settled form. The
% title is set in FRAKTUR here, not antiqua — as are all the German titles in
% this batch (pp. 113, 114, 117, 118).
Modsigelsen i Christendommen, „Kjærlighedens Religion“, seer
% --- p. 111 ---
han afgjørende concentreret i \emph{Modsætningen mellem Tro og Kjærlighed}.

I Kjærligheden seer han nemlig en Aabenbaring af Religionens \emph{skjulte}
Væsen, i Troen det, der constituerer dens \emph{bevidste} Form. Kjærligheden
identificerer Mennesket med Gud og Gud med Mennesket, og derfor Mennesket med
Mennesket; Troen adskiller Gud fra Mennesket, men derfor ogsaa Mennesket fra
Mennesket; thi Gud er intet Andet end Menneskets mystiske Slægtbegreb: Guds
Adskillelse fra Mennesket bliver derfor Menneskets Adskillelse fra Mennesket,
Opløsningen af det fælleds Baand. Idet Troen gjør Gud til et \emph{særeget} Væsen
og opfatter ham som \emph{Person}, sondres Forholdet til Gud fra Forholdet til
Menneskene, og der maa uundgaaelig blive en Conflict mellem de to Forhold,
mellem Religionen og Moralen, og da Forholdet til Gud som det høieste gjør Krav
paa alle Menneskets Kræfter og dets hele Interesse, maa Moralen vige for
Religionen, Kjærligheden tilintetgjøres af Troen. Troen maa hæve sig over den
naturlige Morals Love; der maa gives Handlinger, i hvilke Troen kommer tilsyne
i Forskjel fra og i Modsætning til Moralen; Pligterne mod Gud maae collidere
med de simple menneskelige Pligter, og Troen privilegerer Alt. Selv naar
Christendommen siger: Gud er Kjærlighed, saa bliver der, da Kjærligheden kun
bliver Prædicat, og Sætningen ikke tillige vendes om, saa at den kommer til at
lyde: Kjærligheden er Gud, en Plads for den guddommelige Personlighed som
saadan, og Subjectet: den guddommelige Personlighed, bliver Hovedsagen,
Prædicaterne, de ethiske Bestemmelser, en blot Bisag. Men idet jeg tænker Gud
som Subject, i Adskillelse fra Prædicatet, bliver der et Rum tilovers for
Ukjærligheden, og det bliver uundgaaeligt, at jeg snart giver Subjectet, snart
Prædicatet Overvægten. Denne Modsigelse har Christendommens Historie
tilstrækkelig constateret. --- Den christelige Kjærlighed er, som christelig, en
begrændset Kjærlighed. Men en Kjærlighed, der er begrændset ved Troen, er en
usand Kjærlighed. Den sande Kjærlighed, den frie, universelle Kjærlighed,
kjender ingen anden Lov end
% --- p. 112 ---
sin egen og Fornuftens. Den er Intelligentsens og Naturens universale Lov; den
er intet Andet end Realisationen af Slægtens Enhed ad Sindelagets Vei. Den har
saaledes en Selvgyldighed uafhængig af det Religiøse, behøver ikke at støttes
paa dette. Idet den universale, i sig begrundede Kjærlighed bliver Principet,
forsvinder Christus for Bevidstheden om Slægten, thi Christus er selv kun et
Billede, under hvilket Folkebevidstheden forestillede sig Slægtens Enhed, og
den sande universelle Kjærlighed kan kun grunde sig paa Slægtens, paa
Intelligentsens Enhed, paa Menneskehedens Natur; kun da er den en grundig, i
Principet betrygget, fri Kjærlighed, thi den støtter sig paa Kjærlighedens
Oprindelse, fra hvilken ogsaa Christi Kjærlighed stammede. Denne var selv en
afledet: Christus elskede ikke af sig selv, men i Kraft af Menneskehedens
Natur. Slægtens, Intelligentsens Lov er den, at vi skulle elske Mennesket
\emph{for Menneskets Skyld}; Kjærligheden taaler intet Mellemled.

Skjøndt Christendommen anerkjender Kjærligheden, kommer den dog atter til at
fornægte den, fordi den ikke har givet Kjærligheden fri, ikke opfattet den
absolut, og det har den ikke kunnet gjøre, fordi den er Religion, og
Kjærligheden derfor underkastes Troens Herredømme. Kjærligheden er kun
Christendommens exoteriske Lære, men Troen dens esoteriske; Kjærligheden er kun
dens Moral, men Troen dens Religion.

I Modsigelsen mellem Tro og Kjærlighed ligger Nødvendigheden af at gaae ud over
Christendommen, overhovedet over Religionens eiendommelige Væsen.
\emph{Religionens Indhold og Gjenstand er en fuldstændig menneskelig: Theologiens
Hemmelighed er Anthropologien.}
% p. 112 is a whole page of Feuerbach quotation carrying no quotation marks at
% all; the citation is marked by letterspacing alone and is run into the
% paragraph, so no quotation environment is used. The letterspacing of the
% thesis runs unbroken across the line breaks „menne-skelig“ and
% „Anthropo-logien“. A modern two-word pencil annotation stands in the left
% margin of this leaf; it is the reader's, not the book's, and is not
% transcribed.
Herom har Religionen ingen Bevidsthed, men Opgaven er at hæve Illusionen, at
komme til den Erkjendelse, at Bevidstheden om Gud ikke er Andet end
Bevidstheden om Slægten; at Mennesket kun kan og skal hæve sig over sin
Individualitets Skranker, men ikke over sin Slægts Love og positive
Væsensbestemmelser; at Mennesket ikke kan tænke, ane, forestille sig, føle,
troe, ville, elske og tilbede
% --- p. 113 ---
noget andet Væsen end Menneskeslægtens Væsen --- med Indbegreb af Naturen, der
hører med til Menneskets Væsen, ligesom Mennesket til Naturens --- som det
absolute Væsen.

\emph{Mennesket, der i Religionen i og for sig er det Første, maa sættes og
udtales som det Første.} Alle Menneskets moralske Forhold til Mennesket maae
sættes som i og for sig gyldige, hellige og derved som religiøse. De kunne, som
substantielle Forhold, kun begrundes ved sig selv: en Begrundelse af Moralen
ved Theologien er en Illusion, en Cirkelbevægelse. Det Rette, Sande og Gode har
overalt sin Hellighedsgrund i sig selv, i sin Qualitet. Hvor det er Alvor med
Ret og Sædelighed, der ere de i og for sig guddommelige Magter. Have de ingen
Grund i sig selv, saa gives der heller ikke nogen indre Nødvendighed for dem:
de ere saa priisgivne til Religionens grundløse Vilkaarlighed.

I den selvbevidste Aands Forhold til Religionen gjælder det altsaa kun om
Tilintetgjørelsen af en Illusion, men af en Illusion, der virker
grundfordærvelig paa Menneskeheden, og baade berøver Mennesket Kræfter til det
virkelige Liv og Sandsen for Sandhed og Dyd; thi selv Kjærligheden, i og for
sig det inderligste, sandeste Sindelag, bliver ved Religionen til et blot
tilsyneladende, illusorisk, idet den religiøse Kjærlighed kun elsker Mennesket
for Guds Skyld, altsaa kun tilsyneladende Mennesket, men i Virkeligheden kun
Gud, og i Gud egoistisk sig selv. Men hiin Illusion tilintetgjøres og Sandheden
bliver aabenbar, naar vi kun vende de religiøse Forhold om, naar vi bestandig
opfatte det, Religionen sætter som Middel, som Øiemed; det, den gjør til det
Underordnede, til Bisag og Betingelse, som Hovedsag og Aarsag.
% p. 113: a small raised speck stands immediately before „om“ in „vende de
% religiøse Forhold om“. The same speck recurs at p. 117 (before „denne“ and
% between „Grundlaget“ and „for“) and p. 118 (before „Opfattelsen“). Four
% instances in six pages, all at the same height and all of the same round
% shape: read as set-off / ink specks in this copy, not as type, and therefore
% not transcribed. Both OCR witnesses pick them up (as ‘ or ').
Det er altsaa i Fornuftens, i Menneskelighedens og Sædelighedens Interesse, at
de religiøse Forestillinger opgives; netop hiin Interesse er Atheismens Kilde og
denne Atheisme negerer derfor ikke Gudsforestillingens reale, positive Indhold,
men kun den usande Forestilling om Gud som et særligt personligt Væsen, og hvad
der er uadskilleligt fra denne Forestilling. ---

Hvad der i „\emph{das Wesen der Christenthums}“ er
% p. 113: here the title IS letterspaced (it is not at pp. 110 and 117). The
% signature numeral „8“ at the foot of p. 113 is not transcribed.
% --- p. 114 ---
udviklet om det Religiøse med særligt Hensyn til det Christelige, bliver i
Afhandlingen „\emph{om Religionens Væsen}“ fuldstændiggjort, idet der nærmere
tages Hensyn til de lavere Religionsformer, og idet det almene Forhold mellem
\emph{Naturreligionerne} og \emph{Aandsreligionerne} bestemmes.

Udgangspunktet danner den Sætning, \emph{at Menneskets Afhængighedsfølelse er
Religionens Grund}, Religionen altsaa den bevidste, udtalte
Afhængighedsfølelse. Som den \emph{oprindelige Gjenstand} for denne
Afhængighedsfølelse bestemmes saa \emph{Naturen}, og Naturen sættes ikke blot som
Religionens \emph{oprindelige Gjenstand}, men ogsaa som den \emph{blivende Grund},
den vedvarende, om end skjulte Baggrund for den, som Grunden til Troen paa Guds
\emph{objective} Existents udenfor Tanken og Fornuften, som Grundlaget for de
Egenskaber, der udtrykke det guddommelige Væsens Forskjel fra det menneskelige.
Gud, tænkt som Ophav til Naturen, forestilles vel som et fra Naturen
forskjelligt Væsen, men dette Væsens virkelige Indhold er kun Naturen. Alle
Guds Egenskaber, der gjøre ham til et objectivt, virkeligt Væsen, ere
Egenskaber, der ere abstraherede fra Naturen, der forudsætte den, udtrykke den,
og som altsaa falde bort, naar Naturen falder bort. Vel bliver der da endnu
tilbage et Væsen, et Indbegreb af Egenskaber som Uendelighed, Enhed,
Nødvendighed, Evighed, men dette Væsen er ikke Andet end Naturens
\emph{abstracte} Væsen, der, som saadant, forudsætter dens concrete, sandselige,
virkelige Væsen.

Men naar der nu siges, at Religionens første og væsentlige Gjenstand er
Naturen, saa er dette, vel at mærke, at forstaae saaledes, at Naturen, selv i
Naturreligionerne, \emph{ikke} er Gjenstand \emph{som Natur}, men som \emph{personligt,
levende, fornemmende, menneskeagtigt Væsen}. Mennesket adskiller sig oprindelig
ikke fra Naturen, følgelig heller ikke Naturen fra sig: det gjør derfor de
Fornemmelser, som en Naturgjenstand vækker i det, umiddelbart til
Beskaffenheder ved Gjenstanden. \emph{Naturvæsenet} bliver saaledes til et
\emph{Gemytsvæsen}, ɔ: til et subjectivt, menneskeligt Væsen.
% --- p. 115 ---
Navnlig er det Naturens Foranderlighed, der bringer Mennesket til at opfatte
den som et menneskeligt, vilkaarligt Væsen. Denne Opfattelse aabenbarer sig i
\emph{Offret}, Naturreligionernes væsentligste Act. Og i Offret concentrerer sig
overhovedet hele Religionens Væsen: thi Offrets Grund er
Afhængighedsfølelsen, men dets Resultat og Øiemed Selvfølelsen. \emph{Og saaledes
er Følelsen af Afhængighed af Naturen vel Religionens Grund; men Ophævelsen af
denne Afhængighed er Religionens Øiemed.} Modsigelsen mellem Villen og Kunnen,
Ønsken og Opnaaen, Forestilling og Virkelighed, er Religionens Forudsætning;
det Væsen, i hvilket denne Modsigelse er hævet, er det \emph{guddommelige} Væsen,
og dette Væsen er Gjenstand for Bøn, bestemmeligt ved Bønnen og Offret. Men
idet Naturen bliver Religionens Gjenstand som metamorphoseret til et Gemyttets,
Forestillingens, Phantasiens Væsen, saa bliver Naturreligionen den iøinefaldende
Modsigelse mellem Forestilling og Virkelighed, Indbildning og Sandhed. Denne
Modsigelse undgaaer Religionen kun, idet den gjør sin Gjenstand til en
\emph{usynlig}, overhovedet \emph{usandselig}, til et Væsen, der kun er Gjenstand
for \emph{Troen}, Forestillingen, Phantasien, kortsagt for \emph{Aanden}, altsaa i og
for sig selv et \emph{aandeligt} Væsen. Naar Mennesket fra et blot physisk Væsen
bliver til et \emph{politisk}, overhovedet til et Væsen, der adskiller sig fra
Naturen og concentrerer sig paa sig selv, saa bliver ogsaa dets Gud fra et blot
\emph{physisk} Væsen til et \emph{politisk}, fra Naturen forskjelligt Væsen. Hvor
Mennesket, gjennem Samfundslivet, bliver til et Væsen, der bestemmer sig ved
Grundsætninger, Visdomsregler, Fornuftlove, altsaa et tænkende, forstandigt
Væsen; hvor andre Magter end Naturmagter: politiske, moralske, abstracte
Magter, blive Gjenstande for dets Afhængighedsfølelse, der viser ogsaa Verden,
Naturen, sig for det som et af \emph{Forstand og Villie} afhængigt og bestemt
Væsen. Hvor Mennesket med Villie og Forstand hæver sig over Naturen, bliver
Supranaturalist, der bliver ogsaa Gud et \emph{supranaturalistisk} Væsen,
\emph{Herre over Naturen}, og endelig, hvad der er
% p. 115: „Væsen,“ at the head of the last line is set tight between two spaced
% phrases; confirmed at 600 dpi. The signature „8*“ at the foot is not
% transcribed.
% --- p. 116 ---
Betingelsen for i stræng Forstand at være Herre: \emph{Skaber af Naturen}. Den
menneskelige Bevidstheds og Aands Enhed, der omfatter hele den ydre Verdens
Mangfoldighed, er det, der fører til at statuere \emph{Guds Enhed}. Den egentlige
\emph{Theismus} opstaaer kun der, hvor Mennesket kun refererer Naturen \emph{til sig}
og gjør denne Relation til dens Væsen, altsaa gjør sig til dens sidste Øiemed.
Hvor Naturen har \emph{sit Øiemed} udenfor sig, har den nødvendigviis ogsaa sin
\emph{Grund} og \emph{Begyndelse} udenfor sig; hvor den kun er til \emph{for et andet
Væsen}, der er den nødvendigviis ogsaa til \emph{ved et andet Væsen}, og det et
saadant, der ved Frembringelsen tilsigtede Mennesket som nydende Naturen. Den
Lære, at Gud er Verdens Skaber, har kun sin Grund og Betydning i den Lære, at
Mennesket er Skabelsens Øiemed. Og Naturens Skaber er kun Betingelsen for det
over Naturen hævede Menneskes Skaber, kun han er Gud og Skaber \textit{sensu
eminenti}, ligesom Religionens egentlige Forsyn er det \emph{specielle} Forsyn
\emph{for Mennesket}.
% p. 116: „sensu eminenti“ is the only antiqua in this batch, and it is set
% tight (not letterspaced). Every German title in pp. 106--119 is Fraktur.

Det \emph{aandelige} Væsen, som Mennesket, idet det fra Naturreligionernes Trin er
skredet frem til Aandsreligionernes, sætter over Naturen og lader være dennes
Forudsætning, er ikke Andet end \emph{selve Menneskets aandelige Væsen}, i den
\emph{religiøse} Bevidstheds Opfattelse, men som af den Grund viser sig for
Mennesket som et \emph{andet}, fra Mennesket \emph{forskjelligt}, uforligneligt
Væsen, fordi det bliver gjort til \emph{Aarsag til Naturen}, til Aarsag til
Virkninger, som den menneskelige Aand, den menneskelige Forstand og Villie,
ikke kan frembringe; af den Grund altsaa, at Mennesket med dette aandelige
menneskelige Væsen tillige forbinder Naturens fra det menneskelige Væsen
forskjellige Væsen. Medens Naturens Væsen i denne Forbindelse
menneskeliggjøres, betinger dens Væren i objectiv Sondring fra Mennesket
Objectiveringen af det menneskelige Væsen som guddommeligt.

Troen paa en Gud bliver altsaa enten \emph{Troen paa Naturen}, paa det objective
Væsen, \emph{som et menneskeligt, subjectivt, Væsen}, eller \emph{Troen paa det
menneskelige Væsen som Naturens Væsen}. Hiin Tro er
% --- p. 117 ---
\emph{Naturreligion}, Polytheisme, denne \emph{Aandsreligion}
(„Aand-Menneske-Religion“), Monotheisme. Religionens første Grundsætning lyder:
Jeg er Intet mod Naturen, Alt er Gud mod mig (Fetichismens Standpunkt,
Grundlaget for Polytheismen). Religionens Slutningssætning lyder derimod: Alt
er Intet imod mig, Alt er kun Middel for mig. Men dette kan Naturen kun være,
naar den ikke er ved sig selv, men ved Gud. I \emph{Underet} er Religionens
Øiemed opfyldt paa sandselig populær Maade; Menneskets Herredømme over Naturen,
dets Guddom, er i det en iøinefaldende Sandhed. \emph{Gudmennesket er
Aabenbaringen af Underets hemmelige Tanke: at Mennesket er Gud.}

Gud er et \emph{religiøst} Ord, et religiøst Object og Væsen, et Object, hvis
Tilværelse kun er given med Religionens Tilværelse, hvis Væsen kun er givet med
Religionens Væsen, som altsaa ikke kan existere udenfor Religionen, og i hvilket
der ikke er indeholdt Mere objectivt end subjectivt i Religionen. Gud udtrykker
kun Religionens Væsen: Phantasien og Gemyttet. Religionens Gjenstand,
\emph{saaledes som} den er Gjenstand for den, existerer ikke i Virkeligheden, men
er kun Object for Troen, der forestiller sig det, der ikke er, men attraaes at
være, som værende. --- Som Menneskenes Ønsker ere, saaledes ere deres Guder.
Grækerne havde begrændsede Ønsker og begrændsede Guder, de Christne en
ubegrændset, over al Naturnødvendighed ophøiet, overmenneskelig, transscendent
Gud og absolut ubegrændsede, phantastiske Ønsker. Deres Grundønske er Nydelsen
af en uendelig, ubegrændset, uudsigelig, ubeskrivelig Salighed. ---
\emph{Saligheden og Guddommen ere Eet.} ---

I Afhandlingen om Religionens Væsen er, som det fremgaaer af det Foregaaende,
ikke blot Synskredsen udvidet ved det Hensyn, der er taget til
Naturreligionerne, men Bestemmelsen af det Religiøse er undergaaet en væsentlig
Modification. I „das Wesen der Christenthums“ blev vel Religionens Standpunkt
opfattet som det ensidigt \emph{praktiske}, og Bestemmelsen af „\emph{Gemyttet}“ som
Religionens Organ forberedte de sildigere dragne Consequentser, men der laae
% --- p. 118 ---
dog i Opfattelsen af Religionen som Menneskets indirecte \emph{Selverkjendelse},
naar den sattes i Forbindelse med Opfattelsen af Slægtens, det
Almeenmenneskeliges, Forhold til Individet, og af dens Forhold til Fornuften,
en Tilknytning for et theoretisk, rationelt Element, og dette kom ogsaa
forsaavidt til Gyldighed som Guds Væsen ogsaa fik Bestemmelsen af
„Forstandsvæsen“, om end denne Bestemmelse kun skulde være et yderste
Tilknytningspunkt for den religiøse Bevidsthed, og ikke udtrykke det, der var
Religionens egentlige Kjærne. Nu derimod falder i Opfattelsen af Religionen
Accenten langt mere skarpt og exclusivt paa den Side i Menneskenaturen, der, i
sin Sondring fra det Theoretiske, bliver det \emph{egoistisk Driftagtige, det
Pathologiske, Alogiske.} Som forbindende Overgang fremtræder i Afhandlingen:
\emph{vom Wesen des Glaubens im Sinne Luthers}, Opfattelsen af \emph{Troen} som
Udtryk for \emph{den indirecte Selvkjærlighed}. Men endnu skarpere end i
Afhandlingen om Religionens Væsen, og med endnu fuldstændigere Tilbageskyden af
det rationelle Moment, fremtræder Forandringen i Opfattelsen af det Religiøse i
„\emph{Theogonie}“, og i Feuerbachs sidste Skrift: „\emph{Gottheit, Freiheit und
Unsterblichheit}“.
% PRINTER'S ERROR, transcribed as printed: „Unsterblichheit“ for
% „Unsterblichkeit“. The sort after the „ch“ ligature is unmistakably an h
% (no k-leg) at 600 dpi, and both witnesses read it so. The whole title, and
% „vom Wesen des Glaubens im Sinne Luthers“, are Fraktur, letterspaced; the
% Luther title carries no quotation marks.

I „\emph{Theogonie}“ bliver det, under mange Vendinger, udviklet, hvorledes
\emph{Ønsket}, den irrationelle, egoistiske Attraa, er \emph{Gudernes første Grund},
Kilden til Troen paa dem; at Guderne \emph{ikke} ere Fornuftvæsener, men
væsentlig Udtryk for Attraaen og Ønsket. Gud er Supplementet til Menneskets
mangelfulde Væsen, til dets, i Modsigelse med dets Ønsker, saa endelige,
begrændsede Evne; Visheden om hans Væren hviler paa Selvkjærlighedens Magt, der
ligeoverfor Virkelighedens pinlige Begrændsning fremtvinger en Forestilling om
denne Begrændsnings Ophævelse, der er saa rodfæstet i Menneskets Attraa, at den
uundgaaelig faaer Virkelighed for det. For den guddommelige Magt forsvinder den
pinende Følelse af Begrændsningen. Gud er og kan, hvad Mennesket ikke er eller
kan, men gjerne vilde være og kunne. Ønsket har, som \emph{theogonisk} Magt,
Forrang for Frygten:
% --- p. 119 ---
det er ikke Frygten i og for sig, men Frygtens Ønske: at der Intet skal være at
frygte, der skaber Guder. Men det Ønske, der sammenfatter alle Ønsker, er
\emph{Lyksalighedsønsket}; i dette Ønske er derfor Gudernes Oprindelse,
Religionens Princip og Grundvæsen at søge, og ved Lyksalighedsønskets Natur og
Grændse bestemmes Gudernes Natur og Grændse.

\emph{Ønsket selv som saadant} stammer fra \emph{Mennesket}; men Ønskets
\emph{Gjenstand} stammer fra \emph{den ydre Natur}, fra det ved Sandserne Givne; thi
Mennesket har oprindelig ingen tomme, supranaturalistiske, phantastiske Ønsker;
nei! Gjenstandene for dets Sandser ere ogsaa Gjenstande for dets Ønsker. Netop
derfor ere Guderne som „\emph{Ønskevæsener}“, tillige \emph{Naturvæsener}. Men dette
er ikke saaledes at forstaae, at Guderne som saadanne ere forgudede eller
personificerede Naturkræfter eller Naturlegemer; de ere Udtryk for \emph{Naturens
Virkninger paa Gemyttet og Indbildningskraften}, personificerede,
selvstændiggjorte, objectiverede \emph{Følelser, Affecter}, men Affecter, der ere
bundne til de Naturlegemer, ved hvilke de vækkes eller bevirkes. Saaledes
skinner det \emph{Menneskelige} igjennem i Bestemmelsen af Naturvæsenet, og det
\emph{menneskelige Ønske}, der henter sit Indhold \emph{fra Naturen}, potentseres,
idet Mennesket adskiller sig som det \emph{Høiere} fra Naturen, til det
\emph{supranaturalistiske Ønske}, til \emph{Salighedsønsket}, i hvilket Guddommens
Mysterium først aabenbares. \emph{Salighedsønsket, Selvkjærlighedens intensiveste
Udtryk}, er \emph{Christendommens} Grundønske; den christelige Gud, der opfylder
dette Ønske, er den personificerede Salighed, absolut hævet over Naturens og
Naturnødvendighedens Skranke.

Det, at Menneskene \emph{overalt} have forestillet sig Guder, men som
\emph{menneskelige} Væsener, er det historiske, factiske Beviis for, at Mennesket
i Realiteten forestiller sig det \emph{menneskelige} Væsen som guddommeligt, og at
det menneskelige Væsen er Grundforudsætningen for Forestillingen om en Guddom,
at det er Urvæsenet, der ligger til Grund for
% --- p. 120 ---
\setcounter{footnote}{0}%
Guderne. Som \emph{høieste} Væsen kan jeg kun tilbede hvad der udtrykker
\emph{mit eget} Væsen, men i en Grad og Fuldkommenhed, der fattes mig, men netop
derfor henriver mig til Tilbedelse og Beundring. Gud er vel det
\emph{overmenneskelige}, det uendelige Væsen, men, vel forstaaet! det
\emph{uendelig menneskelige}, det \emph{overmenneskeligt menneskelige} Væsen, ---
et Væsen, der er mere, uendelig mere Menneske end Mennesket selv. Den christelige
Gud er ligesaafuldt som den hedenske et menneskeligt Væsen, kun af en anden Art
end denne, fordi ogsaa den Christne er et Menneske af en anden Art end
Hedningen. ---

Ved dette Resultat slutter dette Skrift sig sammen med de tidligere, og idet det
polemisk tager Hensyn til de forvanskende Omfortolkninger af Religionen, udtaler
det sig imod dem i disse Ord: „Der gives kun een sand, een ærværdig Gud; --- det
er den umiddelbare, selvvirksomme, selv talende, selv lysende, selv lynende, selv
tordnende, selv regnende, den gamle Verdens Gud. Enten denne eller ingen Gud!
Hvor den gamle Verdensanskuelse er forladt, hvor der gives en selvstændig Stats-
og Retslære, en selvstændig Naturlære, der ikke mere i Vinden seer Guds Aande, i
Tordenens Guds Stemme, i Lynet Guds Ild; der er \emph{Atheismen} ikke blot en
videnskabelig, men en sædelig Sandhed og Nødvendighed“. ---

Ved Hjælp af det givne Overblik over Udviklingen af Feuerbachs Theorie af
Religionen, vil det, der i denne Theorie er utilfredsstillende, let kunne gjøres
tydeligt. Der er hos Feuerbach saare Meget, man ikke vil kunne nægte
Anerkjendelse. Der er hos ham en stor Klarhed i det Negative: i Polemiken mod
Theologien og mod de philosophiske Blandingsformer, i hvilke Philosophie og
Christendom søgtes sammensmeltede til en Enhed. Der er hos ham et skarpt Blik for
det, der er \emph{fælleds for alle de positive} Religionsformer, og en skarp
Opfattelse af de Phænomener, der komme frem indenfor de positive Religioners
Omraade\footnote{Dette \emph{Phænomenale, Factiske}, opfatter saa F. som det
\emph{nødvendige} Udtryk for den positive Religions Væsen, uden at lade det
\emph{Spørgsmaal} komme frem, der \emph{indirecte} er besvaret af
\emph{Kierkegaard}, i „Kjærlighedens Gjerninger“ og paa andre Steder: om ikke
dette Factiske, Phænomenale, kan finde sit Correctiv paa selve den positive
Religions Standpunkt, og om det ikke er at forklare ved et ubevidst Tilbagefald
fra det Religiøses Sphære til lavere Livsstadier.}.
% FOOTNOTE p. 120: printed mark *). It runs over onto p. 121 (the catchword-like
% continuation begins with the letterspaced „gaard,“); set here in full at its
% call. „Kierke-gaard“ is letterspaced across the page break, both halves —
% checked at 400 dpi on both pages.
% EMPHASIS p. 120: the run „fælleds for alle de positive“ stops at the line
% break; the following „Re-“/„ligionsformer“ is tight (verified at 400 dpi,
% 2.5x). Transcribed as printed, not regularised.

% --- p. 121 ---
Der er fine og dybt indtrængende Analyser af den positive Religions
Bestemmelser, og af de psychologiske Grundforhold, der opklare de religiøse
Forestillinger. Der er en rigtig Erkjendelse af Religionens væsentlige
Sammenhæng med Menneskenaturen, og af Væsenssammenhængen mellem den religiøse
Bevidstheds Gjenstand og det opfattende Organ. Der er en klar Paaviisning af
Umuligheden af at naae ud over Menneskevæsenets fundamentale Bestemmelser, af
Umuligheden af ved Tanken at kunne negere Tanken. Og dertil kommer, som et
Væsentligt, at alle F.'s Undersøgelser bære Præget af en høi Grad af Redelighed
og Sandhedskjærlighed, af en levende Overbeviisning, og tillige af en energisk
Fremadstræben, der altid paany sætter hans Forskning i Bevægelse og ikke lader
den nøies med engang vundne Resultater. Men trods alt dette har den Feuerbachske
Theorie sine store og væsentlige Mangler. Vi fremhæve de meest afgjørende.

\emph{Grundmanglen} er den, at der hos ham lægges et ensidigt Eftertryk paa
Relionens Fremgaaen \emph{af Menneskevæsenet}, og at det oversees, at
Menneskevæsenet selv er et Sat, et Afledet, og at det Sættende og Sammenfattende
virker igjennem det. Denne Mangel bliver allerede iøinefaldende i F.'s tidligere
Skrifter, som i \emph{Wesen der Christenthums}, hvor dog, ved en Reminiscents fra
Idealismen, „Slægtens“ Begreb endnu accentueredes som det Almene, Væsentlige, i
hvilket Individerne gaae op og som realiserer sig i dem, og hvor Opfattelsen af
\emph{Menneske-Slægtens Forhold til Fornuften og Naturen} lagde det nær at søge
tilbage til et Tilværelsesprincip, der indenfor Væsensenheden med det
Menneskelige havde en Transscendents over det. Men endnu bestemtere maa Manglen
træde frem i de sildigere Skrifter, hvor paa nominalistisk Maade Individerne
% PRINTER'S DEFECT p. 121, l. 16: the page prints „Relionens“ for
% „Religionens“ (the -gi- is simply absent). Both OCR witnesses read it so, and
% the image confirms it at 400 dpi. Transcribed as printed.
% EMPHASIS p. 121: „Slægtens“ inside the quotation marks is NOT letterspaced
% (checked at 400 dpi); the letterspaced runs are „Grundmanglen“, „af
% Menneskevæsenet“, „Wesen der Christenthums“ and „Menneske-Slægtens Forhold
% til Fornuften og Naturen“. In the last two the run carries across the line
% break in both halves.
% --- p. 122 ---
sættes som det egentlig Reale, Slægten kun som et subjectivt gyldigt almindeligt
Udtryk for deres Indbegreb, ligesom Fornuftens Identitet kun opfattes som „et for
Tænkningen og Sproget nødvendigt Udtryk, ved hvilket det Almene fixeres,
selvstændiggjøres, men ved hvilket det ikke maa glemmes, at det kun er et Product
af Tænkningen“. Her kan Religionen kun tænkes som fremgaaende af Individet, og
dens Kilde bliver saa, idet Individet ikke styres af Slægtens ubevidst
gjennemvirkende fornuftige Magt, det kun \emph{Alogiske i In}dividet.

\emph{Det rationelle} Moment i Religionen bliver hos F. paa ensidig Maade trængt
tilbage. Idet Religionens væsentlige Standpunkt bliver det subjectivt-praktiske,
skal Mennesket i Religionen kun opfatte sig fra det Praktiskes Synspunkt, og Alt,
hvad der er Theoriens Gjenstand, skal falde udenfor det religiøse Individs
Bevidsthed om det Menneskelige, og kun i en ved det religiøse Standpunkt
nødvendiggjort Forvandling, eller som et yderste Tilknytningspunkt, der snart
tabes helt af Sigte, finde sin Plads i det Guddommelige. Men i denne Opfattelse
er de forskjellige Evners, de forskjellige Organers, væsentlige Samvirken i
Menneskevæsenet, selv under den sporadiske Udvikling; Nødvendigheden af en
Stræben efter en Tilfredsstillen af \emph{alle} Grundevnerne, miskjendt. Der
indrømmes vel oprindelig, at ogsaa det Theoretiske uvilkaarligt gjør sig
gjældende for Bevidstheden; men denne Indrømmelse tages i Realiteten tilbage ved
de tilføiede Begrændsninger og ved den Betydning, F. tillægger
\emph{Personligheds}bestemmelsen i den religiøse Gudsforestilling. At imidlertid
det Theoretiske \emph{som} Theoretisk maa faae en bestemmende Indflydelse paa den
religiøse Opfattelse af det Absolute, ligger ikke blot deri, at Aanden i
Religionen søger Forsoning, Tilfredsstillelse, \emph{efter sin Heelhed}; men det
fremgaaer ogsaa af den Kjendsgjerning, at det har været muligt for
Menneskeaanden at arbeide sig ud af de første, raaeste religiøse Former, forud
for hvilke endnu ingen selvstændig, ɔ: af det Religiøse uafhængig,
Tankeudvikling har fundet Sted. Paa \emph{høiere} Stadier af Religionens
Udvikling
% EMPHASIS p. 122: two runs stop inside a word at the line break and are kept as
% printed — „Alogiske i In“ (with „dividet.“ tight on the next line) and
% „Personligheds“ (with „bestemmelsen“ tight on the same line). In „det
% Theoretiske som Theoretisk“ only „som“ is letterspaced; both occurrences of
% „Theoretisk“ are tight (verified at 400 dpi, 2x).
% Modern pencil marginalia „29/X 22“ stands after „dividet.“ — not transcribed.
% --- p. 123 ---
\setcounter{footnote}{0}%
er der, trods Opfattelsen af den \emph{guddommelige} Villie som
\emph{vilkaarlig}, en Mulighed for, at et \emph{Tankeindhold} kan lægge sig ind i
de religiøse Former, og i dem hævde det Rationelle en \emph{væsentlig} Betydning,
og dette Tankeindhold, der ligesom fældes ned i hine Former, er deels fremkommet
gjennem en Tankeudvikling, der ikke har sondret sig fra det Religiøse, men er
fremmet ved selve Religionen, deels ifølge en indre Trang \emph{tilegnet} af den
religiøse Bevidsthed fra en af den relativ uafhængig Tankeudvikling.

Den Tilbagetrængen af det Rationelle, der viser sig hos F., hænger sammen med
det, at Betydningen af det \emph{Aprioriskes} Virken i den religiøse Bevidsthed
som rationaliserende og ethicerende oversees. Den ubevidste Magt miskjendes, som
\emph{Menneskets Væsen} udøver over \emph{Individet}, idet det, virkende som
Enhed, behersker det individuelle Ufornuftige og Egoistiske. Men idet denne
Virken af det Aprioriske oversees, accentueres ensidigt ikke blot det
Irrationelle, men ogsaa det \emph{Egoistiske} i Religionerne. Derfor faaer ikke
heller den religiøse Hengivelses Moment sin rette Betydning. F. er ikke blind for
dette Moments Tilstedeværelse i det religiøse Sind, der fordrer uegennyttig,
selvhengivende Kjærlighed: \textit{diligere Deum propter Deum et non propter se
ipsum} (St. Bernhard). Men ligesom de rationelle „metaphysiske“ Bestemmelser af
det Guddommelige kun vare de „yderste Tilknytningspunkter“, ikke de egentlige,
specifisk religiøse Bestemmelser, saaledes skal den selvhengivende Kjærlighed kun
betegne den høieste religiøse Begeistrings Moment, der strax afløses af det
Forhold, i hvilket Subjectet, som adskillende sig fra det guddommelige Object,
forholder sig til sig som Object for Gud, altsaa indirect egoistisk. Og selv hvor
Selvet negerer sig selv i Selvhengivelsen, skal det kun gjøre det fordi det i Gud
har Nydelsen af den fra alle Skranker frigjorte Personlighed; fordi Gud er det
menneskelige Gemyts høieste Selvfølelse og Lystfølelse\footnote{Jfr. hermed F.'s
Opfattelse af Offret: ovenfor S. 115.}.
% ANTIQUA p. 123, the decisive page: „diligere Deum propter Deum et non propter
% se ipsum“ is set in upright antiqua and is TIGHT, not letterspaced (checked at
% 400 dpi against the spaced antiqua of p. 68) — hence \textit{} alone. Inside
% the same Latin sentence „(St. Bernhard)“ is Fraktur and stays Fraktur.
% FOOTNOTE p. 123: printed mark *). „Jfr.“ is a J-word (jævnfør) under the
% house I/J convention.
% EMPHASIS p. 123: „Menneskets Væsen“ and „Individet“ are two separate runs with
% „udøver over“ tight between them. In „det Egoistiske i Religionerne“ only
% „Egoistiske“ is set letterspaced; the single „i“ that follows carries wide
% spaces on both sides but cannot be adjudicated on one sort, and
% „Religionerne“ is plainly tight.

% --- p. 124 ---
Denne ensidige Udskillen af en enkelt Side i det religiøse Forhold bestyrkes hos
F. derved, at, netop gjennem Accentueringen af Personlighedens Form-Bestemmelse,
i Adskillelse fra dens Indhold, den \emph{ensidige Sondring mellem det
Guddommelige og Menneskelige} faaer Overvægten over den \emph{Ligevægt} mellem
Enhed og Sondring, eller rettere: over den \emph{Overvægt} af Enheden over
Sondringen, der oprindelig statueredes ved Bestemmelsen af de \emph{specifisk
religiøse} Prædicater for det Guddommelige. Og den bestyrkes ligeledes derved, at
den \emph{Gradforskjel i Bevidsthedsindholdets Klarhed og Tydelighed}, der finder
Sted under den religiøse Bevidstheds Udvikling, ligesaalidt tilstrækkeligt
fremhæves som \emph{Forholdet mellem den religiøse Bevidsthedsprocesses Momenter}
opfattes correct. Det \emph{mystiske} Element i Troeslivet kommer ikke til sin
Ret i F.s Opfattelse. Men paa dette mystiske Element, der er betinget ved, at i
det religiøse Livsforholds Enhed Bevidsthedsforholdets Momenter \emph{ikke træde
skarpt} ud fra hinanden, beroer netop den umiddelbare religiøse Bevidstheds
Følelse af Enheden med det Guddommelige: den Mulighed, den eier til en
Begeistringens Selvhengivelse, til en Gjenfinden af sig efter sin
Væsensbestemmelse i det Guddommelige. \emph{Det er uundgaaeligt, at Reflexionen},
Sondringens Kilde, gjør sig gjældende paa det religiøse Omraade. Hvor den
menneskelige Tænkning, idet den frigjør sig fra det Religiøse, udvikler sig ud
over det bestemte religiøse Standpunkt og virker tilbage paa dette, eller hvor
Menneskeaanden, paa selve det Religiøses Gebet, ved en spontan Tankeudvikling
begynder at løsne sig fra det tidligere Standpunkt, kan Reflexionen ikke undlade
at gjøre sig gjældende; men det er først hvor Reflexionen faaer en
\emph{væsentlig} Betydning paa det Religiøses Gebet, og hvor den søger at hævde
Religionens Sandhed \emph{i Modsætning til Tænkningens, at Sondringens} Moment
faaer Overvægt i den religiøse Bevidsthed, og at de Consequentser, F. drager,
faae deres Gyldighed.

Med den ensidige Fremhæven af Sondringen hænger
% p. 124 prints „F.s Opfattelse“ without an apostrophe, against „F.'s“ on
% pp. 121 and 123 — kept as printed on both readings.
% EMPHASIS p. 124: the letterspaced runs carry the small words with them —
% „Det er uundgaaeligt, at Reflexionen,“ is spaced throughout, and so are „at“
% in „til Tænkningens, at Sondringens“ and the opening „i“ of „i Modsætning“
% (verified at 400 dpi).
% A modern pencil bracket stands in the left margin beside ll. 27--31 — not
% transcribed.
% --- p. 125 ---
endelig den Mangel sammen hos F., at \emph{Modsætningsforholdet mellem det
Religiøse og det Ethiske ensidig fremhæves}, medens den Enhed, som det mystiske
Element i Religionen begrunder, ensidig skydes tilside. ---

Gjennem denne Angivelse af hvad jeg betragter som Manglerne ved Feuerbachs
Opfattelse af \emph{det positivt Religiøse} vil dens Forskjel fra min være bleven
tilstrækkelig tydelig, og af de i det Foregaaende givne Begrebsudviklinger vil en
væsentlig Differents i vor hele Grundanskuelse fremgaae.

Differentsen fremtræder strax i Bestemmelsen af \emph{Grundbegrebet}. Det bliver
for mig Begrebet om det ideelt-reale Absolute som det Enhedsprincip, der
begrunder Tilværelsen i relativ Selvstændighed; for Feuerbach bliver det, hvor
hans eiendommelige Standpunkt er \emph{fuldt} udviklet, Individerne, de enkelte
sandselige Ting og Væsener, medens det Guddommelige kun bliver en Afspeiling af
den subjective Attraa, det Almene kun et subjectivt uundværligt Udtryk for
Indbegrebet af det Individuelle, og som saadant uden real Gyldighed.

Af denne Differents i Grundbegrebet følger saa som Consequents, at medens paa det
Standpunkt, jeg har angivet, Betingelserne ville være tilstede for at hævde
Muligheden af \emph{Friheden i det} Menneskelige, for at hævde en
\emph{Erkjendelse og en Handlen}, der har det \emph{Almeengyldiges} Præg, altsaa
en sand videnskabelig Erkjendelse og en sand ethisk Handlen, og endelig for at
hævde et \emph{religiøst Forhold} til det gjennemvirkende Absolute, saaledes, at
dette Forhold overalt er i harmonisk Enhed med Forholdet til Mennesket; saa maa
derimod hos Feuerbach den successive Udvikling af hans Grundanskuelse til
Individualisme og Sensualisme føre til, at ikke blot et religiøst Forhold til et
Absolut umuliggjøres, men at ogsaa den det Almenes Grundvold tilintetgjøres, paa
hvilken en sand videnskabelig Erkjendelse maa hvile, og at endog, gjennem
Tilintetgjørelsen af Frihedens Mulighed og af Begrebet om et Ideal for Hand\-%
% EMPHASIS p. 125: „Friheden i det“ is spaced through „det“, with
% „Menneskelige“ tight; in „Erkjendelse og en Handlen, der har det
% Almeengyldiges Præg“ the words „der har det“ are tight, so the page carries
% two runs, not one (verified at 400 dpi).
% --- p. 126 ---
lingen, selve det Ethiske maa undergraves, til Fordeel for hvilket det Religiøse
skulde opoffres.

Naar F. beholder Udtrykkene: \emph{Religion og religiøs}, saa er det kun i
\emph{uegentlig} Betydning som Betegnelse for det, til hvilket Mennesket
forholder sig som til et Høieste, og som ogsaa faaer praktisk Betydning for det.
De betegne ikke for ham hvad de betegne for mig: et absolut Forhold til det
Absolute, et Forhold, der er ved dette Absolute og i hvilket Mennesket forholder
sig til det med sig Væsenslige, men tillige til det Menneskevæsenet og med det
den hele Tilværelse begrundende og i sin evige Enhed sammenfattende høiere
Princip. For mig bliver derfor „Menneskets Selverkjendelse“ i Religionen en
Selverkjendelse \emph{ved og gjennem sit Væsens Grund}, sit absolute Ideal.

\textbf{Den Hegelske Philosophie.} --- Denne Philosophies Opfattelse af
\emph{Livsopgaven} fremgaaer med Nødvendighed af dens Opfattelse af
\emph{Ideen}. Idet Ideen som den „rene Tankes“ høieste Udtryk i sig omfatter al
Tilværelse, er \emph{Enhed af Tanke og Væren}, saa omfatter \emph{Ideens evige
Proces}, i hvilken Naturen og de endelige Aander, hvis Selvstændighed kun er et
Skin, evigt ere optagne som ophævede Momenter, \emph{al} Virkelighed og
\emph{al} Sandhed; al sand Tilværelse er reduceret til Evighedens Element; hvad
der falder udenfor dette i Tidens og Rummets Former, er det \emph{kun Usande},
der \emph{med evig Nødvendighed} ophæves i det Absolute. Derfor blive ogsaa
Menneskelivets Naturformer og dets Naturindhold kun et Uberettiget og Usandt;
Aandens ene absolut sande Væren bliver \emph{Tankens individualitetsløse,
contemplative Liv i Evighedens Element}, i hvilken den veed det Absolute
\emph{som det i dets evige Proces ene Værende}, og veed \emph{sig som Moment i
dette Absolute}, i hvis Viden om sig Menneskets Viden om det Absolute og om sig i
det Absolute er optagen. Ligeoverfor Anskuelsen af den evige metaphysiske Proces,
der ikke er en Stræben hen imod et Maal, men et evigt
% HEAD, p. 126: run-in halbfette (bold) Fraktur head, „Den Hegelske
% Philosophie.“ followed by an em-dash, NOT letterspaced — verified at 400 dpi
% and set exactly like the parallel head „Ludvig Feuerbach.“ on p. 97. A short
% centred rule stands above it; not transcribed. The Indhold instead reads
% „2) den Hegelske Philosophies“. Printed head kept in the body.
% EMPHASIS p. 126: of the three „al“ in „al Virkelighed og al Sandhed; al sand
% Tilværelse“ the first two are letterspaced and the third is tight (checked at
% 400 dpi, 2.5x).
% --- p. 127 ---
Værendes Selvbekræftelse, i hvis Nødvendighed derfor enhver Modsætning evigt
\emph{er hævet}, tabe Personlighedens Virkelighedsproblemer: Problemet om den
endelige Aands Villiesfrihed; om den frie Villies Virkeliggjørelse af Væsenet som
naturligt-aandeligt; om Menneskeaandens Forhold til Tiden, indenfor hvilken den
forholder sig til en Evighedens Opgave, til det Ethisk-Religiøse; om dens
Selvbestemmelse til denne Opgaves Løsning; om dens Forfeilen af Opgaven; om dens
Skyld og dens Forsoning af Skylden, --- deres Betydning og Skarphed. Overfor den
\emph{evige og evigt færdige} Processes Nødvendighed bliver \emph{Endelighedens}
Bevægelse, selv hvor den afslutter sig i det \emph{Onde}, kun en Spøg, der evigt
er taget tilbage i det Absolutes Proces, i hvilken al Alvoren falder. Gjennem
\emph{Tanken om det Absolute} kan Individet tage sig ud af Endelighedens
Modsigelser, og i den individualitetsløse contemplative Enhed med det Absolute
har det \emph{Absolutionen} for al Ufuldkommenhed, og \emph{Forsoningen} med sit
Grund-Væsen. I den naturbestemte Tilværelse skal der leves i den rene
Aandeligheds Abstraction, i den endelige Existents skal der leves kun i
Evighedens Element, saa at Tiden i Tiden skal tages tilbage i Evigheden, saa at
Tidsforholdet ikke faaer Virkelighedsbetydning, ikke bliver Moment i den
ethisk-religiøse Beslutning til Handling og Selvhengivelse. \emph{Tanke er
Handling}, det: at \emph{begribe} det Absolute, den \emph{absolute Handling}.
Livet bliver en Aandelighedens Illusion, \emph{Ethiken} forvandles til abstract
\emph{Metaphysik, Religionen} til en \emph{Begriben af Ideen}.

\emph{Differentsen} mellem denne Opfattelse af Livsopgaven og den ovenfor
angivne, er bestemt ved Differentsen \emph{i Grundbegrebet}. Den tidligere
udviklede Opfattelse af det Absolutes \emph{væsentlige} Forhold til
\emph{Realitets}bestemmelserne, det \emph{væsensgyldige} Reales
\emph{væsentlige Være}former, fører med Nødvendighed til, paa den ene Side at
hævde den \emph{naturbestemte} endelige Aands \emph{relative Selvstændighed}, paa
den anden Side til at lægge et afgjørende Eftertryk paa
\emph{Realitetsbetingelsernes}
% EMPHASIS p. 127: two more runs stop inside the word — „Realitets“ (with
% „bestemmelserne“ tight) and „væsentlige Være“ (with „former,“ tight on the
% next line). „Spøg“ is NOT letterspaced, though „Onde“ just before it is.
% Modern pencil scribbles stand in the right margin and across „skal tages“ —
% not transcribed.
% --- p. 128 ---
Betydning for den ethiske Handlen, overhovedet for Bestemmelsen af den
menneskelige Livsopgave. Denne Differents vil fremtræde i sin concrete Form
gjennem de efterfølgende Udviklinger af centrale ethisk-religiøse Problemer. ---

Fra de kritiske Sammenstillinger gaae vi nu over til disse Udviklinger, i hvilke
væsentlige Begrebsbestemmelser, der i det Foregaaende ere blevne anticiperede,
skulle finde deres Retfærdiggjørelse, og det kun Antydede faae sin concretere
Bestemthed. Efter Sagens Natur bliver Problemet om \emph{Villiens Frihed} det,
hvis Løsning først maa forsøges.

\section*{1. Om Villiens Frihed.}

Problemet om Villiens Frihed fører os ind paa Metaphysikens, Physikens,
Psychologiens og Ethikens Omraade, men maa finde sin definitive Løsning paa det
sidste.

De Begreber, der give de første, \emph{abstracte Tilknyt}ningspunkter for
Frihedens Bestemmelse, er Begrebet om \emph{Væren som Virksomhed} og Begrebet
\emph{Selvhed}. I selve Værens Begreb ligger Virksomhedens Bestemmelse
indesluttet, idet al Væren kun hævder sig som Væren, kun bevarer sin Væren, ved
Virksomhed, og uden denne Virksomhedens Bestemmelse umiddelbart vilde synke
tilbage i Ikke-Væren. Men denne Virksomhed, der er uadskillelig fra Væren,
udtrykker Værens \emph{Eiendommelighed}; gjennem den gjør altsaa enhver Væren sig
gjældende \emph{efter sin Bestemthed}, sætter i Vexelvirkningen et af dens eget
Væsen Fremgaaende, dens eget Væsen Udtrykkende.

I denne Bestemmelse af en eiendommelig \emph{Yttring} af Væsensvirksomhed have vi
den første Antydning af Friheden; men et Skridt nærmere komme vi Frihedens Begreb
gjennem Begrebet \emph{Selvhed}, Subjectivitet. Selvhedens Begreb udtrykker en
Væren for sig gjennem sit Andet; Selvheden, som det \emph{Ideales} Udtryk, er,
efter Begrebsforholdet, det Bestemmende i Forhold til Andetheden og det, der
omsætter
% HEAD, p. 128: centred halbfette (bold) Fraktur display head, „1. Om Villiens
% Frihed.“, wording verified at 400 dpi; a short centred rule stands above it,
% not transcribed. p. 128 carries NO footnote, contrary to the OCR-derived hint
% list — checked by eye.
% EMPHASIS p. 128: „abstracte Tilknyt“ is spaced and „ningspunkter“ tight —
% another run stopping at the line break, kept as printed.
% p. 128, l. 32: a small low mark stands after „Selvhedens Begreb“, set below
% the baseline and unlike the period after „Subjectivitet.“ on the same line;
% read as dirt rather than as a sort, and no punctuation transcribed. The
% sentence runs on.
% --- p. 129 ---
enhver Bestemmethed ved et Andet i Selvbestemmethed. Men ogsaa Selvhedens Begreb
afgiver endnu kun \emph{abstracte} Tilknytningspunkter for Begrebet om
\emph{Villiens Frihed}, det er ikke tilstrækkeligt til Grundlag for dette eller
til at hævde det endelige Existerendes Frihed i Forhold til det Uendelige og
Absolute. Selvheden er Grundbestemmelsen for \emph{alt} Organisk, for \emph{alt}
Levende, altsaa ogsaa for Plantelivet og Dyrelivet. Men selv Dyrelivets høieste,
meest udviklede Former gaae ind under Instinctets Bestemmelse: hiin abstracte,
metaphysiske Bestemmelse af Friheden, der ligger i Selvheden, bliver altsaa i den
reale Virkelighed i det enkelte Dyrelivs Existents omsluttet og skjult af
Nødvendigheden. Den endelige Selvheds Forhold til det Reale, der ikke er sat ved
dette Selv, og ikke tilhører dets Væsen, er ikke umiddelbart bestemt ved det
almene metaphysiske Forhold mellem Selvhed og Andethed. Og hvad Metaphysiken
fastsætter med Hensyn til dette Forhold, der nærmest opfattes som et Forhold
mellem Grundmomenterne i det absolute Princip og i den ved Principet bestemte
\emph{totale} Tilværelse, kan endnu mindre umiddelbart betrygge den endelige,
individuelle Selvhed i dens Forhold til den absolute Subjectivitet.

Gjennem de angivne Tilknytningspunkter naaes der altsaa vel til Momenter i
Frihedens Begreb, men ikke til den concrete Bestemmelse af Friheden som
\emph{Villiens Frihed} eller til en Hævdelse af dette Begrebs Realitet. Forat der
kan findes et sikkrere Grundlag for Bestemmelsen af den menneskelige Villies
Frihed, maa en ny Bestemmelse komme til Selvhedsbegrebet, nemlig
\emph{Væsensuniversalitetens} Bestemmelse. I denne Bestemmelse ligger, som
tidligere paaviist, Bestemmelsen af \emph{indre Uendelighed}, og derigjennem
Muligheden af \emph{en Væren for sig i Selvbevidsthedens og den begribende
Erkjendelses Form}. Selvhedsenergiens Virken bliver, med Væsensuniversaliteten
til Forudsætning, \emph{bevidst} Virken, hvad der er Betingelsen for, at den kan
være en \emph{Villies} Virken, og idet den er Udtryk for Energien i et Væsen med
\emph{indre Uendelighed},
% The signature numeral 9 at the foot of p. 129 (first page of the gathering) is
% not transcribed.
% --- p. 130 ---
viser der sig en Mulighed til, at den kan være en \emph{fri} Villies
Selvbestemmen. Friheden vilde da nærmest udtrykke, at den menneskelige Aand, i
sin Virken gjennem den Væsenets Uendelighed udtrykkende Villie, med Bevidsthed
tager Initiativet for Villien \emph{i sig selv}, og ud fra \emph{sig selv} giver
Villien Indhold. Hvad der i Almindelighed gjaldt om ethvert Værende, at det som
værende maa være virksomt ud fra sin Værens Eiendommelighed; det fik da her den
Form, at Aanden, ifølge sin indre Uendelighed, i sin Selvbestemmen satte sit
Villiesindhold, og, hvor Selvbestemmelsen var den sande, gav Villien et Indhold,
der udtrykte Selvet efter dets Sandhed ɔ: efter dets Væsen, altsaa et
væsensgyldigt, fornuftbestemt Indhold. I denne Bestemthed var da Villien
Væsensvillen ɔ: selvbestemmende ud fra sit almene Væsen. Men dette Væsen, der
sees som selvbestemmende i sin Villen, og som havende sig selv til Indhold for
sin Selvbestemmelse, er det efter sin primitive Begrebsbestemthed selvbevidste og
erkjendende Væsen: i hiin Selvbestemmelsens sande Form maa der derfor være en
Enhed af den theoretiske og den praktiske Aands Grundbestemmelse, og Friheden
blev saaledes det høieste Udtryk for det menneskelige Væsen efter dets Totalitet,
den blev dette totale Væsens sande Virkelighedsform, Erkjendelsens og Villiens
τέλος.

I de givne Bestemmelser er det angivet, der betinger Villiesfrihedens Mulighed,
og Frihedens Forhold til Væsenet er antydet. Men idet det menneskelige Væsen og
den menneskelige Villie ikke blot sees efter det ligesom isolerede Væsens almene
Begreb og efter Villiens Grundforhold til Væsenet, men betragtes under
Realitetsbetingelserne og i Virkelighedens concrete Sammenhæng, reiser der sig en
Tvivl om, at det, der under Abstractionen fra det Concrete viste sig som en
Mulighed, kan hævdes som gyldigt under Virkelighedens reale Forhold. Først naar
denne Tvivl er hævet, faaer det udviklede Begreb af Friheden mere end blot
hypothetisk Gyldighed.

Hvad der vækker Tvivl om Frihedens Virkelighed er paa den ene Side Individets
Forhold til \emph{det Absolute},
% GREEK p. 130: τέλος, antiqua italic with the acute, at the end of the first
% paragraph — the only Greek in pp. 120--133. Neither OCR witness reads it
% (tesseract gives „ré/oc“, ABBYY nothing); read off the image at 400 dpi.
% EMPHASIS p. 130: the run „i sig selv“ takes the preposition with it (its
% flanking spaces match the spaced run), while in „og ud fra sig selv“ the
% words „og ud fra“ are tight.
% --- p. 131 ---
paa den anden dets Forhold til \emph{de reale Existentsbetingelser}. Det første
af disse Forhold vil faae sin egentlige Belysning gjennem Behandlingen af
Problemet om det Ondes Mulighed; det andet bliver her at behandle fra de
forskjellige Synspunkter, fra hvilke det viser sig.

Individets Forhold til sine Existentsbetingelser bliver nemlig deels at see under
Synspunktet af et Forhold til det \emph{Ydre}, deels under Synspunktet af et
Forhold til \emph{Udviklingsbetingelserne}. Under Bestemmelsen af \emph{Forholdet
til det Ydre} sammenfatte vi: Individets Forhold \emph{til det ydre Reale}, i
hvis Causalitetssammenhæng det er indordnet; Forholdet til det \emph{Ydre i
Individet selv} ɔ: til dets \emph{Naturbestemthed}; og Forholdet til \emph{andre
Individer}.

\emph{Det første} Forhold: Forholdet til \emph{det Reales Causalsammenhæng}, er
fortrinsviis blevet accentueret som Tvivlsgrund mod Frihedens Virkelighed, og det
ikke blot fra materialistiske \emph{Systemer}, for hvilke den menneskelige Villie
kun bliver et Led i de materielle Aarsagers og Virkningers mechanisk forbundne
uendelige Række; men ogsaa fra saadanne Systemer, der skille en Numenets Sphære
fra Phænomenets, og indrømme en intelligibel Frihed, men urgere, at Mennesket,
som Phænomen i Tid og Rum, ubetinget er Causalitetslovens Nødvendighed
underkastet, og derfor i sin Villen med Nødvendighed er bestemt ved et
\emph{ydre Motiv}. Skjøndt det her ikke oversees, at den Forudsætning, paa
hvilken overhovedet Nødvendigheden af alle Aarsagers Virkninger beroer, er hver
Tings indre Væsen, og skjøndt der tales om \emph{indre} Incitamenter (i
Organismen), og om en \emph{indre} Tilskyndelse (hos Dyrene, gjennem Instinctet),
fastholdes det dog, idet den intelligible Frihedsact skal have sat den
Grundvillie, der constituerer Væsenet, som en oprindelig, \emph{uforanderlig}
Enhed, at \emph{ethvert} Motiv til en Villiesbevægelse, en bestemt Villiesact, at
hele „Villiens Stof“ skal gives af \emph{Yderverdenen}. Ogsaa det abstracte, blot
tænkte Motiv skal ligesaafuldt være en ydre, Villien bestemmende Aarsag, som det
anskuelige, der bestaaer i et
% EMPHASIS p. 131: „Det første“ is spaced but the following „Forhold:“ is tight
% (verified at 400 dpi), so the sentence carries two runs.
% The signature „9*“ at the foot is not transcribed.
% --- p. 132 ---
realt nærværende Object; i sidste Instants skal det dog altid være et Realt, et
Materielt, forsaavidt det dog altid beroer paa et til en eller anden Tid og paa
en eller anden Maade --- hertil regnes ogsaa Overlevering gjennem Ord og Skrift
fra de fjerneste Tider --- modtaget \emph{ydre} Indtryk (Schopenhauer). For denne
Opfattelse bliver det da ligesaa umuligt som for Materialismen at tænke sig den
menneskelige Villies Frihed; den kan i denne, som Evne til en \emph{absolut
Begynden}, kun see det, der vilde constituere en Conflict med den almene
Verdensorden.

Hvad nu denne saa stærkt accentuerede Tvivlsgrund angaaer, saa maa der strax
vises hen til, at hiin mechaniske Opfattelse af Natursammenhængen, der ogsaa
faaer sin Gyldighed indenfor den charakteriserede Form af Idealismen, lider af
den uundgaaelige Mangel, at en \emph{primitiv} Motor, et \emph{definitivt} Ophav
til den \emph{som Kjendsgjerning} givne Bevægelse, \emph{ikke} kan findes uden en
Inconsequents. \emph{Er al Bevægelse kun som udenfra meddeelt Bevægelse, saa vil
det sige, at den ingensteds er oprindelig tilstede}: en uendelig Forøgelse af
Leddene bliver kun en tautologisk Gjentagelse af det samme Forhold, og \emph{det
bliver en Modsigelse, at anerkjende Bevægelsen, Activiteten, som factisk}, da den
ingen Grund kan have i det Materielle, og ingen udenfor det; thi det, der tænktes
staaende udenfor de materielle Aarsagers Række, vilde, som væsensmodsat det
Materielle, ikke kunne udøve en Virkning paa dette. Skal Bevægelsen,
Virksomheden, i Naturen forklares, ja endog blot anerkjendes som factisk, bliver
det uundgaaelig nødvendigt at gaae tilbage til Begrebet om en
\emph{Totalitetsvirksomhed}, og da at see Totalitetsvirkningen som resulterende
af \emph{Totalitets}leddenes Virken, idet selve Begrebet om en
\emph{Totalitetens} Virken viser hen til \emph{relativ selvstændige individuelle
Virksomhedscentra}. Naar da Begrebet om en \emph{virksom Naturtotalitet} er
traadt i Stedet for Begrebet om en \emph{mechanisk Natursammenhæng}, saa bliver
det ovenfor angivne Forhold mellem det Selvløse og Selvet, nærmere det
selvbevidste Selv, at ac\-%
% EMPHASIS p. 132: „Totalitets“ is spaced with „leddenes“ tight, whereas
% „mechanisk Natursammenhæng“ carries its letterspacing across the line break
% into „hæng,“ — both checked at 400 dpi.
% Modern pencil marginalia „7 X 22“ stands after „Verdensorden.“ — not
% transcribed.
% --- p. 133 ---
centuere, ligesom ogsaa Begrebet om det menneskelige Selvs Centralisation i sig
som relativ Totalitet med indre Uendelighed og derfor med et eiendommeligt
Centrum for bevidst Virksomhed. Det Begreb, der for Mechanismen umiddelbart viste
sig som en Absurditet: Begrebet om en \emph{primitiv Begynden}, maa
\emph{nødvendigviis} komme frem idet det individuelle Væsen, der har indre
Uendelighed, som relativ Totalitet concentrerer, centraliserer sig i sig og
derved, som tagende sig selv i Besiddelse, i sin Constitueren af sig sætter et
nyt Begyndelsespunkt, et eiendommeligt relativt Centrum i Tilværelsens
Totalorganisation. Hvorledes denne den relative Totalitets Virken i Forhold til
det \emph{absolute} Princip bliver at opfatte, vil blive belyst i det
Efterfølgende.

I Selvets Væsen, i dets ideelle Fordobbling, i dets Charakteer som Energie,
ligger Muligheden af en Selvbestemmen, hvis Motiv falder \emph{indenfor} Selvet,
i dets concrete Væsensindhold; i Forholdet mellem det efter sit Grundlag
\emph{evige Selv} og dets Udfoldelse i en \emph{Tidsudvikling} ligger Muligheden
af, at en Bestemmen, der fra Phænomenets Side fremtræder som tidsbestemt
\emph{Forandring}, og som derfor paastodes, at maatte være bevirket ved et
\emph{Ydre}, kan fremgaae af Selvets evige Grund. Og idet Selvet existerer som
indordnet i en Tilværelsens Totalitet, og i hvert Moment modtager Paavirkninger
fra det \emph{Ydre}, saa ophæver ikke heller dette Forhold for det selvbevidste
Selv Muligheden af en fri Selvbestemmen.

Med Hensyn nemlig til Indvirkningerne fra \emph{det ydre Reale}, der, som
Indvirkninger paa Villien, ikke blive mechanisk virkende Aarsager, men
\emph{Motiver}, ligger der for det Første i det menneskelige Væsens
\emph{Universalitet Muligheden til at abstrabere fra det bestemte Motiv}. I Kraft
af Abstractionens indre Uendelighed er Frihedens Mulighed bevaret, om end
foreløbig kun i \emph{Negativitetens Form} som en fri Ikke-Villen af hvad der
udenfra paatrænger sig, og denne Frihed i Negativitetens Form har som sin yderste
Grændse en fri Opgiven af den individuelle Existents forat hævde Friheden. Men
hiin Væsenets Uni\-%
% PRINTER'S DEFECT p. 133, l. 30: the page prints „abstrabere“ for
% „abstrahere“. The sort is a b, not an h — the bowl is closed and there is no
% descending leg, matching the b of „ab-“ in the same word and of „bestemte“ two
% words later, and unlike the h of „Uendelighed“ on the line below (checked at
% 400 dpi, 2x). Transcribed as printed.
% EMPHASIS p. 133: „fri“ in „en fri Selvbestemmen“ is tight, though ABBYY
% renders it spaced; „indenfor“ is spaced but the following „Selvet,“ is not.
% --- p. 134 ---
\setcounter{footnote}{0}%
versalitet, der gjennem Evnen til den uendelige Abstraction begrundede
Frihedens Mulighed i Negativitetens Form, indeslutter tillige en \emph{positiv}
Bestemmelse af Friheden. Idet det ydre Motivs Bestemmen afværges af Villien,
skeer dette i Kraft af \emph{indre ɔ: af Væsenet fremgaaende} Motiver, der danne
Fundamentet for Abstractionens Act, og det i Villiesacten Bestemmende bliver
altsaa en Bestemmelse i \emph{Selvets Væsen}: Villien bliver altsaa bestemt ved
hvad der tilhører det Væsen, den \emph{selv} udtrykker, og idet Væsenet bliver
gjennemsigtigt for Bevidstheden i denne Bestemmen, forvandler den \emph{indre
Væsensnødvendighed sig til bevidst Frihed}. Endelig ligger det i Væsenets
\emph{Universalitet}, at det \emph{ydre Motiv}, hvor en \emph{Væsensenhed}
finder Sted, \emph{kan forvandles til et indre, ɔ:} omsættes til Moment i
Væsensbestemmelsen. Overalt fastholdes det altsaa, \emph{at Villiesacten med
Nødvendighed resulterer af Motiv og Villie}, men idet Motivet bliver et
\emph{indre}, hævdes Friheden, ikke i Betydningen af et grundløst, altsaa
meningsløst, \textit{liberum arbitrium}, men i Betydning af den sande,
væsensbestemte Frihed, den Frihed, gjennem hvilken Villien, idet den bestemmes
ved Motiver, udgaaende fra Væsenet, viser sig som Væsensudtryk, i Enhed med
Væsenet.

Der er i det Foregaaende seet hen til Villiens frie \emph{Beslutning, til den
frie Handlen i det Indre, Frihedens indre Virkeliggjørelse}. Hvad den frie
Villies \emph{ydre Realisation} angaaer, saa bliver denne, naar den kun sees
\emph{i Forhold til den individuelle} Villie, at bestemme som et af
Frihedsbeslutningen væsentlig Uafhængigt. Anderledes viser Forholdet sig naar
Forholdet mellem Villien og dens Udførelse i det Ydre, naar den frie ethiske
Handlens Forhold til Naturens og Historiens Gang og Orden, sees som medieret ved
det \emph{Universelle}, ikke som et udenfra ordnende Tredie (Kant), men som det
immanente Harmoniserende\footnote{Det Spørgsmaal, hvorvidt den frie Villies
Ikke-Realisation i det Ydre, som hængende sammen med en forskyldt Erkjendelsens
Mangel, kan blive Gjenstand for ethisk Ansvar ogsaa fra det Synspunkt, at hiin
Ikke-Realisation kan virke tilbage paa Frihedens \emph{indre} Virkeliggjørelse
og hæmme Frihedsudviklingen, falder udenfor nærværende Undersøgelse.}.
% printed mark: *) --- the only note on p. 134; it runs over onto p. 135 and is
% set here in full at its call. „indre“ in the note's second half is
% letterspaced, „Virkeliggjørelse“ is not (300 dpi).
% ANTIQUA LATIN, p. 134: „liberum arbitrium“ is upright antiqua against the
% Fraktur and is CLOSE-SET, not letterspaced --- unlike the spaced antiqua of
% p. 68 and of „sub specie æterni“ at pp. 144--145 below.
% NAME NOT LETTERSPACED, p. 134: „(Kant)“ is ordinary close-set Fraktur, neither
% antiqua nor spaced. Checked at 4x against the spaced „Universelle“ two words
% earlier in the same sentence. So no \emph{} here, contrary to expectation.
% SPLIT SPACING, p. 134: the run „i Forhold til den individuelle“ stops at the
% line end; the „Villie,“ that completes the phrase on the next line is
% close-set. Transcribed as printed.

% --- p. 135 ---
Idet de frie Villiesacter, der udgaae fra de primitive Virksomhedscentra, træde
ud i ydre Handling, gaaer Handlingen \emph{som Phænomen ind under
Causalitetslovens Nødvendighed}, og af de frie Villiesacters Samvirken
resulterer i det Historiske \emph{den almene Handlings Nødvendighed}. Her findes
en \emph{Virksomhedernes Omsætning} Sted, liig den, der finder Sted mellem
Naturvirksomhederne indbyrdes, og under denne Omsætning bevares Virksomhedens
Maal uforandret.

I det Foregaaende er Villiens \emph{Forhold til det Ydre} seet som Forholdet til
det \emph{udenfor Individet} værende Reale: det enkelte Reale og Totaliteten af
det Reale; det bliver nu at see som et \emph{Forhold til det Ydre i Individet
selv ɔ: til Individets Naturside}. See vi bort fra hvad der vil blive behandlet
under Synspunktet af Indi\emph{videts Forhold til sine Udviklingsbetingelser},
bliver her Forholdet til \emph{Indvirkningerne fra det uigjennemsigtige
Naturlige i Individet}, og Forholdet til de \emph{Hæmmelser for Friheden}, der
kunne ligge i oprindelige Naturbestemtheder: Anlæg, Temperament, o. desl. at
betragte. Hvad hine Indvirkninger angaaer, der hænge sammen med den blivende
Sondring mellem det \emph{Bevidste} og det \emph{Ubevidste i Individets Liv}, og
som kunne betegnes som \emph{Drifttilskyndelser}, naar vi tage Begrebet
\emph{Drift i udstrakt} Betydning, saa bliver med Hensyn til dem Begrebet om
\emph{Væsensenheden} og \emph{Væsenstotaliteten} afgjørende for Frihedens
Hævdelse. Idet det af det \emph{Ubevidste} Fremgaaende fremgaaer af det
\emph{samme} udelte Væsen, som det, der kommer frem i \emph{Bevidsthedslivet},
saa bliver først den deri liggende Nødvendighed --- forsaavidt vi ikke
reflectere paa den muligviis derigjennem virkende udenfra kommende Bestemmethed
--- en \emph{indre Væsensnød}\-%
% PRINTER'S DEFECT, p. 135, line 6: the page prints „Her findes en
% Virksomhedernes Omsætning Sted“, where the idiom requires „finder ... Sted“
% --- and „der finder Sted“ occurs correctly in the very next line, giving a
% control for the r/s ending. Zoomed to 12x: the word ends in a round s, not r.
% Transcribed as printed.
% SPLIT SPACING, p. 135: „Indi-“ at the line end is close-set while
% „videts Forhold til sine Udviklingsbetingelser“ on the next line is spaced.
% Transcribed as printed, cf. „Aands“functioner, p. 6.
% --- p. 136 ---
\emph{vendighed}, altsaa \emph{fri Nødvendighed}. Og idet saa nærmere
\emph{Væsenstotalitetens} og \emph{Væsensuniversalitetens} Bestemmelse gjøres
gjældende, saa bliver det fra Individets Naturside Udgaaende seet som gjennem
den organiske relative \emph{Totalitet optaget i det aandelige} Væsens
Livsenhed, som sammenfattet i dets ideelle Concentration, og den indre
Væsensnødvendighed forklares til det \emph{aandelige} Væsens indre
Væsensnødvendighed ɔ: til den væsensbestemte Frihed. Ligeoverfor den enkelte
Bestemmethed gjør da \emph{Væsenstotaliteten} og \emph{Spontaneiteten} sig
gjældende, og Bevidstheden om Totaliteten bevarer sin Betydning, selv om den
fremtræder i Abstractionens Form.

Hvad det angaaer, der ligger som \emph{Hæmmelser for Friheden} i de individuelle
Anlæg, i Temperamentet o. desl., der som et Naturbestemt, et ved Naturen Givet,
gaaer \emph{forud for} Individets bevidste Aandsliv og forudsættes af dette, saa
bliver, i Forhold til disse individuelle Differentser, hvor stærkt de end kunne
fremtræde, det menneskelige Væsens \emph{blivende Universalitet} at fastholde,
og derigjennem det Almeen-Menneskeliges Gyldighed for \emph{ethvert}
Menneskeindividuum trods Differentsernes Intensitet. Men med det
\emph{Grundmenneskelige} er det \emph{ideelle Forhold til sig selv}, der er
Frihedens Grundbestemmelse, givet, og dette Forhold, der er den menneskelige
Individualitets ideale \textit{prius}, virkeliggjøres for Bevidstheden i den
ideale Concentration (see nedenfor S.\ 139 f.), i hvilken Individet constituerer
sig som Personlighed.

Hvad der gjælder om de naturlige Differentser i Individet i Almindelighed,
gjælder ogsaa om den særegne Form af disse, der træder frem i de \emph{abnorme
physiske} og \emph{psychiske} Tilstande. I Individets Forhold til disse Villien
hæmmende Tilstande ligger deels et Forhold til det Ydre i Individet selv, deels
et Forhold til det Reale udenfor Individet, forsaavidt som en disharmonisk
Indvirkning af dette Reale paa Organismen er betingende hine Tilstande. Med
Hensyn til disse Tilstande bliver det Synspunkt at gjøre gjældende, at det til
et Aands- og Frihedsliv anlagte
% ANTIQUA LATIN, p. 136: „prius“ is upright antiqua, close-set. Tesseract does
% not read it as Fraktur and ABBYY gives „xriu8“, both signals of the roman sort.
% „deels et Forhold“, p. 136: tesseract reads „el“. At 4x the second sort is a
% Fraktur t (short flagged ascender), identical with the t of „til“ and „det“
% on the same line and unlike the l of „Forhold“. Read as „et“.
% p. 136 carries no footnote and no Greek.
% --- p. 137 ---
menneskelige Individs fortsatte Existeren trods disse Tilstande involverer
Muligheden til en Befrielse, der altsaa er tilstede \emph{saalænge Individet
existerer som Individ af denne Art,} som \emph{menneskeligt} Individ. ---

Naar endelig Individets Forhold til det Ydre sees som et Forhold til de
\emph{andre Individer} som handlende, til de andre selvbestemmende Villiers
Virkninger, saa gaaer dette Forhold fra en Side umiddelbart ind under samme
Synspunkt som Individets Forhold til det ydre Reales Paavirkninger, fra en anden
Side viser der sig en \emph{Nuance i For}holdet idet det andet Individ ikke blot
gjennem Væsenslighedens ubevidste Magt, men gjennem sin Erkjendelse har
Muligheden til en Villien langt fastere omspændende Indvirkning end det ydre
Reale. Men Forskjellen bliver dog \emph{kun en Gradforskjel}, der ikke forandrer
Grundforholdet i det Væsentlige. Derimod faaer under Individernes indbyrdes
Forhold det \emph{supplerende Enhedsforhold} mellem dem en væsentlig Betydning
for Frihedens Realisation i Individet; hvad der nærmere vil blive belyst i det
Efterfølgende, under Udviklingen af Problemet om Forsoningen. Det supplerende
Vexelvirkningsforhold mellem Individerne, i hvilket, ligesom ved
Selvbevidsthedens Realisation, Forholdet til det andet Jeg \emph{som Slægtens
Repræsentant} faaer sin Betydning, vil der blive paaviist som Betingelse for
Frihedens Virkeliggjørelse i Individet: en Virkeliggjørelse, der finder Sted i
det \emph{ethiske} Forhold til de aandelige Individualiteter, ved hvilke
Individet integreres, og som viser tilbage til en \emph{Frihedens Forudsætten af
sig selv}, og endelig til dens sidste Grund i det, der sammenfatter de hverandre
integrerende Individualiteter i Aandens Enhed: til dens Grund \emph{i det
absolute Princip}.

Naar saaledes Frihedens Virkelighedsbetydning har viist sig bevaret under
Forholdet til det Ydre, bliver den nu, efter det ovenfor Angivne, at hævde
\emph{i Forholdet til Udviklingsbetingelserne}. Med Hensyn til dette Forhold
gjør et dobbelt Synspunkt sig gjældende, der fører til en dobbelt Aporie. Idet
nemlig Friheden --- i prægnant Be\-%
% SPLIT SPACING, p. 137: two places where the compositor drops the spacing at a
% line break --- „Nuance i For“|„holdet“, and the „som“ between „af denne Art,“
% and „menneskeligt“, which is close-set although both neighbours are spaced.
% Transcribed as printed.
% p. 137 carries no footnote and no Greek.
% --- p. 138 ---
tydning som \emph{Aandens Frihed} --- ikke umiddelbart kan være det
naturbestemte Individ given, men maa \emph{erhverves af Individet}, altsaa faae
en \emph{Vorden} gjennem Individets Tidsudvikling, saa viser Vanskeligheden sig
først, idet der reflecteres paa denne Vorden. Ifølge Aandens Begreb maa Friheden
sættes som den, der skal vorde, fordi den skal erhverves, men idet den kun kan
være som \emph{selverhvervet}, alt\emph{saa maa være ved sig selv, altsaa maa
forudsætte sig selv}, saa synes dette at udelukke en Vorden, da denne som
Frihedens Vorden vilde forudsætte \emph{Ufriheden} som det, ud fra hvilket
Friheden skulde vorde. Og Vanskeligheden viser sig \emph{dernæst} idet
Menneskeaandens Tidsudvikling sees \emph{som en Tilnærmelse til Væsenets Ideal},
idet altsaa Friheden, som Udtryk for og som Virkeliggjørelse af Menneskets sande
Væsen, sees under \emph{Idealets} Synspunkt, som Gjenstand for \emph{en
Approximation}; thi Friheden faaer da Præget af det, der \emph{kun
tilnærmelsesviis} realiseres, og som aldrig i fuld Betydning faaer Virkelighed,
hvilket, hvis kun den fulde omfattende Realisation af Friheden med Rette faaer
Frihedens Navn, vilde være Eet med, \emph{at Friheden overhovedet ikke fik
Virkelighed}. Denne i Approximationen liggende Vanskelighed fremtræder i
bestemtere Form idet der reflecteres paa, at Villiesfrihedens Virkeliggjørelse
har \emph{Erkjendelse} til væsentlig Betingelse, altsaa viser sig som i sin
Fuldkommenhed betinget ved dens Fuldkommenhed, der, naar Erkjendelse tages i sin
concreteste Betydning, kun er et tilnærmelsesviis Opnaaeligt, saa at altsaa
Friheden, under den ovenfor angivne Forudsætning, aldrig vilde faae den
fuldstændige Betingelse for sin Virkeliggjørelse.

Hvad den Vanskelighed angaaer, der ligger i \emph{Frihedens Forudsætten af sig
selv}, saa bliver dens Løsning nærmest at søge i et Begreb, der foresvævede Kant
idet han udviklede sin Lære om den tidløse, \emph{intelligible Frihedsact}; men
som hos ham, ifølge hans abstracte Sondring af Numenets og Phænomenets Sphære,
fik en næsten mythisk Indklædning; nemlig i Begrebet om \emph{det aandelige}
% NAME NOT LETTERSPACED, p. 138: „Kant“, like the „(Kant)“ of p. 134, is
% close-set Fraktur. Both checked at 4--5x. This book letterspaces most names
% („Aristoteles“, p. 139 note; „Statens“ likewise), so the two Kants are a real
% exception, not an oversight of mine.
% SPLIT SPACING, p. 138: „alt“ at the line end is close-set while
% „saa maa være ved sig selv, altsaa maa forudsætte sig selv“ is spaced from the
% next line on. Transcribed as printed.
% p. 138 carries NO footnote --- the OCR-derived hint list is wrong here --- and
% no Greek.
% --- p. 139 ---
\setcounter{footnote}{0}%
\emph{Individs Concentration af sin Udviklingsheelhed i et ideelt Punkt}, der
ophæver Udviklingens Tidsudvorteshed og Nødvendighed. Dette vil med andre Ord
sige: Løsningen bliver at søge i Begrebet om \emph{et Gjennembrud}, i hvilket
\emph{det Evige}, der i det individuelle Selv skjult har bestemt dettes
Udvikling, sammenslutter sig i sin Enhed som sin egen Værens Forudsætning. Og
dette Concentrationens Moment falder sammen med \emph{det ethiske Valgs Moment},
i hvilket Friheden ved det Eviges Gjennembrud først faaer sand Virkelighed. Men
ved Siden af dette afgjørende Begreb faaer --- hvad der allerede ovenfor blev
antydet, --- det enkelte menneskelige Individs Forhold \emph{til de supplerende
Personligheder og til det Guddommelige}\footnote{Ligesom der til hiint Begreb om
en Concentration af Udviklingen i et Evighedens Punkt findes en Parallel hos
\emph{Aristoteles} i Opfattelsen af Formen (εἶδος) som den, der ved Maalet af
den Udvikling, den selv bestemmer, kun maa siges at være \emph{værende, ikke at
være vordet}, saaledes findes der hos ham en Parallel til Tanken om Frihedens
Realisation gjennem Forholdet til de supplerende Individualiteter og til
Principet, i den Bestemmelse, at det Sædeliges Virkeliggjørelse i Individet
forudsætter en \emph{allerede stedfindende actuel Væren af det Sædelige},
ligesom overhovedet al Vorden forudsætter et actuelt Værende. Hiin Bestemmelse
gjør sig gjældende i hans bekjendte Opfattelse af \emph{Statens} Betydning for
den \emph{individuelle} Sædelighed.} sin væsentlige Betydning for Aporiens
Løsning.

Løsningen af den Vanskelighed, der ligger \emph{i det Approximative i Frihedens
Realisation}, maa søges derigjennem, at \emph{et dobbelt} Synspunkt gjøres
gjældende for Opfattelsen af Friheden, idet den paa een Gang sees som
\emph{Opgave og som værende som Væsensbestemmelse}. At Friheden sees fra dette
dobbelte Synspunkt, er begrundet i dens, og overhovedet i det Ideelles, Væsen.
Idet Friheden, som alt Ideelt, er \emph{værende ved sig selv}, er den kun som
\emph{bestandig selvfrembragt}. Ogsaa i det høieste Forhold gjælder det om den
guddommelige Virkelighed, at den, som levende, ideel Virkelighed, er en evig sig
selv frembringende Virkelighed, paa een Gang værende og τέλος for en Virken.
% printed mark: *) --- the only note on p. 139.
% GREEK, p. 139: τέλος in the body and (εἶδος) in the note, both antiqua italic
% with accents and the perispomeni on the iota, read off the image at 4--5x.
% Neither OCR witness renders any of it.
% The tesseract asterisk after „det høieste Forhold“ is show-through, not type.

% --- p. 140 ---
\setcounter{footnote}{0}%
Men i den menneskelige Friheds Forhold træder den Bestemmelse til, i hvilken
Vanskeligheden ligger: at Selvfrembringelsen er en \emph{tidsbestemt}
Selvudvikling. Naar Friheden derfor bliver Opgave for sig selv, maa den som
Opgave blive et Fyldigere end den var som Forudsætning, og, idet Friheden er en
\emph{Evigheds}bestemmelse, og idet den er \emph{et Ideal}, ikke blot for
\emph{Individet}, men for \emph{Menneskeheden}, viser der sig \emph{en uendelig
Udviklingsmulighed}. Men naar dette er statueret, saa bliver det paa den ene Side
at fastholde, at Forudsætningen dog \emph{i sit Væsen er} liig med Opgaven, at
altsaa den Villie, der endnu virkeliggjør sin Frihed, \emph{i Principet er fri},
trods sin relative Abstraction, og paa den anden Side bliver det ovenfor angivne
Synspunkt at fastholde, at den frie Villie i hvert Realisationsmoment
sammenslutter sin Forbigangenhed i et \emph{ideelt Enhedspunkt}, altsaa tager
Tidsudviklingen ideelt \emph{tilbage i et Evighedens Øieblik}, for, i sidste
Instants, at concentrere sig i ideel evig Enhed, som organisk Moment i det
Altomfattende, at concentrere sig til det „at ville Eet“, forstaaet i Evighedens
Betydning. Saaledes sees det Idealet Characteriserende paa hvert Punkt af
Udviklingen.

Naar det altsaa indrømmes, at Frihedens Virkeliggjørelse, idet den maa skee
gjennem en Udvikling henimod et Ideal, der ikke blot er et individuelt, men et
almeenmenneskeligt, ifølge denne Idealets Bestemthed, maa blive en uendelig
Approximation underkastet, saa bliver tillige, --- idet Forholdet til Idealet er
Forholdet til Villiens τέλος, --- det i Udviklingen Virksomme, dens Entelechie,
\emph{selve Idealet som} Udviklingsenergie ɔ: som Villiens actuelle
Væsensbestemmelse\footnote{Jfr.\ forøvrigt hvad der i det Efterfølgende 3, c vil
blive udviklet om Idealets Realisation.}.

Heri ligger saa, at der vel bliver statueret \emph{Grader i Frihed}; men i den
Betydning, at Friheden, idet den constituerer sig, \emph{som} Frihed udfolder sig
fra den relativ abstractere Form til den concretere, og dette bliver atter
nærmere at bestemme saaledes, at den fra den \emph{abstracte Villen}
% printed mark: *) --- the only note on p. 140.
% GREEK, p. 140: τέλος again, antiqua italic with acute. Read off the image.
% SPLIT SPACING, p. 140: three instances --- „Evigheds“|bestemmelse and
% „Realitets“|bestemmelse-type breaks within a single word; „i sit Væsen er“
% stopping at the line end while „liig med Opgaven“ is close-set; and „selve
% Idealet som“ stopping at the line end while „Udviklingsenergie ɔ: som Villiens
% actuelle Væsensbestemmelse“ is close-set throughout. Transcribed as printed.
% --- p. 141 ---
\emph{af Væsenet} --- en Villen, der kun er som \emph{ethisk} Villen ---
udfolder sig til en \emph{concret Væsensvillen}, til en bevidst, besluttet
Villen af Væsenet efter dets fyldige Indhold. Denne gradvise Udvikling, gjennem
hvilken det \emph{i Beslutningen} (i Valget) aabenbarede
Friheds\emph{princip} vinder Concretion, er tillige en Udvikling \emph{i
Intensitet}, idet Frihedens Energie forøges i samme Forhold som Villien bliver
Udtryk for det concrete, for sig selv klare, Væsens Energie, for det sig selv i
sin Fylde bevidste Væsen.

Dette fører os tilbage til den ovenfor særligt udhævede Vanskelighed, der laae
deri, at Friheden forudsatte en, en bestandig Udvikling underkastet,
\emph{Erkjendelse} som sin Betingelse, da kun ved Erkjendelsen Muligheden er
given af Væsenets Klarhed for sig selv. Her bliver der at vise hen til Forholdet
mellem den \emph{aprioriske} Erkjendelse, der indeslutter det Aprioriske i det
Ethiske, og den \emph{empiriske}. Den aprioriske Erkjendelse har i sin
\emph{relative Abstraction en Afsluttethed i sig} og ved sit Forhold til
Menneskets \emph{Væsen en real Gyldighed}, og den er derfor en saadan, paa
hvilken \emph{en fri} Villen kan hvile. Idet jeg, med Bevidsthed om det Godes
Gyldighed, om Pligtens Fordring, vil det Gode, vil Pligten, er min Villen, som
bestemt ved det Aprioriske, altsaa ved mit eget ideelle Væsens Bestemmelser, som
en bevidst Villen af det Væsensgyldige, \emph{en fri} Villen, og denne frie,
væsensgyldige Villen faaer sin, ved ingen ubevægelig Grændse indskrænkede,
Udfyldelse gjennem den Erkjendelse, der, idet den hviler paa det Aprioriske ag i
dette har Selvets og det Reales abstracte Væsen, trænger ind i det empirisk
givne Reale, her altsaa nærmest i det Reale i os selv, i vort eget Væsens
Realitetsbestemmelser og Realitetsforhold, som assimilerende Erkjendelse. Under
denne Erkjendelsesudvikling er altsaa den Erkjendelse \emph{paa ethvert Punkt}
tilstede, der gjør Villien fri, ogsaa hvor Villien forholder sig til det ikke
erkjendte Concrete.

\begin{center}\rule{2cm}{0.4pt}\end{center}
% RULE, p. 141: the page ends short and a short centred SINGLE hairline closes
% the argument below the last line, the rest of the leaf blank. No head follows
% --- p. 142 simply opens a new paragraph --- so this is a section-closing rule
% of the kind at pp. 89, 143 and 201, not a head ornament. Transcribed.
% PRINTER'S DEFECT, p. 141: „paa det Aprioriske ag i dette“ --- „ag“ for „og“.
% Both OCR witnesses read „ag“; at 5x the first sort is unambiguously a Fraktur
% a (flat shouldered top), not the closed round o of „og“ elsewhere on the page.
% Transcribed as printed.
% SPLIT SPACING, p. 141: „Friheds“ is close-set and only „princip“ is spaced,
% within the one word. Transcribed as printed.
% p. 141 carries no footnote and no Greek; the page ends short, closing the
% argument before the new paragraph on p. 142. Its closing rule is set above.

% --- p. 142 ---
Idet i det Foregaaende Frihedens Realitet hævdedes overfor den Tvivl, der
støttede sig paa Individets Forhold til det Ydre, var det afgjørende Begreb
Begrebet om det menneskelige Individs Concentration i sig som en ved
Væsensuniversaliteten bestemt relativ Totalitet med indre Uendelighed, og dette
Begreb fik ogsaa sin væsentlige Betydning i de efterfølgende Udviklinger.
Gjennem hint Begreb om Selvconcentrationen ophævedes Ufriheden i Forholdet til
det Ydre, og omsattes i Selvbestemmelse ud fra Væsensheelheden, i hvilken den
ideelle Aandstotalitet er i Enhed med den organiske Totalitet. Men idet nu
Individet i denne Sammenslutning i sin relative Totalitet sees som bestemmende
sig ved Frihed, saa bliver \emph{Individets} Virken, der udtrykker denne
Selvbestemmen, med det Samme at see som indordnet \emph{i den absolute
Totalitet}. Dette gjælder, hvad enten Friheden opfattes efter sin usande eller
sin \emph{sande} Form. Er Friheden kun den \emph{formelle} Frihed,
\emph{Vilkaarligheden}, den indre Uendeligheds abstracte Udtryk, den
væsensbestemte Friheds \emph{Skin}, den tilsyneladende grundløse Bestemmen, saa
er der her netop bestemte, endelige, om end ikke klart bevidste Motiver
tilstede, der have Sammenhæng med det Reales Totalitet og drage Villiens Virken
ind i denne. Er Friheden derimod den \emph{sande, væsentlige} Frihed, saaledes
at dens „absolute Begynden“ har et væsensbestemt, substantielt, Indhold, saa er
der, i et tidløst Forhold, en Sammenhæng med Tilværelses-Totaliteten og dens
evige Princip, en Bestemmen, der gaaer i Enhed med dettes evige Bestemmen.

Menneskets Frihedsforhold til det Ydre viser saaledes, gjennem Totalitetens
Begreb, tilbage til \emph{Forholdet til det Absolute}. Og kun ved dette Forhold
er det jo, at det individuelle Væsen, der som frit er værende for sig i
Uendelighedens Form --- i hvilken Principets Væren i det reflecterer sig --- kan
have Muligheden til at unddrage sig den ydre Causalitetsnødvendigheds
Herredømme, men saaledes, at det med det Samme, ved selve sin Sætten af en
absolut Begynden, er ordnet ind i den principbestemte Tota\-%
% p. 142 carries NO footnote --- the OCR-derived hint list wrongly lists it ---
% and no Greek. Both of tesseract's commas at „den organiske Totalitet.“ and
% „ind i denne.“ are full stops on the image (checked at 3x); the ABBYY witness
% agrees.
% --- p. 143 ---
litet og undergivet dens Dialektik. Thi naar det Absolute er Menneskeaandens
Princip, saa bliver Menneskets Frihed at see \emph{som begrundet i det}, det
Eneste, der i \emph{absolut Betydning er frit}, og \emph{Friheden i dens
Sandhed} kan da kun være Eet med den \emph{frie absolute Hengivelse} til det
Absolute som den gjennemvirkende Magt. Men naar da Friheden som
\emph{Menneskets} Frihed bliver at see som \emph{afledet}, og dog skal hævdes
som Frihed; naar den bliver at see som værende ved det Absolute, medens den dog
i samme Øieblik, idet den udtrykker Menneskets relative Selvstændighed, viser
sig at indeslutte Muligheden til at skille sig fra det Absolute; --- saa
fremtræde der, gjennem Forholdet til Principet, nye Tvivlsgrunde mod Friheden.

Idet disse Tvivlsgrunde skulle prøves, føre de os ind paa Undersøgelsen af de
andre ovenfor udhævede Problemer. Problemet om den \emph{afledede Frihed}, og om
dens \emph{dobbelte} Virkelighedsform: som hævdende sig selv i
\emph{Selvraadighed}, og som selvstændig sammensluttende sig med Principet i
\emph{Selvhengivelse}, falder væsentligen sammen med Problemerne \emph{om det
Ondes Mulighed og om det Ondes Ophævelse i Forsoningen}.

Problemet om den menneskelige Frihed fører os altsaa, idet de nye Tvivlsgrunde
skulle hæves, atter ind paa det Ethisk-Religiøses Omraade, til hvilket vi
ovenfor vistes hen hvor Individets Selvconcentration viste sig som identisk med
det Eviges Gjennembrud i det Valg, ved hvilket det Ethiske constitueredes i
Individet. Vi vise derfor paa dette Punkt tilbage til den tidligere, ---
anticiperende og derfor foreløbig kun hypothetisk gyldige --- Udvikling af det
Ethiskes Væsen og særligt af Villiens Constitueren af sig som ethisk Villie i
det ethiske Valg, i hvilket hint Gjennembrud skeer, hvori Friheden giver sig
Virkelighed, og, idet den giver sig Virkelighed, forudsætter sig selv som
Væsensbestemmelse og sætter sig som uendelig Opgave.

\begin{center}\rule{2cm}{0.4pt}\end{center}
% RULE, p. 143: the sub-section „1. Om Villiens Frihed“ closes with a short
% centred rule (about half an inch, a tenth of the measure) below the last line,
% the rest of the leaf blank. It is a SINGLE hairline, as at pp. 89 and 141 ---
% not the double rule of pp. 29, 63, 201 and 210. It closes a division and stands
% next to no head, and so IS transcribed. The head „2. Om det Ondes Væsen og
% Mulighed.“ on p. 144 carries no rule of its own.
% p. 143 carries no footnote and no Greek.

% --- p. 144 ---
\section*{2. Om det Ondes Væsen og Mulighed.}
\emph{Det Ondes} Begreb maa, som Begrebet om et \emph{Negativt}, gjøres tydeligt
gjennem Begrebet om det \emph{Positive}, det negerer. \emph{Det Godes} Begreb,
som Begrebet om denne positive Modsætning, maa altsaa danne Udgangspunktet for
Opfattelsen af det Ondes.

\emph{Det Godes} Begreb er allerede i det Foregaaende (S.\ 76) blevet berørt ved
Udviklingen af det Absolutes Idee. Idet det Absolute bestemtes som Tilværelsens
Princip, bestemtes det som det, der sætter den reale Væren som bestemt ved det
Ideelle, og det saaes med det Samme som det \emph{teleologisk} Virkende, og i
denne Virken som Selvøiemed, som Maal for sin subjectiv-ideelle Activitet, og
som i sin evige Selvvirkeliggjørelsesproces omfattende og beherskende de reale
Momenters Totalitet. I disse Bestemmelser af det Absolute laae
Tilknytningspunktet for Opfattelsen af det Absolute som \emph{det Gode} (det
Godes Idee), hvilket Begreb fra det abstractere Metaphysiske udviklede sig
henimod det Metaphysisk-Ethiske.

I Modsætning til det \emph{Gode i dets metaphysiske} Bestemthed bliver det Onde
det \emph{metaphysisk Onde}. Udtrykte nu det Gode i denne Bestemthed
\emph{Ideen}, det \emph{Absolute, som Selvøiemed, som Maal} for sin \emph{evige}
Realisationsproces, saa kan det \emph{Onde}, --- da Ideens Proces omslutter al
Væren, og det Onde derfor kun kan udtrykke Noget, der falder indenfor denne
Proces, --- idet det maa være et Saadant, der staaer i Modsætning til Maalet,
kun faae den almindelige Bestemmelse, at være \emph{Momentet i dets Isolation},
hvilket, som saadant, i Modsætning til det absolute Værende, bliver et blot
\emph{Relativt}, et ved \emph{Negationen} (Determinationen) Bestemt, der fæstner
sig for Tanken i sin Bestemmethed ved et μὴ ὄν.

Men denne abstracteste Almeenbestemmelse af det Onde antager nu en nuanceret
Form efter den forskjellige Opfattelse af det Absolute.

Opfattes det Absolute fra et \emph{abstract Evighedssynspunkt}, saaledes at dets
Proces kun sees \emph{\textit{sub specie}}
% HEAD, p. 144: „2. Om det Ondes Væsen og Mulighed.“ --- bold (halbfette)
% Fraktur, centred, wording verified at 2x. It agrees with the Indhold. No rule
% stands under it; the rule belongs to the foot of p. 143.
% GREEK, p. 144: μὴ ὄν --- grave on the eta, smooth breathing with acute on the
% omicron. Tesseract renders it „uy dv“ and ABBYY „^ o>“; read off the image at
% 4x.
% ANTIQUA LATIN, pp. 144--145: „sub specie æterni“ is antiqua AND letterspaced,
% like the „causa sui“ of p. 68 and unlike the close-set antiqua of pp. 64, 71
% and 134. Encoded \emph{\textit{...}}; it breaks across the page.
% p. 144 carries no footnote.
% --- p. 145 ---
\emph{\textit{æterni}}, saa er det Onde, der fra dette Synspunkt skal sees som
\emph{nødvendigt} Moment i Processen, fordi denne kun er Proces ved Udviklingen
af Momenterne som \emph{saadanne, evigt hævet} ligesom det \emph{evigt} er sat.
Det bliver Ideens Redskab, der fra Ideens Synspunkt ikke kan sees \emph{som
Ondt}, fordi det evigt er draget ind i Processen, men kun som et evigt behersket
Moment. Der kommer da en \emph{Modsigelse} ind, idet det kun kan sees \emph{som
Ondt}, naar det ligesom betragtes nedenfra, fra et Synspunkt, \emph{udenfor og
under} Ideen, og fra dette Synspunkt fastholdes i Isolerthedens Abstraction. Det
metaphysisk Onde, saaledes opfattet, faaer nemlig kun Virkelighed \emph{for den
menneskelige Tanke}, der selv, opfattet fra dette abstracte Standpunkt, ligesom
den menneskelige Existents overhovedet, bliver et \emph{kun Forsvindende}, et i
sin Væsenløshed Ubegribeligt (jfr.\ Spinoza), og det Onde faaer altsaa kun
\emph{en illusorisk Væren for et Forsvindendes og Væsenløsts Tanke}. Og medens
det Onde bliver uvirkeligt, kommer det \emph{metaphysisk Gode}, ifølge
Standpunktets Abstraction, til at lide af den Mangel, at det i sin abstracte
Evighed bliver et abstract Værende, hvorved Øiemedets Bestemmelse bliver
illusorisk.

Skal det Onde blive Mere end et blot forsvindende Skin i en evig Proces, saa
fordres der en Opfattelse af det Absolute, ved hvilken det bliver muligt, at den
\emph{reale Virkelighed hævder en relativ Selvstændighed i Forholdet} til det
Guddommelige.

Ved en saadan Opfattelse maa der da falde et Efter\emph{tryk paa det reale
Endeliges} Begreb, og det Absolutes Proces bliver concretere opfattet, idet den
sees som omfattende det Reale i dets Tidsbestemthed, saaledes at dette i denne
Bestemthed faaer en \emph{væsentlig} Betydning for det Absolute. Her fæstnes
Endeligheden relativt i Kraft af de reale Existentsformers Væsentlighed. Men
ogsaa naar det Endelige saaledes sikkres mod den umiddelbare Forsvinden i det
Evige, og i sin relative Selvstændighed giver en Plads for det Onde, saa vil dog
dette, saalænge det sees \emph{udenfor}
% NAME, p. 145: „(jfr. Spinoza)“ is Fraktur and close-set --- neither antiqua
% nor letterspaced. Checked at 4x. Worth recording, since the book letterspaces
% „Aristoteles“ in the note to p. 139.
% SPLIT SPACING, p. 145: „Efter“ is close-set and only „tryk paa det reale
% Endeliges“ is spaced, the break falling inside the word. Transcribed as
% printed.
% p. 145 carries NO footnote --- the OCR-derived hint list wrongly lists it ---
% and no Greek. The „10“ at the foot is the gathering signature; not
% transcribed.
% --- p. 146 ---
\emph{Villien} eller i Villien som bevirket ved en \emph{Natur}proces, endnu
væsentlig beholde det \emph{metaphysisk} Ondes Charakteer, og være behæftet med
dettes Uvirkelighed. Det Onde bliver da den reale Tilværelses \emph{Endelighed,
Mangel, Ufuldkommenhed}, τὸ μὴ ὄν i det Værende, og sees som saadant i den
blotte \emph{Nødvendigheds} Lys, som begrundet i Ideebevægelsens Lov, men
tillige som behersket ved denne og draget ind i Heelhedens Harmonie, og om det
end, som det „\emph{physisk} Onde“, faaer en concretere Skikkelse end som det
metaphysisk Onde, faaer det, ogsaa naar det opfattes som \emph{Privation}, dog
kun en Væren som Ondt for den Bevidsthed, der, forglemmende Tilværelsens
Nødvendighed, dømmer det Enkelte efter et abstract Almeenbegreb eller efter dets
Forhold til Individet som fornemmende. Og denne det Ondes Væren faaer i sidste
Instants kun Virkelighedsbetydning, idet \emph{Villien giver den Betydning for
Personligheden}.

Den concretere Opfattelse af det Metaphysiske viser saaledes, naar det Onde skal
begribes i sin Virkelighed og i sin fulde Modsætning til det Gode, gjennem det
Reale hen \emph{til Villien og det Ethiske}. Først i den endelige Aands Sphære,
hvor vi faae \emph{en af Villien med Frihed sat Modsætning til det Absolute},
faae vi et Ondt i \emph{egentlig} Betydning: først som det \emph{ethisk} Onde
har det \emph{Onde Virkelighed}.

Ligesom det metaphysisk Ondes Bestemmelse belystes ved det metaphysisk Godes,
saaledes maa ogsaa det \emph{ethisk Ondes Bestemmelse} belyses ved det
\emph{ethisk Godes}. Og idet nu det Godes metaphysiske Bestemmelse skal omsættes
i en ethisk, saa maa Forholdet blive det, at hiin Bestemmelse vedbliver at være
Grundlaget. Ligesom vi i Udviklingen af det Absolutes Begreb (S.\ 74 ff.) saae
den concretere metaphysisk-ethiske Bestemmelse fremgaae af den abstractere
metaphysiske, saaledes maae vi, naar det Gode skal bestemmes som det ethisk
Gode, det menneskeligt Gode, see denne Bestemmelse hvile paa den metaphysiske.
Det Gode maa overalt, for i Sandhed at være det Gode, bære \emph{Ideens}
Bestemmelse i sig: alt Godt er kun Godt ved at udtrykke Ideen. Men det ethisk
Gode i sin \emph{concrete} Bestemthed forholder sig til den \emph{menne}\-%
% GREEK, p. 146: τὸ μὴ ὄν, antiqua italic with accents and breathings; tesseract
% gives „70 uy ôv“ and ABBYY drops it entirely. Read off the image at 4x.
% SPLIT SPACING, p. 146: „Natur“ is spaced and „proces,“ that completes the word
% is close-set. Transcribed as printed.
% A round brown ink blot (about 2 mm) follows „Onde Virkelighed.“ --- a stain in
% this copy, not type; tesseract reads it as „#“. Modern pencil marginalia stand
% in the left margin; not transcribed.
% p. 146 carries no footnote. „Grundlaget.“ is a full stop, not the comma
% tesseract reports.
% --- p. 147 ---
\emph{skelige Villie}, det maa altsaa blive et \emph{Ideens Udtryk gjennem
Villien, et Villiens Ideal}.

Naar vi saaledes lade Forholdet til det Menneskelige, bestemtere: til den
menneskelige Villie, være det, ved hvilket det Gode fra det metaphysisk Gode
bestemmer sig til det ethisk Gode, saa bliver altsaa det Absolute seet som
ethisk Ideal ifølge dette Forhold til Villien som dens Ideal; men idet Forholdet
til den menneskelige Villie, ligesom overhovedet Forholdet til det Reale,
grunder \emph{i selve det Absolutes Væsen}, bliver Bestemmelsen af ethisk Ideal
ikke et ligesom udenfra tilføiet Prædicat for det Absolute, men \emph{oprindelig
Væsensbestemmelse i det}.

Det \emph{ethisk} Godes \emph{første} Bestemmelse bliver da, fra dette
Udgangspunkt, \emph{Virkeliggjørelsen af Ideen (det Godes Idee) ved den frie
menneskelige Villie}. Men denne Bemmelse faaer saa, som det paa et tidligere
Sted (Afsnit II, 3) er blevet paaviist, \emph{det andet Udtryk}, at det
\emph{ethisk} Gode, som det \emph{menneskeligt} Gode, er Virkeliggjørelsen af
\emph{Menneskets Idee}, ɔ: af \emph{Menneskets totale, ideale Væsen}. Denne
Omformning af Udtrykket begrundedes derved, at den \emph{absolute Idee} som den
menneskelige Handlings Norm \emph{er i det menneskelige Ideals Bestemthed},
hvorfor det, at virkeliggjøre Ideen, \emph{for Mennesket} er ligt med at
virkeliggjøre \emph{sit ideale Væsen} gjennem den frie Villiesbeslutning,
gjennem den ved Erkjendelsen for sig klarede Væsensvillen. \emph{Det Absolute}
bliver Øiemed, idet det \emph{menneskelige Væsen} bliver Øiemed, og det har
viist sig, hvorledes Ethiken ikke kunde naae ud over denne Bestemmelse af det
Menneskelige som det Menneskeliges Opgave, naar Opgaven bestemmes
\emph{positivt}, og at den menneskelige Villie, ifølge det menneskelige Individs
Grundforhold til Tilværelsesheelheden og dens Princip, kun kan udtrykke det
Absolute \emph{som saadant i Negativitetens}, i den ubetingede Resignations og
Selvhengivelses Form.

Den specifiske \emph{Realitets}bestemmelse, ved hvilken det metaphysisk Gode
bliver til det \emph{ethisk} Gode, er altsaa \emph{den menneskelige Villie som
Væsensvillen}, ɔ: som be\-%
% PRINTER'S DEFECT, p. 147: „Men denne Bemmelse faaer saa“ --- „Bemmelse“ for
% „Bestemmelse“, a dropped „ste“. Both OCR witnesses read „Bemmelse“ and the
% image at 3x shows B-e-m-m-e-l-s-e with no gap. Transcribed as printed; the
% Trykfeil leaf does not list it.
% ANTIQUA NUMERAL, p. 147: the „II“ of „(Afsnit II, 3)“ is set in antiqua roman
% capitals, not in the Fraktur I/J sort. Left unmarked here, since the house
% \textit{} convention is for Latin words; recorded for a later pass.
% SPLIT SPACING, p. 147: „Realitets“ is spaced and „bestemmelse,“ close-set,
% within the one word. Transcribed as printed.
% p. 147 carries no footnote and no Greek. The „10*“ at the foot is the
% gathering signature; not transcribed.
% --- p. 148 ---
vidst, besluttet Virkeliggjøren af Væsenet. Dette menneskelige Væsen, der i
Kraft af Væsensuniversaliteten er udfoldet til \emph{Selvbevidsthedens} og
\emph{Frihedens} Forsigværen, have vi seet som \emph{naturligt-aandeligt}, som
realiserende sin Idealitetsbestemthed gjennem og paa Grundlaget af en \emph{real
Organisme}, og vi have seet det som \emph{relativ Totalitet} indenfor den
universelle Totalitet. Heri ligger det menneskeligt Godes, det ethisk Godes,
nærmere Bestemmelser. Dette menneskeligt Gode, i hvilket Selvet vil det Absolute,
det Godes Idee, idet det vil sit eget ideale Væsen, er ved og i en Villiesact,
der i det Menneskelige danner Analogien til den teleologiske
Subjectivitetsenergie i Ideens Proces, og idet Villien vil dette Gode, vil den
sit Væsen i det normale Forhold mellem dets Væren som relativ Totalitet og som
Moment i den universelle Totalitet, mellem Aandssiden og Natursiden i det.

Ud fra det saaledes bestemte Begreb af det ethisk Gode dannes nu, gjennem
Modsætningen, Begrebet af det \emph{ethisk Onde}, der kun havde et abstract
Forbillede i det Metaphysiskes Sphære, idet i Ideens Proces Momenternes
nødvendige Sætten som Momenter i deres Relativitet fastholdtes i Abstractionen,
eller, hvad der udtrykker det Samme, idet Negationen, det Ikke-Værende, fæstnede
sig uopløst for den abstraherende Tanke.

Bestemmelsen af det ethisk Onde fører altsaa tilbage til det, der var det
Afgjørende for det Ethiske, \emph{til den frie Villie}, og maa søges ud fra
denne, \emph{ikke} ud fra Natur\emph{bestemtheden} i det Menneskelige, endnu
mindre ud fra noget absolut \emph{udenfor} Mennesket Værende, dette være nu et
\emph{høiere godt} eller et \emph{høiere ondt} Væsen. Den første Bestemmelse for
det bliver da den, at det er et ved \emph{Villien sat Misforhold i det aandelige
Liv}, og hvorledes dette Misforhold bliver at opfatte, ligger implicite i det
ovenfor Udviklede. Var nemlig det Gode bestemt ved det normale Forhold mellem
Menneskets Væren som relativ Totalitet og som Led i den universelle Totalitet, og
mellem Natur- og Aandssiden i det, og er dette sidste Forholds normale Form,
% SPLIT SPACING, p. 148: „Naturbestemtheden“ is close-set in „Natur“ and spaced
% only from „bestemtheden“, within the one word (checked at 300 dpi, 1.75x, on
% the line „ikke ud fra Natur-“). Transcribed as printed, hence
% Natur\emph{bestemtheden}.
% EMPHASIS, p. 148: „Villiesact“ is TIGHT, though ABBYY renders it spaced;
% checked at 300 dpi. „til den frie Villie“ is spaced from the „til“ onwards.
% p. 148 carries no footnote and no Greek.
% --- p. 149 ---
\setcounter{footnote}{0}%
hvad der ikke behøver nogen nærmere Udvikling, betinget ved det førstes normale
Form, saa bliver det Misforhold, der constitueres ved den onde Villen og
udtrykker den, en ved Villien sat Forstyrrelse i disse to Forhold, og,
\emph{oprindelig i det første} af dem; altsaa, efter sin \emph{primitive Form},
et \emph{Misforhold mellem Individets Egenvillie og det Absolutes gjennem
Menneskets Væsensbestemmelse udtalte Fordring til Villien.}

Dette Misforhold kan nu ogsaa udtrykkes paa følgende Maade. Er det Gode det
\emph{Grundvæsentlige}, det \emph{i og for sig Gyldige, der af Villien sættes som
saadant, idet det er det ubetinget Villede,} saa bliver \emph{det Onde i ethisk
Betydning dette Grundvæsentliges, Iogforsigværendes, Modsætning, men saaledes, at
det af Villien sættes som det Iogforsigværende, Absolute, ubetinget
Gyldige}\footnote{I disse Bestemmelser er \emph{Objectivitetens} Moment i det
Gode og det Onde accentueret: de ere bestemte som et ved Villien
\emph{Realiseret, Virkeliggjort}. Men fra Objectivitetens Bestemmelse vises der
netop gjennem dette „ved Villien“ tilbage til \emph{Subjectivitetens} som den
\emph{primitive}, og Objectiviteten bliver en saadan, der bæres af Villien, har
sin Realitet i Villien. Ligesom det ethisk Gode kun er ved denne, saaledes er det
ethisk Onde kun i og ved den onde Villie. Der kan vel tales om et i det Ydre
virkeliggjort Usandt; men det bliver bestandig kun et Ondt ved Relationen til
Villien, til hvilken det forholder sig som Virkning eller Incitament.}. Den onde
Villie bestemmer da ikke blot det Ikke-Iogforsigværende som et Iogforsigværende
og gjør det paa denne Maade til sit Øiemed, men den bestemmer \emph{sig} gjennem
denne Bestemmen. Idet den nemlig vil det Onde som \emph{sit Øiemed}, har den det
til \emph{sin} Bestemmelse. Den er bestemt ved det Forgængelige, Væsenløse,
\emph{Endelige}, idet den selv bestemmer sig for det; men idet den saaledes er
endelig bestemt, har den \emph{Uendelighedens} Bestemmelse i Abstractionens og
Negativitetens Form, og netop deri den indre Modsigelse, \emph{at den ved den
formelle Uendelighed bestemmer sig endeligt,} altsaa giver sig en Bestemthed, der
er incommensurabel med det Bestemmende og derfor atter maa opløses af Villien.
% FOOTNOTE, p. 149: printed mark *). The OCR-derived hint list does not name
% p. 149; the note is there. Its „ved Villien“ is the batch's first quotation
% pair.
% EMPHASIS, p. 149: within the long spaced run „i og for sig Gyldige … ubetinget
% Gyldige“ only „saa bliver“ is close-set; transcribed as printed.
% I/J, p. 149: „Iogforsigværende(s)“ and „Ikke-Iogforsigværende“ take the I of
% the house convention, as at p. 122 („det Iogforsigværende“).
% p. 149 carries no Greek.
% --- p. 150 ---
\setcounter{footnote}{0}%
Misforholdet i Selvet, der constitueres ved det Onde, faaer da ogsaa det Udtryk,
at Personligheden i den onde Villen ved en \emph{Uendelighedsact bestemmer sig
til Endelighed, at den vil sig i Endelighed}\footnote{Denne Bestemmelse gjælder,
hvad der nedenfor vil blive paaviist, ogsaa naar Selvet vil sig efter sin
abstracte, tomme Uendelighed.}, men vil denne Endelighed \emph{som et
Uendeligt.} I Kraft af den \emph{formelle} Frihed vil den altsaa sig, men
\emph{ikke} efter sin \emph{væsentlige, ideale} Bestemmelse, men efter en
\emph{usand}; den vil sig altsaa \emph{i Kraft af Friheden}, men \emph{ikke som
fri}, men, ubevidst, som et Ufrit, Behersket.

I de givne Begrebsbestemmelser af det Onde, der gjælde overalt, hvor Villien er
naaet ud over den blot instinctmæssige Bestemmethed, over den umiddelbare Drifts
Form, og hvor Bevidstheden om den ethiske Fordring er vakt, ligger Muligheden af
\emph{to Hovedformer} af det, der dog ikke simpelthen ere coordinerede, men af
hvilke den ene maa blive den egentlige Grundform. De svare til de to væsentlige
Forhold, hvis normale Form constituerede det Gode: Forholdet mellem Mennesket som
Totalitetsled og Totalitetsprincipet, det Absolute, og Forholdet mellem Natur- og
Aandssiden i Mennesket. Den Form af det \emph{Onde}, der svarer til det første,
\emph{primitive} Forhold, er \emph{Selvraadigheden} (Trodsen); den, der svarer
til det andet, er \emph{Fornedrelsen, Fortabelsen. Trodsen} er
\emph{Grundformen}, fordi Forholdet til det Absolute er det primitive, og fordi
Mennesket \emph{som Aand i enhver} Form af det Onde gjør en \emph{activ} Villen
gjældende, saa at ogsaa \emph{Fornedrelsen}, selv i sit \emph{Maximum} som dyrisk
Aandløshed, kun kan være \emph{som villet}, altsaa som indesluttende
\emph{Trodsen}. Ogsaa naar Villien synker tilbage i relativ Ubevidsthed, bliver
dette Synspunkt gjældende: ikke blot \emph{Tilstanden som Heelhed} er forskyldt
ved Villien, der ligesom ved en enkelt Act har sat denne Tilstand, og derved
bærer Ansvaret for hvad der skeer i den, men i \emph{ethvert} actuelt Udtryk for
Tilstanden er, fordi Mennesket er \emph{Aand}, en \emph{fornyet} Villen, og selv
% FOOTNOTE, p. 150: printed mark *); the hint list does not name p. 150 either.
% EMPHASIS, p. 150: „formelle“ is spaced but the following „Frihed“ is close-set,
% and „primitive“ is spaced but the following „Forhold“ is not (both checked at
% 300 dpi, 1.75x). „sig“ in „vil den altsaa sig“ is close-set, though ABBYY
% renders it spaced.
% p. 150 carries no Greek.
% --- p. 151 ---
den tilsyneladende Liden bliver at see som en Handlen, som \emph{Trodsens
Handlen.} At der i Forholdet mellem det Ondes to Hovedformer er et Dialektisk,
ifølge hvilket, ligesom al Fornedrelse og Svaghed indeslutter Trodsen, omvendt,
al Trods har et Svaghedens, et Fornedrelsens Moment i sig, modsiger ikke det, at
Trodsen er det Ondes primitive Form. Som saadan er den ogsaa bleven opfattet af
den religiøse Bevidsthed, der i Forestillingen om Lucifer og om Syndefaldet søger
det Ondes Kilde i Hovmod og Selvraadighed. I det efterfølgende Afsnit: om det
Ondes Ophævelse, vil Udviklingen derfor væsentlig holde sig til denne det Ondes
Grundform, der altsaa opfattes som værende tilstede i \emph{alt Uethisk}, saa at
dette i \emph{alle} sine Phænomener maa skjule \emph{Trods, Selvraadighed,
Egoisme, Had.}

\emph{Muligheden og Grunden} til det saaledes bestemte ethisk Onde søge vi i
selve den Grundbestemmelse i det menneskelige Væsen, ifølge hvilken det, i sin
Universalitet, som ideelt-realt Led i Tilværelsens organiske Heelhed, er en i
væsensgyldig Realitetssondring bestaaende relativ Totalitet, der har Mulighed til
Bevidsthed, om end i den relative Abstractions Form, om sin indre Uendelighed, og
til en fri Forsigværen. Men denne det \emph{individuelt Ideelles Mulighed til en
uendelig Væren for sig er en i sig fordobblet Mulighed:} den er Muligheden til
den \emph{sande udfyldte Forsigværen i Hengivelsen til og Tilegnelsen af det
gjennemvirkende Universelle,} og Muligheden til en \emph{Afslutten i sig i sin
væsentlige Realitetssondring fra det Absolute, i en indre Uendelighed, der i
Sondringen fra det sande Uendelige bliver en abstract, tom, ved Løsrevetheden
endelig,} og derfor selvmodsigende \emph{Uendelighed}, hvis illusoriske
Udfyldning maa søges i det for den Tilfældige, der kun har \emph{Endelighedens}
Bestemthed. I hiin \emph{første} Forsigværen vil Individet sig som indordnet i
Heelheden, underordnet Principet, det vil sig altsaa som dettes eiendommelige
Organ, og kun saaledes kan det, hvad der ligger i Grundforholdet, ville sig efter
sit Væsen. Men idet
% RULE, p. 151: the page carries NO rule at its foot, and none is expected ---
% the subsection „2. Om det Ondes Væsen og Mulighed“ does not end here but runs
% on into p. 152, and the page breaks in the middle of a sentence.
% p. 151 carries no footnote and no Greek.
% --- p. 152 ---
Individet vil sig som Organ for det Ubetingede, vil det \emph{det Gode}, og dette
er ved den Villie, der er \emph{det Guddommeliges Villen i Individet i sand Enhed
med Individets Villie.} Men i hiin \emph{anden} Forsigværen, i hvilket Individet
vil sig gjennem Egenvillien, vil det sig ikke som Led i Heelheden, men vil sig i
sin Endelighed som et Absolut; det negerer altsaa Principet, \emph{vil det ikke},
men hævder sin Villen af sig imod det, men som den \emph{onde}, og ved denne
onde, det \emph{Gode} negerende Egenvillie negerer det \emph{selve sit Væsen.}

I den \emph{selvsamme} Grundbestemmelse altsaa: i den sig i sin Uendelighed
bevidste individuelle relative Totalitets Bestemmelse, ligger altsaa baade det
\emph{Godes} og det \emph{Ondes} Mulighed: Muligheden \emph{til en Forsigværen i
det Absolute, til Selvhengivelsens Selvbekræftelse,} og Muligheden \emph{til
Isolationens abstracte Forsigværen, til Selviskhedens Selvbekræftelse,} og Selvet
har efter sit Begreb \emph{Valgfrihed i Forhold} til denne dobbelte Mulighed. Men
gjennem de Bestemmelser, der ere givne af det Ondes Væsen og Mulighed, er det
allerede antydet, hvorledes det Onde \emph{i selve sit Væsen} bærer \emph{Spiren}
til en \emph{Selvopløsning}, og gjennem Paaviisningen af den og af dens
Udfoldelse vil \emph{Forsoningens Mulighed} aabenbare sig.

\section*{3. Det Ondes Selvopløsning som Forsoningens Mulighed.}

Ved Udviklingen af dette Punkt gaae vi ud fra de abstractere metaphysisk-ethiske
Begrebsbestemmelser, der saa i de efterfølgende Partier ville faae en concretere
Form.

Den \emph{Selviskhedens Selvbekræftelse}, ved hvilken det \emph{ethisk Onde}
sattes, og som vi have seet under Synspunktet af en \emph{Forrykken} af Aandens
Grundforhold, som en Villen af det Endelige som det Absolute og som en det
Absolutes Villen negerende Villen, viste sig at have sin Mu\-%
% HEAD, p. 152: „3. Det Ondes Selvopløsning som Forsoningens Mulighed.“ ---
% bold (halbfette) Fraktur, set full measure, wording verified at 300 dpi. It
% agrees with the Indhold.
% RULE, p. 152: there is NO printed rule between the close of „2. Om det Ondes
% Væsen og Mulighed“ and this head. The mark standing in that gap is modern
% PENCIL: a bracket (horizontal stroke with a rising tick at the left end and a
% falling tick at the right) with an oblique lead-in stroke, pale grey and quite
% unlike the printed rule of p. 159. Compare the pencil marginalia at pp. 16 and
% 112. Not transcribed.
% p. 152 carries no footnote and no Greek.
% --- p. 153 ---
\setcounter{footnote}{0}%
lighed i Selvets Evne til en \emph{Villen af sig selv efter sin abstracte
Uendelighed.} Men denne sidste Bestemmelse indeslutter et eiendommeligt
dialektisk Forhold til det \emph{sande Uendelige}, det Universelle og Absolute,
og ved dette Forhold bærer den onde Villie sin Negation i sig.

Selviskhedens Selvbekræftelse er ved den \emph{formelle Uendelighed}, i Kraft af
hvilken Selvet afslutter sig i sig som absolut. Men denne Uendelighed har Selvet
som relativ Totalitet kun ved det universelle Totalitetsprincips, det
\emph{Absolutes}, Væren i det. Idet det altsaa, i sin real gyldige Væren som Led
i Totaliteten, bekræfter sig i sin Isolerethed, saa er dets Selvbekræftelse kun
ved en Kraft, \emph{der ikke er dets egen}, fordi den kun har den fra et Høiere.
I sin tilsyneladende største Uafhængighed er Selvet kun virkende ved det
\emph{Absolute}. Idet Villien sætter sig som ond, sætter den sig derfor kun
tilsyneladende udenfor og i Uafhængighed af det \emph{Absolute}: fra et dobbelt
Synspunkt viser sig den vedvarende Afhængighed, gjennem hvilken Selvet bliver i
en Selvmodsigelse. Den viser sig \emph{først} derved, at det villende Selv, der
ved en Uendelighedsact vil sig i sin Isolerethed, i denne Isolerethed sætter sig
i et endeligt Afhængigheds- og Modsætningsforhold til det øvrige af det
Universelle omfattede Endelige\footnote{Dette viser sig fra de forskjellige
Synspunkter, baade naar Ufriheden opfattes som \emph{Driftens} Ufrihed, som
\emph{Vilkaarlighedens} Ufrihed, og som \emph{Handlingens} (Villiesudførelsens)
Ufrihed.}, gaaer ind under de virkende Aarsagers Lov, og bliver, medens det i sin
isolerede Enkelthed er magtesløst ligeoverfor det Endeliges Sum, Led i en ufri
Causalsammenhæng, der har sin sidste Grund i det Absolute, som, fra dette
Synspunkt viser sig som \emph{Naturmagt}, som blind, beherskende, overvældende
\emph{Nødvendighed}, som \emph{Skjæbne}. Saaledes er den tilstræbte Uafhængighed
omsat i Afhængighed, Friheden i Ufrihed, Absolutheden i Relativitet. --- Fra et
\emph{andet} Synspunkt viser den selviske, onde Villies Afhængighed af det
Absolute sig paa følgende Maade. Forat ville sig selv i en Form, der kunde
% FOOTNOTE, p. 153: printed mark *); the hint list does not name p. 153.
% EMPHASIS, p. 153: the spaced run opens at „Villen af sig selv“, not at
% „i Selvets Evne“ (the latter is close-set; checked against ABBYY and the image).
% p. 153 carries no Greek. Tesseract's stray quotation mark after „sande
% Uendelige,“ is show-through, not a printed sort (checked at 3x).
% --- p. 154 ---
svare til Uendeligheden i den Villiesact, der constituerer det Onde, maatte
Selvet drage sig tilbage i en forhærdet Villen af sig selv i sin tomme
Uendelighed. Selvet, der i den onde Villie vil sig selvisk, i Isolation fra
Principet, kan ikke ville sig selv i det harmoniske Forhold af sit Væsens
Momenter, og saaledes ved det Concrete give den abstracte Uendelighed Indhold,
netop fordi Isolationen, Løsrivelsen fra Livsprincipet, maa gjøre den selviske
Villens Uendelighed død og uskikket til indre, levende Udfoldelse. Men i hiin
tomme Uendelighed, i hvilken Selvets Virkelighedsbestemthed forsvinder, kan det
ikke holde sig. Abstractionen bæres kun af det reale Concrete, Abstractionens
Kraft hviler paa det reale Naturgrundlags Kraft, der er virksomt i Villien;
Selvet maa derfor, idet dets Natur som activt sætter det i Bevægelse, fra den
tomme Uendelighed føres tilbage til en \emph{ufri Villen af det concrete
Endelige}, til en vilkaarlig Sætten af en \emph{isoleret} Væsensbestemmelse, og
derfra maa det saa atter, ifølge \emph{Incommensurabiliteten} mellem det
\emph{Villede og Uendeligheds}bestemmelsen, kastes tilbage til den tomme
Uendelighed, der, naar ikke den samme Oscilleren skal gjentage sig rastløs i det
Ubestemte, definitivt kun kan søge at hævde sig \emph{negativt}, gjennem
Forstyrrelsen af det Gode, gjennem en Kamp imod det, i hvilken den egoistiske
Uendelighedsvillie dog bestandig bæres af det reale Naturgrundlag, hvis
Drifttilskyndelser i negativ Form mødes i Tilintetgjørelsesdriften. Men denne
Stræben, i hvilken Selvraadigheden er potentseret, fører ikke blot ind i den
ulige Kamp \emph{udadtil}, men i en Selvets ulige Kamp i \emph{sit eget Indre.}

I Kampen \emph{udadtil} kæmper den onde Villie magtesløs mod den overvældende
Totalsum af de andre Kræfter, og staaer bestandig i dybere Betydning \emph{ene}.
Thi selv under de Ondes Samvirken, under deres gjensidige Benytten af hverandres
Kræfter, kan den Enkelte, der vil sig i Isolationen, og hvis Villies egoistiske
Sondring umuliggjør et Kjærligheds- og Samdrægtigheds-Forhold, under Kampen mod
det Gode aldrig bæres af en Heelhedsvillie, hvorfor selv de udadtil samvirkende
onde Villiers momentane Seir
% SPLIT SPACING, p. 154: „Uendelighedsbestemmelsen“ is spaced through
% „Uendeligheds“ and close-set from „bestemmelsen“ on, within the one word.
% Transcribed as printed.
% EMPHASIS, p. 154: „Sætten af en“ is close-set and only „isoleret“ is spaced
% (checked at 300 dpi, 1.8x, against ABBYY).
% p. 154 carries no footnote and no Greek.
% --- p. 155 ---
maa bære Opløsningens Spire i sig. --- Og Kampen mod det Gode bliver saa tillige
en Selvets Kamp i \emph{sit eget Indre}, gjennem hvilken en \emph{indre
Modsigelse} bliver aabenbar: en Modsigelse i Forholdet til \emph{Selvets Væsen},
og til det i dette Væsen virksomme \emph{Absolute}. Den første Modsigelse træder
frem deri, at Selvet i denne Kamp maa være \emph{deelt}, og \emph{ubevidst maa
ville det Modsatte af, hvad der bevidst villes}, fordi det skal kæmpe mod det,
\emph{der er dets egen sande Natur og Væsen og dets Villies Grund.} Kampen maa
derfor blive frugtesløs, idet den Uendelighed, der bruges til Kampen, faaer
Uendeligheden som \emph{indholdsfyldig} Uendelighed overfor sig, og, idet den har
sit Væsen paa Modsætningens Side, forvandles til et \emph{Endeligt} og
\emph{Væsenløst}, og synker i sin \emph{Afmagt}. I Selvets Kamp mod det, der er
dets Væsen, bliver Potentsationen af den onde Villies Kraft en Potentsation af
\emph{Modstandens} Kraft: den onde Villies Kraftanspændelse bliver derfor en
Kraftslappelse, den høieste Anspændelse deu høieste Afmagt; og Bevidstheden om
denne Afmagt bryder den onde Villie, saa at Væsensvillien kan komme til
Gjennembrud. --- Modsigelsen i Forholdet til det \emph{Absolute} viser sig deri,
at Selvet, da det i sin Villen af det Onde kun er villende ved den Uendelighedens
Kraft, det har fra det Absolute, i sin egoistiske Villen af sig ubevidst maa
ville det Uendelige, ved hvilket hiin Villens Mulighed er, idet dets onde Villen
brydes ved det gjennemvirkende Absolute. Den selviske Villie vil med den Kraft,
der ikke er dens egen, derfor maa den, ubevidst, ville, hvad der ikke er den
selv. Selvet vil, i det Onde, sig selv i sin Realitetssondring, i Uendelighedens
Form, som et Absolut. Denne Uendelighedens Villie, der kun er ved det Absolute,
var det, der, idet Selvet søgte sin Villies Indhold i et Endeligt, fortærede
dette Indhold gjennem Modsigelsen med Uendeligheden, og det var den, der, idet
Selvet saa forholdt sig til selv i Abstractionens Tomhed, fremtvang det negative
Forhold til det sande Uendelige, til det paa een Gang Menneskets Væsen og det
Absolute udtrykkende Gode, i Kampen mod hvilket
% PRINTER'S DEFECT, p. 155, l. 18: the page prints „deu høieste Afmagt“ for
% „den høieste Afmagt“ --- a turned or wrong sort. Both OCR witnesses read „deu“,
% and at 300 dpi / 2.6x the third letter is joined at the FOOT, unlike the „n“ of
% „den høieste“ five words earlier on the same line, whose strokes are joined at
% the head. Transcribed as printed; not on the Trykfeil leaf.
% PRINTER'S DEFECT, p. 155, l. 30: „forholdt sig til selv“ --- „sig“ appears to
% be dropped before „selv“ (expected „til sig selv“). Both witnesses read „til
% selv“ and the image at 2x shows no further word. Transcribed as printed.
% EMPHASIS, p. 155: „Modsigelsen i Forholdet til det“ is close-set and only
% „Absolute“ is spaced (checked at 300 dpi, 1.8x), unlike the parallel clause
% earlier on the page.
% p. 155 carries no footnote and no Greek.
% --- p. 156 ---
\setcounter{footnote}{0}%
Selvet udhuler sig selv, og i den uendelige Anspændelse af Kraften ender i den
absolute Afmagt. \emph{Den Uendelighedens Kraft, ved hvilken det Onde er,} er da
i det \emph{Onde som den det Onde ophævende Kraft;} den er i det Onde som
\emph{det Ondes indre Negation}, som det, der forvandler den onde Villie, der er
en paa en illusorisk Erkjenden, ligesom paa et Feilsyn paa Væsenet, grundet
illusorisk Villen, til en \emph{ubevidst, indirecte Villen af det Gode}, af
Væsenets Sandhed. Gjennem den onde Villies selvopløsende Selvvillen villes det
Gode: Væsenets Sandhed tvinger sig frem gjennem det Væsenløses Selvophævelse. Det
Absolute afspeiler sig, men \emph{som forvandlet} ved det Afspeilende, i
Modsætningens Form, ligesom omvendt, i det Onde. Det er tilstede i
Samvittigheden, ikke som det Forsonende, men som det Fortærende, ikke som det
Kjærlige, men som det Vrede\footnote{I dette Dialektiske i den onde Villen ligger
det \emph{Ondes Straf}, der altsaa deels bliver at see som en \emph{indre}, ɔ:
den, der er given i Væsenets indre Splid og i Forholdet til det Absolute i den
onde Samvittighed, deels som den \emph{ydre}, der ligger i Reactionen af det Gode
og af den forskjelligartede Egoisme mod den egoistiske Villen. Denne ydre Straf
er væsentlig Isolationens Magtesløshed, og det at være hadet. Idet denne ydre
Straf bliver Individet bevidst efter sin Grund, smelter Følelsen af den sammen
med Følelsen af Væsenets indre Splid. At denne indre Splid ogsaa omfatter det i
Fornedrelsen givne Misforhold i Menneskenaturen, er tydeligt af det
Foregaaende.}. Og idet det Absolute virker som det Fortærende, som det
Negerende, som det, der gjennem Modsigelsen og Afmagten fører det efter sin Natur
Tomme og Væsenløse til Opløsning, og gjennem Bevidstheden om Opløsningen og
Afmagten baner Veien for det Godes Gjennembrud i \emph{Angeren}, saa bliver, i
dette Gjennembrud, den i Selvets uendelige Villen af sig selv \emph{actualiserede
Naturkraft:} den reale Naturgrunds \emph{i Villien optagne Magt, tagen i det
Godes Tjeneste}, og giver det Gode \emph{forhøiet Energie}, saaledes at det
\emph{Onde} i Omvendelsen faaer \emph{positiv} Betydning for det \emph{Gode.}
% FOOTNOTE, p. 156: printed mark *); the note carries an id-est mark, „en
% indre, ɔ: den, der er given“. The hint list does not name p. 156.
% EMPHASIS, p. 156: „det Onde i Omvendelsen“ is NOT a spaced run --- only „Onde“
% is spaced, „det“ and „Omvendelsen“ are close-set (checked at 300 dpi, 2.1x).
% Likewise „i det Onde som“ before „det Ondes indre Negation“ is close-set.
% p. 156 carries no Greek.
% --- p. 157 ---
Begrebet om \emph{Realgrundens Actualiseren} gjennem den onde Villie bliver
nærmere at bestemme saaledes. Den onde Villie har som en usand
Uendelighedsvillen til sine Betingelser for det Første en Gjennemvirken af
Principet, det sande Uendelige, og for det Andet en \emph{Energie i det
Naturgrundlag} (den reale relative Totalitet), \emph{der bærer den individuelle
Villie.} Men denne anden Betingelse: Naturgrundlagets Energie, kan, ifølge
Grundforholdet mellem det Reale og Principet, ikke tænkes uden hiin
Gjennemvirken, der altsaa skjult maa være tilstede deri. Naturgrundlagets Energie
gjør sig gjældende i \emph{Driften}, der, potentseret ved Uendeligheden, bliver
\emph{Lidenskab}, en Villen af sin Gjenstand \emph{som absolut}. Driften vil her
uendeligt sig selv, den er et ἄπειρον. Men under denne Drift-Stræben ere Villiens
Form og dens Indhold incommensurable. Deraf følger saa den tidligere paaviste
Vexlen i Villiesindholdet. Men det \emph{Formelle: Lidenskabens Uendelighed},
bliver, og idet nu Energien vendes imod det Gode, forat bekæmpe det, bliver
\emph{Driftvirkningernes Resultant} den for det Godes Tilintetgjørelse
\emph{kæmpende} onde Villies \emph{formelle, negative Uende}lighed. Denne
Villiens Uendelighed er det, der bevares, idet Villiesretningen vendes om: ---
\emph{under Kraftens Forvandling er ogsaa her, ligesom paa det Physiskes Gebet,
Kraftsummen blivende} --- og idet \emph{Angeren}, der er \emph{Væsenets} positive
Hævdelse gjennem den \emph{væsenløse} Villens Negation, bliver seirende, tages
Energien i det Godes Tjeneste, og den er saameget intensivere, som de reale
Momenter ligesom enkeltviis ere gjennemtrængte af Uendeligheden. Ogsaa i denne
Energiens Metamorphose bliver det Absolutes Gjennemvirken at accentuere; men den
er her positiv, i Enhed med Væsensvillien, medens den i den onde Villie var
negativ, som opløsende den væsenløse Villen.

Ifølge det tidligere Udviklede angaaende \emph{det Ondes primitive Form} maa den
paaviste Dialektik i den onde Villie, der directe viser hen til denne primitive
Form, faae Gyldighed for \emph{enhver} phænomenal Form af den onde
% GREEK, p. 157: „ἄπειρον“ --- antiqua italic, smooth breathing with acute on the
% alpha, in „den er et ἄπειρον“. Neither OCR witness reads it (tesseract gives
% „ä7æoov“, the ABBYY layer nothing); read off the image at 2.8x. p. 157 is not
% in any Greek list.
% SPLIT SPACING, p. 157: „Uendelighed“ in „formelle, negative Uendelighed“ is
% spaced through „Uende“ (the end of the line) and close-set in „lighed.“
% Transcribed as printed.
% EMPHASIS, p. 157: „kæmpende“ is spaced while the following „onde Villies“ is
% close-set (checked at 300 dpi, 1.75x); ABBYY misses the spacing of „kæmpende“.
% p. 157 carries no footnote.
% --- p. 158 ---
Villen, fordi hiins Eiendommelighed er gjennemgaaende Bestemmelse i dem alle.
Overalt, selv hvor det for Grundformen Eiendommelige allermeest skjuler sig, som
der, hvor det Onde fremkaldes ved en ydre Paavirkning, hvis Begyndelse gaaer
tilbage til den spæde Alder, eller hvor Tilbøieligheden er nedarvet, bliver der,
netop fordi den onde Villie kun kan være som en \emph{Aandens Yttring}, at sætte
et Punkt, hvor Villien, støttende sig paa en \emph{Bevidsthed}, hvis Klarhedsgrad
gjør \emph{Tilregnelighed} mulig, har \emph{Trodsens} Form, selvraadig vil det
Onde, og ved Selvraadigheden constituerer sig som \emph{ond} Villie.

Men idet det ovenfor Udviklede skal anvendes paa det Ondes enkelte Former,
opstaaer der en Vanskelighed ved saadanne Former, i hvilke Individets Naturkraft
brydes og en energisk Handlens organiske Betingelser destrueres. Ved disse Former
synes der ligesaalidt at kunne være en Tale om en Actualiseren af
Naturgrundlaget, som om en Bevaren af Kraften, og overhovedet synes det, hvor det
Onde har Formen af \emph{Fornedrelse} (Svaghed, Usselhed, Dyriskhed, Fortabthed,
„det, fortvivlet ikke at ville være sig selv“), og hvor Synden fortsættes i
\emph{Sløvhed}, at der gjennem et Liv i saadan \emph{Synd} ikke kan være
realiseret Noget, der kan tjene det Gode, og at en Opreisning maa blive umulig.
Forholdet bliver her at opfatte saaledes. Ved \emph{enhver} uethisk, paa en usand
Driftvirken hvilende Villen er der sat en \emph{Disharmonie} i Væsenet, der til
en vis Grad maa indvirke forstyrrende paa dets Realitetsside: Organismen
(„Døden er Syndens Sold“), og denne \emph{overalt} stedfindende Forstyrrelse er
det, hvis Grad i hine Former forøges saaledes, at den skjulte Sygdom i Selvet
bliver aabenbar, og endog kan antage Charakteren af dyrisk Sløvhed. Her er altsaa
ikke en Modsætning, men en \emph{Gradforskjel}. Under den overalt i det Onde
stedfindende \emph{Disharmonie}, som vi betegnede som en \emph{Sygdom i Selvet},
har Villien dog sin \emph{abstracte Uendelighed} paa Basis af det actualiserede
Reale, fordi Væsenet som Væsen \emph{ikke} kan tilintetgjøres, og idet
Villiesenergien ikke staaer i \emph{directe} Forhold til den
% QUOTATIONS, p. 158: two pairs, „det, fortvivlet ikke at ville være sig selv“
% (the Kierkegaard formula) and („Døden er Syndens Sold“). Both balanced.
% RULE, p. 158: the page carries NO rule at its foot; the subsection runs on
% into p. 159, where a rule does stand.
% p. 158 carries no footnote and no Greek.
% --- p. 159 ---
physiske Kraft, men til den individuelle relative Totalitets Concentration (i
modsat Tilfælde maatte de physisk kraftigste Individer være de villiestærkeste),
saa kan under Disharmonien, selv hvor den antager større Dimensioner,
\emph{Villiesintensiteten} bevares, trods Forstyrrelsen, og dermed
\emph{Angerens} Udgangspunkt være givet. Og selv hvor en Desorganisation af
Villiens organiske Betingelser finder Sted, bliver den ethiske Restitutions
Mulighed tilbage, ligesom Frihedens Mulighed bliver i Selvet saalænge Selvet
existerer som \emph{menneskeligt} Selv. Idet Synden fortsættes i Sløvhed, men dog
som villet (jfr.\ S.\ 150 f.) er der, under den \emph{fortvivlede} Hensynken, i
selve Fortvivlelsen, en \emph{ubevidst Angestens Kamp}, i hvilken det Ondes Magt
undergraves, og af hvilken en frelsende Angers Flamme kan slaae frem. Selv
gjennem et saadant Liv kan da Noget tænkes at være realiseret for det Gode. I
\emph{Fortvivlelsens og Angestens Stigen}, der kun tildækkes under Sløvheden og
Fortabtheden, men dog er tilstede, arbeides der, ligesom i Væsenets Dyb, hen mod
et Gjennembrud, og idet Væsenet bryder igjennem i Angeren, bærer det i sig et
uendeligt Indhold, der ikke har kunnet tilintetgjøres; men som ogsaa kun kan
frigjøres gjennem en langt dybere Lidelse, gjennem en smertelig Udbrænden af det
Onde. Ogsaa her kan da Kraften bevares og Kampen bære sin Frugt i Omvendelsen.

\begin{center}\rule{2cm}{0.4pt}\end{center}
% RULE, p. 159: a short centred SINGLE hairline standing mid-page between
% „…Frugt i Omvendelsen.“ and the summary paragraph below, with no head adjacent
% to it --- a thought-break, as at pp. 19, 28, 42, 62 and 97. It shows a small
% break at its middle. Transcribed. The display head lower on this page carries
% no rule of its own.

Til nærmere Bestemmelse af de i det Foregaaende paaviste Grundforhold i det Onde
bliver endnu at udvikle; \emph{Forholdet mellem Nødvendigheds- og
Frihedsmomentet i det Ondes Virkeliggjørelse og i dets Ophævelse}, samt:
\emph{Forsoningens Væsen, dens Princip og dens Maal.}

\section*{1. Forholdet mellem Nødvendigheds- og Frihedsmomentet i det Onde.}

\emph{a) i Virkeliggjørelsen af det Onde.} --- Den \emph{indifferentistiske}
Opfattelse viser sig strax som uholdbar. Den
% HEAD, p. 159: „1. Forholdet mellem Nødvendigheds- og Frihedsmomentet i det
% Onde.“ --- bold (halbfette) Fraktur, centred over two lines, „Friheds-“ /
% „momentet“ broken at the line end and here joined. Wording verified at 300 dpi.
% The Indhold omits „i det Onde“; the PRINTED head stands in the body.
% HEAD, p. 159: the run-in „a) i Virkeliggjørelsen af det Onde.“ is letterspaced
% regular Fraktur, not bold, followed by an em-dash; hence \emph{} and not
% \textbf{}, unlike the run-in heads of pp. 97 and 126.
% p. 159 carries no footnote and no Greek.
% --- p. 160 ---
er selvmodsigende, idet den lader Handlingens, Valgets, Bestemthed fremgaae af en
i sig tom og ubestemt Villie; den udelukker Ansvarlighed og Tilregnelighed,
derved, at den tilintetgjør Personlighedens Continuitet, og den maa af samme
Grund sætte en ethisk Indvirkning paa Charakteren som umulig. Men fordi
Indifferentismen med dens grundløse \textit{liberum arbitrium} forkastes, sættes
Friheden ikke som absorberet af Nødvendigheden. I Afsnittet om Friheden blev det
viist, hvorledes det menneskelige Individ, ifølge sin Væsensbestemthed som
relativ Totalitet med indre Uendelighed, med Bevidsthed maa concentrere sig i sig
selv, tage sig selv i Besiddelse. Deri, at det \emph{maa} det, er
\emph{Nødvendigheden} i Forholdet sat, men idet det constituerer sig i sin
Væsensbestemthed \emph{ved sig selv}, sætter sig selv ved sig selv, er
Nødvendigheden i denne bevidste Selvconcentrationens Act omsat i \emph{Frihed}. I
Selvets Constitueren af sig som relativ Totalitet ligger saa \emph{et dobbelt
Forhold}: et Selvets Forhold til det Absolute og til sig selv, altsaa \emph{en
dobbelt Mulighed}: Muligheden til Indordningen under den absolute Totalitet og
til Isolationen. Idet det dobbelte Forhold: Forholdet til sig selv og til det
Absolute, er betinget ved \emph{en dobbelt Bevidsthedsact}, maa det sættes i
\emph{forskjellige} Tidsmomenter; men Tidsmomenterne kunne tænkes at grændse
saaledes op til hinanden, at der intet Spatierende bliver, altsaa ingen
Virkelighedssondring af de to Forhold; og de kunne tænkes som i ubestemt Grad
rykkede ud fra hinanden. I denne Tidsbestemmelse, der er betinget ved den
menneskelige Bevidstheds Eiendommelighed, og som faaer Betydning for det efter
sit Væsen Eviges Grundforhold idet disse Forhold blive Gjenstand for
Bevidstheden, ligger Tilknytningen for et \emph{Tilfældighedens} Moment i det
Ondes Realisation. Her faae nemlig de \emph{udenfra} kommende Impulser deres
Spillerum, og det Ondes Mulighed, der er tilstede i ethvert Individ, kan, gjennem
Forholdet til Virkeligheden med dens Impulser, faae en i ubestemmeligt Maal
forhøiet Intensitet, saa at den kan tænkes at staae i Forhold til det Godes
Mulighed som det uendeligt
% ANTIQUA LATIN, p. 160: „liberum arbitrium“ is antiqua and CLOSE-SET, like the
% antiqua of pp. 64, 71 and 134 and unlike the spaced antiqua of pp. 68 and 144;
% encoded \textit{} without \emph{}. Tesseract reads it correctly, the ABBYY
% layer garbles it to „iitm rum arbitcium“.
% PUNCTUATION, p. 160: „…for det efter sit Væsen Eviges Grundforhold idet disse
% Forhold blive Gjenstand for Bevidstheden“ --- no comma is printed after
% „Grundforhold“; both witnesses agree. Transcribed as printed.
% p. 160 carries no footnote and no Greek.
% --- p. 161 ---
Store til Enheden; men idet den gjøres til Virkelighed, saa bliver dog bestandig
Omsætningen fra Mulighed til Virkelighed en \emph{Frihedsact}, bestemtere: en i
Menneskevæsenet begrundet Act af Friheden i dens \emph{Uendelighed}, men i dens
\emph{abstracte} Uendelighed.

Naar der, i Tidsrum af uberegnelig Udstrækning, har været levet et Liv af
Mennesker uden at det i Sandhed Menneskelige i Bevidstheden er kommet til
Udvikling, saa kan det Spørgsmaal opkastes, hvorvidt der under saadanne Forhold
kan tales om det Ondes Virkeliggjørelse. Svaret paa dette Spørgsmaal, der gaaer
ind under det almenere Spørgsmaal om det Ondes tilfældige eller nødvendige
Forhold til \emph{Menneskevæsenet}, vil afhænge af, om den factiske Mangel paa
Aandsudvikling sees som nødvendig og som en virkelig Ikke-Tilstedeværen af det
Principielle i den menneskelige Bevidstheds Indhold. Afgjørelsen heraf ligger i
Forholdet mellem \emph{Væsens}bestemmelsen og de i Intensitet uberegnelige
Paavirkninger og Betingelser. I dette Forhold maa Væsensbestemmelsen fastholdes
som det Overgribende. Er det Liv, der leves, et Liv af \emph{Mennesker}, saa er
det Menneskeliges \emph{Væsensbestemmelse} og den deri liggende \emph{Fordring}
tilstede, og Ikke-Opfyldelsen af denne Fordring constituerer, om end de
modstridende Impulser have været nok saa stærke, Bevidsthedsgraden nok saa svag,
det Ondes Virkelighed.

Sees Friheds- og Nødvendighedsmomentet i Virkeliggjørelsen af det Onde fra
Synspunktet af Forholdet mellem den \emph{individuelle} Villie og det virkende
\emph{Absolute}, saa indeholde de tidligere Udviklinger det til Opfattelsen
Nødvendige. Vi vise her tilbage til hvad der blev udviklet om Forholdet mellem
Idealitetsbestemmelsen og Realitetsbestemmelsen, Evigheds- og Tidsbestemmelsen i
det Absolute, om dette Forholds Betydning for det Endeliges relative
Selvstændighed, om det Absolutes negative Væren i den onde Villie og om denne
Villies Forhold til Væsenets Uendelighed. Anvendes Bestemmelserne: Frihed og
Nødvendighed paa det \emph{Absolute} i dets Forhold til det Onde, saa bliver,
idet i det
% SPLIT SPACING, p. 161: „Væsensbestemmelsen“ is spaced through „Væsens“ and
% close-set from „bestemmelsen“ on, within the one word --- the third such split
% in this batch (cf. pp. 148, 154, 157). Transcribed as printed.
% p. 161 carries no footnote and no Greek. The „11“ at the foot is the gathering
% signature; not transcribed.
% --- p. 162 ---
Absolute Frihed og Nødvendighed maa tænkes i evig Enhed, det Onde paa een Gang
at opfatte som begrundet i en Frihedsact af det Absolute (den absolute
Subjectivitet), og som existerende ifølge en Nødvendighed i dets Væsen. Dette
Friheds- og Nødvendighedsmoments Forhold danner Parallelen til Forholdet mellem
Idealitets- og Realitetsbestemmelsen i Principet.

\emph{b) i Ophævelsen af det Onde.} --- Hvad der i det Foregaaende er udviklet
angaaende Dialektiken i det Onde, dets indre Selvopløsning, gjengiver nærmest et
\emph{Begrebsforhold}, et \emph{Væsensforhold}, en Væsensdialektik, som derfor
nærmest er seet i Abstraction fra Tidsforholdet og under Nødvendighedens Form.
Her derimod bliver det Ondes Ophævelse som et \emph{Virkelighedsforhold} at see i
\emph{Tidsforholdet} og i Forhold til \emph{Friheds}bestemmelsen. Tidsforholdets
væsentlige Betydning er berørt ved Bestemmelsen af det Ondes første
Virkeliggjørelse i Villien. Fra dette Begyndelsespunkt viser Tidens Betydning sig
i en dobbelt Henseende: i Henseende til det i Udviklingen Mangelfulde, til den
voxende Sum af det med Hensyn til Fordringens Opfyldelse Forsømte (den voxende
Skyld); og i Henseende til Udsættelsen eller Udeblivelsen af Villiens Bortvenden
fra det Onde. Det Første vil blive belyst i det Efterfølgende hvor Forsoningen
skal omhandles; det Sidste bliver strax her at belyse.

Den onde Villies indre Modsigelse maa forhindre en Stillestaaen, en Hvile i det
Onde; den maa fremtvinge en rastløs Uro, en hvileløs Fremadstræben. Under denne,
der, hvad enten den faaer Formen af en Omtumlen i Lysten eller af en directe Kamp
mod det Gode, bestandig indeslutter Trodsen, Modstanden mod den ethiske Fordring,
maa den onde Villie mere og mere udhule sig selv; den maa, i sin Negeren af
Væsensfordringen, medens den vil gjøre sig selv absolut, ubevidst stræbe hen mod
en Selvtilintetgjørelse. Men denne Selvtilintetgjørelse kan kun tilnærmes, fordi
den onde Villie, der kæmper mod Væsenet, og vil sætte sig som Væsenets Udtryk,
skjult bevarer Sammenhængen med det,
% HEAD, p. 162: the run-in „b) i Ophævelsen af det Onde.“ is letterspaced regular
% Fraktur, not bold, followed by an em-dash --- set exactly like its companion
% „a) i Virkeliggjørelsen af det Onde.“ on p. 159, hence \emph{} and not
% \textbf{}. Verified at 400 and 800 dpi. The Indhold agrees word for word
% („b) i Ophævelsen af det Onde“, pp. 162--164); no discrepancy to log.
% ANTIQUA ENUMERATOR, p. 162: the „b)“ itself is upright ANTIQUA, not Fraktur ---
% the bowl and the parenthesis are plainly roman beside the Fraktur „Ophævelsen“
% (800 dpi). The same holds of the „a)“ on p. 165 and the „b)“ on p. 167, and of
% the „a)“ on p. 159 already set by the previous batch. Left unmarked here, since
% the house \textit{} convention is for Latin *words* and marking it would break
% with p. 159; recorded for a later pass, like the antiqua „II“ logged at p. 147.
% SPLIT SPACING, p. 162: „Frihedsbestemmelsen“ is spaced through „Friheds“ and
% close-set from „bestemmelsen“ on, within the one word and on one line
% (400 dpi). Transcribed as printed. „Begrebsforhold“, broken „Begrebs-“ /
% „forhold“ at the line end, IS spaced through both halves and is here joined.
% p. 162 carries no footnote and no Greek.
% --- p. 163 ---
der, som Væsen, ikke kan tilintetgjøres, saaledes at den, idet den ubevidst
arbeider hen mod en Selvtilintetgjørelse, --- der, naar Væsenet kunde absorberes
af den onde Villie, vilde være Væsenets Tilintetgjørelse, --- virker med en
Kraft, der er laant fra den kun skjulte Væsensvillie, det Godes Villie, og altid
maa blive som en positiv Rest under den Fordobbling i den onde Villie, i hvilken
Villiens Kraft fortæres, men ved en Kraft, der tilhører Villien. Der bliver
derfor bestandig en Dobbelthed i Forholdet; der bliver bestandig ligesom en
Opblussen af den Flamme, der vil slukkes, og derfor bliver der et bestandig kun
Approximativt i den Stræben, der fra en uendelig Afstand stræber hen imod en
Grændse, ved hvilken Kraften naaer 0, hvor Villien vilde synke hen i
Tilintetgjørelse, i en evig Død. Men i den Villiens \emph{Delthed}, der bliver
Følgen af, at den onde Villie som bevidst actualiseret Villie kæmper mod
Væsensvillien, er der en Splid sat i Væsenet, der er som en \emph{Væsenets
Sygdom}, en Suspension af Integriteten, og Spliden voxer idet den onde Villie
kæmpende udhuler sig selv. I den onde Villies Overvægt bliver Væsenet ligesom
usynligt: dets Kræfter virke negativt i det Onde. Som Væsen kan det ikke tabes,
ikke tilintetgjøres; om end denne dets Tilintetgjørelse, der vilde være Eet med
Villiens Tilintetgjørelse, er det Maal, mod hvilket den onde Villie ubevidst
arbeider hen. Men i \emph{Spliden i Væsenet}, i \emph{Sygdommen}, der truer med
en bestandig stærkere \emph{Opløsning}, er der en \emph{Angest}, og i denne
Syndens Angest er Omslagets Mulighed. I hiin Dobbeltbevægelse i Villien viser
sig, paa de enkelte Punkter, Negationen, Selvets Fortabelse som Væsenets
Forstyrrelse truende for Bevidstheden, og idet den viser sig, viser Væsenet sig
gjennem den i Angesten for sig selv som det, der skal hævdes og restitueres, og
heri ligger Angerens Mulighed. Der kan nu tænkes Mellemstadier med en Oscilleren,
i hvilken Angerens og Tilbagefaldets Momenter vexle; men der bliver et Punkt,
hvor Bevidstheden, under Selvets fortvivlede Arbeiden hen mod en Selvopløsning,
opfyldes af Angesten for den aandelige Død, for Væsenets
% RULE, p. 163: the page carries NO rule at its foot. It ends in mid-sentence
% („for Væsenets“ / „Opløsning“), and the only mark below the last line is the
% gathering signature „11*“, which is not transcribed. Checked at 200 dpi over
% the whole foot of the page.
% ANTIQUA FIGURE, p. 163: „hvor Kraften naaer 0“ --- the zero is an antiqua
% figure (there is no Fraktur nought), read at 400 dpi; ABBYY took it for a
% capital O. Set as the digit 0.
% EMPHASIS, p. 163: „i en evig Død“ is NOT letterspaced, though the parallel
% phrases around it are; checked at 400 dpi. „Spliden i Væsenet, i Sygdommen“
% runs as one spaced phrase across its comma; the second „Angest“ („denne
% Syndens Angest“) is tight where the first is spaced.
% p. 163 carries no footnote and no Greek.
% --- p. 164 ---
Opløsning ved Villien, altsaa ved sig selv, og paa dette Punkt bliver Væsenets
Angestkamp saa stærk ligeoverfor den udhulede onde Villie, der er den delte
Villie, at Forholdet i Muligheden for det Onde vender sig om og bliver at
betegne ved en Formel, der er modsat den, der ovenfor anvendtes: idet det Ondes
Mulighed nu maa sættes som forholdende sig til det Godes som Enheden til det
uendeligt Store. Her er da Omslagets, Gjennembruddets Øieblik, og det er der
ifølge Væsenets \emph{Nødvendighed}, men ved \emph{Frihedens}, den Væsenet og sig
selv hævdende Friheds Beslutning.

To Spørgsmaal maae imidlertid her strax paatænge sig. Det første er det: om ikke
det Tilfælde kan tænkes, at det Onde omfatter \emph{hele} Individets Existents,
og at altsaa Gjennembruddet udebliver for dette Individ? Gjennembruddets Forhold
til Friheden, og Tidsbestemthedens Tilfældighed synes at fordre et bekræftende
Svar. Alligevel vil Svaret vise sig at maatte blive et andet, men dette vil først
kunne begrundes gjennem Udviklingen af Forsoningen og af Forholdet mellem dennes
Factorer. Det \emph{andet} Spørgsmaal er det: om Gjennembruddet, Selvets
Tilbagevenden til en Villen af sin sande Væsensbestemmelse, er Selvets
\emph{egen} Gjerning, eller om det kun er ved et Forhold, i hvilket Selvet staaer
til \emph{et Supplerende}, dette være nu de andre menneskelige Individer, eller
Selvets absolute Princip. Ogsaa dette Spørgsmaal fører os hen til Problemet om
Forsoningen, og ved Udviklingen af dette faaer det ovenfor Berørte: Muligheden af
en Indhenten af det Forsømte, en afgjørende Betydning.

\section*{2) Forsoningen, dens Kilde og Maal.}

Forsoningens Begreb opfatte vi, i Overensstemmelse med det tidligere Udviklede, i
den dobbelte Betydning af Individets Forsoning i \emph{sig selv} ɔ: af den i
Individet til\-%
% HEAD, p. 164: „2) Forsoningen, dens Kilde og Maal.“ --- bold (halbfette)
% Fraktur, centred, on one line. Verified at 600 dpi. TWO discrepancies with the
% book's own Indhold, both real and both left standing: (i) the printed
% enumerator is „2)“ with a closing PAREN, while the Indhold reads „2.“ with a
% full stop --- the reverse of the pattern at pp. 97 and 126, where the body has
% no enumerator and the Indhold supplies „1)“ / „2)“; (ii) the Indhold places the
% entry at „167--201“, but the head is printed here on p. 164, and the Indhold's
% own sub-entry a) is given as 165--167. Neither is in the Trykfeil leaf. The
% comment already standing at the Indhold (p. IV) records the same facts.
% RULE, p. 164: a short centred SINGLE hairline (about one sixth of the measure)
% stands between „…en afgjørende Betydning.“ and the head, i.e. ABOVE the display
% head, as at pp. 97 and 126. Not transcribed, per the standing convention for
% head rules.
% EMPHASIS, p. 164: „Det første er det“ is TIGHT while „Det andet Spørgsmaal“ has
% „andet“ spaced and „Spørgsmaal“ tight --- checked at 400 dpi against each
% other. Another split: „Selvets egen Gjerning“ spaces „egen“ only.
% PRINTER'S DEFECT, p. 164, line 12: „her strax paatænge sig“ --- „paatænge“ for
% „paatrænge“, a dropped „r“. At 600 dpi the „t“ is followed directly by the „æ“
% with no gap and no trace of an „r“. Both OCR witnesses read „paatænge“ /
% „paatcenge“. Transcribed as printed; not in the Trykfeil leaf.
% p. 164 carries no footnote and no Greek.
% --- p. 165 ---
veiebragte indre Harmonie og Samstemmen med Væsenets Fordring, og af Individets
Forsoning \emph{med sit absolute Princip}. Under begge Opfattelser, hvis indre
Sammenhæng tidligere er bleven paaviist, forudsættes her den ethiske Forstyrrelse
som gaaende forud, saa at Forsoningen ikke bliver Udtrykket for et i Bevidstheden
afspeilet uforstyrret harmonisk Forhold, men for den bevidste \emph{Restitution}
af det forstyrrede normale Forhold.

Vi betragte da først:

\begin{center}
\emph{a) det i Forsoningen Virksomme.}
\end{center}

Umiddelbart bliver dette selvfølgelig \emph{selve Individets}, gjennem
\emph{Angeren} frigjorte, i den \emph{ethiske Beslutning} til det Gode til
Gjennembrud komne \emph{væsentlige Villie}, Udtrykket for Individets
Væsensenergie, for Personlighedens intensiveste Activitet. Men Individets ethiske
Beslutning er en Beslutning af den Villie, i hvilken Personligheden vil sig som
indordnet i Totaliteten og som underordnet under dennes Princip. I denne Villie
er derfor fra Begyndelsen af Forholdet til den levende Heelhed, til dens
organiske Led og dens Princip, sat, og dette Forhold, der her sættes, viser sig,
idet det sættes, idet Individet bliver sig bevidst som omsluttet af Heelhedens
Liv, at være \emph{Forudsætning} for hiin Villies Gjennembrud og \emph{væsentlig
Betingelse} for den Villiens Udvikling, ved hvilken Forsoningen bliver
\emph{fuldstændig} som omfattende hele Livets Virkelighed.

Forholdet til Totalitetsleddene, til de \emph{andre Individer}, der i Forening
danne en aandelig Organisation, viser sin Betydning gjennem de supplerende
Individers \emph{ethiske} Indvirkning; gjennem den bevidste og ubevidste
\emph{ethiske Livsmeddelelse}. Indvirkningen træder frem i det endnu ikke ethiske
Individs Opdragelse til sædelig Villen ved de Personligheder, i hvilke
Sædelighedens Aand er virkeliggjort; den træder frem i den fortsatte sædelige
Paavirkning, og, i negativ Form, i det uethiske Individs Bestemmen ved
Afskrækkelse og Straf. Den har sit dybeste, primitiveste og universelleste Udtryk
i \emph{Kjærlighedsforholdet}, i hvilket
% HEAD, p. 165: „a) det i Forsoningen Virksomme.“ --- CENTRED on its own line,
% letterspaced regular Fraktur, NOT bold. Checked at 400 and 800 dpi against the
% bold head of p. 164 four lines above it: the strokes are of the ordinary text
% weight. Hence \begin{center}\emph{…}\end{center} and not \section*{}. The
% Indhold agrees („a) det i Forsoningen Virksomme“, pp. 165--167); the only
% oddity there is the page span, logged at the head of p. 164.
% EMPHASIS, p. 165: the long opening run breaks into four spaced groups with
% „til det Gode til Gjennembrud komne“ tight between the third and the fourth;
% verified at 200 dpi across the whole paragraph. „andre Indi-“ / „vider“ is
% spaced through both halves of the broken word and is here joined.
% p. 165 carries no footnote and no Greek.
% --- p. 166 ---
Væsenets ubevidste Energie gaaer i Enhed med Bevidsthedslivets Energie. Og
ligesom det i Opdragelsen og Afskrækkelsen var Slægten, det Menneskeliges almene,
fornuftbestemte Væsen, der var virksomt gjennem de ethiske Individers Virken,
saaledes træder overhovedet i Kjærlighedsforholdet dette Forhold til
\emph{Slægten} --- Ordet taget i dets \emph{fulde}, aandelig-legemlige Betydning
--- frem som afgjørende. I Kjærlighedsforholdet, for hvilket Kjønskjærligheden
kun er et specielt og i og for sig endnu ikke ethisk Udtryk, er for Individet det
andet Individ Slægtens Repræsentant, den aandelige Totalitets Repræsentant, og
bliver i det ethiske Forhold for Individet ligesom dets personificerede,
objectiverede Samvittighed. Individets Forhold til det saaledes opfattede andet
Individ bliver et \emph{ethisk} Forhold \emph{til sig selv}: i det andet Individs
Bevidsthed har det, gjennem hiin ubevidste Omsætning af Individet til
Repræsentant for Slægten, en bestemmende \emph{Norm} for sin Handlen. Selv den
Illusion, der kan hæfte ved Kjærlighedsforholdet, --- ikke blot ved
Elskovsforholdet ---: den Illusion, ved hvilken Kjærlighedens Gjenstand
idealiseres, har sin ethiske Betydning, naar den ikke bliver en Yttring af den
indirecte Selvkjærlighed; den giver Individet, der i denne Illusion forholder sig
til et Ideal, en \emph{ideal Maalestok} for sin Handlen, en høiere Norm for den.
Og denne Anskuelse af Idealet, af det Almeengyldige, smelter saa sammen med
Anskuelsen af \emph{Personlighedernes Eiendommelighed}, af det Ethiskes
eiendommelige Realisationsformer i de forskjellige Individer, og ved denne
idealiserede ethiske Eiendommelighed, der paa een Gang indeholder et Correctiv
for det i Individet Mangelfulde og en Forvisning om den ethiske Kraft i det
Menneskelige, vise de af Kjærlighedsforholdet omfattede Individer sig i dobbelt
Betydning som supplerende. Saaledes bliver Kjærligheden, hvis Spire vækkes til
Liv i Individet ved den allerede virksomme Kjærlighed, den sande ethiske
Livskilde og derfor en Kilde til Forsoningen.

Gjennem Forholdet til de andre Individer vises der overalt hen til det
Sammenfattende, \emph{Totaliteten}, men
% RULE, p. 166: the page carries NO rule, at its foot or anywhere on it. The type
% block runs the full measure and ends in mid-sentence („Totaliteten, men“ /
% „gjennem dette“). Checked at 200 dpi over the whole page.
% FOOTNOTE, p. 166: NONE. The page is in the OCR-derived hint list of
% note-bearing pages; that hint is wrong here. There is no asterisk in the text,
% no rule above the last line, and no smaller-bodied matter at the foot ---
% checked at 200 dpi and again at 400 dpi over the bottom quarter.
% p. 166 carries no Greek.
% --- p. 167 ---
gjennem dette vises der atter tilbage til \emph{Totalitetens Princip}, og idet
Totalitetens Energie grunder i dette, ligesom de enkelte Individers ethiske
Virken i Totaliteten, bliver Forsoningens Kilde i \emph{sidste Instants} at søge
i \emph{det Absolute}. Dette er allerede i det Væsentlige udviklet i det
Foregaaende. Principets Virken er paaviist i ethvert, positivt og negativt,
individuelt Villiesforhold til det Ethiske. Den er paaviist i den den gode
Villies Selvhengivelse til det Guddommelige, der kun er ved dettes active Iboen i
Villien; den er paaviist i Principets negerende Virken i den onde Villie, den
viser sig ligeledes i det Gjennembrud, i hvilket den gode Villie bryder frem
gjennem det Ondes Opløsning; thi Væsenets Hævdelse i dette Gjennembrud kan kun
være ved en Virken af det, ved hvilket Væsenet er. Den bliver endelig ogsaa at
sætte som det definitivt Virksomme i Restitutionen af det ved den onde Villen
Forstyrrede og Forsømte, i Opnaaelsen af Maalet trods det, at der i et Tidsrum,
der ikke kan gjentages, har fundet en Afvigelse fra Maalet Sted. ---

Men dette sidste Synspunkt fører os fra Betragtningen af det i Forsoningen
Virkende til Betragtningen af:

\begin{center}
\emph{b) det, der skal naaes i Forsoningen.}
\end{center}

Idet vi opfatte Forsoningens Begreb i den ovenfor angivne dobbelte Betydning,
finde vi et sammenfattende Udtryk for Forsoningens Maal i
\emph{Virkeliggjørelsen af Idealet}. Virkeliggjørelsen af Idealet bliver, naar vi
tage Udgangspunktet for Bestemmelsen i Individets Forhold til det Absolute, lig
med Virkeliggjørelsen af \emph{den absolute Villies Fordring}; men denne Fordring
bliver atter, ifølge det tidligere Udviklede, identisk med \emph{Væsenets
Fordring}. Tage vi saa Udgangspunktet fra selve Menneskevæsenet, saa faaer
Idealets Virkeliggjørelse en \emph{dobbelt} Betydning. Den bliver paa den ene
Side at opfatte som Realisationen af det i Individet \emph{væsentlige}
Individuelle, af Individets \emph{eiendommelige} Anlæg, der skulle udvikles til
et ethisk \emph{Personlighedsindhold}; paa den anden Side bliver den
% HEAD, p. 167: „b) det, der skal naaes i Forsoningen.“ --- CENTRED, letterspaced
% regular Fraktur, not bold, exactly like „a)“ on p. 165. It closes with a single
% full stop; tesseract's „Forsoningen.,“ is a misreading, checked at 400 dpi. The
% Indhold agrees word for word („b) det, der skal naaes i Forsoningen“,
% pp. 167--168).
% PRINTER'S DEFECT, p. 167, line 8: „paaviist i den den gode Villies
% Selvhengivelse“ --- the article „den“ is set TWICE, a dittography at the same
% point in the line. Confirmed at 600 dpi; both OCR witnesses read it. Transcribed
% as printed; not in the Trykfeil leaf.
% p. 167 carries no footnote and no Greek.
% --- p. 168 ---
at opfatte som Realisationen af det i Individet \emph{Almene}, af det
\emph{almene ideale Selv i Mennesket}, af \emph{Menneskets Idee}. Men
Realisationen af hint Individuelle og dette Almene maa atter vise sig i
\emph{indre Enhed}, og denne Enhed træder frem i \emph{positiv} Form, idet
Individet i Udfoldelsen af sin Eiendommelighed udtrykker det Almene; den træder
frem i \emph{negativ} Form i Selvhengivelsen, i hvilken Individualitetens Energie
er virksom i Indordningen af Enkeltselvet under Totaliteten, i Resignationen i
Forholdet til denne. I denne Enhed er altsaa Bestemmelsen af Individets
individuelle τέλος og af det almeen-menneskelige, ved \emph{Menneskevæsenet}
bestemte τέλος saaledes forbundne, at det menneskelige Ideal,
\emph{Menneskehedens Idee}, er individualiseret for den Enkelte, og \emph{denne i
den Enkelte levendegjorte Idees Realisation} bliver den ethiske Fordring til den
Enkelte: det \emph{Ideal}, der er \emph{Forsoningens Maal}.

Virkeliggjørelsen af dette Ideal, der bliver den ethiske Villies Maal, faaer nu
en \emph{tredobbelt} Betingelse. Den \emph{første} Betingelse er den, \emph{at
Skylden kan oprettes}, d. v. s. at der, trods det at Idealets Virkeliggjørelse er
\emph{afbrudt} ved den uethiske Handling, og altsaa en, i udvortes Forstand
uerstattelig, Tid er gaaet tabt for Virkeliggjørelsen, dog er en Mulighed til en
Indhenten af det Forsømte, til et Vederlag for den tabte Virketid. Den
\emph{anden} Betingelse, til hvilken der allerede fra en Side er viist hen ved
den første, er den, at det Ethiske kan blive \emph{Totalitetsbestemmelse for
Menneskelivet}, ikke blot tiltrods for hint Tidstab, man ogsaa tiltrods for det,
at Individets ethiske Villen, ifølge de menneskelige Udviklingsbetingelser,
nødvendigviis er begrændset til et bestemt Modenhedsstadium i Individets Liv. Den
\emph{tredie} Betingelse endelig er den, at den individuelle menneskelige
Tilværelse faaer den \emph{Evighedsbestemthed}, der svarer til \emph{Idealets
Begreb}.

\emph{α) Ophævelsen af Skylden.} --- Vi ere ovenfor, fra et rationelt
Udgangspunkt, naaede til Bestemmelsen af den individuelle Villies Incongruents
med den ethiske For\-%
% HEAD, p. 168: the run-in „α) Ophævelsen af Skylden.“ is letterspaced regular
% Fraktur followed by an em-dash, with the Greek enumerator „α)“ in ANTIQUA
% ITALIC (verified at 400 dpi). Set with \emph{} like the run-in heads of pp. 159
% and 162. The Indhold instead reads „α) Skyldens Ophævelse“ (pp. 168--176) ---
% the two nouns in the other order; the PRINTED head stands in the body.
% NOTE ON SCOPE: the dispatch brief for this batch said this head „stands just
% past your range's end“ and should not be set --- but p. 168 lies inside
% pp. 162--175, and the brief's own next clause assigns it to „the pp. 176--189
% batch's predecessor“, which is this batch. It is set here. If it is to be
% dropped, delete the head line and the em-dash, not the sentence.
% RULE, p. 168: a short centred SINGLE hairline stands between „…til Idealets
% Begreb.“ and the α) head, i.e. above a head; not transcribed, as at p. 164.
% GREEK, p. 168: τέλος TWICE, antiqua italic with the acute, in the one sentence
% („Individets individuelle τέλος“ and „ved Menneskevæsenet bestemte τέλος“).
% Neither OCR witness reads either: tesseract gives „7é2oç“ and „7é4oç“, ABBYY
% drops both words entirely, leaving „individuelle og af det almeen-menneskelige,
% ved Menneskevæsenet bestemte saaledes forbundne“. Read off the image at 200 and
% 400 dpi. „Menneskevæsenet“ beside the second one IS letterspaced, which the
% spacing pass missed because the Greek breaks its line model.
% PRINTER'S DEFECT, p. 168: „ikke blot tiltrods for hint Tidstab, man ogsaa
% tiltrods for det“ --- „man“ for „men“. Confirmed at 600 dpi (m-a-n, the „a“
% quite unlike the „e“ of „men“ elsewhere in the line) and read so by both OCR
% witnesses. Transcribed as printed. The Trykfeil leaf lists a „man, l.: men“ at
% Side 119, but that entry belongs to the block for „Problemet om Tro og Viden“,
% not to this volume; this instance is unlisted.
% p. 168 carries no footnote.
% --- p. 169 ---
dring og derved til Begrebet om \emph{Skylden}. Dette Begreb bliver nu nærmere at
udvikle.

Bestemmelsen af Incongruentsforholdet maa faae sit forskjellige Udtryk eftersom
Incongruentsen sees fra \emph{Væsensnødvendighedens} eller fra
\emph{Villiesbestemthedens} Side. Sees den, abstract, fra \emph{Væsenets} Side,
bliver den, hvor den overhovedet indrømmes som virkelig, opfattet i metaphysiske
Bestemmelser, væsentlig svarende til dem, under hvilke den græske Ethik opfattede
den. Det bliver da \emph{Endelighedens} μὴ ὄν, der udpræger sig i
Væsensyttringens \emph{Ufuldkommenhed}: Incongruentsforholdet til \emph{det
Ethiske} bliver bestemt \emph{negativt}, som en \emph{Mangel}. Sees Forholdet
derimod fra \emph{Villiens} Side, gjøre andre Kategorier sig gjældende. Der
bliver da Tale om \emph{Skyld}, endnu bestemtere om \emph{Synd}. Begrebet
\emph{Skyld} kan endnu holdes svævende paa Grændsen af det Metaphysiske og det
Ethiske. Endelighedens, „Enkelttilværelsens Skyld“, der kræver sin Soning, er
Udtryk for Forholdet i den gamle græske Philosophie, der væsentlig betragter
Naturexisteniserne, og i dette Udtryk afspeiler den østerlandske Grundanskuelse
sig. Men i sin strængere og egentlige Betydning bliver Skylden \emph{ethisk}
Kategorie og betegner \emph{Villiens selvforskyldte Misforhold til det Ethiske},
hvad enten nu dette Misforhold er det, at Existentsen endnu ikke, skjøndt
Muligheden var der, er sat under det Ethiskes Bestemmelse, eller at Bruddet med
det Ethiske udtrykkeligt er sat gjennem Villiens villede Selviskhed, der, idet
den gjør sig gjældende i sin negative Uendelighed, sætter en
Endelighedsbestemmelse som det Absolute, altsaa i det Absolutes Sted. Her er
\emph{Villien} accentueret som den, der i \emph{Selvraadighed} paadrager sig
Skyld, og selv om \emph{Erkjendelsens} Betydning for Villiesbeslutningen nok saa
meget betones, saa betones ligesaafuldt \emph{Omsætningen} i Villiesbestemmelse.

\emph{Skyldens} Begreb faaer nu sit bestemtere Udtryk i \emph{Synden}, hvad enten
der nu ved denne Bestemmelse udtrykkelig fremhæves Forholdet til den
\emph{guddommelige Villie}, eller \emph{Væsenets} Affection ved den uethiske
Handling,
% GREEK, p. 169: μὴ ὄν --- antiqua italic, mu, eta with the GRAVE, omicron with
% smooth breathing AND acute (ὄ), nu. Read at 800 dpi. Neither witness has it:
% tesseract gives „uy 0v“, ABBYY „o>“. This page is not in the recon pass's Greek
% list.
% EMPHASIS, p. 169: the quoted „Enkelttilværelsens Skyld“ is set TIGHT inside its
% quotation marks --- checked at 600 dpi against the spaced „Endelighedens“ two
% words earlier. No \emph{} inside the quotation.
% PRINTER'S DEFECT, p. 169: „der væsentlig betragter Naturexisteniserne“ ---
% „Naturexiſteniſerne“ for „Naturexistentserne“: the „t“ of the stem is set as a
% dotted „i“. At 1200 dpi the sort plainly carries a dot and has neither the
% ascender nor the crossbar of a Fraktur „t“; the „ſt“ ligature earlier in the
% same word gives the calibration. Both OCR witnesses read a vowel there
% („Naturexisteniserne“ / „Naturexistentserne“). Transcribed as printed; not in
% the Trykfeil leaf.
% QUOTES, p. 169: one „ and one “, balanced.
% p. 169 carries no footnote.
% --- p. 170 ---
Tabet af dets Integritet. Ved begge disse Bestemmelser kan der føres over paa det
\emph{Paradoxes} Gebet. Opfattes nemlig Forholdet til den guddommelige Villie som
Forhold til en Villie, hvis Lov er den med Erkjendelsens Fornuftlov uensartede, i
sin Vilkaarlighed ubegribelige, saa bliver Synden det Paradoxe, Bevidstheden om
Synden som Synd betinget ved et Troesforhold af paradox Natur, og Syndens
Forsoning bliver betinget ved og Gjenstand for et Troesforhold af samme Art. Og
paa samme Maade forholder det sig med den Opfattelse af Syndens Virkning, der i
Læren om Arvesynden statuerer \emph{Menneskenaturens Fordærvelse} ved den første
syndige Villiesact. Denne Fordærvelse, saaledes som Kirkelæren opfatter den,
bliver det Ubegribelige. Menneskevæsenet skal tænkes som fordærvet ved en af
selve Væsenet, og, bestemtere, af det endnu i Integritetens Tilstand værende
Væsen, fremgaaende Act. Villiesyttringen, der realiserer en Mulighed, der er
anlagt i Væsenet, skal forvandle Væsenet, ved een Act destruere det. Som Parallel
til dette Ubegribelige staaer saa det Ubegribelige i Restitutionen ved et
guddommeligt Under, ved Guds paradoxe Menneskevordelse, ved Gudmenneskets
Selvopoffrelse, der tilegnes gjennem Troen. I denne Opfattelse unddrager Synden
og Forsoningen sig al Erkjendelse. Men Bestemmelsen af \emph{Synd} som omfattende
saavel en Strid \emph{mod den guddommelige Villie}, som en Indvirkning \emph{paa
Væsenet}, kan fastholdes uden at det Rationelles Gebet forlades. Bestemmelsen
Synd kan fastholdes i Betydningen af Overtrædelse af den guddommelige Villies
Fordring, idet den \emph{guddommelige Villies} Begreb opfattes paa den Maade, der
blev angiven under Udviklingen af det Absolutes Begreb. Og Begrebet om \emph{en
Affection af Væsenets Integritet ved Synden} kan fastholdes, idet det opfattes i
Overensstemmelse med hvad der blev udviklet i de foregaaende Afsnit. Ifølge den i
den menneskelige Aand tilstedeværende dobbelte Mulighed: at være og at ville sig
som bevidst indordnet Led i Totaliteten, og at være og at ville sig i Isolationen
efter sin negative Uendelighed,
% READINGS, p. 170, both settled at 600 dpi against tesseract: „fremgaaende Act.
% Villiesyttringen“ takes a FULL STOP, not the comma tesseract read; and „der er
% anlagt i Væsenet“ has NO comma after „i“, where tesseract inserted one.
% p. 170 carries no footnote and no Greek.
% --- p. 171 ---
bliver der for Villien en Mulighed til, ved, som ond, at sætte sig efter sin
negative, abstracte Uendelighed, at \emph{endeliggjøre} sig, og som endeliggjort
er Villien ikke integer, og \emph{Væsensyttringen} Villiens Mangel paa Integritet
berører \emph{Væsenet} ved Spliden, der sættes i dette. Men under denne Væsenets
Splid, under hvilken Integriteten er som suspenderet, er Væsensenergien efter sin
Uendelighed negativt tilstede i den endeliggjorte Villie (jfr.\ de tidligere
Udviklinger). Spliden faaer ikke Magt til at opløse Væsenet; Redintegrationens
Mulighed og Kilde er bevaret i Væsenet selv, og Muligheden bliver til Virkelighed
i Væsenets Selverhverven, i dets Gjennembrud gjennem den onde Villies
Selvopløsning. At denne Selvopløsning kun bevirkes ved Væsenet idet den bevirkes
ved det iboende Absolutes Magt, ligesom Villiens Sætten af sig som ond efter den
usande Uendelighedsabstraction kun var mulig ved det sande Uendeliges Iboen, er
ovenfor paaviist. Men dette Forhold tilintetgjør ikke Væsenets \emph{egen}
Energie; thi det gjennemvirkende Absolute er et saadant, der ifølge Forholdet
mellem dets Idealitets- og Realitetsside lader blive en Plads for det Endeliges
Virken.

Ud fra den givne Bestemmelse af Skyldens og Syndens Begreber, ved hvilken disse
altsaa blive Kategorier i en \emph{rationel} Ethik, bliver da Spørgsmaalet at
opkaste om \emph{Muligheden af en Skyldens Udsoning gjennem en factisk Indhenten
af det under den ethiske Handlen Forsømte}. Denne Mulighed var det, der sattes
som Betingelse for, at det med Skyld behæftede Individs \emph{principielle}
Restitution kunde blive til en \emph{Virkeliggjørelse af Idealet}, til en
\emph{fuld og sand} Forsoning.

Forat en Indhenten af det Forsømte, og derved en ethisk Redintegration af
Individet trods Skylden skal kunne tænkes, fordres det, at i den
\emph{principielle Forvandling} i det Ethiskes Gjennembrud en \emph{Omsætning}
finder Sted \emph{af en gjennem den onde Villen actualiseret Kraft}. En saadan
Actualiseren af en Kraft ved den onde Villen, gjennem denne Villens Uendelighed,
blev ovenfor paaviist,
% READING, p. 171: „i en rationel Ethik“ --- ONE „l“, letterspaced, verified at
% 800 dpi; tesseract's „rationell“ is a misreading and ABBYY agrees with the
% image. „Udsoning“ likewise carries no „ø“.
% EMPHASIS, p. 171: „berører Væsenet ved Spliden“ spaces „Væsenet“ only ---
% „Spliden“ here is set close, unlike the spaced „Spliden i Væsenet“ of p. 163.
% Checked at 400 dpi with the two side by side.
% p. 171 carries no footnote and no Greek.
% --- p. 172 ---
og denne actualiserede Kraft er det, der i det Godes Gjennembrud i Villien tages
i det Godes Tjeneste, saaledes at den bliver en Energie, hvis tidligere
Udvikling, idet gjennem Angeren Kraftretningen vendes om, faaer positiv Betydning
gjennem en forhøiet Villiesvirksomhed. Sees denne Kraftens Actualiseren i Forhold
til \emph{Væsenet}, saa falder Eftertrykket paa Væsenets negative Virken under
Forstyrrelsen, og denne Væsensenergiens skjulte Virken er det, der gjør det
muligt, at, naar Gjennembruddet skeer, den med det Uethiske udfyldte Tid kan faae
positivt Værd for den ethiske Udvikling. Det er saa \emph{Angeren}, som det til
Gjennembrud komne Væsens Reaction mod den selvtilregnede onde Villen, der holder
den ved denne Villen actualiserede Kraft levende som en for det Gode virksom
Kraft. Derfor faaer Angeren en \emph{blivende} Betydning, og den maa --- i
Modsætning til hvad Fichte vilde --- opfattes som bevirkende Kraftforøgelse. ---
Sees Actualiseringen af Kraften i Forhold til \emph{det objectivt Gode og det
Absolutes Villie}, saa falder Eftertrykket paa det Forhold under den onde Villies
Kamp mod det Gode, ifølge hvilket den under denne Kamp \emph{ubevidst} bestemmes
af det \emph{Væsensgyldige}, mod hvilket den kæmper, til en Anspænden af Kraften,
til en Kraftforøgelse i negativ Retning. I selve Kampen er saaledes et Forhold
sat til det Væsensgyldige som det Villien Beherskende, og ved dette Forhold kan,
idet Villiesretningen forvandles i det Godes Gjennembrud ɔ: idet Omvendelsen
skeer, selve Kampen mod det Gode faae positiv Betydning for den ethiske
Udvikling. I Omvendelsen gaae de falske, negative, Yttringer af Villiens Energie,
som forvandlede til positive Kraftbestemmelser, ind i Udviklingen, og det
forudgaaende Liv slutter sig gjennem Angeren sammen med det nærværende til en
continuerlig intensiv Livsenhed. Iøinefaldende Exempler paa et saadant Forhold
fremtræde hvor en energisk Kamp mod det høiere Sande forvandles til en energisk
Tilslutning til det. I selve Kampens Lidenskab er det Sandes Magt aabenbar: ved
selve den lidenskabelige Modstand optages det saaledes i Bevidstheden,
gjennemtrænger
% READING, p. 172: „som det til Gjennembrud komne Væsens Reaction“ --- NO comma
% after „komne“. ABBYY reads one; at 600 dpi the mark is a small round blot lying
% out in the right margin, well clear of the last letter and below the baseline,
% i.e. dirt or show-through, not a comma. The sense („det til Gjennembrud komne
% Væsens Reaction“) takes none either.
% NAME, p. 172: „Fichte“ is Fraktur and NOT letterspaced --- the recon pass's rule
% that names stay Fraktur holds here, and no antiqua name appears in this batch.
% ID-EST MARK, p. 172: „i det Godes Gjennembrud ɔ: idet Omvendelsen skeer“; the
% batch's second, after p. 164.
% p. 172 carries no footnote and no Greek.
% --- p. 173 ---
den saaledes, at i Omslaget Alt er forberedt til en umiddelbar, samlet Virken af
hele Aandens Kraft i den Sandheds Tjeneste, der nu bliver Livsprincip og strax
kan virke som saadant i alle Livets Kræfter. Man erindre kun eet Exempel:
Saulus's Forhold til Christendommen. ---

I det Foregaaende er \emph{Redintegrationen} seet væsentlig i Forhold til
\emph{selve Villien}, men den bliver tillige at see i Forhold til \emph{det ved
det onde Villie Bevirkede}, men saaledes, at dette Synspunkt, hvad det efter
Sagens Natur ikke kan eller maa, ikke abstract holdes ude fra hint.

Det ved den onde Villie Bevirkede bliver deels en Forstyrrelse i selve Individet,
i dets Realitetsside eller i dets aandelige Udvikling; deels en Forstyrrelse i
andre Individers physiske eller aandelige Existents; deels en Forstyrrelse i det
objectivt Almene. Den ved Synden bevirkede Forstyrrelse i \emph{Individets
Realitetsside} bliver i denne Sammenhæng at see som en \emph{Forstyrrelse i de
physiske Livsbetingelser for Villiens Virken}, og dens Betydning at opfatte fra
dette Synspunkt. Med Hensyn hertil er allerede i det Foregaaende, ved Udviklingen
af den Frihedens Mulighed, der bevaredes under Forstyrrelserne og Manglerne i den
physiske Organisation, det Væsentlige angivet. Eftertrykket kommer her, som
overalt i disse Forhold, til at falde paa denne \emph{Frihedens blivende
Mulighed}, i hvilken \emph{Energien til Villiesforvandlingen} er tilstede, og
idet denne Energie er der, bliver den i det Foregaaende charakteriserede
Kraftforvandling mulig trods det Physiskes Mangel, og denne Mangel bliver saa
ethisk Incitament for den Angerens Villie, der, trods den, skal hævde sig i det
Ethiske og er relativt til at hæve ved Villiens Energie. --- Opfattes
Forstyrrelsen i Individet som Forstyrrelse i \emph{dets aandelige Udvikling}, saa
bliver det samme Grundsynspunkt at gjøre gjældende. Frihedens Mulighed bliver at
see som bevaret, og idet den \emph{hele} aandelige Udvikling her sees som
concentreret om den \emph{ethiske} Udvikling, saa ligger i den ethiske Friheds
bevarede Mulighed tillige Muligheden af en væsentlig Restitution af de
Aandsfunctioner, der ere det Ethiskes Betingelser.

% APOSTROPHE, p. 173: „Saulus's Forhold til Christendommen“ --- the printed form
% carries an apostrophe and an „s“, as transcribed.
% EMPHASIS, p. 173: „af hele Aandens Kraft“ in line 2 is TIGHT; the spaced „hele“
% flagged by the spacing pass is the later one, „den hele aandelige Udvikling“.
% SHORT PAGE, p. 173: the type block ends one line short of the measure, at
% „…det Ethiskes Betingelser.“, with white space below and NO rule. p. 174 opens
% a new paragraph --- hence the blank line before the marker.
% p. 173 carries no footnote and no Greek.
% --- p. 174 ---
Sees Forstyrrelsen, der er tilveiebragt ved den onde Villie, som en Forstyrrelse
i \emph{andre Individers physiske} Existents, eller i \emph{deres aandelige}
Livsbetingelser, saa gaaer den, som et \emph{relativt} forstyrrende Indgreb i
denne, ind under samme Synspunkt som de ovenfor omhandlede Forstyrrelser i selve
Individet. For det andet Individ bliver, trods den, Frihedens Mulighed bevaret,
og Forstyrrelsen sætter en \emph{ethisk Opgave} for Individet, hvis Løsnings
Mulighed er sikkret ved hiin Frihedens Mulighed. Er Forstyrrelsen derimod den
\emph{totale} Tilintetgjørelse af den hele physiske Tilværelse, saa vises vi,
ligesom hvor Individet tilintetgjør sin \emph{egen} Existents, med Hensyn til
Spørgsmaalet om den ethiske Opgaves Løsning for Individet hen til Problemet om
Individets Evighed, der vil blive behandlet i det Efterfølgende.

Er Forstyrrelsen i de andre Individer en Forstyrrelse i deres \emph{moralske}
Existents, saa bliver, under alle Nuancer af Forholdet, det andet Individs
\emph{Tilregnelighed} under Paavirkningen at accentuere, saa at Modtagelsen af
denne Paavirkning sees som \emph{villet}, og det Synspunkt for Restitutionen, der
var det oprindelige med Hensyn til det handlende Individs Villie, bliver da ogsaa
at gjøre gjældende med Hensyn til det lidendes.

Er endelig Forstyrrelsen en Forstyrrelse i det \emph{objectivt Almene}, i den
\emph{sædelige Orden}, saa bliver her det Afgjørende hvad der allerede ovenfor er
fremhævet: de gode Kræfters Reaction mod det Forstyrrende, idet dette netop ved
Kampen mod det Gode fremkalder en forøget Kraftudvikling hos dem, der kæmpe for
det Gode, saaledes at i Heelhedens Udvikling en Compensation finder Sted, hvilket
ogsaa gjælder for det Tilfælde, at hele Tidsaldere have en Retning henimod det
Onde; thi da blive Tidsaldere at see i Forhold til Tidsaldere ligesom før
Individer i Forhold til Individer.

Idet Forstyrrelsen i alle disse Forhold sees som Gjenstand for det Individs
Bevidsthed, fra hvilken den er udgaaet, bliver den, naar det principielle
Gjennembrud af den
% EMPHASIS, p. 174: the opening line „Sees Forstyrrelsen, der er tilveiebragt ved
% den onde“ is NOT letterspaced --- the spacing pass flagged „tilveiebragt ved“,
% but at 600 dpi the letters are close and only the word gaps are wide, an
% artefact of a loosely justified first line. Rejected.
% p. 174 carries no footnote and no Greek.
% --- p. 175 ---
gode Villie er skeet, som Gjenstand for Angeren, Incitament for en Villiesvirken
i Individet, der i Forhold til de andre Individer og til den sædelige Orden maa
træde frem ogsaa som en Stræben efter en \emph{ydre} Gjenoprettelse og Erstatning
forsaavidt den staaer i Individets Magt. Men Hovedaccenten falder overalt paa
Restitutionen i \emph{Individernes Indre}.

\begin{center}\rule{2cm}{0.4pt}\end{center}
% RULE, p. 175: a short centred SINGLE hairline standing mid-page between
% „…i Individernes Indre.“ and „De foregaaende Udviklinger…“, with no head
% adjacent --- a thought-break, as at pp. 19, 28, 42, 62, 97 and 159.
% Transcribed. NOTE: the batch pass reported this rule as „not centred, about a
% fifth of the way in from the left margin“. That is wrong. Measured against the
% type measure on the 150 dpi image, its midpoint falls 2.8 % of the measure left
% of centre --- within the scatter of the centred rules elsewhere in the book
% (p. 19 is 0.5 % left, p. 164 is 1.1 % right). No rule in this book is set off
% centre.

De foregaaende Udviklinger have overalt væsentlig taget Hensyn til den Form af
den onde Villie, der er sat som den \emph{primitive: Trodsens,
Selvraadighedens}. Men ifølge denne primitive Forms Forhold til de særskilte
Former af det Onde faaer det ogsaa Anvendelse paa disse. Hvad enten det Onde
fremtræder som Sandselighed, Dyriskhed, Sløvhed, eller som Ladhed og
Uvirksomhed, saa bliver dog i alle, selv de meest passive Former, det at
fastholde, at det Onde, som Villiesbestemmelse, paa ethvert Punkt maa være
villet, at det aldrig er blevet til rent passivt eller fortsættes rent passivt,
men altid bliver til og fortsættes ved et Villiens Samtykke. Derfor er, hvor det
Onde har \emph{Svaghedens} Form i alle dens mulige Nuancer, \emph{Trodsen}
bestandig tilstede, og enhver Villen af det Onde skjuler i sig den onde Villies
primitive Bestemmelse: at stræbe at sætte sig selv i det Absolutes Sted. Den onde
Villen er altsaa overalt en Yttring af en Energie, men af en Energie, der, brudt
i sig selv, arbeider sig selv imod, og derfor bærer Svagheden, Magtesløsheden
skjult i sig, ligesom den i sin Hovmod skjuler Selvhenkastelsen, Retningen nedad,
mod det, der er lavere end Aanden.

Ligesom det Udviklede gjælder om de forskjellige Former af det Onde, saaledes
faaer det ogsaa Anvendelse paa de forskjellige Skikkelser, under hvilke
Villiesforandringen viser sig. Umiddelbart faaer det Anvendelse paa den Form,
hvor Villien til det Gode energisk kommer til Gjennembrud og energisk hævder sig
efter Gjennembruddet i Livets Concretion. Men det faaer ogsaa Anvendelse paa de
Tilfælde,
% p. 175 carries no footnote and no Greek. The batch ends here, in mid-sentence.
% --- p. 176 ---
hvor det Sporadiske i Selvet, ogsaa efter det Ethiskes principielle Gjennembrud,
virker mod dette, og hvor Villien, der i Angeren vender sig bort fra det Onde, er
svag og hæmmet, og hvert Øieblik vakler over imod det Onde. Thi naar ikke denne
Vaklen fører til et principielt Tilbagefald, --- i hvilket Tilfælde det
almindelige Synspunkt for den onde Villies Virken bliver at gjøre gjældende, ---
saa bliver Villiens Oscilleren gjennem Angeren et Incitament for en forhøiet
Villiesanspændelse, gjennem Følelsen af den truende Fare et Incitament for en
inderligere Søgen tilbage til de Forsoningens Kilder, der findes i det
Menneskelige og i det Menneskeliges guddommelige Princip.

Hvor i det enkelte Individ Omvendelsen, det Ethiskes Gjennembrud, slet ikke kan
paavises i noget Moment af Individets Liv, viser Spørgsmaalet om Forsoningen
tilbage til Spørgsmaalet om Individets tidsbestemte Livs Forhold til Individets
Evighedsbestemthed, til dets evige Betydning i Tilværelsens Organisation. Dette
Spørgsmaal vil blive omhandlet i det Efterfølgende.

\begin{center}\rule{2cm}{0.4pt}\end{center}
% RULE, p. 176: a short centred SINGLE hairline standing mid-page between
% „…omhandlet i det Efterfølgende.“ and „Idet Individet som forsonet…“, with no
% head adjacent --- a thought-break, as at pp. 19, 28, 42, 62, 97, 159 and 175.
% Transcribed. p. 176 does NOT close with a rule: its last line is „Aandens Enhed
% til sin Grund.“ and the type block runs the full depth.

Idet Individet som forsonet gjennem sin ethiske Villen forholder sig til det
\emph{Absolute}, viser dette sig afgjørende for det i de \emph{ethiske}
Bestemmelser, hvis Anvendelighed paa det i det Foregaaende (S. 78) blev antydet.
\emph{Det Absolute} sees som \emph{det Gode}, idet det, som gjennemvirkende den
endelige Villie, er denne Villies ethiske τέλος, dens Ideal. Det sees som
\emph{det Hellige} idet dets Energie er bestemt ved Væsensidealitetens Lov. Det
sees som \emph{det Retfærdige} idet det lader Villiens Resultat bestemmes efter
dens Grundforhold til Principet. Det sees endelig under \emph{Kjærlighedens}
Bestemmelse som det, der, idet det meddeler sig og i Væsensmeddelelsen aabenbarer
sig, sammenslutter de enkelte Aander i sin Enhed og sammenfatter dem indbyrdes i
den Kjærlighedens Harmonie, der har Aandens Enhed til sin Grund.
% GREEK, p. 176: τέλος, antiqua italic with the acute and the final sigma, in
% „denne Villies ethiske τέλος, dens Ideal“. Tight, not letterspaced. Neither OCR
% witness reads it (tesseract gives „7é2oc“, the ABBYY layer „rrloy“).
% EMPHASIS, p. 176: the letterspaced runs are „Absolute“ (the preceding „det“ is
% TIGHT), „ethiske“ (but not „Bestemmelser“), „Det Absolute“, „det Gode“, „det
% Hellige“, „det Retfærdige“ (its „det“ stands at the end of the line and IS
% spaced) and „Kjærlighedens“. Checked at 200 and 400 dpi against the tight
% „Væsensidealitetens Lov“ on the same page. spacing.py's flags for „Thi“, „til“,
% „Liv,“, „Livs“ and „Lov.“ are justification noise, not emphasis.
% PENCIL, p. 176: a modern reader has ruled a thin horizontal line through the
% word „lader“ in „idet det lader Villiens Resultat bestemmes“. The word itself
% is unambiguous at 400 dpi and is transcribed; the stroke is marginalia and is
% not, per the standing convention.
% p. 176 carries no footnote.

% --- p. 177 ---
\setcounter{footnote}{0}
I det udviklede ethisk-religiøse Forhold af Individet til det Absolute ligger der
et \emph{rationelt Æquivalent} for det, som den \emph{positive} Religion søger i
Individets Henvendelse til den \emph{personlige} Gud gjennem \emph{Bønnen}: et
Forhold, i hvilket, som andensteds viist\footnote{Jfr. Problemet om Tro og Viden,
S. 95 ff.}, den positive Religions Grundeiendommelighed er concentreret. Der er i
hiint ethisk-religiøse Forhold en \emph{Sindets Henvendelse} til det
gjennemvirkende Absolute for, gjennem Enheden med dette, at naae den manglende
ethiske Kraft. Der er i denne Henvendelse et \emph{Contemplationens} Moment, idet
det Absolute erkjendes efter sit Grundforhold til Menneskevæsenet og
Menneskevillien. Der er i den et \emph{Handlingens} Moment, idet Villien i
Selvhengivelsen frit underordner sig under den guddommelige Villie. Der er en
\emph{Sindets Styrkelse}, en \emph{Følelsens Tilfredsstillelse} i dette Forhold, i
hvilket Væsenets Kilder springe frem fra det Uendeliges Dyb. Der er i det en
\emph{sikker Opnaaelse af det tilstræbte Maal}; thi i selve den fulde
Selvhengivelse er den fulde Selvbekræftelse, en Vinden af det Eviges Kraft, idet
Individet frit taber sig selv til det Evige. Der er endelig i dette Forhold en
Frihed for den Egoisme, der maa blive tilbage, naar der i Bønnen bedes til en
vilkaarlig Villie, selv naar der bedes om Syndsforladelse; thi ogsaa i denne Bøn
er der, naar den bliver Bønnen om Eftergivelsen af en ved den guddommelige
Personligheds vilkaarlige Villie vilkaarligt bestemt Straf, et
Udvorteshedsforhold til det Guddommelige, i hvilket Endeligheden og netop derfor
Egoismen er bevaret.

\section*{β) Det Ethiske som Totalitetsbestemmelse for Menneskelivet.}

Som den \emph{anden} Betingelse for Idealets Realisation blev det ovenfor (S.
168) fordret, at det Ethiske maatte kunne blive \emph{Totalitetsbestemmelse} for
Menneskelivet, ikke blot tiltrods for det ved Synden Forsømte og For\-%
% HEAD, p. 177: „β) Det Ethiske som Totalitetsbestemmelse for Menneskelivet.“ ---
% bold (halbfette) Fraktur, centred, set over two lines („…for“ / „Menneskelivet.“).
% Verified at 300 dpi. The enumerator „β)“ is ANTIQUA ITALIC, as against the
% Fraktur of the rest of the head; the Greek letter is typed directly. The
% Indhold agrees with the printed head.
% RULE, p. 177: a short centred SINGLE hairline stands ABOVE the display head,
% between „…Egoismen er bevaret.“ and the head, as at pp. 97, 126 and 164. Not
% transcribed, per the standing convention for head rules.
% FOOTNOTE, p. 177: printed mark *) --- the only note on the page, below a short
% rule, called after „viist“ and before the comma. „Jfr.“ is a J-word and is set
% with J, per the recorded I/J normalisation.
% EMPHASIS, p. 177: several split runs, all left as printed. „Contemplationens“
% and „Handlingens“ are spaced but the following „Moment“ is TIGHT in both cases,
% checked side by side at 340 dpi. „personlige“ is spaced but „Gud gjennem“ is
% not; „Bønnen“ is. „Totalitetsbestemmelse“ is spaced but „for Menneskelivet“ is
% not. The run „sikker Opnaaelse af det tilstræbte Maal“ begins at „sikker“: the
% preceding „er i det en“ is tight at 600 dpi. „viist“ is NOT letterspaced ---
% spacing.py's 1.20 there is noise.
% p. 177 carries no Greek. The signature numeral „12“ at the foot is not
% transcribed.
% --- p. 178 ---
styrrede, men ogsaa tiltrods for den Tidsbegrændsning for Muligheden af en ethisk
Villen, der er given ved de menneskelige Udviklingsbetingelser. Hvorvidt denne
Fordring kan opfyldes, bliver nu at betragte, nærmest fra det sidste Synspunkt,
da det, der angaaer det første, allerede er omhandlet i det Foregaaende.

I Opfattelsen af den \emph{ethiske Villie} som \emph{Væsens}villen og af den
\emph{ethiske Handling} som Udtryk for \emph{Væsens}bestemmelsen er der viist hen
til den Fordring, at det Ethiske skal være \emph{Totalitets}bestemmelse. Den
ethiske Fordring er, som vort \emph{Væsens} Fordring og som den absolute Villies
Fordring gjennem vort \emph{Væsen}, den \emph{ubetingede} Fordring. Den viser sig
som saadan i en \emph{dobbelt} Betydning: som den, der som et kategorisk Imperativ
har Gyldighed ubetinget for \emph{alle} menneskelige Individer, og som den, der
har Gyldighed ubetinget for \emph{alle} Individets Livsforhold, for den
individuelle menneskelige Livsheelhed. Fordringen om en ethisk Livstotalitet er
altsaa retfærdiggjort ved det Ethiskes Forhold til Væsenet; men forat denne
Fordring skal kunne virkeliggjøres under Individets Tidsudvikling, kræves paa den
ene Side, at det \emph{forethiske Liv}, ɔ: det Liv, der er ført inden den for det
ethiske Valg nødvendige Bevidsthedsudvikling har fundet Sted, kan
\emph{overtages ethisk}; paa den anden Side, at den ethiske Handlen kan faae en
Concretion, svarende til Personlighedens concrete Livsindhold.

At det \emph{forethiske Liv overtages ethisk}, saaledes, at det Livsstadium, der,
ifølge Menneskeaandens psychologiske Udviklingsbetingelser, \emph{maa} ligge
forud for det Ethiske, drages ind under det Ethiske og bliver Moment i en ethisk
Livsheelhed, der har Enhed og Continuitet, vil nærmere sige, at Individet, idet
det i det ethiske Valg constituerer sig som ethisk Personlighed, i denne
Concentration sammenfatter sit forudgaaende Liv under et \emph{nyt Livsprincip},
og idet det i denne Concentration ligesom begynder sit Liv forfra, forvandler den
Udvikling, der viste sig som \emph{forudgaaende Betingelse} for dets \emph{sande},
ɔ: til \emph{Væsenet svarende}
% EMPHASIS, p. 178: three split runs of the kind recorded for pp. 6 and 15, all
% left as printed. (i) „ethiske Villie“ is spaced, the following „som“ tight, and
% of the broken word „Væsens=“ / „villen“ only the first half is spaced. (ii) The
% same word „Væsensbestemmelsen“ stands WHOLE on the next line and is again
% spaced only in „Væsens“ --- so this is not a line-break effect. (iii)
% „Totalitetsbestemmelse“ likewise: „Totalitets“ spaced, „bestemmelse“ tight.
% Also „ubetingede“ is spaced across its own line break but „Fordring“ is not,
% and „dobbelt“ is spaced but „Betydning“ is not. Verified at 340 dpi.
% EMPHASIS, p. 178, last line: in „ɔ: til Væsenet svarende“ the word „til“ is
% TIGHT --- measured at 900 dpi against the tight „for dets“ earlier in the same
% line and the spaced „sande“ between them. The same „til“ recurs tight on p. 179
% between „Existentsform,“ and „Bestemmelse“.
% p. 178 carries no footnote, no Greek and no rule. Two ɔ: marks.
% --- p. 179 ---
\emph{Existentsform}, til \emph{Bestemmelse i denne Existents}, og ved denne
Forvandling, \emph{ethisk} seet, bliver \emph{\textit{causa sui}}. Det, der i det
forethiske Liv var \emph{Naturbestemthed, Anlæg}, bliver da \emph{ethisk
Frihedsopgave}, og derved at den forudgaaende naturbestemte Udvikling, de endnu
ikke ethiserede Kræfters Virken, ethisk opfattes \emph{gjennem Angeren}, bliver
det forethiske Livs og den forethiske Handlens Resultat metamorphoseret, og i
denne Metamorphose indoptaget i det ethiske Livsprincips Sphære, forat
gjennemarbeides ud fra det, saa at den principielle Væsensforvandling udpræges i
den concrete Personligheds ethiske Udfoldelse.

Men idet Personligheden ethisk skal gjennemarbeide sit forethiske Livs Resultat
og som ethisk leve et Liv, heelt igjennem bestemt ved det nye Livsprincip, gjør
Fordringen til \emph{en adæquat Concretion af det Ethiske} sig gjældende. Forat
det Ethiske skal kunne gjennemforme Personlighedens Naturgrundlag, og ud fra det
ethiske Valg forme en Livstotalitet, i hvis enkelte Handlinger det bliver det
bestemmende Princip, fordres der, at det for Bevidstheden udfoldes i en
Concretion, svarende til Livets og Handlingens Concretion, og i denne Concretion,
i hvilken det omfatter et \emph{empirisk Erkjendt}, bevarer det
\emph{Ethisk-Aprioriskes Vished og Nødvendighed}. Idet den \emph{aprioriske}
Bevidsthed om det Gode, om Pligten, om Dyden, i Handlingens Virkelighed skal blive
til en Bevidsthed om dette bestemte Gode, om denne bestemte Pligt, om denne
bestemte Dyd, og idet denne Bevidsthed, gjennem Sammenhængen med hiin, skal bevare
hins \emph{ubetingede Vished}, fordres altsaa, at den \emph{abstracte
Væsensbevidsthed} udfoldes til en \emph{væsentlig Virkelighedsbevidsthed}, ud fra
hvilken der kan handles, saa at ikke i det Concrete Afgjørelsen gjennem en
\emph{umiddelbar æsthetisk Takts relative Bevidstløshed} træder i Stedet for den
af bevidst Reflexion baarne ethiske Beslutning, der ligetil den sidste punktuelle
Tilspidsen til Virkeliggjørelse hviler paa et Alment og for det ethiske Subjects
Bevidsthed hævder sin Gyldighed ved at lade sig føre tilbage til normerende
Almeenbestemmelsers Love.
% LATIN, p. 179: „causa sui“ is antiqua AND letterspaced --- the letters are as
% widely set as in the Fraktur „ethisk“ two words before it, and much wider than
% the tight antiqua of pp. 64, 71, 83, 89 and 91. Encoded \emph{\textit{…}}, as
% „causa sui“ and „causa finalis“ were on p. 68.
% EMPHASIS, p. 179: the run opening the page is broken by a tight „til“ ---
% „Existentsform“ spaced, „til“ tight, „Bestemmelse i denne Existents“ spaced ---
% and is a continuation of the run begun on p. 178; it is set here as two
% \emph{} groups round the tight „til“, which is what the print shows.
% „ligetil“ (spacing.py 1.21) is NOT letterspaced.
% p. 179 carries no footnote, no Greek and no rule; its last word „Love“ is
% followed by a full stop (not the comma tesseract reads), and the paragraph ends
% there. The signature „12*“ at the foot is not transcribed.

% --- p. 180 ---
Men imod en saadan Fordring til den ethiske Bevidsthed stiller sig saa den
\emph{dobbelte} Vanskelighed, at det menneskelige Livs Bevidsthedsbetingelser
medføre det klare Bevidsthedslivs periodiske Afbrydelse ved det relativt og totalt
Ubevidste, og at den menneskelige Erkjendelse, skjøndt ikke hæmmet ved en fast,
ubevægelig Grændse, bevarer en relativ Abstraction. Ved denne dobbelte
Vanskelighed synes den ethiske Livscontinuitet og den Punktualitet i den ethiske
Handling, der betinges ved Erkjendelsens udfoldede Concretion, at blive
forhindret.

Hvad det \emph{sidste} Punkt angaaer, saa ligger der vel i det
Ethisk-Aprioriskes Forhold til de metaphysiske Grundbegreber, i den menneskelige
Naturs Realitetssides Forhold til Væsenets ideelle Enhed, og i Omsætteligheden af
det i Driftens relative Selvstændighed fremtrædende naturligt Reale i ideelle
Væsensmomenter, en Mulighed til en bestandig fremskridende Udvikling til
Concretion af Erkjendelsen af det Ethisk-Aprioriske. Men deri, at
\emph{Erfaringsforholdet} bliver Moment i den Erkjendelse, paa hvilken den ethiske
Handlen hviler, ligger Uundgaaeligheden af, at der i enhver Handlings
Forudsætninger maa blive en partiel Erkjendelsens Ufuldstændighed, en vedblivende,
om end bevægelig, Erkjendelsens Begrændsning. Trods denne Ufuldstændighed og
Begrændsning, der maa indvirke paa Klarheden og Omfanget af den Bevidsthed, ud fra
hvilken der handles, kan dog Handlingen ikke blot i Almindelighed bevare sin
ethiske Sandhed og sin væsentlige Gyldighed som ethisk Handlen, fordi den i
\emph{Principet} er bestemt ved Anerkjendelsen af det Ubetingede, af den absolute
Fordring, ved den ethiske Grundbeslutning, men den kan bevare en
\emph{concretere} Tilsvarenhed til den punktuelle Fordring ved den til den
theoretiske Virkelighedskjendelses paa det Aprioriske begrundede og ligesom
indenfra udvidelige Totalitet svarende \emph{den bevidste Villies Totalitet}, der,
som bestemt ved \emph{det Aprioriske i det Ethiske}, har \emph{real} Gyldighed og
ikke blot indenfra kan udfolde sig til Concretion gjennem Erkjendelsens
fremadskridende Assimilation af det Reale, af Virkelighedsindholdet;
% FOOTNOTE, p. 180: NONE. The page is in the OCR-derived hint list, and the hint
% is wrong: the type block runs the full depth with no rule and no
% smaller-bodied matter at the foot. Checked at 200 dpi over the whole page.
% p. 180 carries no Greek and no rule, and ends in mid-sentence.
% Two ink specks on this page (before „en bestandig“ and between „fremtrædende“
% and „naturligt“) are dirt, not punctuation: at 340 dpi both are round, grey and
% well above the baseline. Not transcribed.
% --- p. 181 ---
men som i sammensluttet Enhed virke i Umiddelbarhedens Form med Genialitetens
Sikkerhed.

Ved denne sidste Bestemmelse vises vi hen til det, der hæver den ovenfor først
berørte Vanskelighed, der laae i det klare Bevidsthedslivs Periodicitet; nemlig
til hvad vi kunde kalde \emph{den ethicerede Naturs Virken}. Ved denne, der har
sit første Udspring fra \emph{Selvets Metamorphose i det ethiske Valg}, vil i det
ved dette Valg constituerede ethiske Liv, trods hiin Periodicitet,
\emph{Totaliteten af det af Væsenets Virksomhed Fremgaaende} kunne have det
\emph{Ethiskes} Charakteer, og det relativt og totalt Ubevidste vil, fra dette
Synspunkt, ligesaalidt ophæve det Ethiske som paa det Theoretiskes Gebet det
uophørlige Vexelforhold mellem Ubevidsthed og Bevidsthed, endog i selve
Bevidsthedsacten i strængere Forstand, ophæver det theoretisk Normale. \emph{I den
ethicerede Naturs Væsensvirken}, der har \emph{Umiddelbarhedens}, men den
\emph{resulterede Umiddelbarheds Form}, ville da \emph{alle} det fremskridende
Livs Acter vare optagne som medvirkende Momenter, efterat være gaaede tilbage fra
\emph{Bevidsthedens Form} til det \emph{Ubevidste}, og i denne Form omsluttede af
Væsenets Enhed.

Denne den ethicerede Naturs Virken, ved hvilken den ethiske Continuitet kan
bevares og Erkjendelsens Ufuldstændighed finde et Supplement, vilde dog endnu,
ligesom hiin Udfoldelse af den ethiske Erkjendelse, lade os blive i det
Approximative og i det ikke i sidste Instants Betryggende, naar vi ikke gik ud
over Individualitetens Grændse. Men vi vises fra den atter tilbage til, hvad der i
Afsnittet om det i Forsoningen Virksomme blev fremhævet: til den fra Begyndelsen
af befrugtende og besjælende, siden opretholdende og klarende Indvirkning, der
udgaaer fra det \emph{sædelige Samfund}, af hvilket Individet aandeligt er
omsluttet, og fra de \emph{ethisk supplerende Individualiteter}. Og definitivt
vises vi tilbage til det, der er det ethiske Samfundslivs og det individuelle
ethiske Livs \emph{absolute Princip}, og det, vi have kaldt den ethicerede Naturs
Virken, see vi da tillige som en Virken af det \emph{universelle Princip},
% RULE, p. 181: the page carries NO rule --- not at its foot, where it ends in
% mid-sentence on „det universelle Princip,“, and not anywhere on it. The type
% block runs the full depth. Checked at 200 dpi.
% EMPHASIS, p. 181: the long run „I den ethicerede Naturs Væsensvirken“ includes
% the opening „I den“, verified at 340 dpi against the tight „der har“ that
% follows it; the „men den“ between „Umiddelbarhedens“ and „resulterede“ is
% tight. „vare optagne“ is printed so („vare“, not „være“) and is left as printed.
% p. 181 carries no footnote and no Greek.
% --- p. 182 ---
til hvilket Selvet træder i Forhold i den ethiske Beslutning og blivende sættes i
Forhold gjennem \emph{Kjærlighedens Selvhengivelse og Væsenstilegnelse}; og i
denne det Absolutes Virken vilde vi da i sidste Instants søge det, der
fuldstændiggjør det Ethiske i Individet.

Ved dette Forhold vises vi da atter tilbage til det tidligere Udviklede angaaende
\emph{det Ethiskes Forhold til det Religiøse}. Idet det Ethiskes Maal bliver den
fulde Virkeliggjørelse af det Menneskevæsen, der som selvbevidst maa blive sig
bevidst efter sit Grundforhold til det Guddommelige, og idet dette Maal kun kan
naaes ved det Guddommeliges Virken, saa kan det Ethiske, i hvis første Tilbliven
det religiøse Forhold var tilstede, ikke opfattes i sin Udvikling uden som overalt
gjennemtrængt af det Religiøse, i Enhed med hvilket det først faaer sin fulde
Virkelighed. \emph{De enkelte ethiske Forhold} blive da, idet de paa
\emph{ethvert} Punkt sees som virkeliggjørende et Personlighedens Forhold til det
Guddommelige, og som blevne til ved dette, at opfatte som \emph{religiøse
Forhold}, og \emph{hele Menneskelivet i dets ethiske Enhed} som \emph{religiøs}
bestemt. I denne Enhed af det Ethiske og det Religiøse har det \emph{Ethiske} sin
\emph{Sandhed}, det \emph{Religiøse} sin \emph{Virkelighed}.

\section*{γ) Det individuelle Menneskelivs til Idealets Begreb svarende Evighedsbestemthed.}

Som den \emph{sidste} Betingelse for Forsoningens fulde Virkeliggjørelse blev det
ovenfor sat, at Individets Liv kan faae den Evighedsbestemthed, der svarer til
Idealets Begreb. Idealet, som Udtryk for det Fuldendte, har, ogsaa naar det sees
under Synspunktet af Virksomhed, Energie --- og saaledes maa Aandens Ideal
opfattes --- det \emph{Eviges} Præg, medens det individuelle Liv stræber hen mod
Idealet gjennem en \emph{Tidsudvikling}, og nærmere gjennem en saadan, der, naar
der sees hen til de organiske Livsbetingelser, viser sig som afsluttet ved en
endelig Tidsgrændse. I denne Indi\-%
% HEAD, p. 182: „γ) Det individuelle Menneskelivs til Idealets Begreb svarende
% Evighedsbestemthed.“ --- bold (halbfette) Fraktur, centred, set over two lines
% („…Begreb“ / „svarende Evighedsbestemthed.“). Verified at 300 dpi. The
% enumerator „γ)“ is ANTIQUA ITALIC against the Fraktur of the head; the Greek
% letter is typed directly. „Idealets“ is an I-word and takes I, per the recorded
% I/J normalisation. The Indhold agrees with the printed head.
% RULE, p. 182: a short centred SINGLE hairline stands ABOVE the display head,
% between „…det Religiøse sin Virkelighed.“ and the head, exactly as on p. 177.
% Not transcribed, per the standing convention for head rules.
% EMPHASIS, p. 182: the closing antithesis is set with TIGHT „sin“ both times ---
% „har det Ethiske sin Sandhed, det Religiøse sin Virkelighed“ spaces only
% „Ethiske“, „Sandhed“, „Religiøse“ and „Virkelighed“, and leaves „det“ and „sin“
% unspaced. Verified at 400 dpi.
% p. 182 carries no footnote and no Greek.
% --- p. 183 ---
videts Tidsbestemthed i dens Forhold til Idealets Evighedsbestemthed ligger
Vanskeligheden.

Til Udjævning af den umiddelbare Incongruents frembyder sig nærmest Tanken om det
tidsbestemte Individs \emph{Udødelighed}, om en fortsat individuel Tilværelse for
det efter Opløsningen af dets reale Organisme. Men naar denne Udødelighed, som en
i det Uendelige fortsat Tilværelse, opfattes saaledes som den opfattes af den
almindelige religiøse Bevidsthed, og saaledes som f. Ex. \emph{Kant} opfattede
den, idet Dualismen mellem Natur og Fornuft, Drift og Frihed, førte ham til at
statuere en uendelig Progres, --- der dog, fra et andet Synspunkt, viste sig for
ham som concentreret i et Evighedens Øieblik, --- saa hjælper den ikke, om den end
tænkes som en fortsat Udvikling, til en Virkeliggjørelse af Idealet, men kun til
en uendelig fortsat Approximation. Man vil naae Evigheden ad \emph{quantitativ}
Vei gjennem en uendelig Tidsudvidelse, og gjennem en Udvikling, der sees under
Tidens Synspunkt, men dette er det Umulige: en saadan fra et bestemt
Begyndelsespunkt uendelig udvidet Tid bliver aldrig Evighedens Correlat, dens
fulde Udfoldelse. Men naar da Læren om en Sjælens Udødelighed som en saadan
uendelig fortsat individuel Tilværelse ikke kan hæve Tvivlen om Idealets
Realisation, saa kunde, fra det Standpunkt, fra hvilket den er fremgaaet,
Eftertrykket lægges paa, at Udødeligheden skulde forestilles som en Væren, der
ikke afficeres ved Endelighedens Existentsbetingelser, ved dens
Udvorteshedssondring og Tidsform, overhovedet som en Væren, der som tilhørende et
\emph{høiere Liv}, ingen Betingelser har tilfælleds med Tidslivets Væren. Herved
vilde dog deels Begrebet om en \emph{Udvikling} falde bort, deels vilde
Individualitetens Væren som \emph{denne} Individualitet, i Bevidsthedscontinuitet
med sin foregaaende Udvikling, blive ubegribelig, og netop paa denne Continuitet
lægges den største Vægt. Men ved hiin Modification vilde imidlertid et Synspunkt
kunne være antydet, der, begrebsmæssigt udviklet, kunde føre til at see Individets
Tidsexistents i en anden Belysning end der, hvor den kun viste sig som en uendelig
Approximation.
% EMPHASIS, p. 183: „Kant“ is letterspaced FRAKTUR, not antiqua --- consistent
% with the standing finding that proper names stay in the Fraktur fount and are
% emphasised by letterspacing instead. Checked at 300 dpi.
% READING, p. 183: after „…tilfælleds med Tidslivets Væren.“ there is NO dash
% before „Herved“, contrary to tesseract, which inserts one. Confirmed at 600 dpi.
% p. 183 carries no footnote, no Greek and no rule. Its last line runs the full
% measure and p. 184 opens without indentation: the paragraph continues across
% the page break, so no blank line stands before the next marker.
% --- p. 184 ---
Der kunde være antydet et Synspunkt, fra hvilket det Evige viser sig som det, der
kan hævde sig som Energie ubehersket af Tidsbestemtheden, som det, der altsaa
ogsaa vil kunne aabenbare sig som det Evige gjennem Tidsbestemtheden, gjøre sig
gjældende under Tidsudviklingen som det, hvad det er efter sit Væsen.

Dette viser os da tilbage til, hvad der tidligere er udviklet om
\emph{Evighedsbetydningen af det ethiske tids}bestemte Liv, og om
\emph{Livsheelhedens Concentration}, under Udviklingen, i et \emph{Evighedens
Moment}. \emph{Det ethiske Liv}, der hviler paa et \emph{Evighedsforhold}, bliver
det, der gjør Mennesketilværelsen til en concret Udfoldelse i Tidsbestemthed af
det Evige. I Livet i Kjærlighed, i den ved Individualitetens Væsensenergie sætte
Selvhengivelse til Principet, leves et Liv med Evighedsbestemthed, der i hvert
Udviklingsmoment concentrerer sig i et ideelt Enhedspunkt, i hvilken
Tidsudviklingen er ideelt taget tilbage i et Evighedens Øieblik. Og idet
Individets Liv i Tiden ophører, idet det som Heelhed sees som sammenfattet i et
Evighedens Moment (jfr. Kant), som concentreret i ideel evig Enhed, saa er denne
individuelle Existents, seet fra det Omfattendes Synspunkt efter sin Betydning som
eiendommeligt Led i det evige Livs Organisation, naaet til sit τέλος, og har
realiseret Evighedsbestemtheden trods Tidsexistentsens Begrændsethed.

Hvad der her er sagt med Hensyn paa et individuelt Liv, seet i Almindelighed som
et \emph{ethisk} Liv, bliver nu, naar den menneskelige Tilværelse overhovedet sees
\emph{fra den universelle Totalitets Synspunkt}, at udvide ikke blot til de
Livsformer, der svagere og ufuldstændigere udtrykke det Ethiske, men ogsaa til de
Livsformer, der synes kun at være udfyldte med et uethisk Indhold og at være
afsluttede eller afbrudte uden noget Gjennembrud af det Ethiske. Naar Synspunktet
tages i den Totalitet, i hvilken de enkelte menneskelige Existentser ere
indbefattede, saa bliver Maalet for den individuelle Existents seet som det, der
\emph{maa} naaes, da \emph{Totalitetens Liv ikke kan} forstyrres, og det, som den
positive Religion udtrykker saaledes, at \emph{ethvert} Væsen naaer
% GREEK, p. 184: τέλος, antiqua italic, tight, in „naaet til sit τέλος“. This
% page is NOT in the recon pass's list of Greek pages, which named only p. 185 in
% this range --- another case of the survey being wrong in both directions.
% Neither OCR witness reads it (tesseract „zé2oç“, ABBYY „rclo?“).
% EMPHASIS, p. 184: the run „Evighedsbetydningen af det ethiske tids-“ stops at
% the line break INSIDE the broken word „tidsbestemte“; „bestemte Liv“ on the next
% line is tight. Transcribed as printed, per the p. 6 / p. 15 precedent, and the
% word is joined. Also: „ethvert“ is spaced but the following „Væsen“ is not
% (checked at 400 dpi at the foot of the page), and „Totalitetens Liv ikke kan“
% is spaced but „forstyrres“ is not.
% READING, p. 184: „Totalitets Synspunkt“ is set as TWO words with a capital S,
% not as the compound „Totalitetssynspunkt“; the word space is narrow because the
% whole phrase is letterspaced. Left as printed. „(jfr. Kant)“ has a lower-case
% „jfr.“ here, unlike the „Jfr.“ of the note on p. 177.
% p. 184 carries no footnote and no rule, and ends in mid-sentence.
% --- p. 185 ---
det Maal, der er det bestemt ved den guddommelige Forudviden, faaer her den Form,
at \emph{ethvert} Individ, som Led i den universelle Tilværelses Organisation, maa
realisere sit τέλος, og i Realisationen indordnes i Evighedens Livstotalitet.

Denne Opfattelse bliver dog ikke at identificere med den, der fra en blot
\emph{Naturnødvendigheds} Synspunkt lader ethvert Individ naae til sit Maal, ɔ:
opfylde sin \emph{Skjæbne}. Individet sees fra dette Synspunkt som i sin
Bestemmethed ved en \emph{ydre} Nødvendighed opnaaende hvad det \emph{kan}
opnaae, og ved denne Opfattelse er Idealets Begreb ophævet. Til Forskjel herfra
lade vi vel det, som Individet \emph{kan} naae, ifølge sin Sammenhæng med
Totaliteten, blive Maalestokken for, hvad det virkeligt naaer; men vi fordobble
Synspunktet saaledes, at det, der \emph{kan} naaes, tillige sees som det, der
\emph{skal} naaes, altsaa som en \emph{Opgave}, der skal realiseres ved
\emph{Frihed}, og saaledes først som et \emph{ethisk Ideal}. Ved dette
\emph{„skal“}, der er \emph{rettet til Villien}, træder altsaa
\emph{Friheds}momentet til Nødvendighedsmomentet og indgaaer i Enhed med dette.
Maalet, der skal naaes, bliver \emph{Idealet}, der realiseres med Frihed, og dette
Ideal er \emph{Individets Væsen efter dets evige Gyldighed}: Gjenstanden for det
ethiske Valg, og det, der realiserer sig gjennem den concrete ethiske Handling.
Men dette Ideal, det \emph{individuelle} Ideal, er tillige, som Kilde til hiint
„skal“, det \emph{almene} Ideal. Dette vil sige: Væsensfordringen til Individet
viser ud over det \emph{Individuelle} til den \emph{Totalitet}, i hvilken
Individet efter sit Grundforhold er ordnet ind, og i hvilken det med Frihed skal
ordne sig ind gjennem Hengivelsen. Saaledes vises der gjennem Frihedsmomentet
atter tilbage til Nødvendigheden; men uden at Friheden forsvinder i denne.

Men naar nu det, at det enkelte Individ, fordi dets Liv er indfattet i
Totaliteten, maa naae sin Bestemmelse, faaer den Betydning, at det med Frihed skal
realisere sit Maal, ved at realisere sin ethiske Opgave, saa bliver det
Spørgsmaalet, hvorledes det, der ovenfor blev sagt om et
% GREEK, p. 185: τέλος, antiqua italic, tight, in „maa realisere sit τέλος“ --- the
% one instance the recon pass predicted for this range. The ABBYY layer drops it
% altogether; tesseract gives „TéAog“.
% EMPHASIS, p. 185: the two occurrences of „skal“ in quotation marks differ, and
% both are left as printed. The FIRST, „Ved dette „skal“, der er rettet til
% Villien“, IS letterspaced inside the quotes; the SECOND, „som Kilde til hiint
% „skal“, det almene Ideal“, is TIGHT. Compared side by side at 600 dpi against
% the tight „dette“ and „hiint“ standing immediately before each. In the first
% sentence the intervening „der er“ is tight and „rettet til Villien“ is spaced
% again. Further splits: „Naturnødvendigheds“ is spaced but „Synspunkt“ is not;
% and of „Frihedsmomentet“ only „Friheds“ is spaced, „momentet“ being tight ---
% a split inside a whole, unbroken word, as on p. 178.
% p. 185 carries no footnote and no rule, and ends in mid-sentence.
% --- p. 186 ---
Liv, der sattes som ethisk uden nærmere Begrændsning, kan faae Anvendelse paa de
forskjellige ufuldkomment ethiske og uethiske Livsformer, og paa Menneskeliv, hvis
Udvikling er afbrudt eller forstyrret.

Skal en saadan Anvendelse blive mulig, maa der fordres et \emph{Gjennembruddets}
og et \emph{Fuldendelsens Øieblik for alle menneskelige Individer}, hvilken end
den tilsyneladende Livsafslutning bliver, og i dette Fuldendelsens Øieblik vil,
hvad enten Gjennembruddet er gaaet forud for det eller falder sammen med det,
Væsenets i det forudgaaede Liv deelte, men bevarede Kraft sluttes sammen til en
udeelt Energie, i hvilken Livsmaalet er naaet. Tidsudstrækningen for dette
Fuldendelsens „Øieblik“ bliver det Uvæsentlige. Hvad psychologiske Erfaringer vise
som muligt i et Øieblik, hvor Livet trues med Tilintetgjørelse: at Individets hele
foregaaende Liv i faa Secunder drages frem for Bevidstheden i klare,
sammenhængende Billeder, ledsagede af en ethisk Dom over det Livsindhold, de
fremstille, kan ogsaa faae sin Anvendelse paa det omhandlede Forhold.
\emph{Dødsøieblikket} som Grændseskjellet mellem to Tilværelsesformer for
Individet, mellem Tidstilværelsens (den phænomenale Livsform) og den evige Værens
Form, vilde da, hvor uendelig lille dets Tidsbestemthed end bliver, kunne blive
Fuldendelsens Tid og give Plads for den gjennemgribende Forvandling. I denne
Fuldendelse vilde da Maalet naaes, hvilken end Villiens Art har været i det
foregaaende Liv; men i den gode Villies forudgaaende Virken vilde
Individualitetens Væsens-Energie have havt direkte, positiv Betydning for
Realisationen af Maalet, i den onde Villies vilde Energien have virket negativt,
men saaledes, at den, gjennem det Dialektiske i Forholdet, indirekte virkede til
Maalets Opnaaelse. Og i den gode som i den onde Villies Virken vilde, som det
Afgjørende og Fuldendende, som det, der i sidste Instants ene kan virkeliggjøre
Idealet, den gjennemvirkende \emph{guddommelige Villie} blive at accentuere.

Men forat denne Anskuelse om et Fuldendelsens Øieblik for Alle ikke skal gjøre
Tidslivets ethiske Betydning
% READING, p. 186: „Individualitetens Væsens-“ / „Energie“ is broken across a
% line, but the second half opens with a CAPITAL E, so the hyphen is a compound
% hyphen and not a line-break artefact; the word is therefore NOT joined as
% „Væsensenergie“ (which is how the same compound is set, whole and lower-case,
% on p. 184) but kept as „Væsens-Energie“. Verified at 340 dpi. This is a genuine
% variation in the printing, not a defect.
% EMPHASIS, p. 186: „Gjennembruddets“ is spaced but the following „og et“ is
% tight; the run resumes at „Fuldendelsens Øieblik for alle menneskelige
% Individer“. „Øieblik“ inside the quotation marks later on the page is NOT
% letterspaced.
% p. 186 carries no footnote, no Greek and no rule, and ends in mid-sentence.
% --- p. 187 ---
illusorisk, saa bliver det \emph{deels} at fastholde, at Forskjellen i Individets
Forhold til den ethiske Væsensfordring gjennem hele Tidsexistentsen maa afspeile
sig i en \emph{væsentlig Bevidsthedsdifferents}, der bestemmer hele Individets
theoretiske og praktiske Forhold; \emph{deels} maa der tænkes en Reflex af
Individets ethiske Forhold i dets Væren som evigt Moment i den totale
Livsorganisation. Dette Sidste vil dog ikke kunne tænkes saaledes, at der i en
Evighedens Existents bliver en forstyrrende Erindring om det Onde, eller en
„metaphysisk Erindring“ om at have kunnet blive til Andet end hvad denne evige
Existents er; --- thi derved vilde den af Totaliteten fordrede Opnaaen af Maalet,
der kun kan være det Godes Maal, være umuliggjort, og den Modsigelse være sat, at
det individuelle Moments Viden af sig i Evighedens Form fik et endeligt Indhold
(Spliden, Syndens Straf). Men det kan, naar der fordres en Evighedsværen for alle
Individer i harmonisk Enhed med det Absolute, kun betegne, at \emph{Intensiteten}
af det Liv, der optages som Moment i den evige Livsfyldes Organisation, udtrykker
\emph{det ethiske Tidslivs Evighedsindhold}.

Fordres der nu herved en fortsat \emph{individuel Tilværelse med
Bevidsthedscontinuitet}? Dette Spørgsmaal har man villet besvare bekræftende idet
man støttede sig paa den \emph{ethiske Forpligtelses Charakteer}. Man slutter:
Individualiteten maa have evig Gyldighed ɔ: Udødelighed som denne bevidste
Individualitet, som denne individuelle Villie, fordi den ethiske Fordring som
ubetinget har evig Gyldighed, og Individet, naar det var forgængeligt, ikke vilde
kunne forpligtes evigt. Men i dette Raisonnement er det overseet, at det, selv om
Individernes bevidste Væren som disse Individer sættes som forgængelig, bliver
Individernes \emph{evige Væsen}, i hvilket Fordringen er begrundet, og at
Fordringen derfor ligefuldt vilde bevare sin \emph{ubetingede} Gyldighed for
Individet som Udtryk for dette Væsen, som Moment i dets Livsproces, og at den
vilde bevare sin Gyldighed for det i ethvert Moment af dets Existents. Væsentligt
bliver det ogsaa at fastholde, at den ethiske For\-%
% READING, p. 187: „Men i dette Raisonnement er det overseet“ --- the faint mark
% between „er“ and „det“ that the ABBYY layer reads as a comma is an ink speck:
% at 500 dpi it is round, grey and sits above the baseline, with none of the
% tail of the true commas elsewhere in the same line. No comma is transcribed.
% EMPHASIS, p. 187: „ubetingede“ is spaced across its own line break but the
% following „Gyldighed“ is tight.
% p. 187 carries no footnote, no Greek and no rule.
% --- p. 188 ---
dring, som en Evighedsfordring til det \emph{tidsbestemte Individ}, bliver en
Fordring, der skal virkeliggjøres i \emph{Tiden} ɔ: i en \emph{bestemt Tid},
forsaavidt den forholder sig til \emph{Individet som Individ}, og at den kun
forholder sig til Tiden overhovedet, og til den Tidens Totalitet sammenfattende
Evighed, idet Individet, medens det er ethisk Selvøiemed, bliver Moment for
Menneskehedens og Tilværelsens Heelhed.

Men naar saaledes den ethiske Fordrings og den ethiske Forpligtelses Art ikke
indeslutter Nødvendigheden af en individuel Udødelighed i \emph{den angivne
Betydning}, saa ligger der dog deri, at Individets Tidsliv ved den faaer en
Evighedsbestemthed i sig, idet Individet forholder sig til en Evighedsopgave, at
der maa sættes en \emph{evig Væsenhed i Individet}. Naar det evige Væsen sees som
Kilde til en Fordring, hvis Opfyldelse faaer Tidens Form og den tidsbestemte
Individualitets Præg, saa kan det Individ, til hvilket det evige Væsens Fordring
er rettet og som virkeliggjør denne Fordring, ikke være et \emph{blot
Forgængeligt}, ikke heller dets Væreform et blot Forsvindende og Betydningsløst.
Dette fører os saa tilbage til Spørgsmaalet om \emph{den Form}, under hvilken den
individuelle Menneskeaands Evighed kan tænkes.

Udgangspunktet for Besvarelsen maa blive \emph{de psychisk-organiske Betingelser
for Individualitetens Existents}. Den reale menneskelige Organisme i sin
Eiendommelighed viser sig, ifølge et Væsensforhold, som den \emph{uundværlige
Betingelse} for Existentsen af den eneste aandelige Individualitet, vi kjende, af
den menneskelige Individualitet i dens eiendommelige Væreform og eiendommelige
Udvikling. Der er altsaa Intet, der berettiger os til at sætte Muligheden af, at
denne individuelle Aandsexistents kan være uden Organismen. Men derfor søge vi dog
ikke Individualitetens \emph{sidste Grund i denne}. Selve Organisationens sidste
Grund kunde jo ikke søges i det Reale, men maatte søges i det Idealitetsprincip,
der, ifølge sit \emph{væsentlige} Forhold til det Reale, sammenfatter de reale
Elementer til Enhed. Og saaledes bliver Individualitetens sidste Grund
% READING, p. 188: „de psychisk-organiske Betingelser“ --- the compound breaks
% across a line here („psychisk=“ / „organiske“), but the hyphen is a real one:
% the same compound is set WITHIN a line on p. 189 („paa Basis af det
% psychisk-organiske“), with a space on either side of the hyphen there, which is
% loose setting and not a suspended hyphen. Kept as „psychisk-organiske“ on both
% pages.
% p. 188 carries no footnote, no Greek and no rule, and ends in mid-sentence.
% --- p. 189 ---
at søge i \emph{Principets Idealitetsbestemthed}, i det sig articulerende Ideelle.
Dette Forhold, ved hvilket der altsaa vises tilbage til et høiere Princip, der
ogsaa begrunder Organismen, indeslutter imidlertid i sin Almindelighed ingen
Garantie for det menneskelige \emph{Individs Udødelighed}. Det evige Princip
aabenbarer sig ogsaa gjennem det individualiserede Selvløse, og dette er dog i sin
Individualitet et Vexlende og Forsvindende. Udødelighedens Garantie maatte da
søges i \emph{Selvheds}bestemmelsen i det menneskelige Individ, og nærmere i
\emph{den Selvhedens Form}, der muliggjør \emph{Friheden}. I Friheden ligger i
Virkeligheden en Garantie for, at Selvet ikke kan tilintetgjøres; i enhver
Frihedsyttring, den være sig positiv eller negativ, aabenbarer Selvet sin
Evighedsbestemthed, viser det sig som det, der ikke kan forsvinde som et
Væsenløst. Men paa den anden Side viser Friheden netop tilbage til det
gjennemvirkende Guddommelige, og den høieste Frihedsact er Selvhengivelsens,
Væsensforeningens med Principet. Spørgsmaalet gjentager sig da om \emph{Formen for
det frie Selvs} Bevaren naar de organiske Individualitetsbetingelser ere
forsvundne, for dets evige Væren efter dette phænomenale Tidslivs Afslutning. Naar
paa den ene Side den Livets Form, der leves i denne Organisme, ikke kan leves naar
den opløses, men paa den anden Side det Selv, der har udfoldet sit Aandsliv paa
Basis af det psychisk-organiske, og som frit har aabenbaret sin
Evighedsbestemthed, ikke kan tilintetgjøres, saa kan det kun tænkes som optaget
\emph{som Moment i det Eviges Liv}, i den Enhed, af hvilken Livsarticulationen
evigt fremgaaer. Men kan der i denne Væren tænkes en \emph{Bevidsthed om en
forudgaaende Udvikling}? Hertil maa svares, at naar Bevidsthedens organiske
Betingelser ere opløste, saa kan Bevidstheden \emph{som saadan} ikke være. Men
hvilken \emph{Metamorphose} kan der da tænkes at foregaae med den? En saadan, at
den individuelle Aand i sin Evighedsconcentration er som et Moment i Guds Viden af
sig i sit evige Aandsrige, i sin evige Livsfyldes Organisation,
% EMPHASIS, p. 189: „Selvhedsbestemmelsen“ is spaced only in „Selvheds“, the
% „bestemmelsen“ being tight, within an unbroken word --- the third instance of
% that split in this batch (cf. pp. 178 and 185). „Formen for det frie Selvs“ is
% spaced but the following „Bevaren“ is not.
% p. 189 carries no footnote, no Greek and no rule, and ends in mid-sentence: the
% batch stops on „i sin evige Livsfyldes Organisation,“ and p. 190 continues it,
% so no blank line stands before the next marker.
% --- p. 190 ---
og --- som det allerede ovenfor blev sagt --- et Moment, hvis Intensitet
udtrykker det ethiske Tidslivs Evighedsindhold, dets evige Betydning i
Livsheelhedens Organisation.

\begin{center}\rule{2cm}{0.4pt}\end{center}
% RULE, p. 190: a short centred rule stands between this paragraph and the next.
% It is not a rule under a display head --- the head follows a further paragraph
% below --- so it is transcribed, as at pp. 19, 28 and 42. Printed as a SINGLE
% thin line (verified at 500 dpi); contrast the DOUBLE rule at the foot of p. 201.

Til klarere Belysning af den i det Foregaaende udviklede Opfattelse af
Forsoningen sammenstille vi den paa den ene Side med den \emph{abstract
metaphysiske} Opfattelse, paa den anden Side med \emph{Christendommens}
Forsoningslære, i hvilken vi see det fuldkomneste Udtryk for den \emph{positive
Religions} Opfattelse.

\begin{center}
1. Den abstract metaphysiske Opfattelse af Forsoningen.
\end{center}
% HEAD, p. 190: „1. Den abstract metaphysiske Opfattelse af Forsoningen.“ ---
% centred on one line in ORDINARY Fraktur: neither halbfette nor letterspaced.
% Verified at 500 dpi against the text lines immediately above and below it and
% against the bold head of p. 192; the strokes are of ordinary text weight and
% the letters are set tight. Hence \begin{center}…\end{center} with neither
% \emph nor \section*. The book's own Indhold reads „1) den abstract
% metaphysiske Opfattelse af Forsoningen“ --- a different enumerator („1)“ for
% the printed „1.“) AND a lower-case „den“ for the printed „Den“. Both
% discrepancies are real and are left standing; the printed head governs.
% EMPHASIS, p. 190: „Christendommens“ is spaced but the „Forsoningslære“ that
% follows it is tight, although the two are one phrase; „positive Religions“ is
% spaced across the line break „posi=“/„tive“ and the following „Opfattelse“ is
% tight. Both verified at 500 dpi.

Som de meest fremtrædende Repræsentanter for denne fremhæve vi \emph{Spinoza} og
\emph{Hegel}. Vi see i denne Sammenhæng bort fra det for den første af dem
Eiendommelige, at Tilværelsen, naar den sees fra de enkelte, begrændsede,
Modificationers Synspunkt, viser sig som en uendelig Kjæde af Aarsager og
Virkninger, behersket af den blinde Naturnødvendigheds Lov, saa at det Ethiske
maa forsvinde i det Physiske, Ethiken blive en Handlingens Naturlære. Vi lægge
derimod Vægten paa den \emph{væsentlige Lighed} mellem de to Tænkeres
Opfattelse, der ligger i den for begge fælleds \emph{abstracte
Evighedsopfattelse}, som har sin Kilde i en af Dualismen ubevidst paavirket
ensidig Idealisme. Denne Lighed, der bliver tydelig, naar hos \emph{Spinoza}
Tilværelsen sees fra \emph{Substantsens}, det Eviges og Uendeliges Synspunkt,
forstyrres ikke ved den tilsyneladende Ligeberettigelse, det Naturlige, Reale, i
Substantsen faaer med det Aandelige og Ideelle, og som kunde vise hen paa et
Correctiv for den abstracte Evighedsopfattelse, og paa en væsentlig Betydning af
det Reales Grundformer for det ethiske Liv. I Virkeligheden faaer nemlig det
\emph{ideelle Mo}ment, der oprindelig dannede Udgangspunktet for Opfattelsen af
% NAMES, p. 190: „Spinoza“ and „Hegel“ are letterspaced Fraktur, not antiqua,
% and the „og“ between them is tight --- hence two \emph groups, not one.
% „Spinoza“ is letterspaced again lower on the page, there broken „Spi=“/„noza“
% with both halves spaced.
% EMPHASIS, p. 190: „ideelle Mo“ is spaced and the „ment“ that completes the
% word on the next line is tight --- the compositor's split-across-the-break
% inconsistency again, transcribed as printed.
Substantsen og var Substantsbegrebets Kilde, Overvæg\-%
% --- p. 191 ---
\setcounter{footnote}{0}%
ten, og ligesom det bliver det overgribende i Forhold til de andre Attributer,
bestemmer det \emph{definitivt} Synspunktet for Tilværelsen som absorberet i
Substantsens Evighed.

\emph{De Consequentser}, som den \emph{abstracte Evighedsopfattelse} og den
\emph{ensidige Idealisme}, af hvilken den fremgaaer, medfører, ere væsentlig
følgende.

Ud fra hiin Opfattelse kan Intet, der foregaaer i Tiden, faae Betydning ɔ:
\emph{væsentlig Betydning}, fordi Tiden kun bliver en \emph{usand}
Existentsform, og ikke staaer i et væsentligt Forhold til det Absolute, eller
til det Tilværende, seet efter dets Sandhed; thi saaledes sees det kun
\textit{sub specie æterni}. Kun naar \emph{Tiden}, som nødvendig Existentsform
for det \emph{væsensgyldige} Reale, anerkjendes som havende \emph{væsentlig}
Betydning for \emph{alt} Værende, kan ogsaa den \emph{ethiske} Existents, som en
Existents i Tiden, der forholder sig til en Evighedsopgave, blive opfattet efter
sin \emph{Virkelighedsbetydning}. Først da kan Tidsforholdet til det Ethiske,
det ethiske Valgs Tidsbestemthed, Idealets Realisation gjennem en Tidstotalitet,
Tabet af Tiden som begrundende en Skyld, og Gjenoprettelsen i Tiden af denne
Skyld faae afgjørende Betydning. Først da kan det blive Opgaven, at see
Evighedsbestemmelsen i det ethiske Individ som realiserende sig i Concretion i
Tidsbestemthed.
% LATIN, p. 191: „sub specie æterni“ is antiqua and TIGHT here (verified at 500
% dpi) --- contrast p. 192, where the same phrase occurs twice on one page, the
% first time in SPACED antiqua and the second time tight. Encoded \textit{…}
% here and \emph{\textit{…}} for the spaced instance.

Ud fra hiin Opfattelse bliver det umuligt at anerkjende \emph{Friheden} i det
Individuelle i dets relative Sondring fra det Absolute --- en Sondring, der
netop er betinget \emph{ved Realitetsbestemmelsernes Væsentlighed for
Principet}. Alt maa drages ind under det Universelles og Eviges metaphysiske
Nødvendighedsproces, i hvilken det Individuelle maa blive et kun Magtesløst og
Forsvindende\footnote{Naar saaledes \emph{Strauß} --- i sin Dogmatik --- søger
at skaffe en Plads for Friheden, for det Onde og for en Forsoning, saa bliver
Bevægelsen rent illusorisk, idet Grundsynspunktet er hiin ensidige Idealismes
abstracte Evigheds- og Nødvendighedsopfattelse.}. Kun naar Realitetsformen
tillægges væsentlig
% FOOTNOTE, p. 191: printed mark *) --- the only note on the page, set below a
% short single rule. „Strauß“ is letterspaced Fraktur (with the ß as printed);
% „Dogmatik“ is tight. „Evigheds- og Nødvendighedsopfattelse“ is a suspended
% compound and keeps its printed hyphen and space.
% NOTE: the OCR-derived hint list gave p. 192 as note-bearing; it is p. 191 that
% carries the note, and p. 192 carries none. Checked by eye on both pages.
% --- p. 192 ---
Gyldighed, og det Reales Betydning for Frihedens Idealitetsbestemmelse indsees,
kan den individuelle Frihed hævdes.

Ud fra hiin Opfattelse kan altsaa \emph{Forsoningen} kun sees som bestemt ved
den \emph{evige} Processes \emph{Nødvendighed}, og den maa sættes som realiseret
i den individuelle \emph{Tankes} Opløftelse over det Endelige til en
\emph{Tænken} af \emph{det Absolute} (Evige) og af sig i det Absolute. I denne
Tænken gaaer Individets \emph{Tanke}, der er dets \emph{Væsens} Udtryk, som
bestemt af det Absolute, op i dette som ophævet Moment, og Tidsbetingelserne og
Tidssondringen blive betydningsløse, uvæsentlige, hvad enten Forholdet opfattes
paa Spinozistisk eller paa Hegelsk Viis. Idet det Væsentlige udelukkende
forholder sig til Evighedsbestemmelsen, der udtrykker den \emph{ene sande
Væren}, og idet den \emph{sande Erkjendelse}, der er Forsoningens \emph{Form},
er \emph{\textit{sub specie æterni}}, saa er der væsentlig \emph{Intet} forsømt
og derfor i Ordets strængeste Betydning ikke Noget at restituere
(redintegrere). I hiin sande, \textit{sub specie æterni} værende, Erkjendelse
forsvinder, idet Tidsforholdet taber sin Betydning, det Livsindhold, der er
udfoldet i en \emph{Tid forud for den sande} Erkjendelse, som et \emph{Skin} med
den forsvindende Tidsform; det ophæves i Evighedens evig nærværende
Virkelighed, i hvilken end ikke en Reflex kan bevares af det endelige
individuelle Liv, der, som \emph{kun} baaret af en Naturdifferents, \emph{kun}
er baaret af et Forsvindende, et Væsenløst.
% LATIN, p. 192: the phrase „sub specie æterni“ stands twice on this page. At
% its first occurrence („er sub specie æterni, saa er der…“) it is antiqua AND
% letterspaced; nine lines later („I hiin sande, sub specie æterni værende“) it
% is tight antiqua. Both verified at 500 dpi on the same page --- the clearest
% instance in the book so far of the antiqua's spacing varying within one page.
% EMPHASIS, p. 192: „Tid forud for den sande“ is spaced and the „Erkjendelse“
% that completes the phrase on the next line is tight (verified at 500 dpi, at
% 200 % magnification: „i en“ is tight, „Tid forud for den sande“ spaced). The
% two occurrences of „kun“ in the last sentence are spaced; the words between
% and after them are tight.

\section*{2. Den christelige Forsoningslære.}
% HEAD, p. 192: „2. Den christelige Forsoningslære.“ --- BOLD (halbfette)
% Fraktur, centred, on one line, verified at 500 dpi; the strokes are plainly
% heavier than the surrounding text and than the head of p. 190 two pages
% earlier. Hence \section*{} here and \begin{center}…\end{center} there. The
% book's own Indhold reads „2) den christelige Forsoningslære“ --- again a
% different enumerator and a lower-case „den“. Left standing; the printed head
% governs.

At den \emph{christelige} Forsoningslære vælges som \emph{Repræsentant} for de
\emph{positive} Religioners Opfattelse af Forsoningen, har sin Berettigelse
deri, at i Christendommen Bevidstheden om Synden, om Spliden i Menneskeaanden og
om dennes Sondring fra det Guddommelige, har naaet en Intensitet og Alvor som i
ingen forudgaaende Religion, og at derfor Læren om Forsoningen i den kunde faae
den rigeste Udvikling og give det skarpe Udtryk for hvad der i de andre
Religioner kun uklart og usikkert antydedes.

% --- p. 193 ---
Der bliver i nærværende Sammenhæng først at tage Hensyn til Christendommens
Opfattelse af Synden og dens Virkning, og derefter til dens Opfattelse af
Forsoningen og af Ophævelsen af Syndens Følger.

\begin{center}
\emph{a) Opfattelsen af Synden og af dens Virkning.}
\end{center}
% HEAD, p. 193: „a) Opfattelsen af Synden og af dens Virkning.“ --- centred on
% its own line, letterspaced regular Fraktur, NOT bold; same treatment as the
% run of centred letterspaced heads at pp. 165 and 167. This head is NOT in the
% table of heads in BATCH-AGENT.md, and neither is its fellow at p. 195; both
% are real and both are set here. Neither appears in the book's own Indhold
% either --- the Indhold jumps straight from „2) den christelige
% Forsoningslære“ (192--201) to „Slutning“ --- so these two are body-only
% heads, with no Indhold reading to compare against.

Afgjørende for Opfattelsen bliver her strax Bestemmelsen af Synden som
\textit{ἀνομία}, (1. Joh. III, 4) som Overtrædelse af Guds Lov, som Ulydighed
mod Guds Bud. Idet nemlig det guddommelige Bud, der i Synden overtrædes,
opfattes som en vilkaarlig guddommelig Villies Bud, saa maa Buddets Indhold vise
sig som det Tilfældige og i dybere Forstand Ligegyldige, og Eftertrykket maa
falde paa Fordringen af en blind Lydighed mod den vilkaarlige Villie (jfr.
Augustin). Overtrædelsen af denne Fordring og af det vilkaarlige Bud sættes da
ikke i Forhold til Menneskets og nærmere til den menneskelige Villies
\emph{Væsens}bestemmelse; kun paa det \emph{Formelle}: paa Uoverensstemmelsen
med den bydende guddommelige Villie, falder Eftertrykket. Og idet denne Villie
staaer som den, der ved Transscendentsen er sondret i Væsen fra den
menneskelige, og idet dens vilkaarlige Bestemmelse, der er udtalt i Buddet, ikke
kan udtrykke en Væsensbestemmelse ved den menneskelige \emph{Villie}, saa kan
Overtrædelsen, som ikke berørende et Væsensforhold, ikke begribes at have en
væsensbestemmende Virkning paa Menneskevillien og Menneskenaturen. Naar derfor
den christelige Lære statuerer en saadan Virkning, saa gaaer den ind under
Bestemmelsen af det Paradoxe.
% GREEK, p. 193: „ἀνομία“, antiqua, the only Greek in this batch (every page of
% pp. 190--201 was read by eye for it). The mark over the initial α is printed
% as a straight slanting stroke, at this scan indistinguishable in shape from
% the acute over the ι; „ἄνομία“ is not a word, so the smooth breathing is set,
% per playbook §2 rule 3. The comma after the Greek is as printed (the book sets
% „som ἀνομία, (1. Joh. III, 4) som Overtrædelse“). „Joh.“ takes J, the word
% being Johannes; „Indhold“, „Idet“, „Individ“ take I throughout.
% EMPHASIS, p. 193: „Væsens“ is spaced and the „bestemmelse“ that completes the
% word is tight --- the same split as at pp. 178, 185 and 189.

Det \emph{Paradoxe} træder allerede frem, hvor Synden, som et Tilfældigt, første
Gang bringes i Forhold til det Menneskelige. Forud for den \emph{første Synd},
der skulde faae afgjørende Betydning for Menneskeslægtens Heelhed, forestilles
nemlig en Urtilstand, i hvilken det endnu syndfrie første Menneske, der var
skabt i Gudsbilledet, havde en intellectuel og sædelig Fuldkommenhed, der gjør
Faldet til et \emph{Ubegribeligt}. Og det \emph{Paradoxe} fortsætter sig saa,
idet i den christelige Kirke, under Kirkelærens Udvikling, den
% --- p. 194 ---
første Synds Virkning bestemmes som den menneskelige Naturs og den menneskelige
Villies Fordærvelse ved Synden --- en Virkning, der maatte sees som en Straf,
paalagt af den guddommelige Vrede. Og det, at Menneskenaturen og
Menneskevillien fordærves ved den første Synd --- det første Menneskes Synd ---
bliver, i Forhold til de \emph{andre In}divider i Slægten, saaledes opfattet, at
Fordærvelsen i Naturen forplanter sig til Menneskeslægtens Heelhed, udbreder sig
fra Slægt til Slægt som Arvesynd, og at Overtrædelsen tilregnes Slægterne som
Arveskyld efter Stamfaderen. Arvesynden bliver den første Synds \emph{Straf}, og
Skylden, der hviler paa Slægten, maa stadig forøges ved Slægtens nye Brøde. Den
Virkning af Synden, der statueres, gaaer ind under det \emph{Paradoxe} saavel
naar den betragtes i Forhold til det Individ, der forestilles som Syndens
Begynder, som naar den sees i Forhold til de andre Individer. Det er et
\emph{Paradox}, at Virkeliggjørelsen af en i Menneskevæsenet nedlagt Mulighed
skal kunne omskabe selve Væsenet, at Villiesyttringen, der er en Væsensyttring,
skal kunne medføre Menneskevæsenets ved Mennesket uoprettelige Fordærvelse. Det
er et \emph{Paradox}, at Synden, den individuelle frie Villies Gjerning, og den
Skyld, den medfører, skal overføres paa de endnu villieløse, kun som Muligheder
existerende Individer, og at den skal forudbestemme deres Handling idet hiin
Slægtens Stamfaders Synd fordærver Slægtens Natur saaledes, at Synden for de
efterfølgende Individer bliver det Uundgaaelige. Og det \emph{Paradoxe} gjør sig
endelig ogsaa gjældende paa den Maade, at, under Forudsætningen af Arvesynden,
af Menneskevæsenets Forstyrrelse, \emph{Bevidstheden om Synden} consequent maa
sættes som manglende i Mennesket, saaledes at den kun kan vækkes ved en
guddommelig \emph{Aabenbaring}; men naar overhovedet Aabenbaringen, som
\emph{Under}, viser sig at høre ind under Paradoxet, saa er det særligt
iønefaldende, at en \emph{saadan} Aabenbaring om Synden til de i Synden
existerende, ved Synden forvandlede Mennesker maa gjøre det.
% p. 194: „deres Handling idet hiin Slægtens Stamfaders Synd“ is printed with no
% punctuation between „Handling“ and „idet“ --- „Handling“ ends a line and „idet“
% opens the next, and both witnesses agree with the image. Transcribed as printed.
% EMPHASIS, p. 194: „andre In“ is spaced and the „divider“ that completes the
% word is tight.

Naar den \emph{positive Religions} Bestemmelser af \emph{Syn\-%
% --- p. 195 ---
dens Begreb} og \emph{Syndens Virkninger} sammenholdes med den \emph{rationelle
Ethiks}, viser Modsætningen sig overalt. Den viser sig i Opfattelsen af Syndens
Væsen, idet Synden for den \emph{positive Religion} bliver Overtrædelsen (eller
Tilbøieligheden til Overtrædelse) af et vilkaarligt guddommeligt Bud, formel
Ulydighed, Brud paa Afhængighedsfølelsen af Almagten; for den \emph{rationelle
Ethik} derimod en Krænkelse af Væsenets Fordring, og netop derfor af det
Absolutes gjennm Menneskevæsenet udtrykte Fordring. Den viser sig i Opfattelsen
af \emph{Syndens Straf}, der for den \emph{positive Religion} bliver bestemt ved
den guddommelige Villies \emph{Vilkaarlighed}, for den \emph{rationelle Ethik}
ved Væsenets Lov, ved Consequentsen af Egoismens Isolation. Den viser sig i
Opfattelsen af \emph{Væsenets Forstyrrelse} ved Synden, der for den
\emph{positive Religion} bliver en ubegribelig \emph{Væsensforvandling}, for den
\emph{rationelle Ethik} en Splid i Væsenet, en Suspension af dets Integritet,
under hvilken Væsenet er tilstede som negativt virkende, som bevarende
Muligheden til Redintegrationen. Den viser sig endelig i
\emph{Syndsbevidstheden}, der efter den \emph{religiøse} Opfattelse consequent
maa vækkes ved en overnaturlig Indvirkning, efter den \emph{rationelle Ethik}
naturligt fremgaaer af Væsenets negative Virken i den onde Villie. Under
\emph{Modsætningen} mod den \emph{rationelle Ethiks} Bestemmelser maae den
\emph{positive Religions alle} gaae ind under det \emph{Antirationelle}, det
\emph{Paradoxe}.
% PRINTER'S DEFECT, p. 195, line 9: „det Absolutes gjennm Menneskevæsenet
% udtrykte Fordring“ --- the page prints „gjennm“, a dropped „e“ in „gjennem“.
% Confirmed at 600 dpi: the six sorts g-j-e-n-n-m stand with no gap and no
% trace of a seventh. Both OCR witnesses read „gjennm“ independently.
% Transcribed as printed; it does not affect the quote balance.

\begin{center}
\emph{b) Opfattelsen af Forsoningen og af Ophævelsen af Syndens Virkninger.}
\end{center}
% HEAD, p. 195: „b) Opfattelsen af Forsoningen og af Ophævelsen af Syndens
% Virkninger.“ --- centred, letterspaced regular Fraktur, NOT bold, set on TWO
% printed lines with the turn after „Ophævelsen“. The turn is not forced here.
% Like the head of p. 193, it is absent from the book's own Indhold.
% Like the head of p. 193 this one is absent from BATCH-AGENT.md's table.

Idet \emph{Synden} og \emph{Syndens Straf} for den \emph{positive Religion}
sammenknyttes ved Forestillingen om den \emph{guddommelige Villie},
\emph{ikke} ved et \emph{Væsensforhold}, bliver Ophævelsen af det ved
\emph{Synden Satte} at see fra et \emph{dobbelt} Synspunkt: som Ophævelse af
\emph{Syndens Virkninger} (\emph{Straffen}), og som Ophævelse af \emph{Syndens
actualiserede Princip}, af den syndige Villiesbestemthed, og disse to
Synspunkter mødes kun forsaavidt, som \emph{Syn\-%
% --- p. 196 ---
den selv}, opfattet som \emph{Arvesynd}, bliver \emph{Syndens Straf}.

Idet \emph{Straffen} er sat ved den \emph{guddommelige Villie}, bliver, ifølge
denne Villies Natur, den \emph{første Betingelse} for \emph{Straffens}
Ophævelse den \emph{objective}: \emph{Guds Tilgivelse}. I Opfattelsen af denne
viser den sig i den christelige Lære en Dobbelthed. Paa den ene Side forestilles
den guddommelige Tilgivelse, som svarende til den vilkaarlige Bestemmen af
Straffen, som \emph{vilkaarlig}, som \emph{uafhængig} af bestemte
Forsoningsbetingelser; paa den anden Side accentueres den guddommelige
\emph{Retfærdighed}, og denne Retfærdighed maa fordre en \emph{adæquat
Erstatning}. Men idet hine to modsatte Synsmaader staae imod hinanden, viser
\emph{hiin} Forestilling om den guddommelige \emph{vilkaarlige Villie} sig som
den \emph{overgribende}, og den bestemmer den modsatte Synsmaade saaledes, at
den guddommelige Retfærdighed tænkes at \emph{lade sig} tilfredsstille ved en
Erstatning, der kun ved den guddommelige Villie bliver \emph{taget for} en
Erstatning. Vilkaarligheden i Valget af Forsoningsmidlet fremhæves ikke blot
stærkt af Scholastikerne, men i de forskjellige Nuancer, Satisfactionstheorien
modtager under Kirkelærens Udvikling, bliver gjennemgaaende det
Vilkaarlighedens Moment, der udtrykkelig anerkjendes i
\emph{Acceptilationstheorien}, bevaret. Forsoningsmidlet, ved hvilket den
guddommelige Villie \emph{lader sig} tilfredsstille, er \emph{den
menneskeblevne Guds Selvhengivelse, Christi Død}; ved denne bringes Forløsning
og Forsoning, og Forsoningen er Menneskevordelsens Øiemed.

Til den \emph{objective} Betingelse for Ophævelsen af Syndestraffen træder saa
den \emph{subjective}: \emph{Troen}, som Tilegnelse af den \emph{objective
Fyldestgjørelse}. Troen som den \emph{retfærdiggjørende Tro} er Troen paa, at
Christi \emph{Død} er fyldestgjørende for \emph{dette} troende Individ. Ved
denne Tro, ikke ved Gjerninger, der, for at have Sandhed, skulle have Troen til
Forudsætning, er det, at Forsoningen vindes.

Ved Troen ophæves den \emph{Syndens Virkning}, der
% --- p. 197 ---
bestaaer i den guddommelige Vrede, og Udslettelsen af Skylden vindes. Men idet
Udslettelsen af Skylden og Eftergivelsen af Straffen indeslutter Ophævelsen af
den Straf, der er \emph{selve Synden}, saa bliver i Troen ogsaa Synden som
\emph{Villiesbestemthed}, som Villiesmagt, seet som ophævet i Principet, og
Begyndelsen til en Handlen i Overensstemmelse med den guddommelige Villie, til
\emph{Helliggjørelse}, som given; og Gjenfødelsens principielle Forvandling
faaer Livsvirkelighed gjennem en \emph{mystisk Forening} med Gud. Her mødes
saaledes de to ovenfor sondrede Synspunkter; men de mødes kun for atter at
skilles.

At de atter skilles, har, bortseet derfra, at Syndens Straf oprindelig er sat
som \emph{vilkaarlig} bestemt, sin Grund i \emph{Troens Forhold til det
Menneskelige} og i dens \emph{Indhold}. Idet Troen, som det vil blive paaviist,
\emph{ikke} har sin \emph{Oprindelse} i det \emph{Menneskelige}, i det Væsen, af
hvilket den syndige Act er fremgaaet, og hvis sande Omdannelse kun kan tænkes
som foregaaende ud fra dets eget Centrum, saa maa den beholde et
Udvorteshedsforhold til det menneskelige Indre, der gjør Omdannelsen
uforklarlig. Og til det samme Resultat føres vi fra det Synspunkt, at Troen, som
retfærdiggjørende Tro, med Hensyn til sit \emph{Indhold}, og derved med Hensyn
til sit \emph{Væsen}, ikke bliver saaledes bestemt, at den kan blive et
virkeligt \emph{Livsprincip}, hvorfor den atter maa komme i et
Udvorteshedsforhold til den sande Handling.

Denne Eiendommelighed ved Troen i Religionens Opfattelse staaer i nøieste
Sammenhæng med, at Religionen paa \emph{dualistisk} Maade lægger \emph{alt Godt}
over i \emph{Gud} og berøver \emph{Mennesket} det, idet den kun i en imaginær
Fortid lader Mennesket have besiddet det Gode som en Gave. Den vil blive
tydelig, naar Forholdet mellem den guddommelige og den menneskelige Virksomhed i
Forsoningen nærmere betragtes.

Naar Menneskenaturen er sat som fordærvet, Gudsbilledet som tabt og Villien som
bunden til Syndens Nødvendighed, saa kan der, consequent, ikke tænkes nogen
Virk\-%
% --- p. 198 ---
somhed som udgaaende fra Menneskenaturen og Menneskevillien, der kan bidrage til
Forsoningen. Al saadan Virksomhed maa da udgaae fra \emph{Gud}; og ligesom
Forsoningens \emph{objective} Betingelse gaves af ham alene, idet han blev
Menneske, døde for Menneskenes Synder og modtog denne Død som Fyldestgjørelse,
saaledes maa ogsaa den \emph{subjective} Betingelse: \emph{Troen}, gives af
\emph{Gud} ved en \emph{Naadesact} af den guddommelige Villie, og gives
\emph{som en Afspeiling af den vilkaarlige Tilgivelse}, som en \emph{Tro} paa,
at \emph{dette} troende Individ er forsonet med Gud. Naar da nogle Mennesker
blive troende, er det ikke ifølge egen Fortjeneste, men ifølge den guddommelige
Naades Valg; naar andre vedblive som Ikke-Troende, altsaa som uforsonede, og
derfor som hjemfaldne til Fordømmelsen, saa er det fordi den guddommelige Naade
unddrages dem, fordi de overlades til deres egne Kræfter. I denne Forestilling
om Naadevalget, der er den \emph{uundgaaelige} Consequents af den christelige
Opfattelse af Grundforholdet, træder Opfattelsen af Guds Villie som den
\emph{absolute Vilkaarlighed} frem i sin største Skarphed, og Modsætningen
mellem den positive Religion og den rationelle Ethik bliver den størst mulige.
Med Rette kalder \emph{Calvin} hiin guddommelige, for Menneskets Tanke
ubegribelige, Beslutning til Saliggjørelse og Fordømmelse
\emph{\textit{et horribile decretum}}, og det klinger som en frygtelig Haan,
naar \emph{Luther} med Hensyn til de ikke Udvalgte fortolker Evangeliets Ord:
„gjører dette, og I skulle blive salige!“ saaledes: „prøver kun derpaa, og I
skulle see, at det ikke er muligt!“, og naar han fra Guds \emph{aabenbarede
Raadslutning} til \emph{alle} Menneskers Salighed, ligesom med en ubevidst
Ironie over \emph{Aabenbaringen}, adskiller en \emph{forborgen Raadslutning} til
Fordømmelse. Og dog ere Calvin og Luther, naar de drage de forfærdeligste
Consequentser, i Overensstemmelse med de Forudsætninger, Religionen giver og,
efter sin Grundsynsmaade, maa give.
% ANTIQUA, p. 198: „et horribile decretum“ is set in antiqua AND letterspaced,
% and the antiqua run begins at the DANISH article „et“ --- verified at 600 dpi
% and 300 % magnification: the „e“ of „et“ is the round roman sort, unlike the
% angular Fraktur „e“ of „melse“ immediately before it, and the „t“ is roman
% too. The article is therefore inside the \textit group, as printed. Contrast
% p. 199, where „hiint horribile decretum“ has „hiint“ in Fraktur and the Latin
% in TIGHT antiqua.
% QUOTES, p. 198: two „…“ pairs, both balanced. The closing mark of the second
% pair stands before its comma, as printed.

I den christelige Lære om Syndestraffens Ophævelse møde vi, ligesom i Læren om
Syndens Væsen og Følger, overalt det \emph{Paradoxe}. Vi møde det i
Forestillingen om \emph{Men\-%
% --- p. 199 ---
neskevordelsen}, om den evige, uendelige Guds Vorden og Leven i den individuelle
Menneskenaturs Begrændsning, under et menneskeligt Livs Tidsudvikling. Vi møde
det i Læren om \emph{Satisfactionen}: i Forudsætningen af Nødvendigheden af en
Fyldestgjørelse af den guddommelige Villie ved en guddommelig Gjerning, og af en
Mulighed af denne Gjernings Antagelse som fyldestgjørende. Vi møde det i
Forestillingen om en \emph{Overførelse} af Retfærdigheden, af det, der, som
Bestemmelse ved den frie Villie, kun kan tænkes som erhvervet ved den frie
Villies egen Virken, medens nu Mennesket forestilles at forholde sig rent
passivt i Forsoningen. Vi møde det endelig i \emph{Vilkaarligheden i den
guddommelige Udvælgelse}, i den al ethisk Kraft tilintetgjørende Forestilling om
en absolut grundløs Bestemmen til Salighed og til Fordømmelse.

Sammenligne vi Forsoningen, saaledes som den bliver Gjenstand for Bevidstheden
paa den \emph{positive Religions} Gebet og indenfor det \emph{rationelt
Ethiskes} Sphære, saa viser der sig for det positivt Religiøses Vedkommende en
Mangel saavel med Hensyn til Forsoningens \emph{Vished} som til dens
\emph{Fuldstændighed}. Hvad \emph{Visheden} angaaer, saa er det indlysende, at
saalænge hiint \textit{horribile decretum} staaer som Baggrund for den religiøse
Bevidsthed, saa kan ingen Hvile findes i en Tro paa den tilveiebragte
Forsoning. Hvad der træder frem for Bevidstheden som Naadevirkninger, kan ikke
med Vished erkjendes som saadanne: Muligheden af en Skuffelse bliver bestandig
tilbage, netop fordi disse Virkninger tænkes udgaaende fra et \emph{andet
Væsen}, der ved \emph{Transscendentsen} er sondret fra det Menneskelige. Og den
Villie, der grundløs udvælger og grundløs fordømmer, kan i hvert Øieblik tage
sin skuffende Naadevirkning tilbage, lade Mennesket synke ned i Fordømmelsen.
Saaledes holdes Angesten bestandig vaagen; der gives ingen Hvile, ingen Fred.
--- Og hvad Forsoningens \emph{Fuldstændighed} angaaer, saa viser Manglen i den
positive Religions Opfattelse sig \emph{først} deri, at Ophævelsen af
\emph{Syndens Straf} atter skilles fra Ophævelsen af \emph{Syndens Grund} i
Individet. Den
% --- p. 200 ---
\emph{Restitution}, der, efter den positive Religions Lære, saaledes som den er
udviklet hos Kirkens største Repræsentanter, i Virkeligheden finder Sted, er den
\emph{vilkaarlig paalagte Strafs Eftergivelse}, altsaa den \emph{forestillede}
Restitution gjennem Forsoningen med den \emph{vilkaarlige} guddommelige Villie,
ifølge den \emph{vilkaarlige} Tilgivelse af Overtrædelsen af det
\emph{vilkaarlige} Bud; men Restitutionen er \emph{ikke Redintegrationen af
Væsenet og Villien}, \emph{ikke} den \emph{virkelige} Restitution,
Restitutionen af Kraften til det Gode, af den Tilstand, der forestilledes som de
første Menneskers. Man sammenligne \emph{Augustins}, \emph{Luthers} og
\emph{Calvins} Udsagn, og man vil finde, at de, trods Menneskets Gjenfødelse som
troende og trods den Troendes mystiske Forening med Gud, sætte den menneskelige
Villie som bestandig hendragen til Synden, som bestandig hildet i det Onde, og
at de kun haabe den sande Helliggjørelse i en \emph{tilkommende} Tilstand, i et
andet Liv, i hvilket Menneskelivets reale Betingelser ere ophævede. Og de gjøre
dette consequent ifølge den ovenfor berørte dualistiske Sondring mellem Gud og
Menneske med Hensyn til det Gode. Paa Grund af den i den \emph{abstracte
Transscendents} liggende Sondring kan Forsoningen med Gud paa den positive
Religions Gebet ikke blive den \emph{fulde} Forsoning i Bevidstheden om
\emph{Væsensenheden} med Principet. \emph{Til denne} Forsoningens Form vises der
kun hen gjennem det \emph{mystiske} Element i Religionen; men
\emph{Reflexionen}, der vaagner i Religionen, trænger det uundgaaelig atter
tilbage ved at give Transscendentsens Sondring Vægt. --- Sees Forsoningen som
Individets Forsoning i \emph{sig selv}, saa viser dens
\emph{Ufuldstændighed} sig ikke blot gjennem den i det Onde hildede Villies Kamp
med den bedre Erkjendelse, men ogsaa i det tidligere Fremhævede, at den
religiøse Bevidsthed, der har sit Støttepunkt i det for Fornuften
Utilgængelige, kun kan leve sit Liv ved at negere det Rationelle, altsaa ved at
negere en Aandens Væsensyttring, og derfor maa bevare en Splid i sit Indre.

Paa det \emph{rationelt Ethiskes} Standpunkt er Forso\-%
% --- p. 201 ---
\setcounter{footnote}{0}%
ningens \emph{Vished} given i Aandens Bevidsthed om sin \emph{Væsensenhed} med
det \emph{iboende} Princip, i Erkjendelsen af den forsonende Kraft som
\emph{Aandens Væsensyttring i dens Enhed med det gjennemvirkende Absolute}. Og
Forsoningens \emph{Fuldstændighed} er muliggjort derved, at Forsoningen ikke
blot bliver Individets \emph{principielle Forsoning med sit Ophav} og i
\emph{sit eget Indre}, men at der, som der i de tidligere Udviklinger er
paaviist, trods det Brud paa det ethiske Livs Udvikling, der foregaaer ved den
onde Handling, og trods Villiens Vaklen og Tilbagefald, er bevaret en Mulighed
til at tilveiebringe en \emph{Redintegration}, at ophæve Discontinuiteten i
Continuitet, at gjøre det Ethiske til Totalitetsbestemmelse for det menneskelige
Liv, og saaledes at realisere den ideale ethiske Fordring\footnote{Baade i den
kirkelige Opfattelse og i de sædvanlige Angreb paa den mangler --- selv hos
\emph{Anselm}, der har været det nærmest --- det Synspunkt, at
Virkeliggjørelsen af Opgaven, af den ideale Fordring, er betinget ved, at en
Continuerlighed tilveiebringes i den ethiske Stræben, eller at det i Tiden
Forsømte atter kan indhentes i Tiden, hvilket synes at umuliggjøres derved, at
ethvert Moment i Tiden har sin Opgave, og derfor ikke synes at kunne afgive
Plads for et tidligere ikke udfyldt eller i negativ Retning udfyldt Moments
positive Udfyldelse og Indordning i Heelheden. Saaledes bliver der f. Ex. hos
\emph{Strauß}, --- hos hvem der kun er Tale om en abstract Nødvendighed i det
enkelte Individs Udvikling, idet hvert enkelt bliver til hvad Heelhedens
Udviklingsnødvendighed fører med sig, hvilket kun opfattes som Frihed fordi
Heelheden kun virkeliggjør sig i Individerne, --- ogsaa kun Tale om, atter i
\emph{Almindelighed} at \emph{kunne handle godt} efterat Individets Erkjendelse
af det Rette har ført Villien i den rigtige Retning, ikke om at handle
\emph{saaledes} godt som der kunde været handlet uden hiin Afbrydelse ved det
Uethiske. I Virkeligheden bliver ved en saadan Opfattelse den enkelte Handling
løsnet fra Væsenet og seet som rent \emph{accidentel}. --- Hvorledes
Vanskeligheden med Hensyn til Tidsforholdet bliver at hæve, er i det
Foregaaende udviklet.}.

\begin{center}\rule{2cm}{0.4pt}\end{center}
% FOOTNOTE, p. 201: printed mark *) --- the only note on the page, set below a
% short single rule. „Anselm“ and „Strauß“ are letterspaced Fraktur. Inside the
% note „Almindelighed“ is spaced, the „at“ that follows it at the end of the
% line is tight, and „kunne handle godt“ on the next line is spaced again ---
% the compositor's split, transcribed as printed.
% RULE, p. 201: the division closes below the footnote with a short centred
% DOUBLE rule --- two thin lines, verified at 600 dpi at 200 % magnification.
% Same weight as the division-closing rules at pp. 4, 29, 63 and 210 (every
% other rule in the book is single). Set here with the single \rule the skeleton
% uses throughout, the printed weight being recorded in this comment.
% p. 201 is a SHORT page: the text block ends at „…den ideale ethiske Fordring*).“
% and the rest of the leaf is the footnote and the closing rule. The „Slutning“
% head is in the skeleton and belongs to the next batch; it is not set here.
% GREEK/FOOTNOTES for the batch: Greek occurs on p. 193 only (ἀνομία); footnotes
% on pp. 191 and 201 only. Every one of pp. 190--201 was read by eye for both.

\chapter*{Slutning.}
\addcontentsline{toc}{chapter}{Slutning}
\markboth{Slutning}{Slutning}

% --- p. 202 ---
\setcounter{footnote}{0}% p. 202 carries one note, printed *)
% The head „Slutning.“ (letterspaced, centred, short rule under it) and the rule
% are already in the skeleton and are not repeated here. p. 202 opens with a
% two-line ornamental initial „A“ (of „Af“); set normally, as at pp. 1, 5, 64, 97.
% ENUMERATION 1.--3., pp. 202--208. In the print the numeral stands at the
% PARAGRAPH INDENT and the following lines run flush to the left margin --- not
% the hanging setting of pp. 1--4, where the numeral stands at the left margin.
% Set here with the same enumitem call as pp. 1--4, for consistency across the
% edition; the divergence in the printing is recorded rather than reproduced.
\begin{enumerate}[label=\arabic*., leftmargin=*, labelsep=1em,
                  topsep=0.35\baselineskip, itemsep=0.35\baselineskip]
\item Af alle de tidligere Udviklinger fremgaaer det som Resultat, at der kan
  blive en væsentlig Enhed mellem den religiøse Bevidsthed og den humane
  Bevidsthed, „Culturbe\-%
vidstheden“, og denne Enhed kan og vil, fordi den er be\-%
grundet i Menneskenaturen, hævde sig trods det, at det Religiøse, hvor det endnu
  har det \emph{positivt Religiøses} Form, træder i \emph{Modsætning} til det
  Humane idet begges Væ\-%
sen begynder at klares for Bevidstheden. I denne Modsætning fastholder den
  religiøse Bevidsthed Dogmet om det, Supra\-%
naturale, det supranaturale Erkjendelsesprincip og det supra\-%
naturale Villiesprincip, medens den humane Bevidsthed fastholder Naturens
  Væsentlighed som Menneskelivets Grund\-%
lag, og den uforvandlede, fuldstændige Menneskenaturs Evne til ud fra sit eget
  Indre at naae til Erkjendelsens og Hand\-%
lingens Sandhed, hvis Kilde Menneskevæsenet gjemmer i sig. Denne Modsætning,
  der, indenfor Christendommen, traadte frem for Bevidstheden ved Middelalderens
  Slutning, klare\-%
des, under den philosophiske Tænknings Udfoldelse, mere og mere i de
  efterfølgende Aarhundreder\footnote{Jfr.\ Problemet om Tro og Viden, S.\ 21
  ff.}, og der udviklede sig, saavel fra realistiske som fra idealistiske
  Forudsætninger, en Opposition mod den christelige Verdens- og Livsanskuelse, og
  Philosophiens Strid mod det Religiøse bestemte enten directe den almindelige
  Bevidsthed eller paavirkede idetmindste indirecte den Udvikling, ved hvilken i
  den moderne Tid denne
% PRINTER'S DEFECT, p. 202: the page prints „Dogmet om det, Supranaturale“, with
% a COMMA between „det“ and „Supranaturale“. At 8x on the 300 dpi render the mark
% is unmistakably a comma --- a wedge descending below the baseline with a tail,
% the same shape as the comma of „naturale,“ on the line below --- and not a
% period, a speck or show-through. Both OCR witnesses register something there
% (tesseract „det.“, ABBYY „deb“). Transcribed as printed. It carries no quote
% marks, so no effect on the quote balance.
% EMPHASIS, p. 202: „positivt Religiøses“ spaced, „Form,“ tight; „Modsætning“
% (after „træder i“) spaced, but the second „Modsætning“ eight words later, in
% „I denne Modsætning“, is TIGHT --- checked side by side at 300 dpi.
% spacing.py's RUN „Aarhundreder *),“ is the footnote mark, not letterspacing:
% „Aarhundreder“ is tight on the image.
% p. 202 carries no Greek.
% --- p. 203 ---
\setcounter{footnote}{0}% p. 203 carries one note, printed *)
Bevidsthed fjernede sig fra det Christelige og fra det positivt Religiøse
  overhovedet.

  Der kunde vel, under Udviklingen af Modsætningen, og medens denne Udvikling
  tildeels førte til aabenbar Kamp, endnu for en Tid \emph{deelviis} bevares en
  Usikkerhed. En enkelt Form af Philosophien kunde\footnote{Jfr.\ Problemet om
  Tro og Viden, S.\ 23.}, idet den ved sit ensidigt idealistiske Princip fik en
  abstract Lighed med den christelige Verdensanskuelse, skuffe sig selv og søge
  en Enhed, hvor dog Modsætningen var det Principielle. Og idet det Christelige
  bestandig ved Siden af sig havde havt en Ud\-%
vikling, der paavirkedes af det, om den end havde sin Rod i et andet Livsprincip
  og bares af den fra Oldtiden nedarvede Anskuelse af Naturen og Menneskelivet,
  som den nye Tid selvstændig gjenerhvervede, saa kunde det, der havde udviklet
  sig samtidig med Christendommen, tildeels under Paavirkning af den, men ud fra
  et andet Princip, opfattes som Noget, der var fremgaaet af selve
  Christendommen, saaledes at denne altsaa blev seet som det Princip, der af sin
  indre Livsfylde kunde udfolde et Virkelighedsindhold, i Videnskab, i Kunst, i
  det sædelige og politiske Livs Former. Og saaledes kunde endnu i vor Tid den
  Anskuelse, der, idet den statuerer en relativ Enhed af det Christelige og det
  Humane, i Christendommen seer Prin\-%
cipet for en Verdensvirkeligheds Udformning i religiøs Aand, staae ligeoverfor
  den, der opfatter Christendommen som kun forholdende sig negativt til det
  Naturlige og til Menneske\-%
livets Udfoldelse i real Virkelighed, og har sine Repræ\-%
sentanter ligesaafuldt indenfor det positivt Religiøse som indenfor det Humane.
  Men der vil ikke kunne bestaae en Uvished om, hvilken af de to modsatte
  Opfattelser der er den sande og den historisk sikkrede, naar man ikke lader sig
  forvirre af de uundgaaelige Inconsequentser, i hvilke en Religion, der sætter
  Verdensfornægtelsen, Naturfornægtelsen, som sit Princip, maa komme ind, naar
  den skal \emph{vedblive at bestaae i den fremmede} Verden, hvis Undergang inden
  en kort Tidsfrist var den oprindelige Forventning.
% „Inconsequentser“: the Fraktur I/J sort, normalised to I per the house
% convention (ABBYY reads it „Jnconsequentser“).
% EMPHASIS, p. 203: only two runs --- „deelviis“, and „vedblive at bestaae i den
% fremmede“, where the following „Verden,“ falls back to tight type in the middle
% of the phrase. Another instance of the compositor's half-executed spacing;
% transcribed as printed. spacing.py's singles („denne“, „har“, „bestaae“,
% „fremmede“ taken separately) are the loose justification of those lines.
% p. 203 carries no Greek.
% --- p. 204 ---
Man analysere blot de christelige Forestillinger om den rent aandelige Gud, om
  Skabelsen af Intet, om Gudmennesket, om Salighedstilstanden og om
  Verdensundergangen, og Chri\-%
stendommens \emph{abstract spiritualistiske} Charakteer, dens \emph{Fornægten af
  det Naturlige} og derigjennem ogsaa af \emph{Menneskelivets Realitetsindhold},
  vil blive utvivlsom. Og idet nu hermed det forenes, som Christendommen maa dele
  med \emph{enhver} positiv Religion: \emph{Henlæggelsen af Sandhedens, den
  theoretiske og praktiske Sand\-%
heds, Princip i et transscendent Væsen}, ligeoverfor hvilket Mennesket staaer som
  det kun endelige, som det, der ved sin Natur er afmægtigt til at finde
  Sandheden, og idet denne Opfattelse sees som skærpet ved den christelige
  Forestilling om Syndighedstilstanden og den første Synds Følger, saa maa
  Modsætningen mellem den christelige og den humane Ver\-%
dens- og Livsanskuelse vise sig som fuldstændig.

  Men under denne Modsætning mellem det positivt Religiøse i dets høieste Form og
  den humane Bevidstheds Princip og Indhold viser sig dog ikke blot Muligheden,
  men Nødvendigheden, af en Forsoning i Aanden, ved hvilken det Humane og det
  Religiøse bringes til Enhed, og til en saadan Forsoning har selve den nyere
  Tids Udvikling ført hen.

  Idet den moderne Bevidsthed blev klar over sin Mod\-%
sætning til den christelig-middelalderlige, kunde Modsætningen til det
  transscendente Supranaturale ved sin Abstraction føre til en
  \emph{Endelighedsopfattelse af Naturen og Menneske\-%
livet}; men denne Opfattelse ender i Consequentser, der maae føre Aanden tilbage
  til Anerkjendelsen af det Religiøse efter dets \emph{væsentlige} Betydning,
  efter dets \emph{Uendelighedsindhold}. Den humane Bevidsthed tvinges ved den
  „reale Virkeligheds“ Endelighed og ved sin egen Endelighed til at søge et
  Uendeligt, og Endeligheden gjemmer i sig Sporet af den Uendelighed, ved hvilken
  den har Væren. Bevidstheden gjenfinder det Uendelige, idet den \emph{Natur},
  der for Mate\-%
rialismen blev et Modsigelsesfuldt og Uforklarligt, viser sig
% p. 204 opens FLUSH, not indented: „Man analysere blot ...“ continues the
% paragraph begun on p. 203, although p. 203 ends on a full stop.
% EMPHASIS, p. 204: the run „Fornægten af det Naturlige“ stops at „og
% derigjennem ogsaa af“, which is tight --- including the „af“ that governs the
% next spaced run „Menneskelivets Realitetsindhold“. Confirmed at 2.4x on the
% 300 dpi render; transcribed as printed, not regularised.
% „Natur“ (in „idet den Natur, der for Materialismen“) is letterspaced; neither
% spacing.py nor either OCR witness flagged it, and it is unambiguous on the
% image.
% p. 204 carries no footnote (the note visible on the image at the foot is
% show-through from the verso of p. 203) and no Greek.
% --- p. 205 ---
for Aanden i \emph{Livets og Idealitetens Lys}, og \emph{idet Bevidstheden
  fordyber sig i sig selv}, og erkjender sit \emph{Forhold} til den
  \emph{levende} Virkeligheds \emph{ideal-reale Prin\-%
cip}. I denne Fordybelse gjenfinder Selvet det \emph{religiøse Forhold},
  gjenfinder det ad \emph{Videns Vei}, og gjenfinder det \emph{uden at opgive sig
  selv} og uden at fornægte en \emph{Grund\-%
virksomhed} af sit Væsen. Den \emph{humane} Bevidsthed bliver, idet den søger ned
  til sit Væsens Grund, i Sandhed \emph{ethisk} Bevidsthed, og som saadan
  \emph{religiøs} Bevidsthed. Først ved dette Enhedsforhold faaer den humane
  Bevidsthed sin \emph{Sandhed}, saavist som Menneskeaanden først i dybeste
  Forstand bliver \emph{sig bevidst}, naar den bliver \emph{sig bevidst som
  grundende i det Guddommelige}, og som frit under\-%
ordnet dette. Og det Religiøse faaer i Enheden med det Humane \emph{Concretion og
  Virkelighed}, og det har i denne sin concrete Virkelighed sin \emph{Sandhed}
  som \emph{Princip for det virkelige menneskelige Livs Totalitet}. I denne Enhed
  bliver den religiøse Forsoning \emph{fuldstændig}, fordi Religionens Liv ikke
  leves ved en Flugt fra Menneskelivets Virkelighed, ved en Fornægten af
  væsentlige Yttringsformer af det: en Fornægten, der altid maa blive
  modsigende, da den kun bliver mulig ved selve Menneskenaturens Kræfter; --- men
  leves i \emph{enhver} væsentlig Form af Aandens Liv, og har sit
  \emph{positive} Udtryk i disse Livsformer, saaledes at det religiøse Forhold
  bliver et \emph{absolut} Forhold, omfattende \emph{hele} Livet, og forbindende
  Menneskeaanden med dens \emph{væ\-%
senslige} Absolute, med \emph{det i Sandhed Uendelige}. Og i denne Enhed er, med
  det \emph{absolute Forhold} og med \emph{For\-%
soningen}, --- Aandens Forsoning med \emph{sit Princip} og i \emph{sit eget
  Indre}, --- \emph{det i de positive Religioner Væsentlige} bevaret, medens den
  historiske Form er opløst, Forestillingstransscendentsen ophævet i det
  Immanentsfor\-%
hold, i hvilket Transscendentsen bliver Moment.

\item Den \emph{væsentlige Opgave} for nærværende Skrift kan altsaa bestemmes som
  den, \emph{at vise hen til den Re\-%
ligiøsitetens Form, som er tilstede}, om end ofte
% DOUBTFUL READING, p. 205: the print gives „dens væ=/senslige Absolute“, not
% „væsentlige“. At 3.2x the fourth letter of the second half is the round s of
% Fraktur, not t --- compare „Livets“ on the line above; both OCR witnesses read
% s as well. The alternative (a compositor's t-for-s, i.e. „væsentlige“) is
% recorded and rejected on the strength of the glyph; Brøchner's own compounds in
% Væsens- (e.g. „Væsensvillen“, p. 22) make the form defensible. As printed.
% EMPHASIS, p. 205: „frit underordnet dette“ is TIGHT, though ABBYY's spacing
% marks suggested otherwise; „som“ between „sin Sandhed“ and „Princip for ...“ is
% likewise tight in the middle of two spaced runs. Both checked at 2.4x.
% „Idealitetens“, not ABBYY's „Identitetens“: clear on the image.
% p. 205 carries no footnote and no Greek.
% --- p. 206 ---
\emph{uklart og halvbevidst, i vor Tids humane Bevidst\-%
hed}, som, klarere eller dunklere, har forkyndt sig i denne Be\-%
vidsthed \emph{fra den nyere Tids Begyndelse}, som har gjort sig gjældende i de
  sociale og sædelige Forhold, i Humani\-%
tetsbestræbelserne, i Aandslivets Form overhovedet, og som kun behøver at
  \emph{frigjøres og klares} gjennem Bevidst\-%
hedens Selvfordybelse, forat vise sig i sin Væsentlighed, som Eet med det i de
  positive Religioner dybest Bevægende og som Eet med Menneskeaandens inderste
  Væsen. Det er \emph{Virkelighedens Religion}, der \emph{ikke} skal stiftes
  eller skabes, men kun paavises som værende og som begrundet i Menneskeaandens
  Udviklingslov; --- der \emph{ikke} skal være en rationaliserende Omdannelse af
  en positiv Religion eller selv en ny positiv Religion, men \emph{Udtrykket for
  Religiøsi\-%
tetens evige Aand som det hele Menneskelivs Livs\-%
princip}. Derfor har det været Opgaven at vise denne Virkelighedens sande og
  evige Religion som den, i hvilken det \emph{totale} Menneskeliv er i Enhed og
  Samklang, og i hvilken derfor den \emph{humane Bevidsthed} kan finde
  \emph{reli\-%
giøs Forsoning}. Løsningen af Problemet om Troens og Videns Forhold er derfor
  søgt ud fra det, der i det udelelige Menneskevæsen har Begrebsprincipatet, men
  saa\-%
ledes, at det i Menneskeaandens væsentlige Yttringsformer Eiendommelige er
  anerkjendt, at Menneskeaandens Virken efter sin concrete Bestemthed er kommen
  til sin Gyldighed. Hvad der i Indledningen til Skriftet: „Problemet om Tro og
  Viden“, blev antydet, har faaet sin nærmere Udførelse. I Menneskeaandens
  \emph{Væsen} er det \emph{religiøse Forholds Nødvendighed} søgt, og den har
  viist sig som Nødvendig\-%
heden af et i Sandhed \emph{absolut Forhold}, af et \emph{sandt
  Uendelighedsforhold}, af et saadant Forhold, i hvilket \emph{den fulde
  Forsoning: Forsoningen indenfor Virke\-%
ligheden}, kan findes, idet Menneskeaanden, i den \emph{abso\-%
lute Hengivelse til Principet}, bliver i \emph{indre Enhed med sig}. I dette
  Forhold ere \emph{Religion, Videnskab og Sædelighed} seete i indre Enhed. Ud
  fra Anerkjen\-%
% p. 206 is set almost throughout in letterspaced type; the tight passages are
% the exception. „Væsen“ in „I Menneskeaandens Væsen“ is spaced while the
% following „er det“ is tight (checked at 2.6x). The small dot to the right of
% „Enhed“ on the „indre Enhed“ line is a speck, not a hyphen (tesseract read one).
% p. 206 carries no footnote and no Greek.
% --- p. 207 ---
delsen af Menneskeaandens Naturgrund som et ved For\-%
holdet til det absolute Princip Væsentligt, er der viist hen til en \emph{ethisk
  Handlen}, der faaer et \emph{Virkelighedsind\-%
hold} ved denne Naturgrund, og som, idet den faaer det \emph{Religiøses Præg,
  udfylder det Religiøse med Vir\-%
kelighed}, og nærmere, med en \emph{tankebestemt, fornuft\-%
mæssig Virkelighed}.

\item Det Standpunkt, der i dette Skrift er gjort gjæl\-%
dende, sætter altsaa \emph{ikke et abstract Metaphysisk} i den \emph{religiøse
  Livsvirkeligheds} Sted. Det finder i selve Principet Garantien for det Reales
  Væsentlighed, og ved at accentuere denne skaffer det Plads for et
  \emph{Frihedsforhold} i Modsætning til det abstract metaphysiske
  Nødvendigheds\-%
forhold, og sikkrer det Reales nødvendige Væreforms væsent\-%
lige Betydning for det ethisk-religiøse Forhold; og herved, og ved at lade dette
  Forhold omfatte en concret Livstotalitet, hævder det det netop som et
  \emph{Virkelighedsforhold}. Her\-%
imod strider det ikke, at der overalt gaaes tilbage til meta\-%
physiske Bestemmelser; thi det Metaphysiske bevarer sin \emph{al\-%
mene} Gyldighed; men det maa kun bestemmes saaledes, at det viser hen \emph{til
  det Reale}, til den endelige \emph{Virkelighedens Livsbetingelser}.

  Det Standpunkt, der er gjort gjældende, kan ikke heller med Rette, som Nogle
  have villet, betegnes som et \emph{negativt}, og som saadant bringes sammen med
  de Standpunkter, fra hvilket de Angreb paa det Christelige ere gjorte, der
  tilhøre vor Tid og den nærmest foregaaende. Dets Forhold til den betydeligste
  Repræsentant for den negative Retning i vor Tid, til \emph{Feuerbach}, er
  ovenfor udførlig bleven fremstillet. Gjennem denne Fremstilling, og gjennem
  alle de tidligere Udviklinger overhovedet, vil det have viist sig, at det kun
  forholder sig negativt til de positive Religioner idet det ne\-%
gerer \emph{Endeligheds}bestemtheden i dem, men at det gjennem denne Negation
  netop stræber at hævde \emph{det væsentlig Religiøses Uendelighed}. Fra dette
  Synspunkt bliver Standpunktet netop at betegne som \emph{positivt}: det
  Negative
% „til den endelige Virkelighedens Livsbetingelser“: the print reads
% „Virkelighedens“ --- letterspaced --- with no comma; both witnesses agree.
% As printed.
% LETTERSPACING ACROSS A LINE BREAK, p. 207: in „ne=/gerer
% Endeligheds bestemtheden“ the compositor spaces „Endeligheds“ and then
% falls back to tight type for „bestemtheden“ inside the same word. At 2.6x.
% Transcribed as printed, as at pp. 6, 15, 43 and 205.
% „fra hvilket de Angreb ... ere gjorte“: as printed; the antecedent
% („Standpunkter“) is plural, but „hvilket“ is a word in the language and the
% burden of proof is on the claim of a defect (playbook §2). Not emended.
% „Feuerbach“ is set in letterspaced FRAKTUR, not antiqua --- the book's rule for
% proper names. No antiqua and no Greek anywhere on p. 207. No footnote.
% --- p. 208 ---
bliver Bestemmelse for et Saadant, der, idet det forholder sig polemisk ud fra et
  \emph{ensidigt} Udgangspunkt, \emph{kun kan nedbryde}, og for et Saadant, der
  \emph{bevidst fastholder Endelighedsbestemtheden}. ---

  Skulde nu dog Nogen fra det negative Forhold, jeg ind\-%
tager til Endelighedsmomentet i de positive Religioner, medens jeg stræber at
  hævde det i Religionerne Væsentlige, Evige og Uendelige, ville tage Anledning
  til, paa Syko\-%
phanters Viis at henkaste de paa Ukyndigheden velberegnede, og for Fanatismen
  altid velkomne, Beskyldninger for Atheisme, ulmende Had til Christendommen, og
  hvad det nu ellers er, der er Stikordene, saa skal det fra min Side kun blive
  besvaret med den Taushed, der er Foragtens. Jeg har engang i et andet Skrift
  udtalt mig om mit personlige For\-%
hold til den positive Religion, og en Udtalelse af den Art er det meningsløst at
  gjentage.
\end{enumerate}

\begin{center}\rule{2cm}{0.4pt}\end{center}
% A short centred rule stands here on p. 208, between the close of the numbered
% argument and the book's valedictory paragraphs. It is a SINGLE hairline, as at
% pp. 89 and 143 --- not the double rule that closes the divisions. The
% enumeration is closed above it because nothing after it is numbered.

Hermed afslutter jeg saa dette Skrift og overgiver det til deres Dom, der ere
istand til med Indsigt at følge en videnskabelig Undersøgelse, og at kunne
underkaste den en \emph{videnskabelig} Kritik. En saadan Kritik vil bestandig
være mig velkommen, og hvor fjernt det Standpunkt, fra hvilket den udgaaer, end
maatte ligge fra mit, skal jeg altid være rede til at anerkjende, hvad Kritiken
maatte bringe af vir\-%
kelig Gyldigt og Sandt. Men ogsaa kun med en saadan Kritik, der hviler paa en
virkelig Forstaaelse af, hvad jeg har givet, som har Sagen selv til sit Øiemed,
og som kæm\-%
per med Grunde, Videnskaben maa anerkjende, vil jeg kunne indlade mig. Naar en
nylig udkommen, saakaldet „\emph{Reli\-%
gionsphilosophie}“ istedenfor at modbevise mit Beviis mod Muligheden af en Videns
Selvbegrændsning, af Videns Anerkjenden af et med den \emph{absolut} uensartet
\emph{Princips Gyldighed}, indskrænker sig til at paavise det Factum:
% p. 208 carries no footnote and no Greek. „Ukyndigheden“ (ABBYY's
% „Myndigheden“ is wrong) and „velkomne“ (tesseract's „velflomne“ is wrong),
% both settled on the image.
% --- p. 209 ---
at Videnskaben erklærer den positive Religions Princip for \emph{uforeneligt} med
sit --- (netop hvad jeg har urgeret!) --- men ikke gjør nogetsomhelst alvorligt
Forsøg paa at begrunde, at Viden maa anerkjende dette Princip \emph{som gyl\-%
digt}, eller at den overhovedet \emph{kan} erkjende Noget som faldende
\emph{udenfor} dens Grændser; naar denne „Religionsphilo\-%
sophie“ ved et simpelt Postulat gjør Differentsen mellem ab\-%
stract og concret Viden til en \emph{absolut} Modsætning mellem to Arter af
Viden, og ved lignende Postulater vil naae til, at kun den personlige
Tilfredsstillelse, den positive Religion giver, er den sande; og naar den endelig
først trivialiserer og saa bortescamoterer den „absolute Uensartethed“, der
skulde give den dens Charakteer; --- saa er en Forhandling med en saadan omvendt
i Luften svævende Videnskab overflødig; der er simpelthen at vise tilbage til den
Kritik, jeg har givet af den i mit tidligere Skrift; der behøver ikke at føies
Noget dertil.

Fra \emph{Theologiens} Side er der hidtil ikke gjort noget Forsøg paa at
imødegaae hvad der var Hovedpunkterne i min Kritik af den: min \emph{Benægtelse
af Theologiens Ret som Videnskab og af dens Gudsbegrebs Gyl\-%
dighed}. Jeg retter en Opfordring til dem blandt vore Theo\-%
loger, der have Sands for og Evne til en videnskabelig Under\-%
søgelse, til at gjøre et saadant Forsøg, og da til at gjøre det, ikke ad
\emph{indirecte} Vei, ved at paavise mulige Mangler og Mod\-%
sigelser i min Theorie, men ad \emph{directe} Vei, ved at modbevise min Kritiks
Grunde. Ad hiin indirecte Vei vilde intet endeligt Resultat kunne naaes; thi om
det end blev godt\-%
gjort, at jeg havde feilet nok saa meget i hvad jeg \emph{posi\-%
tivt} har opstillet, saa kunde det dog ikke blive til nogen Fordeel for
Theologien, saalænge det ikke var paaviist, at hine Feil i det Positive stode i
indre Sammenhæng med min Kritik; men selve denne Paaviisning vilde jo forvandle
det indirecte Angreb til et directe. Kun ved en directe Kritik af min Kritik kan
et Resultat naaes, og det fore\-%
kommer mig, at Sagen har Vigtighed nok til at der bør
% „Religionsphilosophie“ is set TIGHT inside its quotation marks here, whereas on
% p. 208 the same title is letterspaced inside the marks. As printed on each page.
% The signature numeral „14“ at the foot of p. 209 is not transcribed, per the
% book's convention.
% „til et directe.“ ends the sentence with a full stop (tesseract reads a comma);
% „Fordeel“, not tesseract's „Vordeel“. Both settled on the image.
% p. 209 carries no footnote and no Greek.
% --- p. 210 ---
stræbes derefter. Hvad der nylig af en af vor Theologies fortrinligste
Repræsentanter er blevet bebreidet en vis theologisk Retning, at den, istedenfor
at modbevise alvorlige Angreb, bortsnakker eller borttier dem, burde ikke kunne
gjøres til Anke mod Theologien overhovedet, ogsaa mod den Theologie, der
repræsenteres af Mænd, som virkelig ere Videnskabs\-%
mænd, om de end ikke ere det \emph{som} Theologer.

\begin{center}\rule{2cm}{0.4pt}\end{center}
% p. 210 is the LAST page of the book's text and is a short page: seven lines of
% type, then the closing rule. The rule is a DOUBLE rule --- two thin parallel
% lines, confirmed at 4x on the 300 dpi render --- as at pp. 4, 29, 63 and 201,
% and unlike every other rule in the book, which is single. Set with the same
% single \rule the skeleton uses for every short rule in this book.
% The mark between „dem,“ and „burde“ (tesseract reads a dash) is a dust speck:
% tiny, round, at mid x-height, with none of the weight of type.
% EMPHASIS, p. 210: „som“ alone, in „ere det som Theologer“, is letterspaced;
% „Theologer.“ is tight. The only emphasis on the page.
% p. 210 carries no footnote and no Greek. Printed p. 211, the Trykfeil leaf,
% is already in the skeleton and is not transcribed here.

% =====================================================================
%  TRYKFEIL  (printed p. 211 = PDF 222)
%  As printed, including the two blocks that belong to the earlier
%  volumes. The leaf carries no folio. Heading letterspaced and centred.
% =====================================================================
\chapter*{Trykfeil:}
\addcontentsline{toc}{chapter}{Trykfeil}
\markboth{Trykfeil}{Trykfeil}

% --- p. 211 ---
\begin{center}
\emph{Trykfeil:}
\end{center}

\noindent Side 62, Lin. 24: fra det, læs: for det.

\begin{center}\rule{2cm}{0.4pt}\end{center}

\begin{center}
i „Problemet om Tro og Viden“.
\end{center}

\noindent
\begin{tabbing}
Side \=\ 114,\ \ \=\ \ 26: \=\kill
Side \> 14, \> Lin.\ \ 7: \> sat, læs: satte.\\
--- \> 29, \> --- \ \ 2 f.\ n: \> findes. l.: finder.\\
--- \> 42, \> --- \ 20: \> afgivende, l.: afvigende.\\
--- \> 51, \> --- \ \ 1: \> om sine, l.: og sine.\\
--- \> 71, \> --- \ 22: \> Indhold, l.: Forhold.\\
--- \> 114, \> --- \ 16: \> dobbeltlydige, l.: dobbelttydige.\\
--- \> 119, \> --- \ 22: \> man, l.: men.\\
--- \> 176, \> --- \ 26: \> viser, l.: vises.\\
--- \> 223, \> --- \ 15: \> som at, l.: som i at.
\end{tabbing}

\begin{center}\rule{2cm}{0.4pt}\end{center}

\begin{center}
i „Et Svar til Prof.\ R.\ Nielsen“ (2det Oplag).
\end{center}

\noindent
\begin{tabbing}
Side \=\ 22,\ \ \=\ \ 14: \=\kill
Side \> 6, \> Lin.\ 14: \> fører, l.: føres.\\
--- \> 22, \> --- \ \ 6: \> lønner sig, l.: lønner.
\end{tabbing}

\end{document}
